%=============================================================================
%  chapter4.tex -- Does the Model Work?
%  Uses macros from thesis_main.tex Part 2. Defines no notation of its own.
%
%  COORDINATOR NOTES (delete before submission):
%
%  1. thesis_main.tex carries an inline skeleton for this chapter at the
%     \chapter{Does the Model Work?} block. This file reproduces the labels
%     ch:validation, sec:gate_philosophy, sec:gates_abc, sec:gate_e,
%     sec:fujimoto, sec:balmer_equivalence, sec:errors_found and
%     sec:validation_summary, so replacing the skeleton with %=============================================================================
%  chapter4.tex -- Does the Model Work?
%  Uses macros from thesis_main.tex Part 2. Defines no notation of its own.
%
%  COORDINATOR NOTES (delete before submission):
%
%  1. thesis_main.tex carries an inline skeleton for this chapter at the
%     \chapter{Does the Model Work?} block. This file reproduces the labels
%     ch:validation, sec:gate_philosophy, sec:gates_abc, sec:gate_e,
%     sec:fujimoto, sec:balmer_equivalence, sec:errors_found and
%     sec:validation_summary, so replacing the skeleton with %=============================================================================
%  chapter4.tex -- Does the Model Work?
%  Uses macros from thesis_main.tex Part 2. Defines no notation of its own.
%
%  COORDINATOR NOTES (delete before submission):
%
%  1. thesis_main.tex carries an inline skeleton for this chapter at the
%     \chapter{Does the Model Work?} block. This file reproduces the labels
%     ch:validation, sec:gate_philosophy, sec:gates_abc, sec:gate_e,
%     sec:fujimoto, sec:balmer_equivalence, sec:errors_found and
%     sec:validation_summary, so replacing the skeleton with %=============================================================================
%  chapter4.tex -- Does the Model Work?
%  Uses macros from thesis_main.tex Part 2. Defines no notation of its own.
%
%  COORDINATOR NOTES (delete before submission):
%
%  1. thesis_main.tex carries an inline skeleton for this chapter at the
%     \chapter{Does the Model Work?} block. This file reproduces the labels
%     ch:validation, sec:gate_philosophy, sec:gates_abc, sec:gate_e,
%     sec:fujimoto, sec:balmer_equivalence, sec:errors_found and
%     sec:validation_summary, so replacing the skeleton with \input{chapter4}
%     breaks no existing cross-reference.
%
%  2. OVERLAP, RESOLVED BY HOMES. Three passages exist in both chapters:
%       chapter6 sec:open_system   vs  Section~\ref{sec:attacks}   (home: ch4)
%       chapter6 sec:r1_deficit    vs  Section~\ref{sec:fujimoto}
%       chapter6 sec:nmax          vs  Section~\ref{sec:convergence}
%     The Chapter 6 versions carry the SCOPE consequence; the Chapter 4
%     versions carry the EVIDENTIAL standing. On the r_1 deficit: the table
%     is bundled-n (backlog I.1, 10 Sep 2026; findings_10 ADDENDUM I), the
%     l-closure was tested and does NOT account for the gap, and the deficit
%     DOES reach a_3 and a_4; its consequence for Delta, the cap and S is
%     quantified in the Status paragraph of sec:fujimoto
%     (validation/fujimoto_r1_consequence/). Do not reintroduce the earlier
%     "does not propagate" sentence.
%
%     [17-18 Sep 2026] verify_bundling_psm20.py HAS now been executed
%     (sec:bundling_check in chapter 6, outputs/bundling/); an earlier
%     version of this comment said the "never executed" sentence was still
%     true and must not be corrected. It was corrected. The n_max
%     convergence result carried in Section~\ref{sec:convergence} is now
%     stamped (verify_nmax_downward_scan.py, 18 Sep 2026); the working-note
%     values (ADDENDUM D.8) it replaced had the opposite sign at the
%     benchmark and the cold corner, and the text records that. Its \todo asking for provenance on the
%     4.9-6.2x terminal-shell excess is also closed there, by the internal
%     terminal-vs-interior measurement (a_p at [49,0]: 9.4x at p=8, 10.7x at
%     p=9, 6.3x at p=10), which does not bracket it (band 6.3 to 10.7 against
%     4.9 to 6.2) but lies within a factor 2.2 of it and shares the density
%     trend; stamped 21 Sep 2026 by verify_terminal_vs_interior.py.
%
%  3. chapter3.tex:222-225 asserts a severity for the conservation check that
%     Section~\ref{sec:conservation_severity} measures and refutes. Either
%     soften line 224 or let this chapter carry the correction; do not leave
%     both standing as written.
%=============================================================================

\chapter{Does the Model Work?}
\label{ch:validation}

Chapter~\ref{ch:theory} built a $43\times43$ rate matrix, eliminated the fast
variables, and extracted from the result a closed-form statement about how a
line ratio responds to its reservoir. Chapter~\ref{ch:results} will use that
machinery to make quantitative claims about a diagnostic in use on real
machines. Between the two sits a question that no amount of further algebra
answers: why should anyone believe a number this model produces?

The answer is not a single verdict. It is a graded one, because the
checks that have been run on this model are not of equal worth, and several of
the ones that look most impressive are worth almost nothing. This chapter sorts
them.

\medskip
\noindent\textbf{What this chapter takes as given.} The state space, the atomic
data and their sources, and the benchmark of the excitation rates against an
independent R-matrix-with-pseudostates (RMPS) calculation are established in
Chapter~\ref{ch:atoms}, Sections~\ref{sec:atomic_uncertainty}
and~\ref{sec:atoms_corrections}, and are not repeated. The construction of
$\Lmat$, the conservation identity, the quasi-steady-state (QSS) elimination,
and the two-channel decomposition of the excited populations come from
Chapter~\ref{ch:theory}. This chapter asks only what the checks run on that
construction are entitled to establish.

\medskip
\noindent\textbf{What this chapter does not do.} It reports no result about
divertor plasmas. Every number below is about the instrument, not about the
physics the instrument is pointed at. Keeping those separate matters here more
than anywhere else in the thesis: a model that predicts something interesting
is not thereby a correct model, and a check that uses the interesting
prediction as its evidence has verified nothing.

%-----------------------------------------------------------------------------
\section{What a check can and cannot show}
\label{sec:gate_philosophy}
% QUESTION: what distinguishes a check that could have failed from one that
% could not, and how is that distinction measured rather than asserted?

\subsection{Three things that get called validation}

Three activities are routinely reported under one heading, and they license
different conclusions.

\emph{Verification} asks whether the equations that were intended got solved
correctly. Conservation residuals, detailed balance, matrix identities,
convergence with tolerance, agreement between two independent code paths: all
of these compare the model against itself or against arithmetic. They can
detect a transposed index. They cannot detect a rate coefficient that is wrong
by a factor of three, because a uniformly wrong dataset satisfies every one of
them.

\emph{Validation} asks whether the model reproduces something established
independently of it. Published coefficients from another group's code, measured
cross sections, a tabulated limit. Only this class of check can reach the
atomic physics.

\emph{Analysis} asks what the model then says. That is Chapter~\ref{ch:results},
and none of it appears here.

\subsection{Severity, and how to measure it}

The useful property of a check is not whether it passes. It is the size of the
set of wrong models it would reject. Call that its severity.

The failure mode is the one a mass balance on a continuous stirred-tank
reactor shows. Compute the outlet stream as inlet minus consumption, then
verify that the balance closes. It closes to machine
precision, always, for any rate constant whatever, because the outlet was
constructed from the balance being tested. The check is real arithmetic and it
is entirely uninformative about the kinetics. A model can be wired correctly
and be wrong.

Severity is measurable, and two operations measure it.

\begin{enumerate}
\item \textbf{Provenance inspection.} Ask whether the quantity under test was
derived from the quantity doing the testing. If the de-excitation coefficients
were computed from the excitation coefficients through detailed balance, then a
detailed-balance check on them is arithmetic, not physics. Chapter~\ref{ch:atoms}
already states this for the rate dataset (Section~\ref{sec:two_energies}); this
chapter carries it through to the matrix.

\item \textbf{Fault injection.} Break the model deliberately, in a way that a
real mistake could plausibly break it, and see whether the check notices. A
check that survives a tenfold error in a rate coefficient has told you that a
tenfold error in that coefficient is invisible to it.
\end{enumerate}

Fault injection is the sharper of the two, because it produces a number rather
than an argument. It is used in Section~\ref{sec:conservation_severity} on the
conservation check, and its result is the least comfortable finding in this
chapter.

\subsection{The three grades used below}

Every check in this chapter is placed in one of three grades, and the grade is
stated wherever the check is reported.

\begin{description}
\item[Tautological.] Cannot fail except on an arithmetic bug, because the
quantity tested was constructed from the relation being tested. It confirms
that the code implements the construction. It is not evidence about the
physics, and reporting it as though it were makes the validation weaker, not
stronger, by displacing the checks that carry information.

\item[Weak.] Could fail in principle, but the tolerance is so much looser than
the measured value that only a gross error would trip it. A gate that passes
with a factor of $150$ in hand has told you the answer is not absurd.

\item[Severe.] A plausible wrong version of this model fails it. Either the
check has actually fired at some point in the project's history, or a specific
corruption is known to trip it.
\end{description}

Four quantities currently printed by this project's own scripts as validation
are tautological by this definition, and Table~\ref{tab:ladder} names them:
Gate~A, the column-sum conservation residual, the $\tanh$ bound, and Gate~E.
Sections~\ref{sec:conservation_severity} to~\ref{sec:severe_checks} give the
measurements behind each.

%-----------------------------------------------------------------------------
\section{The conservation check is exact by construction}
\label{sec:conservation_severity}
% QUESTION: does the column-sum identity detect a wrong rate coefficient?

Section~\ref{sec:conservation} established that each column of $\Lmat$ must sum
to the ionisation leak from the level that column represents, and reported the
measured residual of Eq.~\eqref{eq:conservation_measured} as
$3.61\times10^{-11}$, a relative residual (Section~\ref{sec:conservation}). The
$3.61\times10^{-11}$ is the maximum over every column at every one of the $400$
grid points; the fault-injection run below has its own baseline, $2.238\times
10^{-12}$, at the single point it was run at. The claim attached to it there is that ``an error in a
single off-diagonal element, a transposed index, or a missing back-reaction all
break this immediately.''

That claim is testable, and it is false.

\subsection{Four injected faults}

Four faults were introduced into the assembly, one at a time, and the residual
recomputed after each. The first three are the exact failure modes the claim
names. Table~\ref{tab:faultinject} gives the outcome.

\begin{table}[tbp]
  \centering
  \caption[Fault injection into the conservation check]{Four deliberate faults
  injected into the assembly of $\Lmat$, with the resulting change in the
  matrix and the conservation residual of
  Eq.~\eqref{eq:conservation_measured} recomputed after each. Three faults
  large enough to move individual matrix elements by up to
  $2.7\times10^{11}\,\si{\per\second}$ leave the residual unchanged at the
  baseline value; only an inconsistency between the Einstein coefficients $A$
  and the total radiative loss rates $\gamma$ is detected. Source:
  \texttt{outputs/CHANGE\_REPORT.md}~\S4.2.}
  \label{tab:faultinject}
  \begin{tabular}{lrl}
    \toprule
    injected fault & $\max|\Delta \Lmat|$ (\si{\per\second}) & residual \\
    \midrule
    none (baseline)                     & $0$                  & $2.238\times10^{-12}$ \\
    $K_{\rm exc}(1s\to2p)$ $\times$ ten   & $6.9\times10^{5}$    & $2.238\times10^{-12}$ \\
    all de-excitation deleted           & $1.9\times10^{11}$   & $2.171\times10^{-12}$ \\
    excitation array transposed         & $2.7\times10^{11}$   & $1.837\times10^{-12}$ \\
    $A_{pq}$ halved, $\gamma_p$ left alone & $3.1\times10^{8}$   & $3.63\times10^{1}$ \textbf{caught} \\
    \bottomrule
  \end{tabular}
\end{table}

A tenfold error in the single most important excitation coefficient in the
model, the one that supplies the entire ground-fed channel of
Section~\ref{sec:R_depends_on_b}, changes the residual in no digit. Deleting
every de-excitation process in the matrix, which is not a subtle mistake,
changes it in the third digit and leaves it three orders of magnitude below any
tolerance. Transposing the whole excitation input array, the classic
index-order mistake, does the same.

\subsection{Why, and what the check is actually worth}

The reason is in the assembly, not in the physics. The diagonal of $\Lmat$ is
built as minus the column sum of the same off-diagonal array that the check
then sums. Whatever numbers are in that array, correct or corrupted, the column
sum returns the ionisation leak. The identity of Eq.~\eqref{eq:conservation} is
imposed by the construction rather than tested by it. This is the reactor mass
balance of Section~\ref{sec:gate_philosophy} exactly.

The one fault that was caught is the one that breaks the construction. In the
notation of Eq.~\eqref{eq:level_balance}, $A_{pq}$ is the Einstein coefficient
for spontaneous emission from level $p$ to level $q$, in \si{\per\second}, and
$\gamma_p = \sum_{q<p} A_{pq}$ is the total radiative loss rate from $p$, also in
\si{\per\second}. The individual $A_{pq}$ enter the off-diagonal elements while
their sum $\gamma_p$ enters the diagonal, so halving one without the other makes
the two inconsistent, and the residual rises from $2.238\times10^{-12}$ to
$36.3$, thirteen orders of magnitude. The check raises rather than warns, it
costs nothing to run, and on this evidence that is the width of what it covers.

The rate magnitudes are reached only by Section~\ref{sec:atomic_uncertainty}'s
comparison against RMPS and by the external comparisons of
Section~\ref{sec:fujimoto}, and by nothing internal.

\paragraph{The grid point of the table.} Three of the four
$\max|\Delta\Lmat|$ values in Table~\ref{tab:faultinject} scale with $\nel$, so
the table fixes its own grid point. Sweeping all $400$ grid points and matching on those three values
identifies $[23,5]$, the benchmark, with a worst disagreement of
\SI{1.63}{\percent}, the rounding of the printed table to two significant
figures; every residual in Table~\ref{tab:faultinject} then reproduces to the
digits shown (\texttt{verify\_fault\_injection.py}, stamped in
\texttt{validation/fault\_injection/}). The script rebuilds $\Lmat$ through the
pipeline's own assembly rather than patching the stored matrix, and the rebuild
reproduces \texttt{L\_grid} exactly before anything is injected, which makes the
faults faults in the model rather than in a re-implementation of it.

%-----------------------------------------------------------------------------
\section{The gate suite, read one gate at a time}
\label{sec:gates_abc}
% QUESTION: for each automated gate, what quantity does it compare, against
% what tolerance, and what would have to be wrong for it to fail?

Five automated gates run over the grid in \texttt{src/validation/validate\_gates.py}
and write \texttt{validation/gate\_summary.txt}. Their headline outcome is four
passes and one failure. Read as a headline it says little, because the five
differ in severity by a wide margin and two of them cannot fail. This section
takes them one at a time; Gate~D, the failure, gets
Section~\ref{sec:gate_d} to itself.

\subsection{Gate A: detailed balance, and where its physics content actually
lives}

Gate~A takes every excitation rate coefficient $K_{\rm exc}(i\to j)$ in the
dataset, in \si{\cubic\centi\metre\per\second}, and predicts the reverse
coefficient from the Boltzmann relation between forward and reverse rates in a
Maxwellian electron distribution at temperature $\Te$,
\begin{equation}
  K_{\rm deexc}(j\to i)
  = K_{\rm exc}(i\to j)\,\frac{g_i}{g_j}\,
    \exp\!\left(\frac{E_j - E_i}{k_B\Te}\right),
  \label{eq:gateA}
\end{equation}
then compares the prediction against the stored reverse coefficient at three
temperatures spanning the grid. Here $g_i$ is the dimensionless statistical
weight of level $i$, $E_j-E_i$ is the excitation energy in
\si{\electronvolt}, and $k_B$ is Boltzmann's constant. The relation holds only
for a Maxwellian distribution, which Section~\ref{sec:maxwellian} established
and which is assumed throughout. The gate passes
if the maximum relative error is below \SI{0.01}{\percent}
(\texttt{validate\_gates.py}:118). Measured: below \SI{0.01}{\percent} over all
$819$ pairs at all three temperatures.

\textbf{Grade: tautological.} The stored de-excitation coefficients were
produced from the excitation coefficients through Eq.~\eqref{eq:gateA} in the
first place (\texttt{compute\_K\_CCC.py}, function \texttt{detailed\_balance}),
as Section~\ref{sec:two_energies} states. Gate~A confirms that the same formula
was applied twice with the same energies and weights. That is a wiring check,
and Chapter~\ref{ch:atoms} already flagged it as one.

The physics content of detailed balance in this project sits somewhere else and
is currently invisible in the outputs. \texttt{src/parsers/qc\_ccc.py}:87--168
tests whether Bray's \emph{raw} convergent-close-coupling (CCC) cross sections,
which this project did not produce, satisfy microscopic reversibility on their
own. The ratio of the two directions has mean $0.9995$ with standard deviation
$0.0032$, and \SI{98.7}{\percent} of comparisons fall within \SI{1}{\percent}
(\texttt{ccc\_qc\_report.png}). A \SI{0.05}{\percent} mean bias with
\SI{0.32}{\percent} scatter on external data is a genuine external check, and it
is the one to lead with.

\subsection{One energy ladder on both sides, which is severe}

There is a second detailed-balance result, and unlike Gate~A it could have
failed.

A model of this kind carries energies twice: once in the Boltzmann factors that
relate forward and reverse collisional rates, and once in the Saha factors that
relate ionisation to three-body recombination. If two different energy tables
were used, or one table were indexed differently from the other, the model would
still conserve particles, still satisfy Gate~A, and still produce plausible
populations. Nothing internal would notice.

Testing ionisation against three-body recombination through the Saha relation
at every level, with CODATA constants on the Saha side, gives a ratio of
$0.999997$ for every one of the $43$ states at every temperature, identical
across all $2150$ entries to $1.4\times10^{-15}$
(\texttt{verify\_saha\_balance.py}, \texttt{validation/saha\_balance/}); the
residual is the rate module's rounded values of $h$, $m_e$ and the
electron-volt against CODATA, and nothing else. The point is not the
closeness to unity, which detailed balance imposes. It is the identity of the residual across
levels whose Saha exponents span the full binding-energy ladder of hydrogen,
from $\mathrm{e}^{13.6\,\mathrm{eV}/k_B\Te}$ for the ground state down to
$\mathrm{e}^{0.06\,\mathrm{eV}/k_B\Te}$ for the top shell. At
$\Te = \SI{1}{\electronvolt}$ those two factors differ by
$\mathrm{e}^{13.54}$, a factor of $7.6\times10^{5}$. Two energy ladders that
differed would produce a residual drifting across that range; the measured
residual does not drift, and its common value is the eight-figure rounding of
the stored state table. \textbf{Grade: severe.} The check could have failed and
did not.

The same audit gives the collisional side over the $819$ excitation pairs at
each of three temperatures, $2457$ comparisons in total: maximum deviation
$7.31\times10^{-9}$, median $7.5\times10^{-10}$, again consistent with the state
table's rounding.

\subsection{Gate B: the coronal limit}

At low enough density an excited level is populated by electron impact from the
ground state and empties radiatively before a second collision reaches it, so
its population per unit ground-state density and per unit electron density
becomes independent of $\nel$. Gate~B tests that plateau: it computes
$n_{2p}/(\ngs\nel)$ and $n_{3d}/(\ngs\nel)$ at the two lowest densities on the
grid and asks whether they have converged.

The tolerance is \SI{50}{\percent} per temperature
(\texttt{validate\_gates.py}:175) and the gate passes if \SI{70}{\percent} of
temperatures pass (line 181). Measured: all $50$ temperatures pass, worst
convergence $0.0034$.

\textbf{Grade: weak.} Passing with a factor of $150$ between the measured
convergence and the tolerance means the gate rejects only a gross error. It
does establish that the low-density end of the grid is in the coronal regime,
which Chapter~\ref{ch:results} relies on when it reads the low-density column
as ionising, and that is worth having. It is not evidence about rate
magnitudes.

\subsection{Gate C: what the Saha gate tests, which is not the Saha limit}

Gate~C is named for the high-density limit in which collisions dominate
radiation and populations approach their Saha--Boltzmann values. It computes
three quantities per temperature: whether $n_{2p}/\ngs$ increases monotonically
with density, whether it stays below the Saha value (with \SI{5}{\percent}
slack), and whether it approaches that value, defined as reaching $10^{-4}$ of
it.

The pass flag uses the first two and \textbf{excludes the third}
(\texttt{validate\_gates.py}:250; the \texttt{approaching} column is computed at
line 239, stored, printed at line 261, and never consulted). The exclusion is
not what makes the gate weak, because the excluded column would pass: the
measured approach fraction runs $0.0123$ to $0.0155$, comfortably above the
$10^{-4}$ the criterion asks for. What makes it weak is that the criterion is
$150$ times slacker than the measurement. Reaching $10^{-4}$ of Saha is not
approaching Saha, and neither is reaching \SI{1.5}{\percent}.

The measurement itself is the useful part. At
$\nel = 10^{15}\,\si{\per\cubic\centi\metre}$, the top of the grid, the $2p$
population sits at about \SI{1.5}{\percent} of its Saha--Boltzmann value, so
that level is nowhere near local thermodynamic equilibrium (LTE) anywhere on
this grid. Only $2p$ was tested, so this is a statement about $2p$ and not about
the whole distribution. It is the expected behaviour for an optically thin
plasma whose ion population is an externally imposed reservoir rather than a
Saha partner, and Chapter~\ref{ch:results} depends on it: it is why the two
supply channels of Section~\ref{sec:R_depends_on_b} never merge, and why
$\ffrac{3}-\ffrac{4}$ does not vanish at high density.

\textbf{Grade: weak, and mis-named.} What Gate~C tests is monotonicity in
density and a Saha ceiling. Both would fail on a sign error in the recombination
source, which is worth something. Neither tests the limit the gate is named for.

%-----------------------------------------------------------------------------
\section{Gate E: the timescale hierarchy}
\label{sec:gate_e}
% QUESTION: is the two-clock structure of Chapter 3 present everywhere on the
% grid, and how demanding is the test that says so?

Gate~E diagonalises $\Lmat$ at every grid point, orders the negative
eigenvalues, and forms $\Msep = \tauslow/\taurel = |\lambda_1|/|\lambda_0|$. It
passes if $\Msep > 1$ at \SI{95}{\percent} of points
(\texttt{validate\_gates.py}:409--412).

Measured over all $400$ points: $\Msep$ ranges from $86.77$ at
$[49,7]$ ($\Te = \SI{10}{\electronvolt}$,
$\nel = \SI{e15}{\per\cubic\centi\metre}$) to $1.72928\times10^{9}$ at the cold
corner; $\tauslow$ spans \SI{75.4}{\nano\second} to \SI{67.2}{\second} and
$\taurel$ spans \SI{0.87}{\nano\second} to \SI{38.9}{\nano\second}
(project working note; Table~\ref{tab:ladder}). The gate asks for $\Msep>1$. The smallest value
anywhere is $86.8$.

\textbf{Grade: tautological as posed, and the real content is elsewhere.} The
threshold sits between two and seven orders of magnitude below the measured
values. No matrix this construction has produced comes within a factor of $86$
of failing, and none plausibly could, because the fast radiative rates and the
slow ionisation rate that set the two eigenvalues are separated by construction
of the physics rather than by a numerical accident.
The reportable quantity is the range, not the pass, together with the fact that
three independent implementations produce it bit-for-bit: \texttt{qss\_analysis.py},
Gate~E itself, and an unconditional recomputation stored in
\texttt{timescales\_unfiltered\_CHECK.npz}. That agreement is what caught the
eigenvalue-filter error of Section~\ref{sec:errors_found}, and a check that has
actually fired is worth more than a threshold nothing approaches.

Two properties of these numbers must travel with them, and both come from
Chapter~\ref{ch:theory}. $\tauslow$ is boundary-dependent, which
Section~\ref{sec:attacks} quantifies. And the slow eigenvalue is not an
equilibration rate: measured against the ground-state ionisation rate
coefficient $K_{\rm ion}(1s)$, in \si{\cubic\centi\metre\per\second}, it
satisfies $|\lambda_0| = 2.286\,K_{\rm ion}(1s)\,\nel$ across the grid
(project working note; Table~\ref{tab:ladder}), so $\tauslow$ is the collisional--radiative
ionisation lifetime of the ground-state reservoir and not the time for anything
to come to equilibrium.

\subsection{The gates do not stop anything}

No gate raises on failure. Each prints \texttt{PASS} or \texttt{FAIL} and
returns a flag; \texttt{main} sets \texttt{all\_pass = False}, prints a warning,
and exits normally. \textbf{The script returns exit code zero whether or not
Gate~D fails.} The residual gates elsewhere in the project do raise; this one
does not, and that asymmetry is a defect.

%-----------------------------------------------------------------------------
\section{The checks that could have failed}
\label{sec:severe_checks}
% QUESTION: which checks in this project would a plausible wrong model fail?

Everything so far has been a demotion. This section is the other half, and it
is short by design: a small number of severe checks is worth more than a long
list of consistency tests, and these are the ones the thesis leans on.

\subsection{The superposition residual}
\label{sec:superposition}

The entire mechanism of Chapter~\ref{ch:results} rests on one structural claim:
that the excited populations $\nF$ of Section~\ref{sec:qss} split exactly into a
ground-fed and a recombination-fed part,
$\nF = \bm{a}\,\ngs + \bm{c}\,\nion$, where the network responses $\bm{a}$ and
$\bm{c}$ of Section~\ref{sec:R_depends_on_b} are dimensionless and independent
of the two reservoir densities. If that split is not exact, the ground-fed
fractions $\ffrac{p}$ of Section~\ref{sec:ground_fed_fraction} are not well
defined and the closed form for the diagnostic error does not exist.

The check solves the full system three times at each grid point and each step
direction, once with each reservoir alone and once with both, and measures the
relative residual of the sum against the joint solve. Across the $784$
grid-point and step-direction pairs it covers, out of the $800$ the grid admits,
the largest residual is $3.075\times10^{-14}$
(\texttt{plateau\_gridmap.txt}, line 19).

\textbf{Grade: severe}, for a specific reason. Feeding the check a corrupted
input in which the $n=3$ and $n=4$ entries are swapped, which is the mistake
that would most easily go unnoticed in a two-shell observable, trips it
(project working note; Table~\ref{tab:ladder}). It is one of only two checks in the
project reported to catch that swap.

\subsection{The reduced-versus-full ratio test}

The second is its companion. The observable $\Robs = n_3/n_4$ can be computed
two ways: from the full $43$-state steady-state solve, and from the reduced
two-channel form of Eq.~\eqref{eq:R_of_b} with the ground-fed fractions
substituted in. The two routes share no code path beyond the matrix itself. At
the benchmark point $\Te = \benchTe$, $\nel = \benchne$, an independent
re-implementation obtained $\Robs = 0.777768$ by three routes against
Chapter~\ref{ch:theory}'s recorded $0.7778$
(project working note; Table~\ref{tab:ladder}). A $3\leftrightarrow4$ swap breaks the
agreement.

\subsection{Pinned hashes on the inputs}

The most common way for a computation like this to become wrong is not an
algebra error. It is a stale array: a matrix regenerated after a correction,
combined with an analysis output produced before it. This project has three
retractions in its history and at least one of them had that shape.

\texttt{src/validation/preflight.py} pins the SHA-256 hashes, fixed-length
fingerprints that change if any byte of a file changes, of the two canonical
inputs and compares them byte for byte before any verification script
runs. Recomputed for this chapter:
\begin{center}
\begin{tabular}{ll}
  \toprule
  file & SHA-256 (first 16 hex digits) \\
  \midrule
  \texttt{data/processed/cr\_matrix/L\_grid.npy} & \texttt{2d92b58e1693107d} \\
  \texttt{data/processed/cr\_matrix/S\_grid.npy} & \texttt{7822f536590c76ba} \\
  \bottomrule
\end{tabular}
\end{center}
Both match the pinned baselines at \texttt{preflight.py}:37--38 exactly, so the
matrix on disk at the time of writing is the one the pins were taken from; the
statement is only as good as the discipline of running the preflight before
each analysis.

\textbf{Grade: severe on the data, incomplete on the scripts.} Of the six
scripts the same preflight claims to check, \texttt{verify\_ch3\_groupB.py} and
\texttt{verify\_grid\_coverage.py} do not exist on disk under those names; a
missing script is downgraded to a non-blocking advisory
(\texttt{preflight.py}:258--262) while the summary still prints ``All checks
passed''. The hash check is sound; the sentence after it overstates what ran.

\subsection{Guards that recompute rather than trust}

\texttt{src/analysis/make\_ch5\_figures.py} does not plot the stored results. It
rebuilds every plotted quantity from $\Lmat$ and the recombination source
$\Svec$ of Eq.~\eqref{eq:source} and compares against
the stored comma-separated tables at a relative tolerance of $10^{-9}$, raising
with the offending grid point, column, and matrix hash if any row disagrees
(lines 393--432). It reads the temperature step size out of the file header
rather than redeclaring it, so a figure cannot be drawn against a step
convention different from the one that produced the data. It also verifies that
the untrapped run of the radiation-trapping sweep reproduces the canonical
matrix exactly, without which no trapped number in Chapter~\ref{ch:limits} would
be interpretable as a difference.

The same script carries a negative control: it computes the correlation between
$\Msep$ and the diagnostic error both raw and after controlling for temperature
and density, and refuses to draw the figure if the controlled correlation fails
to change sign (Section~\ref{sec:M_does_not_predict} draws the conclusion).

\subsection{Two checks that fired, and one that was withdrawn}

Three episodes say more than the pass/fail table, because in each the check did
something.

The Lyman-$\alpha$ line-centre cross-section gate, an independent implementation
in \texttt{escape\_factor.py} against the project's own Einstein coefficients,
fired at \SI{1156}{\percent} on its first run, located a $4\pi$ factor, and now
agrees to \SI{0.059}{\percent} against a raising \SI{1}{\percent} tolerance
(Section~\ref{sec:opacity}).

\texttt{verify\_fujimoto\_table41.py} carries a deliberately broken variant with
proton-impact $\ell$-mixing removed. It returns a recombination-fed population
coefficient $r_0(2) = 104.5$, and the script flags it itself as unphysical on
the ground that a level cannot exceed Saha equilibrium. Production gives
$0.7392$. A check that fires on a known-bad input is the definition of a severe
one.

Third, a semigroup propagator identity,
$\exp(\Lmat t_1)\exp(\Lmat t_2) = \exp(\Lmat(t_1+t_2))$ for a matrix exponential
computed by the same routine, returned $0.000\times10^{0}$ by construction and
was \emph{withdrawn by the author} (Section~\ref{sec:bridge}), the judgement
this chapter applies to Gate~A, applied earlier and unprompted.

\subsection{The bound that is a theorem}

Section~\ref{sec:bound} derives an upper bound on the diagnostic error of the
form $\tanh(\Delta/4)$, where $\Delta$ is the separation between the two shells'
switching points. A script reports that the bound is honoured at all $248$
points tested, and that line reads as validation.

It is not. Here $\Delta$ is the dimensionless separation between the two shells'
switching points, in units of the logarithm of the reservoir strength. Given
$\bm{a},\bm{c}\ge0$, which Section~\ref{sec:R_depends_on_b} proves from the
M-matrix structure of $-\Lmat_{FF}$, the bound is a theorem: it cannot fail
except on an arithmetic bug. Fed corrupted input it still passes. Swapping
shells $3\leftrightarrow4$ sends $\Delta\to-\Delta$ and leaves both $|\Delta|$
and the logarithmic change in $\Robs$ that it bounds unchanged, so the gate
cannot see it. Swapping $c_3\leftrightarrow c_4$ also
passes, but for a different reason: it moves $|\Delta|$ from $1.945614$ to
$0.939493$, a factor $2.07$, and the gate still passes because it compares
$|f_3-f_4|$ against $\tanh(|\Delta|/4)$ evaluated from the same corrupted
$\bm{c}$. The bound is satisfied by whatever $\bm{a},\bm{c}\ge0$ it is
handed.
\textbf{Grade: tautological.} The script's own docstring calls it a wiring
check; the line it prints does not, and the printed line is what a reader sees.

The bound is a result, and a useful one, because it holds for any atomic dataset
whatever. It is not evidence that this atomic dataset is right.

%-----------------------------------------------------------------------------
\section{Conditioning, precision, and state-space convergence}
\label{sec:conditioning}
% QUESTION: could the answers be artifacts of finite-precision arithmetic or of
% where the level ladder was cut off?

Two numerical questions sit underneath every result in this thesis. The
elimination of Section~\ref{sec:qss} inverts a $42\times42$ matrix whose entries
span many orders of magnitude, which invites the question of whether the inverse
means anything. And the level ladder was cut at $n_{\max}=15$, which invites the
question of what the discarded levels would have changed.

\subsection{The inverse is well behaved, and here is how well}

The relevant quantity is the condition number of $\Lmat_{FF}$, the block of
$\Lmat$ coupling the $42$ excited levels to each other
(Section~\ref{sec:qss}), and it is the factor by which a relative perturbation
of the matrix can be amplified in the solution.
Computed from the canonical $\Lmat$ at all $400$ grid points by
\texttt{src/validation/verify\_operator\_conditioning.py} and stamped to
\texttt{validation/operator\_conditioning/operator\_conditioning.csv}: minimum
$1.4821\times10^{3}$ at $[0,0]$ ($\Te = \SI{1.000}{\electronvolt}$), maximum
$1.7436\times10^{5}$ at $[33,7]$ ($\Te = \SI{4.715}{\electronvolt}$), median
$2.1044\times10^{4}$. In double precision, which carries about sixteen
significant digits, a condition number of $10^{5}$ costs five of them. Eleven
remain, and every conclusion in this thesis turns on the first three.

The range was first asserted in Chapter~\ref{ch:theory} on an unscripted
working note, which is not a check; the stamped script reproduces it to
\SI{0.14}{\percent} and \SI{0.21}{\percent}. The number was right the whole
time and was not evidence until it could be re-run.

Three further measurements bound the arithmetic. The eigenvalue condition
number is a different quantity from the matrix condition number above: it is
$1/|\bm{y}^H\bm{x}|$ for a left and right eigenvector pair, and it measures how
far a single eigenvalue moves under a perturbation of the matrix. For the two
eigenvalues the thesis uses it is of order unity
($1.71$ and $1.98$ at the benchmark point, $3.26$ and $2.59$ at the cold
corner), despite a largest absolute column sum of
$2.0\times10^{14}\,\si{\per\second}$; that large norm is a
statement about the fast radiative modes and does not propagate into
$\lambda_0$ or $\lambda_1$. Recomputing the two supply channels in single
precision moves them by $1.3\times10^{-5}$, so double precision is nowhere near
a cliff. And tightening the integrator changes nothing that matters: a relative
tolerance of $10^{-9}$ against $10^{-12}$ shifts the result by
$1.20\times10^{-8}$, and computing the propagator by a dense matrix exponential
against a Krylov action differs by $1.13\times10^{-10}$.

The eigenvalue condition numbers and the single-precision test cover three of
$400$ points; the tolerance comparisons cover two. The condition numbers and the
displacement are recomputed under the definition $1/|\bm{y}^H\bm{x}|$ with
unit-norm left and right eigenvectors by \texttt{verify\_residual\_claims.py}
(\texttt{validation/residual\_claims/}). They are point measurements, not grid
scans, and are quoted as such.

\subsection{Where the ladder was cut}
\label{sec:convergence}

The quantity that has to converge is not a population but the difference in
ground-fed fraction between the two Balmer shells, $\ffrac{3}-\ffrac{4}$, since
Chapter~\ref{ch:results} shows the diagnostic error to be proportional to it.

Rebuilding $\Lmat$ with successively lower $n_{\max}$, from $15$ down to $9$,
with the loss channels of every removed state returned to the diagonal of the
survivors, and recomputing $\ffrac{3}-\ffrac{4}$ at each truncation's own
$u^{\mathrm{CRE}}$ (\texttt{verify\_nmax\_downward\_scan.py},
\texttt{validation/nmax\_downward\_scan/}) gives increments of one sign at the
three named points, with the last measured ratio $d(15)/d(14)$ equal to
$0.88$ at the benchmark, $0.72$ at the cold corner $[0,0]$ and $0.33$ at the
ridge $[15,3]$; at the $245$ of the $280$ defended points where the increments
are one-signed and geometric the ratio has median $0.67$. Summing the geometric
tail with that ratio gives an extrapolated truncation error of
\SI{-0.46}{\percent} at the benchmark, \SI{+1.6}{\percent} at the cold corner
and \SI{+0.005}{\percent} at the ridge; holding $u$ at its $n_{\max}=15$ value
instead gives \SI{-0.38}{\percent}, \SI{+1.4}{\percent} and
\SI{+0.005}{\percent}, so the sign does not depend on that choice. A historical
project value with no producing script, not used as evidence here and superseded
by \texttt{verify\_nmax\_downward\_scan.py}, records
\SI{+0.9}{\percent}, \SI{-1.4}{\percent} and \SI{-0.55}{\percent} with a
common ratio of about $0.75$; the stamped run reverses both signs and the
conclusion that rested on them. Truncation makes
$\ffrac{3}-\ffrac{4}$ slightly too large at the benchmark and too small at the
cold corner, so it is not a bias that can be corrected by a single factor.
Over the defended points where the tail can be summed the extrapolated error
has median \SI{+0.002}{\percent} and runs from \SI{-0.82}{\percent} to
\SI{+1.03}{\percent}, smaller than the spread between the two step conventions
of Section~\ref{sec:step_convention} at every such point. At the $35$ defended
points of the density column $\nel = \SI{5.2e13}{\per\cubic\centi\metre}$ the
increments change sign among $n_{\max} = 12$ to $15$ and no tail can be
summed; the last measured increment there is below \SI{0.041}{\percent} of
$\ffrac{3}-\ffrac{4}$ at every point.

\paragraph{A trap in this test.} The first attempt dropped the top shells from
the state vector without removing their loss channels from the diagonal of the
surviving states, which then lost population to states that no longer existed,
and $\ffrac{3}-\ffrac{4}$ jumped. The stamped scan repeats the mistake
deliberately: dropping $n=15$ alone without the correction moves
$\ffrac{3}-\ffrac{4}$ by \SI{11.5}{\percent} at the benchmark,
\SI{12.6}{\percent} at the cold corner and \SI{5.0}{\percent} at the ridge,
where the correct truncation to $n_{\max}=14$ moves it by \SI{0.06}{\percent},
\SI{0.64}{\percent} and \SI{0.01}{\percent}; dropping six shells without it
gives \SI{17.9}{\percent}, \SI{25.9}{\percent} and \SI{13.3}{\percent} (a
working note recorded \SIrange{22}{37}{\percent}). The jump is an artifact of
the truncation procedure, not a dependence on $n_{\max}$. Removing state $k$
requires adding $\sum_{k\ \rm dropped} L_{ki}$ back to $L_{ii}$ for every
surviving $i$, as the production matrix does. Anyone repeating this test will
reproduce the spurious number first.

\paragraph{The terminal shell is over-populated, and by a measured amount.}
Truncating the ladder at $n_{\max}$ removes, from the balance of the terminal
shell, the collisional excitation out of it to $n > n_{\max}$ and the
de-excitation and cascade into it from those levels (cascade is downward, so
what is lost is a way out, not something ``to cascade out to''). The loss removed is not small. Continuing the $12\to13$, $13\to14$
and $14\to15$ excitation coefficients geometrically, the missing $15\to16$
channel would carry $1.30$ times the whole retained loss out of $n=15$ at the
benchmark and $1.20$ times at $[0,4]$, between $1.20$ and $1.35$ times at
every grid point, while radiative decay is below $10^{-6}$ of that loss and
ionisation is $7$ to \SI{10}{\percent} of it
(\texttt{validation/terminal\_shell\_budget/}). The top shell sits within
$6\times10^{-5}$ of Saha--Boltzmann, so for the total population the removed
exchange is nearly balanced and the sign of the bias is decided channel by
channel: the ground-fed channel, whose net flux through the top of the ladder
is upward, loses only an exit and can only be biased high, and it is the
ground-fed coefficient $r_1(15)$ that the external comparison finds high while
$r_0(15)$ agrees (Section~\ref{sec:fujimoto}). Comparing each shell's ground-fed coefficient $a_p$ as terminal
shell (operator rebuilt with $n_{\max}=p$, the loss channels into the removed states kept on the diagonal) with its
value in the full $n_{\max}=15$ operator gives at the Fujimoto point $[49,0]$ (\SI{10}{\electronvolt},
\SI{e12}{\per\cubic\centi\metre}) excesses of $9.4$ at $p=8$, $10.7$ at $p=9$ and $6.3$ at $p=10$, then $4.4$,
$3.0$, $2.1$ and $1.4$ at $p=11$ to $14$ (\texttt{verify\_terminal\_vs\_interior.py},
\texttt{validation/terminal\_vs\_interior/}). At the benchmark $p=8$ to $10$ give $14.2$, $12.3$ and $6.5$; over
the $280$ defended points the $p=8$ excess spans $8.2$ to $17.7$ and the $p=10$ excess $5.4$ to $7.5$. The excess
is confined to the ground-fed channel: the shell's total population and recombination-fed coefficient move under
\SI{1}{\percent} at the benchmark and at most \SI{32}{\percent} and \SI{30}{\percent} at the low-density named
points. A historical project value with no producing script, not used as evidence here and superseded by
\texttt{verify\_terminal\_vs\_interior.py}, records $9.4$, $10.7$ and
$6.3$ with neither quantity nor grid point named; the stamped run reproduces them to \SI{0.4}{\percent} as $a_p$
at $[49,0]$. The absence of a monotone trend over $p=8$ to $10$ holds there and at $47$ of the $280$ points;
from $p=9$ upward the excess falls at every grid point, at the benchmark from $p=8$. The external excess at the
terminal shell, $4.9$ to $6.2$ against published values (Section~\ref{sec:fujimoto}), lies outside the $6.3$ to
$10.7$ band, not inside it as the working note states, but within a factor $2.2$ of every value in it, and both rise
with density ($4.9$ to $6.2$ externally, $9.4$ to $17.7$ at $p=8$ internally at \SI{10}{\electronvolt}). That
agreement in magnitude and in the sign of the density trend makes truncation a sufficient explanation for the
$n=15$ discrepancy without invoking the atomic data. The fall with $p$ cannot be extrapolated to $p=15$: the
reference for $p=14$ is itself a ladder with one shell above it, so the internal measurement neither predicts nor
excludes $4.9$ to $6.2$ there.
No result in this thesis rests on the
top shell: the observable is built from $n=3$ and $n=4$, both fully resolved in
orbital angular momentum $\ell$ and eleven shells below the cut at
$n_{\max}=15$.

\paragraph{What has now been checked.} Above $n=8$ the shells are bundled and
their internal $\ell$-distribution is assumed statistical rather than computed.
The script that tests that assumption, \texttt{verify\_bundling\_psm20.py}, had
never been executed; two of the three defects found on its first run, the
synthesised temperature grid and the identically zero $\ell$-mixing array for
the bundled indices, were predicted here. Section~\ref{sec:bundling_check}
reports all three defects and the result.

%-----------------------------------------------------------------------------
\section{Is the shell ratio the line ratio?}
\label{sec:balmer_equivalence}
% QUESTION: the analysis computes n_3/n_4, but a spectrometer measures
% Halpha/Hbeta. What does the substitution cost?

Every analytic result in Chapter~\ref{ch:theory} is written for the shell ratio
$\Robs = n_3/n_4$, because the two-channel decomposition applies to shell
populations. A spectrometer does not measure shell populations. It measures the
$H_\alpha$ and $H_\beta$ line intensities, each of which is a sum over the
electric-dipole-allowed sublevel transitions, weighted by their Einstein
coefficients. Section~\ref{sec:radiative} established that these are not
proportional: $4f$ holds \SI{43.75}{\percent} of the $n=4$ shell by statistical
weight and contributes nothing to $H_\beta$, since $4f\to2p$ would change $\ell$
by two and is forbidden.

The substitution therefore has to be measured, not assumed.

The line intensities are built as $A$-weighted sums over resolved channels,
$H_\alpha$ from $3s\to2p$, $3p\to2s$ and $3d\to2p$, and $H_\beta$ from
$4s\to2p$, $4p\to2s$ and $4d\to2p$, with populations taken from the solve.
Nothing statistical is assumed anywhere in that construction: searching the two
producing scripts for statistical-weight expressions returns no matches, and the
$A$-value lookup raises unless exactly one row matches each transition.

Two measurements follow. The $4f$ fraction of the $n=4$ shell runs from
$0.4361$ to $0.4375$ across the grid, against the statistical value
$14/32 = 0.4375$; proton-impact $\ell$-mixing drives the shell to a statistical
distribution to better than \SI{0.3}{\percent} everywhere, with the departure
concentrated at $\nel = \SI{e12}{\per\cubic\centi\metre}$ where $\ell$-mixing
weakens and radiative decay competes. That fraction is a property of the full
equilibrium population; the quantity that sets the line-versus-shell difference
below is the $\ell$-distribution within each supply channel separately, which
is not statistical at the low-density edge. And the ratio of the transient error
computed on the line ratio to the same error computed on the shell ratio runs
from $0.9482$ to $0.9999$ over the $784$ point-and-direction pairs, below $1$ at
every one: \textbf{worst case \SI{5.2}{\percent}, at the heating step from
$9.54$ to \SI{10.0}{\electronvolt} at the lowest density}
(\texttt{verify\_weighted\_census.py}, stamped to
\texttt{validation/weighted\_census/}). At
$[0,4]$, the point carrying the largest error anywhere, the two agree to
\SI{0.06}{\percent} ($0.386683$ against $0.386903$). Because the ratio is below
$1$ everywhere, the shell census of Chapter~\ref{ch:results} is conservative
for the line observable: counted on the line ratio, $44$ instead of $45$ of the
$448$ pairs above \SI{2}{\electronvolt} exceed \SI{10}{\percent} at
\SI{100}{\micro\second}, with the worst case $0.1746$ instead of $0.1748$ at
the same point. The \SI{5.2}{\percent} scales inversely with the proton
$\ell$-mixing rate: with the within-shell $\ell$-changing rates scaled by a
common factor and the column sums of the operator preserved, the corner
departure is \SI{10.0}{\percent} at half the rate, \SI{2.6}{\percent} at double
and \SI{1.1}{\percent} at five times, while the benchmark line-to-shell ratio
stays between $0.9987$ and $0.9999$ (\texttt{validation/weighted\_census/},
$\ell$-mixing sensitivity block and \texttt{weighted\_census\_lmix.csv}); and
it sits on the grid boundary.

That \SI{0.06}{\percent} is a single-point figure and must not be quoted as a
grid-wide one. The grid-wide statement is the \SI{5.2}{\percent}.

\subsection{And is the plasma transparent to its own Balmer light?}

The shell-to-line substitution assumes the emitted photons escape. For the
Lyman series that assumption fails at the cold edge and
Chapter~\ref{ch:limits} treats it at length. For the Balmer lines themselves it
closes here, because it bears on the observable rather than on the matrix.

Section~\ref{sec:opacity} defines the optical depth $\tau$ and the escape factor
and treats the Lyman series, where transparency fails. For the Balmer lines the
absorbers are the $n=2$ atoms, $n_{2s}+n_{2p}$ in collisional--radiative
equilibrium with $\nion=\nel$, the Doppler width is taken at $\Tn=\Te$, and the
multiplet oscillator strengths are $f(2\to3)=0.641$ and
$f(2\to4)=0.119$~\cite{WieseFuhr2009}. The line-centre optical depth of
$H_\alpha$ across a \SI{10}{\centi\metre} slab is then $4.9\times10^{-7}$ at the
cold corner, $8.2\times10^{-4}$ at $[0,4]$ and $0.214$ at $[0,7]$, the single
worst cell on the grid; resolving the absorbers into $2s$ and $2p$ with the
model's own $A$ coefficients changes these by at most \SI{4}{\percent}. At
$[0,7]$ the population escape factor for a photon born at the slab centre
($\kappa_0 D/2 = 0.107$) is $0.928$ for $H_\alpha$ and $0.990$ for $H_\beta$, so
the observed ratio is lowered by \SI{6.3}{\percent} in that one cell; over the
full path ($\tau=0.214$) the figures are $0.861$ and \SI{12}{\percent}. Above
\SI{2}{\electronvolt} the largest $H_\alpha$ optical depth is $0.010$ at
$[15,7]$ and the ratio moves by less than \SI{0.7}{\percent} anywhere
(\texttt{verify\_balmer\_optical\_depth.py},
\texttt{validation/balmer\_optical\_depth/}). The Balmer lines are optically
thin over the region where this thesis makes quantitative claims. A historical
project value with no producing script, not used as evidence here and superseded
by \texttt{verify\_balmer\_optical\_depth.py}, records
$6.3\times10^{-7}$, $1.0\times10^{-3}$ and $0.263$ for the same three cells;
those values are $1.23$ to $1.28$ times the stamped ones and no convention
tested recovers them, and the escape factor above $0.85$ and ratio shift of at
most about \SI{10}{\percent} recorded with them are half-slab figures attached
to a full-slab optical depth. The conclusion holds with either set of values.

%-----------------------------------------------------------------------------
\section{Two tests of the reduction, and their expected outcomes}
\label{sec:attacks}
% QUESTION: the two objections most likely to destroy the result were tested;
% what was the falsifier in each case, and did it appear?

The checks so far ask whether the model computes what it intends. This section
asks something harder: whether two specific structural choices manufacture the
result. In each case the size of effect that would have destroyed the claim was
written down before the test was run, which is what makes the outcome
informative rather than reassuring.

\subsection{Does the open boundary manufacture the slow clock?}

$\Lmat$ describes an open system: atoms that ionise leave the accounting, the
recombination feed $\Svec\nion$ enters as a constant source rather than a matrix
element, and there is no forty-fourth state for ions to accumulate in. Since the
persistence claim of Section~\ref{sec:persistence} rests entirely on $\tauslow$
exceeding the duration of the disturbance, the first objection is that
$\tauslow$ is an artifact of the open boundary.

\textbf{The falsifier, named first.} Write $\taudrive$ for the duration of the
disturbance, fixed at \SI{100}{\micro\second} throughout this thesis as a
representative edge-localised-mode (ELM) time
(Section~\ref{sec:persistence}). At $[0,4]$, the worst point on the grid,
$\tauslow/\taudrive \approx 2300$. If closing the system shortened $\tauslow$ by
that factor, the slow clock would fall to $\taudrive$ itself and the persistence
would vanish. A factor of $2300$ refutes the claim.

Two grid labels appear in the table for the same worst case: the error is
evaluated at $[0,4]$, but the operator governing the decay after the step is the
post-step one, $\Lmat[1,4]$, and every timescale in this thesis names the
operator it belongs to (Chapter~\ref{ch:results}).

The manifold was closed by appending an ion state,
\begin{equation}
  \Lmat_c =
  \begin{pmatrix}
    \Lmat & \Svec \\
    -\operatorname{colsum}(\Lmat) & -\textstyle\sum\Svec
  \end{pmatrix},
  \label{eq:closure_ch4}
\end{equation}
so that every ionising atom lands in the new state instead of being deleted.
Each column of $\Lmat_c$ sums to zero to machine precision, and the construction
reproduces the closed-system value at the cold corner recorded in this
project's working notes, \SI{1.6226}{\second}, to five digits, which confirms it is the same
closure rather than a new one. The measured factors are given in
Table~\ref{tab:closure}.

\begin{table}[tbp]
  \centering
  \caption[Effect of closing the ion boundary on the slow timescale]{The slow
  timescale $\tauslow$ computed with the ion treated as an external reservoir
  (open) and with a forty-fourth ion state appended (closed), at four grid
  points. Over the whole grid the ratio has minimum $1.00$, median $1.00$ and
  maximum $41.4$: the boundary condition matters only where $\tauslow$ already
  exceeds the disturbance duration by three orders of magnitude. Every entry is
  reduced from \texttt{validation/ion\_closure/ion\_closure\_summary.csv}
  (\texttt{verify\_ion\_closure.py}) by
  \texttt{verify\_headline\_provenance.py}; the open column is the float64
  eigensolver, which at the cold corner differs from the artifact's 40-digit
  path by one part in $10^{4}$.}
  \label{tab:closure}
  \begin{tabular}{lrrr}
    \toprule
    point & $\tauslow$ open & $\tauslow$ closed & factor \\
    \midrule
    cold corner $[0,0]$ & \SI{67.23}{\second} & \SI{1.6226}{\second} & $41.4$ \\
    post-step operator $\Lmat[1,4]$ & \SI{0.23324}{\second} & \SI{0.015173}{\second} & $15.4$ \\
    ridge $[15,3]$ & \SI{2.027}{\milli\second} & \SI{1.999}{\milli\second} & $1.01$ \\
    benchmark $[23,5]$ & \SI{22.728}{\micro\second} & \SI{22.708}{\micro\second} & $1.00$ \\
    \bottomrule
  \end{tabular}
\end{table}

The shape of the distribution carries the argument: the factor is large only
where $\tauslow$ already exceeds $\taudrive$ by about a thousand, so the error
has already saturated, and it is unity everywhere $\tauslow$ approaches
$\taudrive$. Recomputing the entire persistence census with closed-system
$\tauslow$ moves the breakdown count from $202$ to $200$ of $680$ and the worst
case from $0.38682$ to $0.38563$, about one part in $325$; restricted to
$\Te \ge \SI{2}{\electronvolt}$, where Chapter~\ref{ch:results} makes
quantitative claims, the count is $45$ before and $45$ after.

The falsifier did not appear. A factor of $2300$ was needed; the measured factor
at that operator is $15$.

\textbf{The caveat, which must not be softened.} The closure adds an ion
reservoir; it does not close the electron balance, since $\nel$ is held fixed
while the ion population moves, and it adds no transport. It answers the
boundary condition on the atomic manifold and nothing else:
Section~\ref{sec:transport} states the objection it cannot answer, and
Section~\ref{sec:quasineutrality} the one it makes sharper rather than softer.

\subsection{Is the single-slow-mode estimate a bound on the time-averaged error?}

Chapter~\ref{ch:results} evaluates the error at the partial-equilibrium plateau
and builds from it the single-slow-mode estimate of Eq.~\eqref{eq:elm_bound},
which multiplies the plateau error $\epsplat$ of
Section~\ref{sec:eps_definitions} by the fraction of $\taudrive$ over which an
exponential decay on $\tauslow$ keeps it alive. It was first written as a lower
bound, on the assumption that the observable cannot decay faster than the
slowest mode. That assumption is testable, and an estimate quoted as though it
were the answer is a claim about accuracy that has to be measured.

Propagating the full $43$-state system through the step by direct
eigen-decomposition and integrating the error over \SIrange{0}{100}{\micro\second}
gives the comparison in Table~\ref{tab:lowerbound}.

\begin{table}[tbp]
  \centering
  \caption[Accuracy of the single-slow-mode estimate]{The plateau error, the
  true error averaged over \SIrange{0}{100}{\micro\second} from direct
  eigen-propagation of the full $43$-state system, and the single-slow-mode
  estimate of Eq.~\eqref{eq:elm_bound}, at five grid points. The averages are
  integrated at $400$ time samples per decade; at $1600$ per decade the five
  change by at most $4.7\times10^{-6}$ relative, and the $[0,4]$ average by
  $2.7\times10^{-10}$. Integration on a uniform \SI{25}{\nano\second} mesh, which
  under-resolves the nanosecond rise, reads low by $6\times10^{-5}$ to
  $5\times10^{-4}$ relative. The estimate reproduces
  the true average to better than \SI{0.8}{\percent} at the five points. The
  value below unity at $[0,4]$ is real: a deficit of $2.2\times10^{-5}$ that
  survives the denser grid. Source: \texttt{verify\_trajectory\_census.py},
  \texttt{validation/trajectory\_census/}.}
  \label{tab:lowerbound}
  \begin{tabular}{lrrrr}
    \toprule
    point & $\epsplat$ & true \SI{100}{\micro\second} average & estimate, Eq.~\eqref{eq:elm_bound} & true/estimate \\
    \midrule
    $[0,4]$ worst case & $0.386903$ & $0.386811$ & $0.386820$ & $0.99998$ \\
    $[15,3]$ ridge     & $0.180728$ & $0.175283$ & $0.174812$ & $1.00270$ \\
    $[15,5]$           & $0.103550$ & $0.072185$ & $0.071722$ & $1.00646$ \\
    $[23,5]$ benchmark & $0.063612$ & $0.011796$ & $0.011712$ & $1.00722$ \\
    $[0,0]$            & $0.083004$ & $0.083084$ & $0.083004$ & $1.00096$ \\
    \bottomrule
  \end{tabular}
\end{table}

At the five points tested the estimate sits within \SI{0.8}{\percent} of the
directly propagated average, and at $[0,4]$ it sits above it, so it is an
estimate and not a bound. The accuracy at these five points does not
generalise: over the $448$ pairs above \SI{2}{\electronvolt} the ratio of true
average to estimate departs from unity by up to \SI{2.2}{\percent} at
\SI{100}{\micro\second} and \SI{3.8}{\percent} at \SI{506}{\micro\second}
(Section~\ref{sec:persistence}). What matters here is that the assumption
underneath it was tested rather than asserted.

The reason is structural, and it is a property the shell ratio has and another
observable in this same project does not. The spectrum of $\Lmat$ has one slow
eigenvalue and then a band, with nothing between: at $\Lmat[1,4]$ the three
slowest eigenvalues are $-4.29$, $-1.60\times10^{8}$ and
$-3.44\times10^{8}\,\si{\per\second}$, a gap of $3.7\times10^{7}$. An observable
that depends on the excited populations only through their ratio cannot decay
faster than $\tauslow$: once the manifold is quasi-steady the excited vector is
a function of $u$ alone, so
$\ln\Robs - \ln\Robs^{\mathrm{CRE}} = (S/u^{*})\,\delta u + O(\delta u^{2})$
with $\delta u \propto \mathrm{e}^{-t/\tauslow}$, and the leading term
vanishes only where $S = \ffrac{3}-\ffrac{4} = 0$. It does not: $S$ is positive
at all $400$ grid points, with minimum $0.046$ at $[44,7]$, and the slow
eigenvector's projection onto $\ln\Robs$ equals $S/u^{*}$ to
$1.2\times10^{-4}$ at the benchmark and $1.6\times10^{-8}$ at $[0,4]$, the
deviation being the first-order quasi-steady correction
$\lambda_{\mathrm{slow}}/\lambda_{2}$ itself. It exceeds \SI{1}{\percent} only
at the two hot, dense points $[48,7]$ and $[49,7]$, where the separation is
$87$-fold and which the window criterion excludes from the analysed set
(\texttt{validation/slow\_mode\_projection/}). That decay on an intermediate
mode was the failure mode being hunted, and it does occur for the
ground-normalised error measure, which does not cancel the reservoir and whose
$1/\mathrm{e}$ time is \SI{2.58}{\second} against
$\tauslow = \SI{67.2}{\second}$ at the cold corner $[0,0]$. The estimate is
accurate for the ratio and not for that measure; why it is nevertheless not a
bound, the denominator relaxing on the same $\tauslow$, is set out in
Section~\ref{sec:persistence}.

%-----------------------------------------------------------------------------
\section{External comparison: Fujimoto's population coefficients}
\label{sec:fujimoto}
% QUESTION: does an independently computed collisional-radiative model, built
% decades earlier from different atomic data, agree with this one?

Everything to this point compares the model against itself or against
arithmetic. Only comparison against an independent calculation reaches the
atomic physics, and this project has one such comparison in depth: against
Fujimoto's collisional--radiative treatment of atomic
hydrogen~\cite{Fujimoto2004}, computed decades earlier, with different atomic
data, by a different method.

\subsection{Timescales, at a point that lies exactly on this grid}

Fujimoto works a numerical example at $\Te = \SI{10}{\electronvolt}$ and
$\nel = \SI{e18}{\per\cubic\metre}$, which is
$\SI{e12}{\per\cubic\centi\metre}$ and lands exactly on grid point $[49,0]$. No
interpolation is involved.

\begin{center}
\begin{tabular}{lrrr}
  \toprule
  quantity & Fujimoto App.~4B & this model, $[49,0]$ & ratio \\
  \midrule
  excited-state response & ${\sim}\SI{1e-7}{\second}$ & \SI{3.4057e-8}{\second} & $0.34$ \\
  ground-state depletion & ${\sim}\SI{1e-4}{\second}$ & \SI{1.7657e-4}{\second} & $1.77$ \\
  \bottomrule
\end{tabular}
\end{center}

The two-clock structure appears in both, separated by three orders of
magnitude in both, with the fast time short and the slow time long in the same
sense.

\textbf{What this comparison is not.} Fujimoto's response time is
$\max_p t_{\rm tr}(p)$, the \SI{63}{\percent} rise time of the slowest-responding
level, which is not $1/|\lambda_1|$. The two are different functionals of the
same spectrum, so agreement to a factor of a few is the meaningful comparison
and exact agreement would be suspicious. Quoting these ratios as though they
were a precision test would overstate them.

One qualification on the independence of this check, found on reading the
printed appendix rather than a summary of it. Fujimoto's Figs.~4B.1 and 4B.2 are
captioned as quoted from Sawada and Fujimoto (1994)
\cite{SawadaFujimoto1994}, so the appendix and the paper
Chapter~\ref{ch:conclusions} discusses are the same calculation, not two.
\textbf{The comparison is therefore independent of this thesis but not two
independent checks}, and it is counted once. The appendix also states the same
construction the paper does, that $n(1)$ is held constant, which matters for
Chapter~\ref{ch:conclusions} and is taken up there.

A second, independent route corroborates the same physics. Fujimoto's boundary
level, the principal quantum number above which collisions dominate radiative
decay, sits at $p_G \approx 4$--$5$ at this density. Tracing the boundary
descent in this model puts it at $5g$ at the low-density end of the grid, which
is the crossing Chapter~\ref{ch:atoms} deferred to this chapter
(Section~\ref{sec:radiative}). Two independently computed models place the same
structural feature within one shell of each other, at one density.

\subsection{The recombining coefficients agree}

Fujimoto's Table~4.1 tabulates the two population coefficients of the
decomposition that Section~\ref{sec:R_depends_on_b} derives: $r_0(p)$, how
strongly level $p$ is fed from the ion continuum, and $r_1(p)$, how strongly it
is fed from the ground state. Both are dimensionless, and both are normalised to
the Saha--Boltzmann population of level $p$, which is why $r_0$ is of order unity
while $r_1$ at low density is of order $10^{-5}$, and why a value of $r_0$ above
one is unphysical.

The recombining coefficients agree to \SIrange{0.1}{1}{\percent} for $n \ge 4$.
The discriminating case is $r_0(2)$, where cascade and $\ell$-mixing structure
matter most and a hydrogenic approximation would go wrong: this model gives
$0.7392$ against the tabulated $0.730$, a difference of \SI{1.3}{\percent}, a
test the model could have failed and did not (the broken variant of
Section~\ref{sec:severe_checks} drives the same quantity to $104.5$).

\subsection{The ionising coefficients at low $p$ disagree, and the target is
not verified}

At $\nel = \SI{e12}{\per\cubic\centi\metre}$ this model's $r_1$ sits a factor
$8.3$ low at $p=3$ and $4.4$ low at $p=4$ against Fujimoto's Table~4.1(b).
Those are precisely the two shells whose difference is the mechanism of
Chapter~\ref{ch:results}, so the discrepancy cannot be dismissed as peripheral.

Reading it as a failure of the model is not supportable, for six reasons, and
the status is an open comparison whose \emph{target} is unverified.

\begin{enumerate}
\item \textbf{The model's $r_1$ is reproduced by hand.} In the coronal limit
$r_1(2)$ follows from the single excitation rate coefficient
$K_{\rm exc}(1s\to n{=}2) = \SI{9.2796e-9}{\cubic\centi\metre\per\second}$ and
gives $1.37\times10^{-5}$, against the full $43$-state solve's
$1.66\times10^{-5}$. The two agree to \SI{21}{\percent}, which is the accuracy a
coronal estimate that ignores cascades should have. Whatever produces a factor
of $8$ is not inside the solve.

\item \textbf{The atomic stack passes three independent external checks.} The
radiative recombination coefficient into the ground state is within
\SI{2}{\percent} of the scaled standard value; the sum over excited states,
$\SI{2.08e-13}{\cubic\centi\metre\per\second}$, sits \SI{10}{\percent} below the
standard case-B coefficient $\alpha_B \approx
\SI{2.31e-13}{\cubic\centi\metre\per\second}$, the accepted total for capture
into every excited level, and that \SI{10}{\percent} shortfall is exactly the
$n>15$ truncation of Section~\ref{sec:convergence}; and the ground-state ionisation
coefficient at \SI{10}{\electronvolt}, $\SI{4.864e-9}{\cubic\centi\metre\per\second}$,
sits \SI{10.4}{\percent} below ADAS SCD96 read in its low-density limit,
$\SI{5.43e-9}{\cubic\centi\metre\per\second}$ at
$\nel = \SI{5e7}{\per\cubic\centi\metre}$ \cite{Summers2006, ADAS}, where the
effective coefficient reduces to ionisation out of the ground state because
every excited level decays before it can be ionised. The comparison is at the
low-density limit deliberately: SCD96 varies by \SI{0.9}{\percent} over a
factor four in density there, so the limit is reached. \textbf{This
\SI{10}{\percent} is not the truncation of the item above.} Direct ionisation
from $1s$ is a single cross section and does not know where the level ladder
stops; the difference is between this model's ionisation cross-section source
and the one behind SCD96, and it is reported rather than reconciled.

\item \textbf{The tabulated asymptote demands an excitation coefficient fifteen
times the accepted value.}
Reproducing Fujimoto's quoted coronal asymptote requires
$K_{\rm exc}(1s\to n{=}2) = \SI{1.44e-7}{\cubic\centi\metre\per\second}$ at
\SI{11.03}{\electronvolt}, fifteen times the accepted value. No error in this
model's implementation produces that number; only a different input does.

\item \textbf{Three candidate explanations are eliminated quantitatively.}
Truncation moves $r_1(3)$ by \SI{1.2}{\percent} between $n_{\max}=10$ and
$n_{\max}=15$, which is a factor $8$ short. Lyman trapping at an escape factor
of $0.1$ brings $r_1(2)$ to $0.92$ of the table but leaves $r_1(3)$ at $0.26$ of
it and drives $r_0(2)$ to $8.2$, a population above Saha equilibrium and
therefore unphysical. The third candidate, the $\ell$~closure, is eliminated in
the item below.

\item \textbf{The $\ell$~closure is a real difference, and it does not account
for the gap.} Table~4.1(b) is indexed by $p = 2, 3, 4, 5, 7, 10, 15$, principal
quantum numbers, and carries no $\ell$ label: the published coefficients are
shell quantities computed with the $\ell$~substructure bundled, while this model
resolves $\ell$ below $n=9$. That difference is real, and it is not the
explanation. Imposing the source's own
closure directly, by bundling this model's operator under statistical
$\ell$~weights as $\Lmat_{\rm bundle} = \mathsf{P}\,\Lmat\,\mathsf{R}$ with
$\mathsf{P}_{ki} = 1$ for $i$ in shell $k$ and $\mathsf{R}_{ik} = g_i/g_k$,
moves $r_1(3)$ by \SI{0.42}{\percent}: the factor $8.29$ becomes $8.26$
(\texttt{verify\_fujimoto\_bundle.py}, stamped in
\texttt{validation/fujimoto\_bundle/}). At
$p=2$ the agreement becomes \emph{worse}, from a factor $10.8$ low to a factor
$11.8$ low, and $r_0(2)$ moves from $0.7392$, which agrees with the tabulated
$0.730$, to $0.6459$, which does not. The reason is structural: the bundling
projector annihilates the $\ell$-mixing operator identically, $\max|\mathsf{P}
\mathsf{B} \mathsf{R}| = 8.5\times10^{-7}$ against a matrix scale of
$1.0\times10^{11}$, so bundling cannot express what $\ell$-mixing does; and the
production model is already at the statistical-$\ell$ limit for $p \ge 3$, its
sublevel populations sitting within \SI{5}{\percent} of $g_{n\ell}/g_n$. The
$\ell$-resolved calculation was therefore already the like-for-like comparison.
Removing proton $\ell$-mixing altogether moves $r_1(2)$ by a factor $118$, but
that is the opposite limit from a statistical closure rather than a proxy for
it, and the two do not bracket the tabulated value.

\item \textbf{The transcription and the row labels are correct.} Both were
checked against the printed table. The $\log_{10}\nel = 18$ row reads
$0.730$, $0.835$, $0.947$ and $0.983$ for $r_0$ at $p = 2, 3, 4, 5$, and
$1.79\times10^{-4}$, $5.72\times10^{-5}$, $1.66\times10^{-5}$ and
$5.08\times10^{-6}$ for $r_1$, matching the values used here exactly. The row
labels are $\log_{10}(\nel/\si{\per\cubic\metre})$, which the table establishes
internally: $r_0(2)$ reaches $0.981$ at $\log_{10}\nel = 21$, and Griem's
criterion places local thermodynamic equilibrium for $n=2$ at
\SI{1.74e16}{\per\cubic\centi\metre}, which is $10^{21}\,\si{\per\cubic\metre}$
to within the criterion's own precision. Read as \si{\per\cubic\centi\metre} the
onset would sit five orders above Griem, which is impossible; and the rows run
from $\log_{10}\nel = 12$ to $24$, a span of $10^{6}$ to
$10^{18}\,\si{\per\cubic\centi\metre}$ read as \si{\per\cubic\metre} and no
plasma at all at the top read as \si{\per\cubic\centi\metre}. A one-decade
offset in the density rows is not a tenable reading either.
\end{enumerate}

\textbf{Status.} The $r_1$ deficit at low $p$ is an open external
disagreement, and this thesis does not explain it; the three candidate causes
above each fail to account for a factor of $8$. The $\ell$-closure test carries
a severity check: bundling an operator whose $\ell$~distribution is deliberately
non-statistical moves $r_1(2)$ by a factor $129$, so the null result at $p=3$ is
a measurement and not a blind instrument. What the section does \emph{not}
establish is that the disagreement leaves the quantity
Chapter~\ref{ch:results} uses untouched.
Equation~\eqref{eq:fujimoto_correspondence} makes
$a_p = r_1(p)\,\Zsb{p}/\Zsb{1}$: $r_1$ is the ground-fed coefficient from which
$\ffrac{p}$ is built, so a deficit in $r_1(3)$ and $r_1(4)$ is a deficit in
$a_3$ and $a_4$ themselves. Bundling is not the protection, since neither shell
passes through a bundled coefficient; the only protection is that $\Delta$ and
the cap depend on the ratio $a_3/a_4$, in which a deficit common to both shells
cancels. At \SI{e12}{\per\cubic\centi\metre} the deficit is not common, $8.3$
at $p=3$ against $4.4$ at $p=4$; at \SI{e15}{\per\cubic\centi\metre} it nearly
is, $1.8$ against $1.9$. Dividing the model's $a_3$ and $a_4$ by those ratios,
so that they become what Fujimoto's coefficients would give, and recomputing
at every grid point (\texttt{verify\_fujimoto\_r1\_consequence.py},
\texttt{validation/fujimoto\_r1\_consequence/}) moves $\Delta$ by $+0.64$
under the first pair and by $-0.08$ under the second. Over the $280$ points
above \SI{2}{\electronvolt} the cap then moves by \SI{+29}{\percent} at the
median under the first and \SI{-4}{\percent} under the second; the crest
$u^{\mathrm{peak}}$ falls by a factor $6.0$ and $1.9$; and the sensitivity at
the operating point, $\ffrac{3}-\ffrac{4}$ evaluated at $u^{\mathrm{CRE}}$,
moves by \SI{+30}{\percent} and \SI{+11}{\percent} at the median, with ranges
across the points of $-58$ to \SI{+524}{\percent} and $-38$ to
\SI{+68}{\percent}. Both exceed the \SIrange{8}{9}{\percent} that the R-matrix
excitation substitution of Chapter~\ref{ch:atoms} produces. The disagreement
is therefore direct evidence about the coefficients Chapter~\ref{ch:results}
uses; what is open is whether the deficit lies in this model or in the
tabulation, and the scenario applies a deficit measured at
\SI{10}{\electronvolt} at every temperature, which is its stated assumption.
An unexplained disagreement in an absolute population coefficient is reported
here as unexplained, with its stake stated.

\textbf{What the benchmark does establish.} The recombination-fed channel
agrees, and because $r_0$ and $r_1$ carry the same $p$-dependent factor $Z(p)$,
a normalisation error large enough to explain the $r_1$ deficit is excluded by
$r_0$ agreeing. That is a pass, and it should be read as one.

\textbf{What it does not establish, and two limits on it.} It does not fix the
absolute normalisation of the ground-fed channel against an independent
calculation, because no $\ell$-resolved tabulation was available to compare
against. Table~4.1(b) is computed at $T_\mathrm{e} = \SI{1.28e5}{\kelvin}$,
which is \SI{11.03}{\electronvolt}, while this grid reaches
\SI{10}{\electronvolt}, so the comparison sits at the grid edge, with the
tabulated temperature \SI{10.3}{\percent} above it. The companion
Table~4.1(a) is at \SI{e3}{\kelvin} =
\SI{0.0862}{\electronvolt}, two orders below this grid, so there is no second
usable case. The external check is one temperature at one edge.

%-----------------------------------------------------------------------------
\section{Gate D fails, and the failure is in the gate}
\label{sec:gate_d}
% QUESTION: the one gate that compares against an external atomic database
% disagrees at every point. What does that mean, and what does it not mean?

The Atomic Data and Analysis Structure (ADAS) distributes effective
stage-to-stage rate coefficients for hydrogen: SCD, the effective ionisation
coefficient, and ACD, the effective recombination coefficient~\cite{Summers2006}.
These are the coefficients a transport code uses when it does not resolve
excited states, and they are the most widely used external reference this model
could be compared against.

Gate~D forms the model's counterpart by taking the steady-state population
$n_p^{\rm ss}$ of each level at each grid point, weighting it by that level's
ionisation rate coefficient $Q_p$ of Eq.~\eqref{eq:level_balance}, and dividing
by the total neutral population:
\begin{equation}
  {\rm SCD}_{\rm model}
  = \frac{\sum_p Q_p\,n_p^{\rm ss}}{\sum_p n_p^{\rm ss}},
  \label{eq:scd_model}
\end{equation}
then compares against ADAS interpolated to the same temperature within a
factor-of-two density band. The gate passes a point if the ratio
$\eta = {\rm SCD}_{\rm model}/{\rm SCD}_{\rm ADAS}$ lies between $0.5$ and $2.0$.

\begin{center}
\texttt{Gate D: FAIL  (0\% of points)} \hfill \texttt{validation/gate\_summary.txt}
\end{center}

Measured, $\eta$ ranges from $5.73$ to $8346$ over the full grid, and no point
agrees within a factor of two. The extremes are not typical: over most of the
grid the model sits high by a factor between $9$ and $65$, and that systematic
offset peaks near $\Te \approx \SI{2}{\electronvolt}$ (these are the gate's
own numbers, under its linear read of the table; the read is discussed
below). The disagreement has a
shape, which is what makes it diagnosable rather than random.

\subsection{The diagnosis, and the test that settles it}
\label{sec:gate_d_diagnosis}

Two readings were open. One of them has now been tested and it is the
explanation.

The first is that Eq.~\eqref{eq:scd_model} is not the quantity ADAS tabulates.
ADAS defines SCD as the \emph{ionising} coefficient, the ionisation flux
attributable to the ground-state population. The sum in
Eq.~\eqref{eq:scd_model} runs over the full steady state, which includes the
recombination-fed Rydberg levels, the $\nzero = \bm{c}\,\nion$ channel of
Section~\ref{sec:R_depends_on_b}. Ionisation out of
those levels is a recombination-driven flux and belongs under ACD, not SCD.
Since the high levels have small ionisation potentials and large ionisation
coefficients, moving them to the wrong side of the ledger inflates the model's
SCD by a large factor, and inflation is what is measured. The test that would
settle it is to recompute Eq.~\eqref{eq:scd_model} over the ground-fed channel
$\none$ alone, and it has now been run
(\texttt{diagnose\_gate\_d.py}).

The two channels separate exactly, with no approximation:
$\nvec_E = \nzero + \none$ to a relative residual below $10^{-8}$ at every grid
point, which the script checks before using either. Rebuilding the coefficient
on the ground-fed channel alone,
\begin{equation}
  {\rm SCD}_{\rm model} = Q_{1s}
  + \sum_{p \neq 1s} Q_p\,\frac{n^{(1)}_p}{\ngs},
  \label{eq:scd_groundfed}
\end{equation}

\noindent and comparing against the same ADAS table with the same factor-of-two
tolerance moves the agreement from $0$ of $400$ points to $400$ of $400$. The
ratio $\eta$ falls from a range of $5.73$ to $8346$ to a range of $0.702$ to
$1.004$, with median $0.828$; above \SI{2}{\electronvolt} it runs from $0.721$
to $0.937$, median $0.831$, and at the benchmark it is $0.768$
(\texttt{validation/reservoir\_vs\_adas/}). The read matters. ADAS96
tabulates seven temperatures between $1$ and \SI{10}{\electronvolt}, and
\texttt{diagnose\_gate\_d.py}, which first made this comparison, interpolates
linearly in value between them (\texttt{adas\_at}, lines 114 to 124); across
$1$ to \SI{1.5}{\electronvolt}, where SCD rises by nearly two orders of
magnitude, that overstates the table by up to a factor $6.65$. That read gives
$323$ of $400$, a range of $0.130$ to $0.931$ and a median of $0.696$, with a deficit deepening toward low temperature; read logarithmically
in both axes, as the table's spacing requires, the tail disappears, the row
median being $0.72$ at \SI{1}{\electronvolt} and $0.90$ by
\SI{1.2}{\electronvolt}. The diagnosis script is left in place; every ADAS
number in this section comes from the logarithmic read.

The mechanism is visible directly. At collisional--radiative equilibrium the
recombination-fed channel carries between \SI{89.5}{\percent} and
\SI{99.99}{\percent} of $\sum_p Q_p n_p$, and its share rises monotonically with
density, from \SI{97.8}{\percent} at $\SI{e12}{\per\cubic\centi\metre}$ to
\SI{99.99}{\percent} at $\SI{e15}{\per\cubic\centi\metre}$ at
\SI{1}{\electronvolt}. Equation~\eqref{eq:scd_model} is therefore reporting the
recombination channel with a small ground-fed correction, which is why $\eta$
tracks $\nel$: three-body recombination carries an extra factor of $\nel$ and
the mislabelled flux carries it into the coefficient.

\paragraph{The recombination half, built and compared for the first time.}
The same decomposition supplies the coefficient the gate never computed. The net
flux into the neutral ground state carried by the recombination-fed channel,
per unit $\nel\nion$, is
\begin{equation}
  {\rm ACD}_{\rm model}
  = \frac{S_{1s}\,\nion + \sum_{p \neq 1s} L_{1s,p}\,n^{(0)}_p}{\nel\,\nion},
  \label{eq:acd_groundfed}
\end{equation}

\noindent direct radiative recombination into $1s$ plus everything that cascades
or is collisionally transferred down into it. Against ACD96 this agrees at
\emph{all} $400$ points, with $\eta$ between $0.894$ and $1.009$ and median
$0.982$; at the benchmark the model gives
$\SI{2.788e-13}{\cubic\centi\metre\per\second}$ against ADAS's
$\SI{2.883e-13}{}$, a ratio of $0.967$.

\textbf{That is the strongest external check in this thesis.} It is an absolute
coefficient, compared against an independent database built from independent
atomic data, agreeing to better than \SI{10}{\percent} at the benchmark and
within a factor of two everywhere. It was available all along and was not taken,
because the half of the gate that would have taken it was never written.

The second reading was that the two quantities are not defined on the same
system at all. ADAS coefficients describe a closed atom-plus-ion system; this
model holds the ion as a fixed external reservoir with no forty-fourth state
(Section~\ref{sec:two_references}). A stage-to-stage coefficient extracted from
an open system need not equal one extracted from a closed one, and
Section~\ref{sec:attacks} has already shown that the boundary condition moves
$\tauslow$ by up to a factor of $41$. \textbf{The ACD result argues against this
being the dominant effect}: ACD is extracted from the same open system through
the same fixed-ion reservoir and agrees with ADAS to within \SI{3.3}{\percent}
at the benchmark, \SI{1.8}{\percent} at the grid median and \SI{11}{\percent}
at every point; an open-system artifact large enough to produce $\eta$ of
$8346$ in SCD would not leave ACD intact.

\paragraph{The residual SCD deficit, stated rather than dismissed.} Rebuilding
SCD on the ground-fed channel does not make the agreement perfect, and the
remaining disagreement is systematic in a way worth reporting: the model's
ionisation sits below ADAS at all but one of the $400$ points, by up to
\SI{30}{\percent} and by \SI{17}{\percent} at the median, with the maximum
$\eta = 1.004$ at $[4,7]$. Under the linear read of the table the gap closes
steeply with temperature, from a median of $0.302$ over $1$ to
\SI{1.5}{\electronvolt} to $0.817$ over $5$ to \SI{10}{\electronvolt}, a shape
that reads as the signature of truncation; the steep part is the linear read,
and under the logarithmic read the deficit is between $0.72$ and $0.94$ at every
temperature above \SI{2}{\electronvolt}. A deficit of that size is still what
truncation would produce, since ADAS96 carries levels above $n_{\max} = 15$
and the levels nearest the continuum ionise most easily, but its temperature
shape no longer singles that cause out (the $n \ge 12$ flux share the diagnosis
script printed against the linear read is not carried forward). The test that
would settle the
cause is the $n_{\max}$ convergence study of Chapter~\ref{ch:limits} rather
than this one.

\paragraph{The reservoir itself, against the same table.} ADAS's two
coefficients fix the equilibrium reservoir: at ionisation balance
$n_{1s}\,\mathrm{SCD} = \nion\,\mathrm{ACD}$, so
$u^{\mathrm{CRE}} = \mathrm{ACD}/\mathrm{SCD}$, and the model obeys the same
identity with its own ground-fed SCD and ACD to $4\times10^{-13}$ at every grid
point. That makes the quantity Chapter~\ref{ch:results} is built on directly
comparable, and the comparison has two parts
(\texttt{verify\_reservoir\_vs\_adas.py},
\texttt{validation/reservoir\_vs\_adas/}). The \emph{level} is high:
$u^{\mathrm{CRE}}_{\mathrm{model}}/u^{\mathrm{CRE}}_{\mathrm{ADAS}}$ has median
$1.18$ over the $400$ points, runs from $0.89$ at $[4,7]$ to $1.40$ at $[0,1]$,
and is $1.26$ at the benchmark, structured about equally in $\Te$ and in
$\nel$. The model's ground reservoir per ion sits a fifth above ADAS's, which
moves the operating point beneath the cap of Section~\ref{sec:bound} without
touching the cap. The \emph{displacement}, which is what the plateau error is
built from, agrees. Over the $392$ one-interval temperature steps the model's
$\Delta\ln u$ has median magnitude $0.26$ and differs from ADAS's by a median
$0.014$; above \SI{2}{\electronvolt} the median difference is $0.011$ and the
worst $0.050$, at $[15,0]$; at the benchmark step the two are $-0.270$ and
$-0.258$. The reservoir gain $G$ this thesis tabulates is therefore within
about \SI{5}{\percent} of the one ADAS implies, at the median over the
defended range, and the refuter set before the run, a difference above $0.1$
at any interior step above \SI{2}{\electronvolt}, appeared at none of $264$.
Every number in this paragraph comes from the logarithmic read; on the linear
read of \texttt{diagnose\_gate\_d.py}, which overstates SCD by up to a factor
$6.65$ over $1$ to \SI{1.5}{\electronvolt}, the refuter would have appeared at
$9$ of the $264$ steps.

\subsection{Three further defects in the gate itself}

The gate is also incompletely implemented, and this must be said alongside the
failure rather than instead of it.

The recombination half does not exist \emph{in the gate}. The ACD table is
loaded once, at line 299 of \texttt{validate\_gates.py}, and the variable
holding it is never read again anywhere in the file; only the ionisation
comparison runs, despite a docstring that describes both. That half is computed
in \texttt{diagnose\_gate\_d.py} (Section~\ref{sec:gate_d_diagnosis}).

The interpolation is wrapped in a catch-all exception handler that sets the
comparison value to not-a-number and continues, so a failure to find the right
density band is indistinguishable from a point that was legitimately skipped.

And the gate's status is recorded three incompatible ways: \texttt{README.md}
line 119 calls it documented, the docstring at line 47 predicts a pass above
$\Te = \SI{5}{\electronvolt}$, and \texttt{gate\_summary.txt} reports failure
at \SI{0}{\percent} of points. The correct entry is the measured one: Gate~D
fails.

\subsection{What is claimed, and what is not}

\textbf{Gate~D itself still fails, and it has not been changed.}
\texttt{validate\_gates.py} is untouched: no tolerance was widened, no subset
was selected, and \texttt{gate\_summary.txt} still records
\SI{0}{\percent}. What has changed is that the failure now has a demonstrated
cause rather than two undemonstrated candidates. The gate compares a mixed
quantity against a table that keeps the two channels apart by design, which is
a defect in the comparison and not in the matrix.

Deciding what Gate~D should test belongs with whoever owns the gate. What is
claimed here is narrower: rebuilt on the two channels separately, the
coefficients agree with ADAS within a factor two at all $400$ points for SCD,
the model $0$ to \SI{30}{\percent} low with a median ratio of $0.83$, consistent
with the $n_{\max}$ truncation but not uniquely so, and at all $400$ for ACD
with a median ratio of $0.98$.

The failure does not cost the thesis its only comparison of absolute effective
coefficients against an external database; that comparison exists, hidden by
the definition. The failure does not touch the structural results
of Chapter~\ref{ch:theory}, which are statements about whatever $\Lmat$ is
supplied and would hold for the ADAS-consistent matrix equally. It does not
touch the level-resolved comparison of Section~\ref{sec:fujimoto}, which
compares the same decomposition against tabulated coefficients rather than
against a contraction of $42$ coefficients into one. And the contraction is the
reason the two comparisons carry different weight: many incorrect $\bm{a}$ would
contract to the correct SCD, so a passing Gate~D would have been weak evidence
even had it passed.

%-----------------------------------------------------------------------------
\section{Errors found by our own audit}
\label{sec:errors_found}
% QUESTION: what did this project get wrong, how was each error found, and how
% large was it?

Section~\ref{sec:atoms_corrections} reports two implementation errors that this
project found in itself, with their sizes, and observed that both were invisible
to every consistency test the pipeline ran. This section states what each
changes here, adds two more, and keeps the four together rather than in an
appendix: a model whose authors found and bounded four of its own errors is
better characterised than one reporting none, and in three of the four cases
the bound is itself a result.

\subsection{The $\ell$-mixing Coulomb logarithm}

Proton-impact $\ell$-mixing redistributes population between sublevels of the
same $n$ and, as Section~\ref{sec:lmixing} showed, is the dominant loss channel
for the metastable $2s$ level. The rate coefficient follows Badnell's
formulation~\cite{Badnell2021}, whose Eq.~9 carries a factor
\begin{equation}
  F(U_m) = \frac{\sqrt{\pi}}{2}U_m^{-3/2}\operatorname{erf}(\sqrt{U_m})
           - \frac{\mathrm{e}^{-U_m}}{U_m} + E_1(U_m),
  \label{eq:FUm}
\end{equation}
in which $\operatorname{erf}$ is the error function, $E_1$ the exponential
integral, and $U_m$ the dimensionless ratio of the smallest energy transfer the
Debye cutoff permits to $k_B\Te$; $F(U_m)$ is the Coulomb logarithm of the
interaction, which diverges as $U_m \to 0$.

The implementation set $F = 1$, justified in its docstring as the low-density
limit $U_m \to 0$; the limit was inverted, and
Section~\ref{sec:atoms_corrections} gives the finding, the factor of three to
seven by which the rates were low, and the measured effect on the spectrum.
Evaluated with the local Debye length at the benchmark point $[23,5]$,
$F = 6.62$ for $2p$--$2s$, $5.71$ for $3d$--$3p$, $4.26$ for $5g$--$5f$ and
$2.89$ for $8k$--$8i$ (\texttt{validation/lmix\_cutoffs/}); with the
production table's frozen density of \SI{e14}{\per\cubic\centi\metre} the
$2p$--$2s$ value is $6.95$ (Section~\ref{sec:atoms_corrections}). A derivation
note records $6.65$, $5.73$, $4.29$ and $2.92$ for the same four pairs.

The insensitivity of the slow eigenvalues to that error, which
Section~\ref{sec:atoms_corrections} measures, is specific to a process already
far faster than $\taurel$; it is not a general licence to be careless about
fast rates, and the $\ell$-mixing scan below is what establishes the range over
which it holds.

\subsection{And $\ell$-mixing is saturated}

The same insensitivity was then tested directly rather than argued. Scaling
every $\ell$-mixing rate in the matrix over a range of $\times0.1$ to $\times10$
moves $\ffrac{3}-\ffrac{4}$ by less than \SI{0.3}{\percent} at the benchmark
point and at the ridge, and by less than \SI{5}{\percent} at the cold corner.
Only switching $\ell$-mixing off entirely changes anything. Rebuilding the
matrix with a density-consistent Debye cutoff in place of a frozen one moves
$\ffrac{3}-\ffrac{4}$ by at most \SI{0.42}{\percent}.

The claim that the frozen Debye density varies the cutoff factor by about
\SI{30}{\percent} across the grid is wrong: the module's own functions give
$\times3.7$ for $n=2$ and at least $\times50$ for $n=8$
(Section~\ref{sec:lmixing}); the conclusion survives because the observable does
not depend on the input over that range.

\subsection{The eigenvalue magnitude filter}

The spectrum of $\Lmat$ was extracted by a routine containing the line
\texttt{eigs = eigs[eigs < -1.0]}, which discards every eigenvalue smaller in
magnitude than $1\,\si{\per\second}$. Where the slow eigenvalue is smaller than
that, at the cold end of the grid, the filter deleted it and promoted the whole
ladder by one step, so $\tauslow$, $\taurel$ and $\Msep$ were wrong together
while all three remained finite, negative and correctly ordered. It was found by
recomputing the spectrum unconditionally and comparing bit-for-bit against the
stored arrays, the three-implementation agreement of Section~\ref{sec:gate_e}.
Section~\ref{sec:atoms_corrections} gives the extent of the damage, the ladder
shift that confirmed the mechanism, the preserved filtered outputs and the
removal of the literal from two dormant files. What belongs here is the shape
of the wrong answer: at the cold corner the filter turned
$\Msep = 1.729\times10^{9}$ into $\Msep = 1.467$, and the grid minimum from
$86.77$ into $1.341$, values that read as a physics finding, since
near-breakdown of timescale separation in the coldest, thinnest plasma is
exactly what a reader would expect to see. The later correction of the two
dormant files changes no number in this thesis: at the benchmark the two filters
return the same $\tauslow$ to every digit, and they differ only where the slow
mode is slower than $1\,\si{\per\second}$.

\subsection{The central quantity was misnamed}

The largest of the four errors changed no number and inverted a sentence.

The quantity that had been reported as the quasi-steady-state error was
defined as the difference between the tabulated ratio and the true
one, with the tabulated ratio evaluated from a two-argument function of $\Te$
and $\nel$. Section~\ref{sec:two_references} shows that the excited-state ratio
is a function of three arguments, not two, the third being the state of the
ground-state reservoir. Fixing the first two therefore silently imposes
equilibrium ionisation balance, so what was measured was the distance to the
equilibrium answer, not the error of the quasi-steady-state closure.

Integrating the full system and evaluating both errors along the same trajectory
separates them:

\begin{center}
\begin{tabular}{lrr}
  \toprule
  point & error against equilibrium & error of the QSS closure \\
  \midrule
  benchmark $[23,5]$ & $6.34\times10^{-2}$ & $8.66\times10^{-6}$ \\
  worst case $[0,4]$ & $3.869\times10^{-1}$ & $6.73\times10^{-9}$ \\
  \bottomrule
\end{tabular}
\end{center}

They differ by four to eight orders of magnitude, most at the point carrying the
largest error. The physics was intact and every number reproduced; what was
wrong was the name and the sentence built on it, and Chapter~\ref{ch:results}
states the consequence, that the approximation everyone worries about is not
the one that fails. The error was found by an adversarial reading of the thesis text
itself, not of the code, and it is why Section~\ref{sec:eps_definitions}
defines four error measures separately.

\subsection{A statement in Chapter~\ref{ch:theory} that measurement contradicted}
\label{sec:promise_broken}

The statement, in Section~\ref{sec:source_quadratic}, was that at
$\nel = \SI{e12}{\per\cubic\centi\metre}$ three-body recombination supplies
under \SI{10}{\percent} of the feed into any level, and it was used to propose
testing the radiative and three-body datasets separately at the two ends of the
density grid. Measured against the matrix, $30$ of $36$ levels exceed
\SI{10}{\percent} three-body feed at both \SI{1.00}{\electronvolt} and
\SI{2.95}{\electronvolt}, and $26$ of $36$ at \SI{10}{\electronvolt}. Into the
bundled shells it is essentially the whole feed.

The density grid therefore does not separate the two recombination datasets, and
\textbf{no such separate test was run or can be}; Section~\ref{sec:source_quadratic}
now states the measured fractions and withdraws the promise. What remains is the
gap: the radiative and three-body datasets enter this model only in combination,
nothing in this thesis tests them apart, and a test that could would need a
density below the bottom of this grid, outside the conditions the thesis is
about.

%-----------------------------------------------------------------------------
\section{Reproduction, and what a document cannot certify}
\label{sec:reproducibility}
% QUESTION: can the numbers in this thesis be regenerated, and by whom?

\subsection{Regeneration from the canonical matrix}

Every quantitative result in Chapters~\ref{ch:theory} and~\ref{ch:results} was
regenerated from $\Lmat$ and compared against the stored output. The timescale
and step-error grids reproduce bit-identically, including the values that come
from numerical integration and not from linear algebra. The plateau and persistence
tables reproduce byte for byte, $680$ of $680$ rows, with $784$ superposition
residuals at exactly zero difference. The nineteen points corrected by the
eigenvalue filter reproduce exactly, including the ladder-shift identity.

This is reproduction in the narrow sense: the same code, run again, on the same
inputs, gives the same answers. It excludes a stale array and it excludes a
non-deterministic solver. It says nothing about whether the algebra is right.

\subsection{Independent re-derivation}

The stronger check is re-derivation. All six central results of
Chapter~\ref{ch:theory} were derived from the linear collisional--radiative
system alone by a party working from the physical statement of the problem,
implemented from scratch, with the predicted values written down before any
project file was opened.

In the table below, $\bdep$ is the ground-state departure coefficient of
Eq.~\eqref{eq:departure}, the ratio of the actual ground-state population to its
Saha--Boltzmann value, and $\bdep^{\rm CRE}$ is the value it takes in
collisional--radiative equilibrium (CRE), the steady state the lookup tables
assume.

\begin{center}
\begin{tabular}{lrr}
  \toprule
  quantity & \texttt{chapter3.tex} & independent re-derivation \\
  \midrule
  $\bdep^{\rm CRE}$ at the benchmark        & $956$                        & $956.509$ \\
  switching points $c_m/a_m$                & $2.73\times10^{3}$, $1.91\times10^{4}$ & $2727.3$, $19085.5$ \\
  peak $\bdep$                              & $7.2\times10^{3}$            & $7214.7$ \\
  $\ffrac{3}$, $\ffrac{4}$, difference       & $0.260$, $0.048$, $0.212$    & $0.259652$, $0.047725$, $0.211927$ \\
  maximum of $\ffrac{3}-\ffrac{4}$          & $0.4514$                     & $0.451357$ \\
  $\Robs$ at the benchmark                  & $0.7778$                     & $0.777768$ \\
  \bottomrule
\end{tabular}
\end{center}

The timescales reproduce to four digits, $\tauslow = \SI{22.73}{\micro\second}$,
$\taurel = \SI{2.277}{\nano\second}$, $\Msep = 9981.9$.

Two further invariances were measured in the same pass and belong with the
result they protect. The shell separation $\Delta$ that sets the bound of
Section~\ref{sec:bound} is invariant under rescaling the ground-to-ion
normalisation, verified exactly under factors of $10^{10}$ and $10^{-7}$, while
the individual channel fractions are not; and $\Delta$ varies by
\SI{0.07}{\percent} across five different $\ell$-weightings of the shell sum, so
the bound is not an artifact of how sublevels were weighted.

\subsection{A prediction that was falsified, and improved the claim}

The same re-derivation made one prediction that turned out to be wrong, and it
is more informative than the six that were right.

The prediction was that $\ffrac{3}-\ffrac{4}$ must vanish at high density, on
the reasoning that in local thermodynamic equilibrium both supply channels
produce the same Boltzmann distribution, so the two shells would draw
identically and the difference would close. It does not happen. At
$\Te = \benchTe$, $\Delta$ doubles over the lowest part of the grid and then
stops moving: $0.979$ at $\SI{e12}{\per\cubic\centi\metre}$, $1.946$ at
$\SI{1.4e14}{\per\cubic\centi\metre}$, $1.939$ at
$\SI{e15}{\per\cubic\centi\metre}$
(\texttt{validation/fujimoto\_r1\_consequence/}). Above about
$\SI{7e12}{\per\cubic\centi\metre}$ it is flat, and the bound $\tanh(\Delta/4)$
built on it spans only $0.394$ to $0.450$, a \SI{14}{\percent} range, across the
whole of the rest of the grid.

The reason is the open boundary again, this time in the model's favour: the
recombination source is imposed externally rather than tied to the ion density
by a Saha relation, so the two channels cannot merge by construction, and
Gate~C's \SI{1.5}{\percent} of Saha says the grid never approaches the limit in
which they would. The consequence for the density maximum of the diagnostic
error, a position effect rather than a span effect, is
Table~\ref{tab:position_effect} in Section~\ref{sec:density_mechanism}.

\subsection{One item demoted: the May/August cross-validation}

A cross-validation between two runs four months apart, on different matrices,
with different observables and different numerical methods, was recorded in this
project's register as verified and described as the strongest single check it
had. The agreement is close: $\tauslow$ to \SI{0.064}{\percent}, $\taurel$ to
\SI{0.200}{\percent}, $\Msep$ to \SI{0.264}{\percent}, $\epsstep$ to
\SI{0.014}{\percent}. Predicting how long the error stays above
\SI{10}{\percent} from the plateau-then-decay picture gives
\SI{12.478}{\micro\second} against the earlier run's full integration,
\SI{12.4794}{\micro\second}.

Two corrections apply, and both weaken it.

The apparent agreement to eight parts in $10^{5}$ on the duration is
cancellation. The inputs agree only to \SIrange{0.06}{0.6}{\percent}, and
carrying the earlier run's own internal numbers through gives
\SI{12.406}{\micro\second}, off by \SI{0.59}{\percent}. \textbf{The figure is
agreement at the \SI{0.6}{\percent} level}, which is still strong.

More seriously, \textbf{the earlier run has never been reproduced against the
canonical matrix}. This project's own derivation record says of these same
numbers that they are a reported result, not independently re-run, and should be
treated as strong documented evidence rather than as self-verified. A
cross-validation whose earlier half cannot be re-executed is a citation, not a
check. It is reported here as documented evidence and is not counted on the
ladder of Section~\ref{sec:validation_summary}. Under the ground-truth ordering
this thesis works to, physics over mathematics over code over documents,
a document is the weakest of the four, and it is where this particular claim
currently rests.

\subsection{Two provenance defects that are not repaired}

Two arrays of timescales, $\Msep$, $\tauslow$ and $\taurel$, are written to the
same three file paths by two different scripts. Whichever ran last wins, and six
downstream readers cannot tell whose numbers they are holding. This is exactly
the condition under which the eigenvalue filter propagated, since the filtered
and unfiltered versions of the same array have the same name.

Two things bound what that costs here, and both were checked rather than
assumed. The two writers are not two calculations: each sorts the eigenvalues of
the same canonical $\Lmat$ and keeps the negative ones, neither applies the
filter that caused the earlier incident, and the arrays on disk agree with the
independently produced \texttt{verify\_timescales.py} artifact to $2\times
10^{-16}$. And none of the six readers feeds this document: every figure and
table they write is absent from it. The benchmark $\Msep = 9982$ quoted in
Chapter~\ref{ch:theory} and in the abstract is not read from these arrays at
all; it is $\tauslow/\taurel$ computed from $\Lmat$ by
\texttt{verify\_timescales.py} and recorded in
\texttt{validation/spectrum\_testpoint.csv}. The defect is therefore a latent
hazard in the repository rather than an ambiguity in any number printed here,
and it is left in place and named rather than quietly repaired.

An audit script exists to detect single-writer violations, and it reports an
unfixed failure of this kind. It is also incomplete. Its test for a write reads the file mode from the second
argument or a keyword, so the idiom \texttt{path.open("w")} passes through it
unrecorded. At the tip of \texttt{main} on 23~August that idiom had four users;
in the tree this thesis is built
from it has $45$ call sites in $29$ scripts, $42$ of them writing
comma-separated output, among them the producer of the plateau map that
fifteen places in Chapter~\ref{ch:results} read from
(\texttt{verify\_writer\_census.py}, \texttt{validation/writer\_census/}). The
audit should not be cited as evidence that the single-writer property holds.
The census confirms that the property fails for exactly the three timescale
arrays, each written by both \texttt{qss\_analysis.py} and
\texttt{validate\_gates.py}, and holds for the plateau map, which has one
writer.

%-----------------------------------------------------------------------------
\section{What this chapter establishes}
\label{sec:validation_summary}
% QUESTION: after grading every check, what is the model licensed to be used
% for, and where?

Table~\ref{tab:ladder} collects the checks in the order a reader should weight
them, with the grade of Section~\ref{sec:gate_philosophy} attached to each.

\begin{table}[tbp]
\centering
\caption[The validation ladder, graded by severity and provenance]{Every check
this chapter grades, with what it tests, its severity, whether its result can be
regenerated, and the measured outcome. \emph{Tautological} means the check
cannot fail except on an arithmetic bug, because the quantity tested was
constructed from the relation being tested; \emph{weak} means it can fail but
only for an error far larger than any plausible one; \emph{severe} means a
plausible wrong version of this model fails it. The provenance column is the
more uncomfortable of the two: \emph{stamped} means a script in the repository
regenerates the number into an artifact under \texttt{validation/};
\emph{notes} means the number is recorded in this project's working files with
no producing script, so it cannot be re-executed and cannot fail anything today.
None of the checks graded severe are in that condition, and the one
\emph{notes} entry left is the $\tanh$ bound, which is graded tautological. The
counts stated in this caption are checked against the table's own rows by
\texttt{verify\_ladder\_counts.py}, which exits non-zero if they drift apart. Of the eighteen
internal checks, twelve are severe, four are tautological and two are weak; of the five external comparisons, two
pass, one passes only at the higher densities, one fails, and one is open.}
\label{tab:ladder}
\footnotesize
\begin{tabular}{p{3.9cm}p{1.8cm}p{1.3cm}p{5.6cm}}
\toprule
check & grade & record & measured \\
\midrule
Atomic-data provenance (Ch.~\ref{ch:atoms})  & severe & stamped & every coefficient traces to a named dated source; tallies exact \\
Excitation rates against RMPS (Ch.~\ref{ch:atoms}) & severe & stamped & \SI{81.4}{\percent} within \SI{20}{\percent} for $n_{\rm upper}\le4$, mean \SI{13.0}{\percent}; $n=5$ disagrees by factors of $2$--$5$, localised to the top shell of the RMPS basis \\
Raw cross sections, microscopic reversibility & severe & figure & ratio $0.9995 \pm 0.0032$, \SI{98.7}{\percent} within \SI{1}{\percent} \\
One energy ladder on both sides & severe & stamped & Saha ratio $0.999997$ for all $43$ states, spread $1.4\times10^{-15}$ \\
Superposition of the two channels & severe & stamped & $3.075\times10^{-14}$ over $784$ pairs; catches a $3\leftrightarrow4$ swap \\
Reduced against full $\Robs$ & severe & stamped & $0.777768$ by three routes; catches the same swap \\
Deliberate fault injection & severe & stamped & three of four injected faults invisible to the conservation residual; the fourth moves it by thirteen orders \\
Line-weighted against shell ratio & severe & stamped & ratio $0.948$ to $1.000$ over $784$ pairs, below unity at every one; census $44$ against $45$ \\
Escape-factor cross-check & severe & stamped & $\sigma_0$ against an independent module, \SI{0.059}{\percent} apart against a \SI{1}{\percent} raising tolerance \\
Pinned input hashes & severe & stamped & both SHA-256 match byte for byte \\
Recomputation guards on the figures & severe & code & every plotted row rebuilt from $\Lmat$ at $10^{-9}$, raises \\
Truncation convergence of $\ffrac{3}-\ffrac{4}$ & severe & stamped & increments one-signed and geometric at $245$ of $280$ defended points, ratio median $0.67$; extrapolated \SI{-0.46}{\percent} at the benchmark, \SI{+1.6}{\percent} at the cold corner, within \SIrange{-0.82}{+1.03}{\percent} on the defended set. Tests one quantity, not the state space \\
Gate~A, detailed balance & tautological & stamped & $<\SI{0.01}{\percent}$, but $K_{\rm deexc}$ was derived from $K_{\rm exc}$ \\
Column-sum conservation & tautological & stamped & $3.61\times10^{-11}$; see the fault-injection row for what it cannot see \\
The $\tanh$ bound & tautological & notes & holds at all $248$ points, and holds on corrupted input \\
Gate~E, timescale hierarchy & tautological & stamped & $\Msep \ge 86.8$ against a threshold of $1$ \\
Gate~B, coronal limit & weak & stamped & worst convergence $0.0034$ against a tolerance of $0.5$ \\
Gate~C, Saha ceiling & weak & stamped & monotone and below Saha; approach fraction \SI{1.5}{\percent}, excluded from the flag \\
\midrule
Fujimoto App.~4B timescales & external, passes & stamped & $0.34\times$ and $1.77\times$ at a grid point his example lands on \\
ACD96 against this model's recombination coefficient & external, passes & stamped & all $400$ points within \SI{11}{\percent}; $\eta$ from $0.894$ to $1.009$, median $0.982$, \SI{3.3}{\percent} at the benchmark (logarithmic read) \\
Fujimoto Table~4.1, $r_0$ & external, passes at $\lg\nel = 21$ & stamped & \SIrange{0.1}{1}{\percent} for $n\ge4$ at the higher density; \SIrange{12}{14}{\percent} at $p=3,4$ at $\lg\nel = 18$; $r_0(2) = 0.7392$ against $0.730$ \\
Fujimoto Table~4.1(b), $r_1$ & external, \textbf{open} & stamped & $8.3\times$ low at $p=3$. Truncation, Lyman trapping and the $\ell$~closure each tested and each eliminated; unexplained (\S\ref{sec:fujimoto}) \\
Gate~D, SCD96 against ADAS & \textbf{FAIL} & stamped & $\eta \in [5.73,\ 8346]$, $0$ of $400$ within a factor $2$. Rebuilt on the ground-fed channel alone and read logarithmically it reaches $400$ of $400$, $\eta$ from $0.702$ to $1.004$, median $0.828$; the gate itself is unchanged \\
\bottomrule
\end{tabular}
\end{table}

\subsection{What the model may be used for}

Three statements are supported by the evidence above, and they are the ones
Chapter~\ref{ch:results} relies on.

\textbf{The equations are solved correctly.} The superposition residual, the
independent re-derivation of all six central results, the bit-for-bit
regeneration, and the grid-wide conditioning measurement together establish that
the linear algebra does what Chapter~\ref{ch:theory} says it does, to a
precision many orders of magnitude finer than any quantity the thesis reports.

\textbf{The input rates are traceable and benchmarked where the observable
lives; one derived coefficient there is not.} The rates feeding $n \le 4$, which is where $H_\alpha$ and $H_\beta$
originate, agree with an independent R-matrix calculation at a mean absolute
error of \SI{13.0}{\percent}, and the two transitions this thesis leans on
hardest, $1s\to2p$ and $2p\to3d$, agree to between \SI{2.4}{\percent} and
\SI{13.3}{\percent} across the temperature range. Above $n=4$ the atomic data are
less certain and the truncation is measurable. Section~\ref{sec:convergence}
tests one component of the mechanism, the sensitivity $\ffrac{3}-\ffrac{4}$,
and finds it converged to within \SI{1.03}{\percent} at the defended points
where the geometric tail can be summed, in a stamped run
(\texttt{validation/nmax\_downward\_scan/}); under the same truncation
$u^{\mathrm{CRE}}$ is converged to within \SI{0.13}{\percent} at the $257$
defended points where its increments are geometric, $\Delta$ to within
\SI{0.57}{\percent} and $\tauslow$ to within \SI{1.6}{\percent} at all $280$,
and $\Delta\ln u$ across the defended heating steps moves by at most
$2.3\times10^{-5}$ between $n_{\max}=14$ and $15$ against a median step of
$0.20$ (Appendix~\ref{app:numerics}). What is measured
for the reservoir is a scan rather than a truncation: weakening the ionisation
and three-body exchange of the top six shells eightfold moves $\epsplat$ by
\SI{0.7}{\percent} at the median over the defended pairs
(Section~\ref{sec:ion_recomb}). All of that concerns the rates going in. The
coefficients that come out at the two shells the observable uses are a separate
matter, and one of them is in open disagreement: the ground-fed $r_1(3)$ and
$r_1(4)$ of this model sit factors $8.3$ and $4.4$ below Fujimoto's tabulation
at the lowest density, which is a deficit in $a_3$ and $a_4$ themselves and so
in the sensitivity built from them (Section~\ref{sec:fujimoto}). Benchmarked
inputs are therefore not the same thing as a benchmarked susceptibility, and
this thesis can claim the first but not the second.

\textbf{The structural results are robust to the choices that could have
manufactured them.} Closing the ion boundary changes the persistence census by
one part in $325$. Scaling $\ell$-mixing over two decades moves the observable
by under \SI{0.3}{\percent} where the thesis makes claims. The single-slow-mode
estimate of Eq.~\eqref{eq:elm_bound} reproduces the true time-averaged error to
within \SI{2.2}{\percent} at \SI{100}{\micro\second} over the $448$ pairs above
\SI{2}{\electronvolt}.

\subsection{What it may not be used for}

\textbf{No effective coefficient from this model should be compared against a
database without first separating its two supply channels.} Gate~D as written
fails at every point on the grid because it compares a coefficient that mixes
both channels against a table that separates them (Section~\ref{sec:gate_d});
rebuilt channel by channel it agrees within a factor two at all $400$ points
for SCD and ACD, and the model's total effective ionisation coefficient remains
a quantity ADAS does not tabulate.

\textbf{No claim rests on the terminal shell}, whose population is high by a
factor consistent with truncation, or on the $\ell$-distribution inside the
bundled shells, which is assumed rather than measured.
Section~\ref{sec:bundling_check} tests that assumption: $\ell$-mixing licenses
it over most of the grid, and at the two highest densities below
\SI{3}{\electronvolt}, where the locally evaluated mixing loses to collisional
loss, detailed balance among shells within $10^{-4}$ of Saha--Boltzmann
equilibrium licenses it instead; the adjacent-shell rates themselves are not
$\ell$-blind, and no grid point escapes both criteria.

\textbf{Nothing here validates the physics that is absent.} Every check in this
chapter tests a $43$-state atomic hydrogen model against itself, against
arithmetic, and against other atomic hydrogen models. None of them tests
transport, molecular channels, or radiation trapping, because none of those
appears in $\Lmat$ at all. A model can pass every check in
Table~\ref{tab:ladder} and still be the wrong model for a divertor.
Chapter~\ref{ch:limits} is where that question is asked, and it is asked with
numbers rather than with hedges.

\medskip
The instrument is now characterised: what it computes, how precisely, on which
region of the grid, and where the evidence for it runs out. What it says is the
subject of Chapter~\ref{ch:results}.

%     breaks no existing cross-reference.
%
%  2. OVERLAP, RESOLVED BY HOMES. Three passages exist in both chapters:
%       chapter6 sec:open_system   vs  Section~\ref{sec:attacks}   (home: ch4)
%       chapter6 sec:r1_deficit    vs  Section~\ref{sec:fujimoto}
%       chapter6 sec:nmax          vs  Section~\ref{sec:convergence}
%     The Chapter 6 versions carry the SCOPE consequence; the Chapter 4
%     versions carry the EVIDENTIAL standing. On the r_1 deficit: the table
%     is bundled-n (backlog I.1, 10 Sep 2026; findings_10 ADDENDUM I), the
%     l-closure was tested and does NOT account for the gap, and the deficit
%     DOES reach a_3 and a_4; its consequence for Delta, the cap and S is
%     quantified in the Status paragraph of sec:fujimoto
%     (validation/fujimoto_r1_consequence/). Do not reintroduce the earlier
%     "does not propagate" sentence.
%
%     [17-18 Sep 2026] verify_bundling_psm20.py HAS now been executed
%     (sec:bundling_check in chapter 6, outputs/bundling/); an earlier
%     version of this comment said the "never executed" sentence was still
%     true and must not be corrected. It was corrected. The n_max
%     convergence result carried in Section~\ref{sec:convergence} is now
%     stamped (verify_nmax_downward_scan.py, 18 Sep 2026); the working-note
%     values (ADDENDUM D.8) it replaced had the opposite sign at the
%     benchmark and the cold corner, and the text records that. Its \todo asking for provenance on the
%     4.9-6.2x terminal-shell excess is also closed there, by the internal
%     terminal-vs-interior measurement (a_p at [49,0]: 9.4x at p=8, 10.7x at
%     p=9, 6.3x at p=10), which does not bracket it (band 6.3 to 10.7 against
%     4.9 to 6.2) but lies within a factor 2.2 of it and shares the density
%     trend; stamped 21 Sep 2026 by verify_terminal_vs_interior.py.
%
%  3. chapter3.tex:222-225 asserts a severity for the conservation check that
%     Section~\ref{sec:conservation_severity} measures and refutes. Either
%     soften line 224 or let this chapter carry the correction; do not leave
%     both standing as written.
%=============================================================================

\chapter{Does the Model Work?}
\label{ch:validation}

Chapter~\ref{ch:theory} built a $43\times43$ rate matrix, eliminated the fast
variables, and extracted from the result a closed-form statement about how a
line ratio responds to its reservoir. Chapter~\ref{ch:results} will use that
machinery to make quantitative claims about a diagnostic in use on real
machines. Between the two sits a question that no amount of further algebra
answers: why should anyone believe a number this model produces?

The answer is not a single verdict. It is a graded one, because the
checks that have been run on this model are not of equal worth, and several of
the ones that look most impressive are worth almost nothing. This chapter sorts
them.

\medskip
\noindent\textbf{What this chapter takes as given.} The state space, the atomic
data and their sources, and the benchmark of the excitation rates against an
independent R-matrix-with-pseudostates (RMPS) calculation are established in
Chapter~\ref{ch:atoms}, Sections~\ref{sec:atomic_uncertainty}
and~\ref{sec:atoms_corrections}, and are not repeated. The construction of
$\Lmat$, the conservation identity, the quasi-steady-state (QSS) elimination,
and the two-channel decomposition of the excited populations come from
Chapter~\ref{ch:theory}. This chapter asks only what the checks run on that
construction are entitled to establish.

\medskip
\noindent\textbf{What this chapter does not do.} It reports no result about
divertor plasmas. Every number below is about the instrument, not about the
physics the instrument is pointed at. Keeping those separate matters here more
than anywhere else in the thesis: a model that predicts something interesting
is not thereby a correct model, and a check that uses the interesting
prediction as its evidence has verified nothing.

%-----------------------------------------------------------------------------
\section{What a check can and cannot show}
\label{sec:gate_philosophy}
% QUESTION: what distinguishes a check that could have failed from one that
% could not, and how is that distinction measured rather than asserted?

\subsection{Three things that get called validation}

Three activities are routinely reported under one heading, and they license
different conclusions.

\emph{Verification} asks whether the equations that were intended got solved
correctly. Conservation residuals, detailed balance, matrix identities,
convergence with tolerance, agreement between two independent code paths: all
of these compare the model against itself or against arithmetic. They can
detect a transposed index. They cannot detect a rate coefficient that is wrong
by a factor of three, because a uniformly wrong dataset satisfies every one of
them.

\emph{Validation} asks whether the model reproduces something established
independently of it. Published coefficients from another group's code, measured
cross sections, a tabulated limit. Only this class of check can reach the
atomic physics.

\emph{Analysis} asks what the model then says. That is Chapter~\ref{ch:results},
and none of it appears here.

\subsection{Severity, and how to measure it}

The useful property of a check is not whether it passes. It is the size of the
set of wrong models it would reject. Call that its severity.

The failure mode is the one a mass balance on a continuous stirred-tank
reactor shows. Compute the outlet stream as inlet minus consumption, then
verify that the balance closes. It closes to machine
precision, always, for any rate constant whatever, because the outlet was
constructed from the balance being tested. The check is real arithmetic and it
is entirely uninformative about the kinetics. A model can be wired correctly
and be wrong.

Severity is measurable, and two operations measure it.

\begin{enumerate}
\item \textbf{Provenance inspection.} Ask whether the quantity under test was
derived from the quantity doing the testing. If the de-excitation coefficients
were computed from the excitation coefficients through detailed balance, then a
detailed-balance check on them is arithmetic, not physics. Chapter~\ref{ch:atoms}
already states this for the rate dataset (Section~\ref{sec:two_energies}); this
chapter carries it through to the matrix.

\item \textbf{Fault injection.} Break the model deliberately, in a way that a
real mistake could plausibly break it, and see whether the check notices. A
check that survives a tenfold error in a rate coefficient has told you that a
tenfold error in that coefficient is invisible to it.
\end{enumerate}

Fault injection is the sharper of the two, because it produces a number rather
than an argument. It is used in Section~\ref{sec:conservation_severity} on the
conservation check, and its result is the least comfortable finding in this
chapter.

\subsection{The three grades used below}

Every check in this chapter is placed in one of three grades, and the grade is
stated wherever the check is reported.

\begin{description}
\item[Tautological.] Cannot fail except on an arithmetic bug, because the
quantity tested was constructed from the relation being tested. It confirms
that the code implements the construction. It is not evidence about the
physics, and reporting it as though it were makes the validation weaker, not
stronger, by displacing the checks that carry information.

\item[Weak.] Could fail in principle, but the tolerance is so much looser than
the measured value that only a gross error would trip it. A gate that passes
with a factor of $150$ in hand has told you the answer is not absurd.

\item[Severe.] A plausible wrong version of this model fails it. Either the
check has actually fired at some point in the project's history, or a specific
corruption is known to trip it.
\end{description}

Four quantities currently printed by this project's own scripts as validation
are tautological by this definition, and Table~\ref{tab:ladder} names them:
Gate~A, the column-sum conservation residual, the $\tanh$ bound, and Gate~E.
Sections~\ref{sec:conservation_severity} to~\ref{sec:severe_checks} give the
measurements behind each.

%-----------------------------------------------------------------------------
\section{The conservation check is exact by construction}
\label{sec:conservation_severity}
% QUESTION: does the column-sum identity detect a wrong rate coefficient?

Section~\ref{sec:conservation} established that each column of $\Lmat$ must sum
to the ionisation leak from the level that column represents, and reported the
measured residual of Eq.~\eqref{eq:conservation_measured} as
$3.61\times10^{-11}$, a relative residual (Section~\ref{sec:conservation}). The
$3.61\times10^{-11}$ is the maximum over every column at every one of the $400$
grid points; the fault-injection run below has its own baseline, $2.238\times
10^{-12}$, at the single point it was run at. The claim attached to it there is that ``an error in a
single off-diagonal element, a transposed index, or a missing back-reaction all
break this immediately.''

That claim is testable, and it is false.

\subsection{Four injected faults}

Four faults were introduced into the assembly, one at a time, and the residual
recomputed after each. The first three are the exact failure modes the claim
names. Table~\ref{tab:faultinject} gives the outcome.

\begin{table}[tbp]
  \centering
  \caption[Fault injection into the conservation check]{Four deliberate faults
  injected into the assembly of $\Lmat$, with the resulting change in the
  matrix and the conservation residual of
  Eq.~\eqref{eq:conservation_measured} recomputed after each. Three faults
  large enough to move individual matrix elements by up to
  $2.7\times10^{11}\,\si{\per\second}$ leave the residual unchanged at the
  baseline value; only an inconsistency between the Einstein coefficients $A$
  and the total radiative loss rates $\gamma$ is detected. Source:
  \texttt{outputs/CHANGE\_REPORT.md}~\S4.2.}
  \label{tab:faultinject}
  \begin{tabular}{lrl}
    \toprule
    injected fault & $\max|\Delta \Lmat|$ (\si{\per\second}) & residual \\
    \midrule
    none (baseline)                     & $0$                  & $2.238\times10^{-12}$ \\
    $K_{\rm exc}(1s\to2p)$ $\times$ ten   & $6.9\times10^{5}$    & $2.238\times10^{-12}$ \\
    all de-excitation deleted           & $1.9\times10^{11}$   & $2.171\times10^{-12}$ \\
    excitation array transposed         & $2.7\times10^{11}$   & $1.837\times10^{-12}$ \\
    $A_{pq}$ halved, $\gamma_p$ left alone & $3.1\times10^{8}$   & $3.63\times10^{1}$ \textbf{caught} \\
    \bottomrule
  \end{tabular}
\end{table}

A tenfold error in the single most important excitation coefficient in the
model, the one that supplies the entire ground-fed channel of
Section~\ref{sec:R_depends_on_b}, changes the residual in no digit. Deleting
every de-excitation process in the matrix, which is not a subtle mistake,
changes it in the third digit and leaves it three orders of magnitude below any
tolerance. Transposing the whole excitation input array, the classic
index-order mistake, does the same.

\subsection{Why, and what the check is actually worth}

The reason is in the assembly, not in the physics. The diagonal of $\Lmat$ is
built as minus the column sum of the same off-diagonal array that the check
then sums. Whatever numbers are in that array, correct or corrupted, the column
sum returns the ionisation leak. The identity of Eq.~\eqref{eq:conservation} is
imposed by the construction rather than tested by it. This is the reactor mass
balance of Section~\ref{sec:gate_philosophy} exactly.

The one fault that was caught is the one that breaks the construction. In the
notation of Eq.~\eqref{eq:level_balance}, $A_{pq}$ is the Einstein coefficient
for spontaneous emission from level $p$ to level $q$, in \si{\per\second}, and
$\gamma_p = \sum_{q<p} A_{pq}$ is the total radiative loss rate from $p$, also in
\si{\per\second}. The individual $A_{pq}$ enter the off-diagonal elements while
their sum $\gamma_p$ enters the diagonal, so halving one without the other makes
the two inconsistent, and the residual rises from $2.238\times10^{-12}$ to
$36.3$, thirteen orders of magnitude. The check raises rather than warns, it
costs nothing to run, and on this evidence that is the width of what it covers.

The rate magnitudes are reached only by Section~\ref{sec:atomic_uncertainty}'s
comparison against RMPS and by the external comparisons of
Section~\ref{sec:fujimoto}, and by nothing internal.

\paragraph{The grid point of the table.} Three of the four
$\max|\Delta\Lmat|$ values in Table~\ref{tab:faultinject} scale with $\nel$, so
the table fixes its own grid point. Sweeping all $400$ grid points and matching on those three values
identifies $[23,5]$, the benchmark, with a worst disagreement of
\SI{1.63}{\percent}, the rounding of the printed table to two significant
figures; every residual in Table~\ref{tab:faultinject} then reproduces to the
digits shown (\texttt{verify\_fault\_injection.py}, stamped in
\texttt{validation/fault\_injection/}). The script rebuilds $\Lmat$ through the
pipeline's own assembly rather than patching the stored matrix, and the rebuild
reproduces \texttt{L\_grid} exactly before anything is injected, which makes the
faults faults in the model rather than in a re-implementation of it.

%-----------------------------------------------------------------------------
\section{The gate suite, read one gate at a time}
\label{sec:gates_abc}
% QUESTION: for each automated gate, what quantity does it compare, against
% what tolerance, and what would have to be wrong for it to fail?

Five automated gates run over the grid in \texttt{src/validation/validate\_gates.py}
and write \texttt{validation/gate\_summary.txt}. Their headline outcome is four
passes and one failure. Read as a headline it says little, because the five
differ in severity by a wide margin and two of them cannot fail. This section
takes them one at a time; Gate~D, the failure, gets
Section~\ref{sec:gate_d} to itself.

\subsection{Gate A: detailed balance, and where its physics content actually
lives}

Gate~A takes every excitation rate coefficient $K_{\rm exc}(i\to j)$ in the
dataset, in \si{\cubic\centi\metre\per\second}, and predicts the reverse
coefficient from the Boltzmann relation between forward and reverse rates in a
Maxwellian electron distribution at temperature $\Te$,
\begin{equation}
  K_{\rm deexc}(j\to i)
  = K_{\rm exc}(i\to j)\,\frac{g_i}{g_j}\,
    \exp\!\left(\frac{E_j - E_i}{k_B\Te}\right),
  \label{eq:gateA}
\end{equation}
then compares the prediction against the stored reverse coefficient at three
temperatures spanning the grid. Here $g_i$ is the dimensionless statistical
weight of level $i$, $E_j-E_i$ is the excitation energy in
\si{\electronvolt}, and $k_B$ is Boltzmann's constant. The relation holds only
for a Maxwellian distribution, which Section~\ref{sec:maxwellian} established
and which is assumed throughout. The gate passes
if the maximum relative error is below \SI{0.01}{\percent}
(\texttt{validate\_gates.py}:118). Measured: below \SI{0.01}{\percent} over all
$819$ pairs at all three temperatures.

\textbf{Grade: tautological.} The stored de-excitation coefficients were
produced from the excitation coefficients through Eq.~\eqref{eq:gateA} in the
first place (\texttt{compute\_K\_CCC.py}, function \texttt{detailed\_balance}),
as Section~\ref{sec:two_energies} states. Gate~A confirms that the same formula
was applied twice with the same energies and weights. That is a wiring check,
and Chapter~\ref{ch:atoms} already flagged it as one.

The physics content of detailed balance in this project sits somewhere else and
is currently invisible in the outputs. \texttt{src/parsers/qc\_ccc.py}:87--168
tests whether Bray's \emph{raw} convergent-close-coupling (CCC) cross sections,
which this project did not produce, satisfy microscopic reversibility on their
own. The ratio of the two directions has mean $0.9995$ with standard deviation
$0.0032$, and \SI{98.7}{\percent} of comparisons fall within \SI{1}{\percent}
(\texttt{ccc\_qc\_report.png}). A \SI{0.05}{\percent} mean bias with
\SI{0.32}{\percent} scatter on external data is a genuine external check, and it
is the one to lead with.

\subsection{One energy ladder on both sides, which is severe}

There is a second detailed-balance result, and unlike Gate~A it could have
failed.

A model of this kind carries energies twice: once in the Boltzmann factors that
relate forward and reverse collisional rates, and once in the Saha factors that
relate ionisation to three-body recombination. If two different energy tables
were used, or one table were indexed differently from the other, the model would
still conserve particles, still satisfy Gate~A, and still produce plausible
populations. Nothing internal would notice.

Testing ionisation against three-body recombination through the Saha relation
at every level, with CODATA constants on the Saha side, gives a ratio of
$0.999997$ for every one of the $43$ states at every temperature, identical
across all $2150$ entries to $1.4\times10^{-15}$
(\texttt{verify\_saha\_balance.py}, \texttt{validation/saha\_balance/}); the
residual is the rate module's rounded values of $h$, $m_e$ and the
electron-volt against CODATA, and nothing else. The point is not the
closeness to unity, which detailed balance imposes. It is the identity of the residual across
levels whose Saha exponents span the full binding-energy ladder of hydrogen,
from $\mathrm{e}^{13.6\,\mathrm{eV}/k_B\Te}$ for the ground state down to
$\mathrm{e}^{0.06\,\mathrm{eV}/k_B\Te}$ for the top shell. At
$\Te = \SI{1}{\electronvolt}$ those two factors differ by
$\mathrm{e}^{13.54}$, a factor of $7.6\times10^{5}$. Two energy ladders that
differed would produce a residual drifting across that range; the measured
residual does not drift, and its common value is the eight-figure rounding of
the stored state table. \textbf{Grade: severe.} The check could have failed and
did not.

The same audit gives the collisional side over the $819$ excitation pairs at
each of three temperatures, $2457$ comparisons in total: maximum deviation
$7.31\times10^{-9}$, median $7.5\times10^{-10}$, again consistent with the state
table's rounding.

\subsection{Gate B: the coronal limit}

At low enough density an excited level is populated by electron impact from the
ground state and empties radiatively before a second collision reaches it, so
its population per unit ground-state density and per unit electron density
becomes independent of $\nel$. Gate~B tests that plateau: it computes
$n_{2p}/(\ngs\nel)$ and $n_{3d}/(\ngs\nel)$ at the two lowest densities on the
grid and asks whether they have converged.

The tolerance is \SI{50}{\percent} per temperature
(\texttt{validate\_gates.py}:175) and the gate passes if \SI{70}{\percent} of
temperatures pass (line 181). Measured: all $50$ temperatures pass, worst
convergence $0.0034$.

\textbf{Grade: weak.} Passing with a factor of $150$ between the measured
convergence and the tolerance means the gate rejects only a gross error. It
does establish that the low-density end of the grid is in the coronal regime,
which Chapter~\ref{ch:results} relies on when it reads the low-density column
as ionising, and that is worth having. It is not evidence about rate
magnitudes.

\subsection{Gate C: what the Saha gate tests, which is not the Saha limit}

Gate~C is named for the high-density limit in which collisions dominate
radiation and populations approach their Saha--Boltzmann values. It computes
three quantities per temperature: whether $n_{2p}/\ngs$ increases monotonically
with density, whether it stays below the Saha value (with \SI{5}{\percent}
slack), and whether it approaches that value, defined as reaching $10^{-4}$ of
it.

The pass flag uses the first two and \textbf{excludes the third}
(\texttt{validate\_gates.py}:250; the \texttt{approaching} column is computed at
line 239, stored, printed at line 261, and never consulted). The exclusion is
not what makes the gate weak, because the excluded column would pass: the
measured approach fraction runs $0.0123$ to $0.0155$, comfortably above the
$10^{-4}$ the criterion asks for. What makes it weak is that the criterion is
$150$ times slacker than the measurement. Reaching $10^{-4}$ of Saha is not
approaching Saha, and neither is reaching \SI{1.5}{\percent}.

The measurement itself is the useful part. At
$\nel = 10^{15}\,\si{\per\cubic\centi\metre}$, the top of the grid, the $2p$
population sits at about \SI{1.5}{\percent} of its Saha--Boltzmann value, so
that level is nowhere near local thermodynamic equilibrium (LTE) anywhere on
this grid. Only $2p$ was tested, so this is a statement about $2p$ and not about
the whole distribution. It is the expected behaviour for an optically thin
plasma whose ion population is an externally imposed reservoir rather than a
Saha partner, and Chapter~\ref{ch:results} depends on it: it is why the two
supply channels of Section~\ref{sec:R_depends_on_b} never merge, and why
$\ffrac{3}-\ffrac{4}$ does not vanish at high density.

\textbf{Grade: weak, and mis-named.} What Gate~C tests is monotonicity in
density and a Saha ceiling. Both would fail on a sign error in the recombination
source, which is worth something. Neither tests the limit the gate is named for.

%-----------------------------------------------------------------------------
\section{Gate E: the timescale hierarchy}
\label{sec:gate_e}
% QUESTION: is the two-clock structure of Chapter 3 present everywhere on the
% grid, and how demanding is the test that says so?

Gate~E diagonalises $\Lmat$ at every grid point, orders the negative
eigenvalues, and forms $\Msep = \tauslow/\taurel = |\lambda_1|/|\lambda_0|$. It
passes if $\Msep > 1$ at \SI{95}{\percent} of points
(\texttt{validate\_gates.py}:409--412).

Measured over all $400$ points: $\Msep$ ranges from $86.77$ at
$[49,7]$ ($\Te = \SI{10}{\electronvolt}$,
$\nel = \SI{e15}{\per\cubic\centi\metre}$) to $1.72928\times10^{9}$ at the cold
corner; $\tauslow$ spans \SI{75.4}{\nano\second} to \SI{67.2}{\second} and
$\taurel$ spans \SI{0.87}{\nano\second} to \SI{38.9}{\nano\second}
(project working note; Table~\ref{tab:ladder}). The gate asks for $\Msep>1$. The smallest value
anywhere is $86.8$.

\textbf{Grade: tautological as posed, and the real content is elsewhere.} The
threshold sits between two and seven orders of magnitude below the measured
values. No matrix this construction has produced comes within a factor of $86$
of failing, and none plausibly could, because the fast radiative rates and the
slow ionisation rate that set the two eigenvalues are separated by construction
of the physics rather than by a numerical accident.
The reportable quantity is the range, not the pass, together with the fact that
three independent implementations produce it bit-for-bit: \texttt{qss\_analysis.py},
Gate~E itself, and an unconditional recomputation stored in
\texttt{timescales\_unfiltered\_CHECK.npz}. That agreement is what caught the
eigenvalue-filter error of Section~\ref{sec:errors_found}, and a check that has
actually fired is worth more than a threshold nothing approaches.

Two properties of these numbers must travel with them, and both come from
Chapter~\ref{ch:theory}. $\tauslow$ is boundary-dependent, which
Section~\ref{sec:attacks} quantifies. And the slow eigenvalue is not an
equilibration rate: measured against the ground-state ionisation rate
coefficient $K_{\rm ion}(1s)$, in \si{\cubic\centi\metre\per\second}, it
satisfies $|\lambda_0| = 2.286\,K_{\rm ion}(1s)\,\nel$ across the grid
(project working note; Table~\ref{tab:ladder}), so $\tauslow$ is the collisional--radiative
ionisation lifetime of the ground-state reservoir and not the time for anything
to come to equilibrium.

\subsection{The gates do not stop anything}

No gate raises on failure. Each prints \texttt{PASS} or \texttt{FAIL} and
returns a flag; \texttt{main} sets \texttt{all\_pass = False}, prints a warning,
and exits normally. \textbf{The script returns exit code zero whether or not
Gate~D fails.} The residual gates elsewhere in the project do raise; this one
does not, and that asymmetry is a defect.

%-----------------------------------------------------------------------------
\section{The checks that could have failed}
\label{sec:severe_checks}
% QUESTION: which checks in this project would a plausible wrong model fail?

Everything so far has been a demotion. This section is the other half, and it
is short by design: a small number of severe checks is worth more than a long
list of consistency tests, and these are the ones the thesis leans on.

\subsection{The superposition residual}
\label{sec:superposition}

The entire mechanism of Chapter~\ref{ch:results} rests on one structural claim:
that the excited populations $\nF$ of Section~\ref{sec:qss} split exactly into a
ground-fed and a recombination-fed part,
$\nF = \bm{a}\,\ngs + \bm{c}\,\nion$, where the network responses $\bm{a}$ and
$\bm{c}$ of Section~\ref{sec:R_depends_on_b} are dimensionless and independent
of the two reservoir densities. If that split is not exact, the ground-fed
fractions $\ffrac{p}$ of Section~\ref{sec:ground_fed_fraction} are not well
defined and the closed form for the diagnostic error does not exist.

The check solves the full system three times at each grid point and each step
direction, once with each reservoir alone and once with both, and measures the
relative residual of the sum against the joint solve. Across the $784$
grid-point and step-direction pairs it covers, out of the $800$ the grid admits,
the largest residual is $3.075\times10^{-14}$
(\texttt{plateau\_gridmap.txt}, line 19).

\textbf{Grade: severe}, for a specific reason. Feeding the check a corrupted
input in which the $n=3$ and $n=4$ entries are swapped, which is the mistake
that would most easily go unnoticed in a two-shell observable, trips it
(project working note; Table~\ref{tab:ladder}). It is one of only two checks in the
project reported to catch that swap.

\subsection{The reduced-versus-full ratio test}

The second is its companion. The observable $\Robs = n_3/n_4$ can be computed
two ways: from the full $43$-state steady-state solve, and from the reduced
two-channel form of Eq.~\eqref{eq:R_of_b} with the ground-fed fractions
substituted in. The two routes share no code path beyond the matrix itself. At
the benchmark point $\Te = \benchTe$, $\nel = \benchne$, an independent
re-implementation obtained $\Robs = 0.777768$ by three routes against
Chapter~\ref{ch:theory}'s recorded $0.7778$
(project working note; Table~\ref{tab:ladder}). A $3\leftrightarrow4$ swap breaks the
agreement.

\subsection{Pinned hashes on the inputs}

The most common way for a computation like this to become wrong is not an
algebra error. It is a stale array: a matrix regenerated after a correction,
combined with an analysis output produced before it. This project has three
retractions in its history and at least one of them had that shape.

\texttt{src/validation/preflight.py} pins the SHA-256 hashes, fixed-length
fingerprints that change if any byte of a file changes, of the two canonical
inputs and compares them byte for byte before any verification script
runs. Recomputed for this chapter:
\begin{center}
\begin{tabular}{ll}
  \toprule
  file & SHA-256 (first 16 hex digits) \\
  \midrule
  \texttt{data/processed/cr\_matrix/L\_grid.npy} & \texttt{2d92b58e1693107d} \\
  \texttt{data/processed/cr\_matrix/S\_grid.npy} & \texttt{7822f536590c76ba} \\
  \bottomrule
\end{tabular}
\end{center}
Both match the pinned baselines at \texttt{preflight.py}:37--38 exactly, so the
matrix on disk at the time of writing is the one the pins were taken from; the
statement is only as good as the discipline of running the preflight before
each analysis.

\textbf{Grade: severe on the data, incomplete on the scripts.} Of the six
scripts the same preflight claims to check, \texttt{verify\_ch3\_groupB.py} and
\texttt{verify\_grid\_coverage.py} do not exist on disk under those names; a
missing script is downgraded to a non-blocking advisory
(\texttt{preflight.py}:258--262) while the summary still prints ``All checks
passed''. The hash check is sound; the sentence after it overstates what ran.

\subsection{Guards that recompute rather than trust}

\texttt{src/analysis/make\_ch5\_figures.py} does not plot the stored results. It
rebuilds every plotted quantity from $\Lmat$ and the recombination source
$\Svec$ of Eq.~\eqref{eq:source} and compares against
the stored comma-separated tables at a relative tolerance of $10^{-9}$, raising
with the offending grid point, column, and matrix hash if any row disagrees
(lines 393--432). It reads the temperature step size out of the file header
rather than redeclaring it, so a figure cannot be drawn against a step
convention different from the one that produced the data. It also verifies that
the untrapped run of the radiation-trapping sweep reproduces the canonical
matrix exactly, without which no trapped number in Chapter~\ref{ch:limits} would
be interpretable as a difference.

The same script carries a negative control: it computes the correlation between
$\Msep$ and the diagnostic error both raw and after controlling for temperature
and density, and refuses to draw the figure if the controlled correlation fails
to change sign (Section~\ref{sec:M_does_not_predict} draws the conclusion).

\subsection{Two checks that fired, and one that was withdrawn}

Three episodes say more than the pass/fail table, because in each the check did
something.

The Lyman-$\alpha$ line-centre cross-section gate, an independent implementation
in \texttt{escape\_factor.py} against the project's own Einstein coefficients,
fired at \SI{1156}{\percent} on its first run, located a $4\pi$ factor, and now
agrees to \SI{0.059}{\percent} against a raising \SI{1}{\percent} tolerance
(Section~\ref{sec:opacity}).

\texttt{verify\_fujimoto\_table41.py} carries a deliberately broken variant with
proton-impact $\ell$-mixing removed. It returns a recombination-fed population
coefficient $r_0(2) = 104.5$, and the script flags it itself as unphysical on
the ground that a level cannot exceed Saha equilibrium. Production gives
$0.7392$. A check that fires on a known-bad input is the definition of a severe
one.

Third, a semigroup propagator identity,
$\exp(\Lmat t_1)\exp(\Lmat t_2) = \exp(\Lmat(t_1+t_2))$ for a matrix exponential
computed by the same routine, returned $0.000\times10^{0}$ by construction and
was \emph{withdrawn by the author} (Section~\ref{sec:bridge}), the judgement
this chapter applies to Gate~A, applied earlier and unprompted.

\subsection{The bound that is a theorem}

Section~\ref{sec:bound} derives an upper bound on the diagnostic error of the
form $\tanh(\Delta/4)$, where $\Delta$ is the separation between the two shells'
switching points. A script reports that the bound is honoured at all $248$
points tested, and that line reads as validation.

It is not. Here $\Delta$ is the dimensionless separation between the two shells'
switching points, in units of the logarithm of the reservoir strength. Given
$\bm{a},\bm{c}\ge0$, which Section~\ref{sec:R_depends_on_b} proves from the
M-matrix structure of $-\Lmat_{FF}$, the bound is a theorem: it cannot fail
except on an arithmetic bug. Fed corrupted input it still passes. Swapping
shells $3\leftrightarrow4$ sends $\Delta\to-\Delta$ and leaves both $|\Delta|$
and the logarithmic change in $\Robs$ that it bounds unchanged, so the gate
cannot see it. Swapping $c_3\leftrightarrow c_4$ also
passes, but for a different reason: it moves $|\Delta|$ from $1.945614$ to
$0.939493$, a factor $2.07$, and the gate still passes because it compares
$|f_3-f_4|$ against $\tanh(|\Delta|/4)$ evaluated from the same corrupted
$\bm{c}$. The bound is satisfied by whatever $\bm{a},\bm{c}\ge0$ it is
handed.
\textbf{Grade: tautological.} The script's own docstring calls it a wiring
check; the line it prints does not, and the printed line is what a reader sees.

The bound is a result, and a useful one, because it holds for any atomic dataset
whatever. It is not evidence that this atomic dataset is right.

%-----------------------------------------------------------------------------
\section{Conditioning, precision, and state-space convergence}
\label{sec:conditioning}
% QUESTION: could the answers be artifacts of finite-precision arithmetic or of
% where the level ladder was cut off?

Two numerical questions sit underneath every result in this thesis. The
elimination of Section~\ref{sec:qss} inverts a $42\times42$ matrix whose entries
span many orders of magnitude, which invites the question of whether the inverse
means anything. And the level ladder was cut at $n_{\max}=15$, which invites the
question of what the discarded levels would have changed.

\subsection{The inverse is well behaved, and here is how well}

The relevant quantity is the condition number of $\Lmat_{FF}$, the block of
$\Lmat$ coupling the $42$ excited levels to each other
(Section~\ref{sec:qss}), and it is the factor by which a relative perturbation
of the matrix can be amplified in the solution.
Computed from the canonical $\Lmat$ at all $400$ grid points by
\texttt{src/validation/verify\_operator\_conditioning.py} and stamped to
\texttt{validation/operator\_conditioning/operator\_conditioning.csv}: minimum
$1.4821\times10^{3}$ at $[0,0]$ ($\Te = \SI{1.000}{\electronvolt}$), maximum
$1.7436\times10^{5}$ at $[33,7]$ ($\Te = \SI{4.715}{\electronvolt}$), median
$2.1044\times10^{4}$. In double precision, which carries about sixteen
significant digits, a condition number of $10^{5}$ costs five of them. Eleven
remain, and every conclusion in this thesis turns on the first three.

The range was first asserted in Chapter~\ref{ch:theory} on an unscripted
working note, which is not a check; the stamped script reproduces it to
\SI{0.14}{\percent} and \SI{0.21}{\percent}. The number was right the whole
time and was not evidence until it could be re-run.

Three further measurements bound the arithmetic. The eigenvalue condition
number is a different quantity from the matrix condition number above: it is
$1/|\bm{y}^H\bm{x}|$ for a left and right eigenvector pair, and it measures how
far a single eigenvalue moves under a perturbation of the matrix. For the two
eigenvalues the thesis uses it is of order unity
($1.71$ and $1.98$ at the benchmark point, $3.26$ and $2.59$ at the cold
corner), despite a largest absolute column sum of
$2.0\times10^{14}\,\si{\per\second}$; that large norm is a
statement about the fast radiative modes and does not propagate into
$\lambda_0$ or $\lambda_1$. Recomputing the two supply channels in single
precision moves them by $1.3\times10^{-5}$, so double precision is nowhere near
a cliff. And tightening the integrator changes nothing that matters: a relative
tolerance of $10^{-9}$ against $10^{-12}$ shifts the result by
$1.20\times10^{-8}$, and computing the propagator by a dense matrix exponential
against a Krylov action differs by $1.13\times10^{-10}$.

The eigenvalue condition numbers and the single-precision test cover three of
$400$ points; the tolerance comparisons cover two. The condition numbers and the
displacement are recomputed under the definition $1/|\bm{y}^H\bm{x}|$ with
unit-norm left and right eigenvectors by \texttt{verify\_residual\_claims.py}
(\texttt{validation/residual\_claims/}). They are point measurements, not grid
scans, and are quoted as such.

\subsection{Where the ladder was cut}
\label{sec:convergence}

The quantity that has to converge is not a population but the difference in
ground-fed fraction between the two Balmer shells, $\ffrac{3}-\ffrac{4}$, since
Chapter~\ref{ch:results} shows the diagnostic error to be proportional to it.

Rebuilding $\Lmat$ with successively lower $n_{\max}$, from $15$ down to $9$,
with the loss channels of every removed state returned to the diagonal of the
survivors, and recomputing $\ffrac{3}-\ffrac{4}$ at each truncation's own
$u^{\mathrm{CRE}}$ (\texttt{verify\_nmax\_downward\_scan.py},
\texttt{validation/nmax\_downward\_scan/}) gives increments of one sign at the
three named points, with the last measured ratio $d(15)/d(14)$ equal to
$0.88$ at the benchmark, $0.72$ at the cold corner $[0,0]$ and $0.33$ at the
ridge $[15,3]$; at the $245$ of the $280$ defended points where the increments
are one-signed and geometric the ratio has median $0.67$. Summing the geometric
tail with that ratio gives an extrapolated truncation error of
\SI{-0.46}{\percent} at the benchmark, \SI{+1.6}{\percent} at the cold corner
and \SI{+0.005}{\percent} at the ridge; holding $u$ at its $n_{\max}=15$ value
instead gives \SI{-0.38}{\percent}, \SI{+1.4}{\percent} and
\SI{+0.005}{\percent}, so the sign does not depend on that choice. A historical
project value with no producing script, not used as evidence here and superseded
by \texttt{verify\_nmax\_downward\_scan.py}, records
\SI{+0.9}{\percent}, \SI{-1.4}{\percent} and \SI{-0.55}{\percent} with a
common ratio of about $0.75$; the stamped run reverses both signs and the
conclusion that rested on them. Truncation makes
$\ffrac{3}-\ffrac{4}$ slightly too large at the benchmark and too small at the
cold corner, so it is not a bias that can be corrected by a single factor.
Over the defended points where the tail can be summed the extrapolated error
has median \SI{+0.002}{\percent} and runs from \SI{-0.82}{\percent} to
\SI{+1.03}{\percent}, smaller than the spread between the two step conventions
of Section~\ref{sec:step_convention} at every such point. At the $35$ defended
points of the density column $\nel = \SI{5.2e13}{\per\cubic\centi\metre}$ the
increments change sign among $n_{\max} = 12$ to $15$ and no tail can be
summed; the last measured increment there is below \SI{0.041}{\percent} of
$\ffrac{3}-\ffrac{4}$ at every point.

\paragraph{A trap in this test.} The first attempt dropped the top shells from
the state vector without removing their loss channels from the diagonal of the
surviving states, which then lost population to states that no longer existed,
and $\ffrac{3}-\ffrac{4}$ jumped. The stamped scan repeats the mistake
deliberately: dropping $n=15$ alone without the correction moves
$\ffrac{3}-\ffrac{4}$ by \SI{11.5}{\percent} at the benchmark,
\SI{12.6}{\percent} at the cold corner and \SI{5.0}{\percent} at the ridge,
where the correct truncation to $n_{\max}=14$ moves it by \SI{0.06}{\percent},
\SI{0.64}{\percent} and \SI{0.01}{\percent}; dropping six shells without it
gives \SI{17.9}{\percent}, \SI{25.9}{\percent} and \SI{13.3}{\percent} (a
working note recorded \SIrange{22}{37}{\percent}). The jump is an artifact of
the truncation procedure, not a dependence on $n_{\max}$. Removing state $k$
requires adding $\sum_{k\ \rm dropped} L_{ki}$ back to $L_{ii}$ for every
surviving $i$, as the production matrix does. Anyone repeating this test will
reproduce the spurious number first.

\paragraph{The terminal shell is over-populated, and by a measured amount.}
Truncating the ladder at $n_{\max}$ removes, from the balance of the terminal
shell, the collisional excitation out of it to $n > n_{\max}$ and the
de-excitation and cascade into it from those levels (cascade is downward, so
what is lost is a way out, not something ``to cascade out to''). The loss removed is not small. Continuing the $12\to13$, $13\to14$
and $14\to15$ excitation coefficients geometrically, the missing $15\to16$
channel would carry $1.30$ times the whole retained loss out of $n=15$ at the
benchmark and $1.20$ times at $[0,4]$, between $1.20$ and $1.35$ times at
every grid point, while radiative decay is below $10^{-6}$ of that loss and
ionisation is $7$ to \SI{10}{\percent} of it
(\texttt{validation/terminal\_shell\_budget/}). The top shell sits within
$6\times10^{-5}$ of Saha--Boltzmann, so for the total population the removed
exchange is nearly balanced and the sign of the bias is decided channel by
channel: the ground-fed channel, whose net flux through the top of the ladder
is upward, loses only an exit and can only be biased high, and it is the
ground-fed coefficient $r_1(15)$ that the external comparison finds high while
$r_0(15)$ agrees (Section~\ref{sec:fujimoto}). Comparing each shell's ground-fed coefficient $a_p$ as terminal
shell (operator rebuilt with $n_{\max}=p$, the loss channels into the removed states kept on the diagonal) with its
value in the full $n_{\max}=15$ operator gives at the Fujimoto point $[49,0]$ (\SI{10}{\electronvolt},
\SI{e12}{\per\cubic\centi\metre}) excesses of $9.4$ at $p=8$, $10.7$ at $p=9$ and $6.3$ at $p=10$, then $4.4$,
$3.0$, $2.1$ and $1.4$ at $p=11$ to $14$ (\texttt{verify\_terminal\_vs\_interior.py},
\texttt{validation/terminal\_vs\_interior/}). At the benchmark $p=8$ to $10$ give $14.2$, $12.3$ and $6.5$; over
the $280$ defended points the $p=8$ excess spans $8.2$ to $17.7$ and the $p=10$ excess $5.4$ to $7.5$. The excess
is confined to the ground-fed channel: the shell's total population and recombination-fed coefficient move under
\SI{1}{\percent} at the benchmark and at most \SI{32}{\percent} and \SI{30}{\percent} at the low-density named
points. A historical project value with no producing script, not used as evidence here and superseded by
\texttt{verify\_terminal\_vs\_interior.py}, records $9.4$, $10.7$ and
$6.3$ with neither quantity nor grid point named; the stamped run reproduces them to \SI{0.4}{\percent} as $a_p$
at $[49,0]$. The absence of a monotone trend over $p=8$ to $10$ holds there and at $47$ of the $280$ points;
from $p=9$ upward the excess falls at every grid point, at the benchmark from $p=8$. The external excess at the
terminal shell, $4.9$ to $6.2$ against published values (Section~\ref{sec:fujimoto}), lies outside the $6.3$ to
$10.7$ band, not inside it as the working note states, but within a factor $2.2$ of every value in it, and both rise
with density ($4.9$ to $6.2$ externally, $9.4$ to $17.7$ at $p=8$ internally at \SI{10}{\electronvolt}). That
agreement in magnitude and in the sign of the density trend makes truncation a sufficient explanation for the
$n=15$ discrepancy without invoking the atomic data. The fall with $p$ cannot be extrapolated to $p=15$: the
reference for $p=14$ is itself a ladder with one shell above it, so the internal measurement neither predicts nor
excludes $4.9$ to $6.2$ there.
No result in this thesis rests on the
top shell: the observable is built from $n=3$ and $n=4$, both fully resolved in
orbital angular momentum $\ell$ and eleven shells below the cut at
$n_{\max}=15$.

\paragraph{What has now been checked.} Above $n=8$ the shells are bundled and
their internal $\ell$-distribution is assumed statistical rather than computed.
The script that tests that assumption, \texttt{verify\_bundling\_psm20.py}, had
never been executed; two of the three defects found on its first run, the
synthesised temperature grid and the identically zero $\ell$-mixing array for
the bundled indices, were predicted here. Section~\ref{sec:bundling_check}
reports all three defects and the result.

%-----------------------------------------------------------------------------
\section{Is the shell ratio the line ratio?}
\label{sec:balmer_equivalence}
% QUESTION: the analysis computes n_3/n_4, but a spectrometer measures
% Halpha/Hbeta. What does the substitution cost?

Every analytic result in Chapter~\ref{ch:theory} is written for the shell ratio
$\Robs = n_3/n_4$, because the two-channel decomposition applies to shell
populations. A spectrometer does not measure shell populations. It measures the
$H_\alpha$ and $H_\beta$ line intensities, each of which is a sum over the
electric-dipole-allowed sublevel transitions, weighted by their Einstein
coefficients. Section~\ref{sec:radiative} established that these are not
proportional: $4f$ holds \SI{43.75}{\percent} of the $n=4$ shell by statistical
weight and contributes nothing to $H_\beta$, since $4f\to2p$ would change $\ell$
by two and is forbidden.

The substitution therefore has to be measured, not assumed.

The line intensities are built as $A$-weighted sums over resolved channels,
$H_\alpha$ from $3s\to2p$, $3p\to2s$ and $3d\to2p$, and $H_\beta$ from
$4s\to2p$, $4p\to2s$ and $4d\to2p$, with populations taken from the solve.
Nothing statistical is assumed anywhere in that construction: searching the two
producing scripts for statistical-weight expressions returns no matches, and the
$A$-value lookup raises unless exactly one row matches each transition.

Two measurements follow. The $4f$ fraction of the $n=4$ shell runs from
$0.4361$ to $0.4375$ across the grid, against the statistical value
$14/32 = 0.4375$; proton-impact $\ell$-mixing drives the shell to a statistical
distribution to better than \SI{0.3}{\percent} everywhere, with the departure
concentrated at $\nel = \SI{e12}{\per\cubic\centi\metre}$ where $\ell$-mixing
weakens and radiative decay competes. That fraction is a property of the full
equilibrium population; the quantity that sets the line-versus-shell difference
below is the $\ell$-distribution within each supply channel separately, which
is not statistical at the low-density edge. And the ratio of the transient error
computed on the line ratio to the same error computed on the shell ratio runs
from $0.9482$ to $0.9999$ over the $784$ point-and-direction pairs, below $1$ at
every one: \textbf{worst case \SI{5.2}{\percent}, at the heating step from
$9.54$ to \SI{10.0}{\electronvolt} at the lowest density}
(\texttt{verify\_weighted\_census.py}, stamped to
\texttt{validation/weighted\_census/}). At
$[0,4]$, the point carrying the largest error anywhere, the two agree to
\SI{0.06}{\percent} ($0.386683$ against $0.386903$). Because the ratio is below
$1$ everywhere, the shell census of Chapter~\ref{ch:results} is conservative
for the line observable: counted on the line ratio, $44$ instead of $45$ of the
$448$ pairs above \SI{2}{\electronvolt} exceed \SI{10}{\percent} at
\SI{100}{\micro\second}, with the worst case $0.1746$ instead of $0.1748$ at
the same point. The \SI{5.2}{\percent} scales inversely with the proton
$\ell$-mixing rate: with the within-shell $\ell$-changing rates scaled by a
common factor and the column sums of the operator preserved, the corner
departure is \SI{10.0}{\percent} at half the rate, \SI{2.6}{\percent} at double
and \SI{1.1}{\percent} at five times, while the benchmark line-to-shell ratio
stays between $0.9987$ and $0.9999$ (\texttt{validation/weighted\_census/},
$\ell$-mixing sensitivity block and \texttt{weighted\_census\_lmix.csv}); and
it sits on the grid boundary.

That \SI{0.06}{\percent} is a single-point figure and must not be quoted as a
grid-wide one. The grid-wide statement is the \SI{5.2}{\percent}.

\subsection{And is the plasma transparent to its own Balmer light?}

The shell-to-line substitution assumes the emitted photons escape. For the
Lyman series that assumption fails at the cold edge and
Chapter~\ref{ch:limits} treats it at length. For the Balmer lines themselves it
closes here, because it bears on the observable rather than on the matrix.

Section~\ref{sec:opacity} defines the optical depth $\tau$ and the escape factor
and treats the Lyman series, where transparency fails. For the Balmer lines the
absorbers are the $n=2$ atoms, $n_{2s}+n_{2p}$ in collisional--radiative
equilibrium with $\nion=\nel$, the Doppler width is taken at $\Tn=\Te$, and the
multiplet oscillator strengths are $f(2\to3)=0.641$ and
$f(2\to4)=0.119$~\cite{WieseFuhr2009}. The line-centre optical depth of
$H_\alpha$ across a \SI{10}{\centi\metre} slab is then $4.9\times10^{-7}$ at the
cold corner, $8.2\times10^{-4}$ at $[0,4]$ and $0.214$ at $[0,7]$, the single
worst cell on the grid; resolving the absorbers into $2s$ and $2p$ with the
model's own $A$ coefficients changes these by at most \SI{4}{\percent}. At
$[0,7]$ the population escape factor for a photon born at the slab centre
($\kappa_0 D/2 = 0.107$) is $0.928$ for $H_\alpha$ and $0.990$ for $H_\beta$, so
the observed ratio is lowered by \SI{6.3}{\percent} in that one cell; over the
full path ($\tau=0.214$) the figures are $0.861$ and \SI{12}{\percent}. Above
\SI{2}{\electronvolt} the largest $H_\alpha$ optical depth is $0.010$ at
$[15,7]$ and the ratio moves by less than \SI{0.7}{\percent} anywhere
(\texttt{verify\_balmer\_optical\_depth.py},
\texttt{validation/balmer\_optical\_depth/}). The Balmer lines are optically
thin over the region where this thesis makes quantitative claims. A historical
project value with no producing script, not used as evidence here and superseded
by \texttt{verify\_balmer\_optical\_depth.py}, records
$6.3\times10^{-7}$, $1.0\times10^{-3}$ and $0.263$ for the same three cells;
those values are $1.23$ to $1.28$ times the stamped ones and no convention
tested recovers them, and the escape factor above $0.85$ and ratio shift of at
most about \SI{10}{\percent} recorded with them are half-slab figures attached
to a full-slab optical depth. The conclusion holds with either set of values.

%-----------------------------------------------------------------------------
\section{Two tests of the reduction, and their expected outcomes}
\label{sec:attacks}
% QUESTION: the two objections most likely to destroy the result were tested;
% what was the falsifier in each case, and did it appear?

The checks so far ask whether the model computes what it intends. This section
asks something harder: whether two specific structural choices manufacture the
result. In each case the size of effect that would have destroyed the claim was
written down before the test was run, which is what makes the outcome
informative rather than reassuring.

\subsection{Does the open boundary manufacture the slow clock?}

$\Lmat$ describes an open system: atoms that ionise leave the accounting, the
recombination feed $\Svec\nion$ enters as a constant source rather than a matrix
element, and there is no forty-fourth state for ions to accumulate in. Since the
persistence claim of Section~\ref{sec:persistence} rests entirely on $\tauslow$
exceeding the duration of the disturbance, the first objection is that
$\tauslow$ is an artifact of the open boundary.

\textbf{The falsifier, named first.} Write $\taudrive$ for the duration of the
disturbance, fixed at \SI{100}{\micro\second} throughout this thesis as a
representative edge-localised-mode (ELM) time
(Section~\ref{sec:persistence}). At $[0,4]$, the worst point on the grid,
$\tauslow/\taudrive \approx 2300$. If closing the system shortened $\tauslow$ by
that factor, the slow clock would fall to $\taudrive$ itself and the persistence
would vanish. A factor of $2300$ refutes the claim.

Two grid labels appear in the table for the same worst case: the error is
evaluated at $[0,4]$, but the operator governing the decay after the step is the
post-step one, $\Lmat[1,4]$, and every timescale in this thesis names the
operator it belongs to (Chapter~\ref{ch:results}).

The manifold was closed by appending an ion state,
\begin{equation}
  \Lmat_c =
  \begin{pmatrix}
    \Lmat & \Svec \\
    -\operatorname{colsum}(\Lmat) & -\textstyle\sum\Svec
  \end{pmatrix},
  \label{eq:closure_ch4}
\end{equation}
so that every ionising atom lands in the new state instead of being deleted.
Each column of $\Lmat_c$ sums to zero to machine precision, and the construction
reproduces the closed-system value at the cold corner recorded in this
project's working notes, \SI{1.6226}{\second}, to five digits, which confirms it is the same
closure rather than a new one. The measured factors are given in
Table~\ref{tab:closure}.

\begin{table}[tbp]
  \centering
  \caption[Effect of closing the ion boundary on the slow timescale]{The slow
  timescale $\tauslow$ computed with the ion treated as an external reservoir
  (open) and with a forty-fourth ion state appended (closed), at four grid
  points. Over the whole grid the ratio has minimum $1.00$, median $1.00$ and
  maximum $41.4$: the boundary condition matters only where $\tauslow$ already
  exceeds the disturbance duration by three orders of magnitude. Every entry is
  reduced from \texttt{validation/ion\_closure/ion\_closure\_summary.csv}
  (\texttt{verify\_ion\_closure.py}) by
  \texttt{verify\_headline\_provenance.py}; the open column is the float64
  eigensolver, which at the cold corner differs from the artifact's 40-digit
  path by one part in $10^{4}$.}
  \label{tab:closure}
  \begin{tabular}{lrrr}
    \toprule
    point & $\tauslow$ open & $\tauslow$ closed & factor \\
    \midrule
    cold corner $[0,0]$ & \SI{67.23}{\second} & \SI{1.6226}{\second} & $41.4$ \\
    post-step operator $\Lmat[1,4]$ & \SI{0.23324}{\second} & \SI{0.015173}{\second} & $15.4$ \\
    ridge $[15,3]$ & \SI{2.027}{\milli\second} & \SI{1.999}{\milli\second} & $1.01$ \\
    benchmark $[23,5]$ & \SI{22.728}{\micro\second} & \SI{22.708}{\micro\second} & $1.00$ \\
    \bottomrule
  \end{tabular}
\end{table}

The shape of the distribution carries the argument: the factor is large only
where $\tauslow$ already exceeds $\taudrive$ by about a thousand, so the error
has already saturated, and it is unity everywhere $\tauslow$ approaches
$\taudrive$. Recomputing the entire persistence census with closed-system
$\tauslow$ moves the breakdown count from $202$ to $200$ of $680$ and the worst
case from $0.38682$ to $0.38563$, about one part in $325$; restricted to
$\Te \ge \SI{2}{\electronvolt}$, where Chapter~\ref{ch:results} makes
quantitative claims, the count is $45$ before and $45$ after.

The falsifier did not appear. A factor of $2300$ was needed; the measured factor
at that operator is $15$.

\textbf{The caveat, which must not be softened.} The closure adds an ion
reservoir; it does not close the electron balance, since $\nel$ is held fixed
while the ion population moves, and it adds no transport. It answers the
boundary condition on the atomic manifold and nothing else:
Section~\ref{sec:transport} states the objection it cannot answer, and
Section~\ref{sec:quasineutrality} the one it makes sharper rather than softer.

\subsection{Is the single-slow-mode estimate a bound on the time-averaged error?}

Chapter~\ref{ch:results} evaluates the error at the partial-equilibrium plateau
and builds from it the single-slow-mode estimate of Eq.~\eqref{eq:elm_bound},
which multiplies the plateau error $\epsplat$ of
Section~\ref{sec:eps_definitions} by the fraction of $\taudrive$ over which an
exponential decay on $\tauslow$ keeps it alive. It was first written as a lower
bound, on the assumption that the observable cannot decay faster than the
slowest mode. That assumption is testable, and an estimate quoted as though it
were the answer is a claim about accuracy that has to be measured.

Propagating the full $43$-state system through the step by direct
eigen-decomposition and integrating the error over \SIrange{0}{100}{\micro\second}
gives the comparison in Table~\ref{tab:lowerbound}.

\begin{table}[tbp]
  \centering
  \caption[Accuracy of the single-slow-mode estimate]{The plateau error, the
  true error averaged over \SIrange{0}{100}{\micro\second} from direct
  eigen-propagation of the full $43$-state system, and the single-slow-mode
  estimate of Eq.~\eqref{eq:elm_bound}, at five grid points. The averages are
  integrated at $400$ time samples per decade; at $1600$ per decade the five
  change by at most $4.7\times10^{-6}$ relative, and the $[0,4]$ average by
  $2.7\times10^{-10}$. Integration on a uniform \SI{25}{\nano\second} mesh, which
  under-resolves the nanosecond rise, reads low by $6\times10^{-5}$ to
  $5\times10^{-4}$ relative. The estimate reproduces
  the true average to better than \SI{0.8}{\percent} at the five points. The
  value below unity at $[0,4]$ is real: a deficit of $2.2\times10^{-5}$ that
  survives the denser grid. Source: \texttt{verify\_trajectory\_census.py},
  \texttt{validation/trajectory\_census/}.}
  \label{tab:lowerbound}
  \begin{tabular}{lrrrr}
    \toprule
    point & $\epsplat$ & true \SI{100}{\micro\second} average & estimate, Eq.~\eqref{eq:elm_bound} & true/estimate \\
    \midrule
    $[0,4]$ worst case & $0.386903$ & $0.386811$ & $0.386820$ & $0.99998$ \\
    $[15,3]$ ridge     & $0.180728$ & $0.175283$ & $0.174812$ & $1.00270$ \\
    $[15,5]$           & $0.103550$ & $0.072185$ & $0.071722$ & $1.00646$ \\
    $[23,5]$ benchmark & $0.063612$ & $0.011796$ & $0.011712$ & $1.00722$ \\
    $[0,0]$            & $0.083004$ & $0.083084$ & $0.083004$ & $1.00096$ \\
    \bottomrule
  \end{tabular}
\end{table}

At the five points tested the estimate sits within \SI{0.8}{\percent} of the
directly propagated average, and at $[0,4]$ it sits above it, so it is an
estimate and not a bound. The accuracy at these five points does not
generalise: over the $448$ pairs above \SI{2}{\electronvolt} the ratio of true
average to estimate departs from unity by up to \SI{2.2}{\percent} at
\SI{100}{\micro\second} and \SI{3.8}{\percent} at \SI{506}{\micro\second}
(Section~\ref{sec:persistence}). What matters here is that the assumption
underneath it was tested rather than asserted.

The reason is structural, and it is a property the shell ratio has and another
observable in this same project does not. The spectrum of $\Lmat$ has one slow
eigenvalue and then a band, with nothing between: at $\Lmat[1,4]$ the three
slowest eigenvalues are $-4.29$, $-1.60\times10^{8}$ and
$-3.44\times10^{8}\,\si{\per\second}$, a gap of $3.7\times10^{7}$. An observable
that depends on the excited populations only through their ratio cannot decay
faster than $\tauslow$: once the manifold is quasi-steady the excited vector is
a function of $u$ alone, so
$\ln\Robs - \ln\Robs^{\mathrm{CRE}} = (S/u^{*})\,\delta u + O(\delta u^{2})$
with $\delta u \propto \mathrm{e}^{-t/\tauslow}$, and the leading term
vanishes only where $S = \ffrac{3}-\ffrac{4} = 0$. It does not: $S$ is positive
at all $400$ grid points, with minimum $0.046$ at $[44,7]$, and the slow
eigenvector's projection onto $\ln\Robs$ equals $S/u^{*}$ to
$1.2\times10^{-4}$ at the benchmark and $1.6\times10^{-8}$ at $[0,4]$, the
deviation being the first-order quasi-steady correction
$\lambda_{\mathrm{slow}}/\lambda_{2}$ itself. It exceeds \SI{1}{\percent} only
at the two hot, dense points $[48,7]$ and $[49,7]$, where the separation is
$87$-fold and which the window criterion excludes from the analysed set
(\texttt{validation/slow\_mode\_projection/}). That decay on an intermediate
mode was the failure mode being hunted, and it does occur for the
ground-normalised error measure, which does not cancel the reservoir and whose
$1/\mathrm{e}$ time is \SI{2.58}{\second} against
$\tauslow = \SI{67.2}{\second}$ at the cold corner $[0,0]$. The estimate is
accurate for the ratio and not for that measure; why it is nevertheless not a
bound, the denominator relaxing on the same $\tauslow$, is set out in
Section~\ref{sec:persistence}.

%-----------------------------------------------------------------------------
\section{External comparison: Fujimoto's population coefficients}
\label{sec:fujimoto}
% QUESTION: does an independently computed collisional-radiative model, built
% decades earlier from different atomic data, agree with this one?

Everything to this point compares the model against itself or against
arithmetic. Only comparison against an independent calculation reaches the
atomic physics, and this project has one such comparison in depth: against
Fujimoto's collisional--radiative treatment of atomic
hydrogen~\cite{Fujimoto2004}, computed decades earlier, with different atomic
data, by a different method.

\subsection{Timescales, at a point that lies exactly on this grid}

Fujimoto works a numerical example at $\Te = \SI{10}{\electronvolt}$ and
$\nel = \SI{e18}{\per\cubic\metre}$, which is
$\SI{e12}{\per\cubic\centi\metre}$ and lands exactly on grid point $[49,0]$. No
interpolation is involved.

\begin{center}
\begin{tabular}{lrrr}
  \toprule
  quantity & Fujimoto App.~4B & this model, $[49,0]$ & ratio \\
  \midrule
  excited-state response & ${\sim}\SI{1e-7}{\second}$ & \SI{3.4057e-8}{\second} & $0.34$ \\
  ground-state depletion & ${\sim}\SI{1e-4}{\second}$ & \SI{1.7657e-4}{\second} & $1.77$ \\
  \bottomrule
\end{tabular}
\end{center}

The two-clock structure appears in both, separated by three orders of
magnitude in both, with the fast time short and the slow time long in the same
sense.

\textbf{What this comparison is not.} Fujimoto's response time is
$\max_p t_{\rm tr}(p)$, the \SI{63}{\percent} rise time of the slowest-responding
level, which is not $1/|\lambda_1|$. The two are different functionals of the
same spectrum, so agreement to a factor of a few is the meaningful comparison
and exact agreement would be suspicious. Quoting these ratios as though they
were a precision test would overstate them.

One qualification on the independence of this check, found on reading the
printed appendix rather than a summary of it. Fujimoto's Figs.~4B.1 and 4B.2 are
captioned as quoted from Sawada and Fujimoto (1994)
\cite{SawadaFujimoto1994}, so the appendix and the paper
Chapter~\ref{ch:conclusions} discusses are the same calculation, not two.
\textbf{The comparison is therefore independent of this thesis but not two
independent checks}, and it is counted once. The appendix also states the same
construction the paper does, that $n(1)$ is held constant, which matters for
Chapter~\ref{ch:conclusions} and is taken up there.

A second, independent route corroborates the same physics. Fujimoto's boundary
level, the principal quantum number above which collisions dominate radiative
decay, sits at $p_G \approx 4$--$5$ at this density. Tracing the boundary
descent in this model puts it at $5g$ at the low-density end of the grid, which
is the crossing Chapter~\ref{ch:atoms} deferred to this chapter
(Section~\ref{sec:radiative}). Two independently computed models place the same
structural feature within one shell of each other, at one density.

\subsection{The recombining coefficients agree}

Fujimoto's Table~4.1 tabulates the two population coefficients of the
decomposition that Section~\ref{sec:R_depends_on_b} derives: $r_0(p)$, how
strongly level $p$ is fed from the ion continuum, and $r_1(p)$, how strongly it
is fed from the ground state. Both are dimensionless, and both are normalised to
the Saha--Boltzmann population of level $p$, which is why $r_0$ is of order unity
while $r_1$ at low density is of order $10^{-5}$, and why a value of $r_0$ above
one is unphysical.

The recombining coefficients agree to \SIrange{0.1}{1}{\percent} for $n \ge 4$.
The discriminating case is $r_0(2)$, where cascade and $\ell$-mixing structure
matter most and a hydrogenic approximation would go wrong: this model gives
$0.7392$ against the tabulated $0.730$, a difference of \SI{1.3}{\percent}, a
test the model could have failed and did not (the broken variant of
Section~\ref{sec:severe_checks} drives the same quantity to $104.5$).

\subsection{The ionising coefficients at low $p$ disagree, and the target is
not verified}

At $\nel = \SI{e12}{\per\cubic\centi\metre}$ this model's $r_1$ sits a factor
$8.3$ low at $p=3$ and $4.4$ low at $p=4$ against Fujimoto's Table~4.1(b).
Those are precisely the two shells whose difference is the mechanism of
Chapter~\ref{ch:results}, so the discrepancy cannot be dismissed as peripheral.

Reading it as a failure of the model is not supportable, for six reasons, and
the status is an open comparison whose \emph{target} is unverified.

\begin{enumerate}
\item \textbf{The model's $r_1$ is reproduced by hand.} In the coronal limit
$r_1(2)$ follows from the single excitation rate coefficient
$K_{\rm exc}(1s\to n{=}2) = \SI{9.2796e-9}{\cubic\centi\metre\per\second}$ and
gives $1.37\times10^{-5}$, against the full $43$-state solve's
$1.66\times10^{-5}$. The two agree to \SI{21}{\percent}, which is the accuracy a
coronal estimate that ignores cascades should have. Whatever produces a factor
of $8$ is not inside the solve.

\item \textbf{The atomic stack passes three independent external checks.} The
radiative recombination coefficient into the ground state is within
\SI{2}{\percent} of the scaled standard value; the sum over excited states,
$\SI{2.08e-13}{\cubic\centi\metre\per\second}$, sits \SI{10}{\percent} below the
standard case-B coefficient $\alpha_B \approx
\SI{2.31e-13}{\cubic\centi\metre\per\second}$, the accepted total for capture
into every excited level, and that \SI{10}{\percent} shortfall is exactly the
$n>15$ truncation of Section~\ref{sec:convergence}; and the ground-state ionisation
coefficient at \SI{10}{\electronvolt}, $\SI{4.864e-9}{\cubic\centi\metre\per\second}$,
sits \SI{10.4}{\percent} below ADAS SCD96 read in its low-density limit,
$\SI{5.43e-9}{\cubic\centi\metre\per\second}$ at
$\nel = \SI{5e7}{\per\cubic\centi\metre}$ \cite{Summers2006, ADAS}, where the
effective coefficient reduces to ionisation out of the ground state because
every excited level decays before it can be ionised. The comparison is at the
low-density limit deliberately: SCD96 varies by \SI{0.9}{\percent} over a
factor four in density there, so the limit is reached. \textbf{This
\SI{10}{\percent} is not the truncation of the item above.} Direct ionisation
from $1s$ is a single cross section and does not know where the level ladder
stops; the difference is between this model's ionisation cross-section source
and the one behind SCD96, and it is reported rather than reconciled.

\item \textbf{The tabulated asymptote demands an excitation coefficient fifteen
times the accepted value.}
Reproducing Fujimoto's quoted coronal asymptote requires
$K_{\rm exc}(1s\to n{=}2) = \SI{1.44e-7}{\cubic\centi\metre\per\second}$ at
\SI{11.03}{\electronvolt}, fifteen times the accepted value. No error in this
model's implementation produces that number; only a different input does.

\item \textbf{Three candidate explanations are eliminated quantitatively.}
Truncation moves $r_1(3)$ by \SI{1.2}{\percent} between $n_{\max}=10$ and
$n_{\max}=15$, which is a factor $8$ short. Lyman trapping at an escape factor
of $0.1$ brings $r_1(2)$ to $0.92$ of the table but leaves $r_1(3)$ at $0.26$ of
it and drives $r_0(2)$ to $8.2$, a population above Saha equilibrium and
therefore unphysical. The third candidate, the $\ell$~closure, is eliminated in
the item below.

\item \textbf{The $\ell$~closure is a real difference, and it does not account
for the gap.} Table~4.1(b) is indexed by $p = 2, 3, 4, 5, 7, 10, 15$, principal
quantum numbers, and carries no $\ell$ label: the published coefficients are
shell quantities computed with the $\ell$~substructure bundled, while this model
resolves $\ell$ below $n=9$. That difference is real, and it is not the
explanation. Imposing the source's own
closure directly, by bundling this model's operator under statistical
$\ell$~weights as $\Lmat_{\rm bundle} = \mathsf{P}\,\Lmat\,\mathsf{R}$ with
$\mathsf{P}_{ki} = 1$ for $i$ in shell $k$ and $\mathsf{R}_{ik} = g_i/g_k$,
moves $r_1(3)$ by \SI{0.42}{\percent}: the factor $8.29$ becomes $8.26$
(\texttt{verify\_fujimoto\_bundle.py}, stamped in
\texttt{validation/fujimoto\_bundle/}). At
$p=2$ the agreement becomes \emph{worse}, from a factor $10.8$ low to a factor
$11.8$ low, and $r_0(2)$ moves from $0.7392$, which agrees with the tabulated
$0.730$, to $0.6459$, which does not. The reason is structural: the bundling
projector annihilates the $\ell$-mixing operator identically, $\max|\mathsf{P}
\mathsf{B} \mathsf{R}| = 8.5\times10^{-7}$ against a matrix scale of
$1.0\times10^{11}$, so bundling cannot express what $\ell$-mixing does; and the
production model is already at the statistical-$\ell$ limit for $p \ge 3$, its
sublevel populations sitting within \SI{5}{\percent} of $g_{n\ell}/g_n$. The
$\ell$-resolved calculation was therefore already the like-for-like comparison.
Removing proton $\ell$-mixing altogether moves $r_1(2)$ by a factor $118$, but
that is the opposite limit from a statistical closure rather than a proxy for
it, and the two do not bracket the tabulated value.

\item \textbf{The transcription and the row labels are correct.} Both were
checked against the printed table. The $\log_{10}\nel = 18$ row reads
$0.730$, $0.835$, $0.947$ and $0.983$ for $r_0$ at $p = 2, 3, 4, 5$, and
$1.79\times10^{-4}$, $5.72\times10^{-5}$, $1.66\times10^{-5}$ and
$5.08\times10^{-6}$ for $r_1$, matching the values used here exactly. The row
labels are $\log_{10}(\nel/\si{\per\cubic\metre})$, which the table establishes
internally: $r_0(2)$ reaches $0.981$ at $\log_{10}\nel = 21$, and Griem's
criterion places local thermodynamic equilibrium for $n=2$ at
\SI{1.74e16}{\per\cubic\centi\metre}, which is $10^{21}\,\si{\per\cubic\metre}$
to within the criterion's own precision. Read as \si{\per\cubic\centi\metre} the
onset would sit five orders above Griem, which is impossible; and the rows run
from $\log_{10}\nel = 12$ to $24$, a span of $10^{6}$ to
$10^{18}\,\si{\per\cubic\centi\metre}$ read as \si{\per\cubic\metre} and no
plasma at all at the top read as \si{\per\cubic\centi\metre}. A one-decade
offset in the density rows is not a tenable reading either.
\end{enumerate}

\textbf{Status.} The $r_1$ deficit at low $p$ is an open external
disagreement, and this thesis does not explain it; the three candidate causes
above each fail to account for a factor of $8$. The $\ell$-closure test carries
a severity check: bundling an operator whose $\ell$~distribution is deliberately
non-statistical moves $r_1(2)$ by a factor $129$, so the null result at $p=3$ is
a measurement and not a blind instrument. What the section does \emph{not}
establish is that the disagreement leaves the quantity
Chapter~\ref{ch:results} uses untouched.
Equation~\eqref{eq:fujimoto_correspondence} makes
$a_p = r_1(p)\,\Zsb{p}/\Zsb{1}$: $r_1$ is the ground-fed coefficient from which
$\ffrac{p}$ is built, so a deficit in $r_1(3)$ and $r_1(4)$ is a deficit in
$a_3$ and $a_4$ themselves. Bundling is not the protection, since neither shell
passes through a bundled coefficient; the only protection is that $\Delta$ and
the cap depend on the ratio $a_3/a_4$, in which a deficit common to both shells
cancels. At \SI{e12}{\per\cubic\centi\metre} the deficit is not common, $8.3$
at $p=3$ against $4.4$ at $p=4$; at \SI{e15}{\per\cubic\centi\metre} it nearly
is, $1.8$ against $1.9$. Dividing the model's $a_3$ and $a_4$ by those ratios,
so that they become what Fujimoto's coefficients would give, and recomputing
at every grid point (\texttt{verify\_fujimoto\_r1\_consequence.py},
\texttt{validation/fujimoto\_r1\_consequence/}) moves $\Delta$ by $+0.64$
under the first pair and by $-0.08$ under the second. Over the $280$ points
above \SI{2}{\electronvolt} the cap then moves by \SI{+29}{\percent} at the
median under the first and \SI{-4}{\percent} under the second; the crest
$u^{\mathrm{peak}}$ falls by a factor $6.0$ and $1.9$; and the sensitivity at
the operating point, $\ffrac{3}-\ffrac{4}$ evaluated at $u^{\mathrm{CRE}}$,
moves by \SI{+30}{\percent} and \SI{+11}{\percent} at the median, with ranges
across the points of $-58$ to \SI{+524}{\percent} and $-38$ to
\SI{+68}{\percent}. Both exceed the \SIrange{8}{9}{\percent} that the R-matrix
excitation substitution of Chapter~\ref{ch:atoms} produces. The disagreement
is therefore direct evidence about the coefficients Chapter~\ref{ch:results}
uses; what is open is whether the deficit lies in this model or in the
tabulation, and the scenario applies a deficit measured at
\SI{10}{\electronvolt} at every temperature, which is its stated assumption.
An unexplained disagreement in an absolute population coefficient is reported
here as unexplained, with its stake stated.

\textbf{What the benchmark does establish.} The recombination-fed channel
agrees, and because $r_0$ and $r_1$ carry the same $p$-dependent factor $Z(p)$,
a normalisation error large enough to explain the $r_1$ deficit is excluded by
$r_0$ agreeing. That is a pass, and it should be read as one.

\textbf{What it does not establish, and two limits on it.} It does not fix the
absolute normalisation of the ground-fed channel against an independent
calculation, because no $\ell$-resolved tabulation was available to compare
against. Table~4.1(b) is computed at $T_\mathrm{e} = \SI{1.28e5}{\kelvin}$,
which is \SI{11.03}{\electronvolt}, while this grid reaches
\SI{10}{\electronvolt}, so the comparison sits at the grid edge, with the
tabulated temperature \SI{10.3}{\percent} above it. The companion
Table~4.1(a) is at \SI{e3}{\kelvin} =
\SI{0.0862}{\electronvolt}, two orders below this grid, so there is no second
usable case. The external check is one temperature at one edge.

%-----------------------------------------------------------------------------
\section{Gate D fails, and the failure is in the gate}
\label{sec:gate_d}
% QUESTION: the one gate that compares against an external atomic database
% disagrees at every point. What does that mean, and what does it not mean?

The Atomic Data and Analysis Structure (ADAS) distributes effective
stage-to-stage rate coefficients for hydrogen: SCD, the effective ionisation
coefficient, and ACD, the effective recombination coefficient~\cite{Summers2006}.
These are the coefficients a transport code uses when it does not resolve
excited states, and they are the most widely used external reference this model
could be compared against.

Gate~D forms the model's counterpart by taking the steady-state population
$n_p^{\rm ss}$ of each level at each grid point, weighting it by that level's
ionisation rate coefficient $Q_p$ of Eq.~\eqref{eq:level_balance}, and dividing
by the total neutral population:
\begin{equation}
  {\rm SCD}_{\rm model}
  = \frac{\sum_p Q_p\,n_p^{\rm ss}}{\sum_p n_p^{\rm ss}},
  \label{eq:scd_model}
\end{equation}
then compares against ADAS interpolated to the same temperature within a
factor-of-two density band. The gate passes a point if the ratio
$\eta = {\rm SCD}_{\rm model}/{\rm SCD}_{\rm ADAS}$ lies between $0.5$ and $2.0$.

\begin{center}
\texttt{Gate D: FAIL  (0\% of points)} \hfill \texttt{validation/gate\_summary.txt}
\end{center}

Measured, $\eta$ ranges from $5.73$ to $8346$ over the full grid, and no point
agrees within a factor of two. The extremes are not typical: over most of the
grid the model sits high by a factor between $9$ and $65$, and that systematic
offset peaks near $\Te \approx \SI{2}{\electronvolt}$ (these are the gate's
own numbers, under its linear read of the table; the read is discussed
below). The disagreement has a
shape, which is what makes it diagnosable rather than random.

\subsection{The diagnosis, and the test that settles it}
\label{sec:gate_d_diagnosis}

Two readings were open. One of them has now been tested and it is the
explanation.

The first is that Eq.~\eqref{eq:scd_model} is not the quantity ADAS tabulates.
ADAS defines SCD as the \emph{ionising} coefficient, the ionisation flux
attributable to the ground-state population. The sum in
Eq.~\eqref{eq:scd_model} runs over the full steady state, which includes the
recombination-fed Rydberg levels, the $\nzero = \bm{c}\,\nion$ channel of
Section~\ref{sec:R_depends_on_b}. Ionisation out of
those levels is a recombination-driven flux and belongs under ACD, not SCD.
Since the high levels have small ionisation potentials and large ionisation
coefficients, moving them to the wrong side of the ledger inflates the model's
SCD by a large factor, and inflation is what is measured. The test that would
settle it is to recompute Eq.~\eqref{eq:scd_model} over the ground-fed channel
$\none$ alone, and it has now been run
(\texttt{diagnose\_gate\_d.py}).

The two channels separate exactly, with no approximation:
$\nvec_E = \nzero + \none$ to a relative residual below $10^{-8}$ at every grid
point, which the script checks before using either. Rebuilding the coefficient
on the ground-fed channel alone,
\begin{equation}
  {\rm SCD}_{\rm model} = Q_{1s}
  + \sum_{p \neq 1s} Q_p\,\frac{n^{(1)}_p}{\ngs},
  \label{eq:scd_groundfed}
\end{equation}

\noindent and comparing against the same ADAS table with the same factor-of-two
tolerance moves the agreement from $0$ of $400$ points to $400$ of $400$. The
ratio $\eta$ falls from a range of $5.73$ to $8346$ to a range of $0.702$ to
$1.004$, with median $0.828$; above \SI{2}{\electronvolt} it runs from $0.721$
to $0.937$, median $0.831$, and at the benchmark it is $0.768$
(\texttt{validation/reservoir\_vs\_adas/}). The read matters. ADAS96
tabulates seven temperatures between $1$ and \SI{10}{\electronvolt}, and
\texttt{diagnose\_gate\_d.py}, which first made this comparison, interpolates
linearly in value between them (\texttt{adas\_at}, lines 114 to 124); across
$1$ to \SI{1.5}{\electronvolt}, where SCD rises by nearly two orders of
magnitude, that overstates the table by up to a factor $6.65$. That read gives
$323$ of $400$, a range of $0.130$ to $0.931$ and a median of $0.696$, with a deficit deepening toward low temperature; read logarithmically
in both axes, as the table's spacing requires, the tail disappears, the row
median being $0.72$ at \SI{1}{\electronvolt} and $0.90$ by
\SI{1.2}{\electronvolt}. The diagnosis script is left in place; every ADAS
number in this section comes from the logarithmic read.

The mechanism is visible directly. At collisional--radiative equilibrium the
recombination-fed channel carries between \SI{89.5}{\percent} and
\SI{99.99}{\percent} of $\sum_p Q_p n_p$, and its share rises monotonically with
density, from \SI{97.8}{\percent} at $\SI{e12}{\per\cubic\centi\metre}$ to
\SI{99.99}{\percent} at $\SI{e15}{\per\cubic\centi\metre}$ at
\SI{1}{\electronvolt}. Equation~\eqref{eq:scd_model} is therefore reporting the
recombination channel with a small ground-fed correction, which is why $\eta$
tracks $\nel$: three-body recombination carries an extra factor of $\nel$ and
the mislabelled flux carries it into the coefficient.

\paragraph{The recombination half, built and compared for the first time.}
The same decomposition supplies the coefficient the gate never computed. The net
flux into the neutral ground state carried by the recombination-fed channel,
per unit $\nel\nion$, is
\begin{equation}
  {\rm ACD}_{\rm model}
  = \frac{S_{1s}\,\nion + \sum_{p \neq 1s} L_{1s,p}\,n^{(0)}_p}{\nel\,\nion},
  \label{eq:acd_groundfed}
\end{equation}

\noindent direct radiative recombination into $1s$ plus everything that cascades
or is collisionally transferred down into it. Against ACD96 this agrees at
\emph{all} $400$ points, with $\eta$ between $0.894$ and $1.009$ and median
$0.982$; at the benchmark the model gives
$\SI{2.788e-13}{\cubic\centi\metre\per\second}$ against ADAS's
$\SI{2.883e-13}{}$, a ratio of $0.967$.

\textbf{That is the strongest external check in this thesis.} It is an absolute
coefficient, compared against an independent database built from independent
atomic data, agreeing to better than \SI{10}{\percent} at the benchmark and
within a factor of two everywhere. It was available all along and was not taken,
because the half of the gate that would have taken it was never written.

The second reading was that the two quantities are not defined on the same
system at all. ADAS coefficients describe a closed atom-plus-ion system; this
model holds the ion as a fixed external reservoir with no forty-fourth state
(Section~\ref{sec:two_references}). A stage-to-stage coefficient extracted from
an open system need not equal one extracted from a closed one, and
Section~\ref{sec:attacks} has already shown that the boundary condition moves
$\tauslow$ by up to a factor of $41$. \textbf{The ACD result argues against this
being the dominant effect}: ACD is extracted from the same open system through
the same fixed-ion reservoir and agrees with ADAS to within \SI{3.3}{\percent}
at the benchmark, \SI{1.8}{\percent} at the grid median and \SI{11}{\percent}
at every point; an open-system artifact large enough to produce $\eta$ of
$8346$ in SCD would not leave ACD intact.

\paragraph{The residual SCD deficit, stated rather than dismissed.} Rebuilding
SCD on the ground-fed channel does not make the agreement perfect, and the
remaining disagreement is systematic in a way worth reporting: the model's
ionisation sits below ADAS at all but one of the $400$ points, by up to
\SI{30}{\percent} and by \SI{17}{\percent} at the median, with the maximum
$\eta = 1.004$ at $[4,7]$. Under the linear read of the table the gap closes
steeply with temperature, from a median of $0.302$ over $1$ to
\SI{1.5}{\electronvolt} to $0.817$ over $5$ to \SI{10}{\electronvolt}, a shape
that reads as the signature of truncation; the steep part is the linear read,
and under the logarithmic read the deficit is between $0.72$ and $0.94$ at every
temperature above \SI{2}{\electronvolt}. A deficit of that size is still what
truncation would produce, since ADAS96 carries levels above $n_{\max} = 15$
and the levels nearest the continuum ionise most easily, but its temperature
shape no longer singles that cause out (the $n \ge 12$ flux share the diagnosis
script printed against the linear read is not carried forward). The test that
would settle the
cause is the $n_{\max}$ convergence study of Chapter~\ref{ch:limits} rather
than this one.

\paragraph{The reservoir itself, against the same table.} ADAS's two
coefficients fix the equilibrium reservoir: at ionisation balance
$n_{1s}\,\mathrm{SCD} = \nion\,\mathrm{ACD}$, so
$u^{\mathrm{CRE}} = \mathrm{ACD}/\mathrm{SCD}$, and the model obeys the same
identity with its own ground-fed SCD and ACD to $4\times10^{-13}$ at every grid
point. That makes the quantity Chapter~\ref{ch:results} is built on directly
comparable, and the comparison has two parts
(\texttt{verify\_reservoir\_vs\_adas.py},
\texttt{validation/reservoir\_vs\_adas/}). The \emph{level} is high:
$u^{\mathrm{CRE}}_{\mathrm{model}}/u^{\mathrm{CRE}}_{\mathrm{ADAS}}$ has median
$1.18$ over the $400$ points, runs from $0.89$ at $[4,7]$ to $1.40$ at $[0,1]$,
and is $1.26$ at the benchmark, structured about equally in $\Te$ and in
$\nel$. The model's ground reservoir per ion sits a fifth above ADAS's, which
moves the operating point beneath the cap of Section~\ref{sec:bound} without
touching the cap. The \emph{displacement}, which is what the plateau error is
built from, agrees. Over the $392$ one-interval temperature steps the model's
$\Delta\ln u$ has median magnitude $0.26$ and differs from ADAS's by a median
$0.014$; above \SI{2}{\electronvolt} the median difference is $0.011$ and the
worst $0.050$, at $[15,0]$; at the benchmark step the two are $-0.270$ and
$-0.258$. The reservoir gain $G$ this thesis tabulates is therefore within
about \SI{5}{\percent} of the one ADAS implies, at the median over the
defended range, and the refuter set before the run, a difference above $0.1$
at any interior step above \SI{2}{\electronvolt}, appeared at none of $264$.
Every number in this paragraph comes from the logarithmic read; on the linear
read of \texttt{diagnose\_gate\_d.py}, which overstates SCD by up to a factor
$6.65$ over $1$ to \SI{1.5}{\electronvolt}, the refuter would have appeared at
$9$ of the $264$ steps.

\subsection{Three further defects in the gate itself}

The gate is also incompletely implemented, and this must be said alongside the
failure rather than instead of it.

The recombination half does not exist \emph{in the gate}. The ACD table is
loaded once, at line 299 of \texttt{validate\_gates.py}, and the variable
holding it is never read again anywhere in the file; only the ionisation
comparison runs, despite a docstring that describes both. That half is computed
in \texttt{diagnose\_gate\_d.py} (Section~\ref{sec:gate_d_diagnosis}).

The interpolation is wrapped in a catch-all exception handler that sets the
comparison value to not-a-number and continues, so a failure to find the right
density band is indistinguishable from a point that was legitimately skipped.

And the gate's status is recorded three incompatible ways: \texttt{README.md}
line 119 calls it documented, the docstring at line 47 predicts a pass above
$\Te = \SI{5}{\electronvolt}$, and \texttt{gate\_summary.txt} reports failure
at \SI{0}{\percent} of points. The correct entry is the measured one: Gate~D
fails.

\subsection{What is claimed, and what is not}

\textbf{Gate~D itself still fails, and it has not been changed.}
\texttt{validate\_gates.py} is untouched: no tolerance was widened, no subset
was selected, and \texttt{gate\_summary.txt} still records
\SI{0}{\percent}. What has changed is that the failure now has a demonstrated
cause rather than two undemonstrated candidates. The gate compares a mixed
quantity against a table that keeps the two channels apart by design, which is
a defect in the comparison and not in the matrix.

Deciding what Gate~D should test belongs with whoever owns the gate. What is
claimed here is narrower: rebuilt on the two channels separately, the
coefficients agree with ADAS within a factor two at all $400$ points for SCD,
the model $0$ to \SI{30}{\percent} low with a median ratio of $0.83$, consistent
with the $n_{\max}$ truncation but not uniquely so, and at all $400$ for ACD
with a median ratio of $0.98$.

The failure does not cost the thesis its only comparison of absolute effective
coefficients against an external database; that comparison exists, hidden by
the definition. The failure does not touch the structural results
of Chapter~\ref{ch:theory}, which are statements about whatever $\Lmat$ is
supplied and would hold for the ADAS-consistent matrix equally. It does not
touch the level-resolved comparison of Section~\ref{sec:fujimoto}, which
compares the same decomposition against tabulated coefficients rather than
against a contraction of $42$ coefficients into one. And the contraction is the
reason the two comparisons carry different weight: many incorrect $\bm{a}$ would
contract to the correct SCD, so a passing Gate~D would have been weak evidence
even had it passed.

%-----------------------------------------------------------------------------
\section{Errors found by our own audit}
\label{sec:errors_found}
% QUESTION: what did this project get wrong, how was each error found, and how
% large was it?

Section~\ref{sec:atoms_corrections} reports two implementation errors that this
project found in itself, with their sizes, and observed that both were invisible
to every consistency test the pipeline ran. This section states what each
changes here, adds two more, and keeps the four together rather than in an
appendix: a model whose authors found and bounded four of its own errors is
better characterised than one reporting none, and in three of the four cases
the bound is itself a result.

\subsection{The $\ell$-mixing Coulomb logarithm}

Proton-impact $\ell$-mixing redistributes population between sublevels of the
same $n$ and, as Section~\ref{sec:lmixing} showed, is the dominant loss channel
for the metastable $2s$ level. The rate coefficient follows Badnell's
formulation~\cite{Badnell2021}, whose Eq.~9 carries a factor
\begin{equation}
  F(U_m) = \frac{\sqrt{\pi}}{2}U_m^{-3/2}\operatorname{erf}(\sqrt{U_m})
           - \frac{\mathrm{e}^{-U_m}}{U_m} + E_1(U_m),
  \label{eq:FUm}
\end{equation}
in which $\operatorname{erf}$ is the error function, $E_1$ the exponential
integral, and $U_m$ the dimensionless ratio of the smallest energy transfer the
Debye cutoff permits to $k_B\Te$; $F(U_m)$ is the Coulomb logarithm of the
interaction, which diverges as $U_m \to 0$.

The implementation set $F = 1$, justified in its docstring as the low-density
limit $U_m \to 0$; the limit was inverted, and
Section~\ref{sec:atoms_corrections} gives the finding, the factor of three to
seven by which the rates were low, and the measured effect on the spectrum.
Evaluated with the local Debye length at the benchmark point $[23,5]$,
$F = 6.62$ for $2p$--$2s$, $5.71$ for $3d$--$3p$, $4.26$ for $5g$--$5f$ and
$2.89$ for $8k$--$8i$ (\texttt{validation/lmix\_cutoffs/}); with the
production table's frozen density of \SI{e14}{\per\cubic\centi\metre} the
$2p$--$2s$ value is $6.95$ (Section~\ref{sec:atoms_corrections}). A derivation
note records $6.65$, $5.73$, $4.29$ and $2.92$ for the same four pairs.

The insensitivity of the slow eigenvalues to that error, which
Section~\ref{sec:atoms_corrections} measures, is specific to a process already
far faster than $\taurel$; it is not a general licence to be careless about
fast rates, and the $\ell$-mixing scan below is what establishes the range over
which it holds.

\subsection{And $\ell$-mixing is saturated}

The same insensitivity was then tested directly rather than argued. Scaling
every $\ell$-mixing rate in the matrix over a range of $\times0.1$ to $\times10$
moves $\ffrac{3}-\ffrac{4}$ by less than \SI{0.3}{\percent} at the benchmark
point and at the ridge, and by less than \SI{5}{\percent} at the cold corner.
Only switching $\ell$-mixing off entirely changes anything. Rebuilding the
matrix with a density-consistent Debye cutoff in place of a frozen one moves
$\ffrac{3}-\ffrac{4}$ by at most \SI{0.42}{\percent}.

The claim that the frozen Debye density varies the cutoff factor by about
\SI{30}{\percent} across the grid is wrong: the module's own functions give
$\times3.7$ for $n=2$ and at least $\times50$ for $n=8$
(Section~\ref{sec:lmixing}); the conclusion survives because the observable does
not depend on the input over that range.

\subsection{The eigenvalue magnitude filter}

The spectrum of $\Lmat$ was extracted by a routine containing the line
\texttt{eigs = eigs[eigs < -1.0]}, which discards every eigenvalue smaller in
magnitude than $1\,\si{\per\second}$. Where the slow eigenvalue is smaller than
that, at the cold end of the grid, the filter deleted it and promoted the whole
ladder by one step, so $\tauslow$, $\taurel$ and $\Msep$ were wrong together
while all three remained finite, negative and correctly ordered. It was found by
recomputing the spectrum unconditionally and comparing bit-for-bit against the
stored arrays, the three-implementation agreement of Section~\ref{sec:gate_e}.
Section~\ref{sec:atoms_corrections} gives the extent of the damage, the ladder
shift that confirmed the mechanism, the preserved filtered outputs and the
removal of the literal from two dormant files. What belongs here is the shape
of the wrong answer: at the cold corner the filter turned
$\Msep = 1.729\times10^{9}$ into $\Msep = 1.467$, and the grid minimum from
$86.77$ into $1.341$, values that read as a physics finding, since
near-breakdown of timescale separation in the coldest, thinnest plasma is
exactly what a reader would expect to see. The later correction of the two
dormant files changes no number in this thesis: at the benchmark the two filters
return the same $\tauslow$ to every digit, and they differ only where the slow
mode is slower than $1\,\si{\per\second}$.

\subsection{The central quantity was misnamed}

The largest of the four errors changed no number and inverted a sentence.

The quantity that had been reported as the quasi-steady-state error was
defined as the difference between the tabulated ratio and the true
one, with the tabulated ratio evaluated from a two-argument function of $\Te$
and $\nel$. Section~\ref{sec:two_references} shows that the excited-state ratio
is a function of three arguments, not two, the third being the state of the
ground-state reservoir. Fixing the first two therefore silently imposes
equilibrium ionisation balance, so what was measured was the distance to the
equilibrium answer, not the error of the quasi-steady-state closure.

Integrating the full system and evaluating both errors along the same trajectory
separates them:

\begin{center}
\begin{tabular}{lrr}
  \toprule
  point & error against equilibrium & error of the QSS closure \\
  \midrule
  benchmark $[23,5]$ & $6.34\times10^{-2}$ & $8.66\times10^{-6}$ \\
  worst case $[0,4]$ & $3.869\times10^{-1}$ & $6.73\times10^{-9}$ \\
  \bottomrule
\end{tabular}
\end{center}

They differ by four to eight orders of magnitude, most at the point carrying the
largest error. The physics was intact and every number reproduced; what was
wrong was the name and the sentence built on it, and Chapter~\ref{ch:results}
states the consequence, that the approximation everyone worries about is not
the one that fails. The error was found by an adversarial reading of the thesis text
itself, not of the code, and it is why Section~\ref{sec:eps_definitions}
defines four error measures separately.

\subsection{A statement in Chapter~\ref{ch:theory} that measurement contradicted}
\label{sec:promise_broken}

The statement, in Section~\ref{sec:source_quadratic}, was that at
$\nel = \SI{e12}{\per\cubic\centi\metre}$ three-body recombination supplies
under \SI{10}{\percent} of the feed into any level, and it was used to propose
testing the radiative and three-body datasets separately at the two ends of the
density grid. Measured against the matrix, $30$ of $36$ levels exceed
\SI{10}{\percent} three-body feed at both \SI{1.00}{\electronvolt} and
\SI{2.95}{\electronvolt}, and $26$ of $36$ at \SI{10}{\electronvolt}. Into the
bundled shells it is essentially the whole feed.

The density grid therefore does not separate the two recombination datasets, and
\textbf{no such separate test was run or can be}; Section~\ref{sec:source_quadratic}
now states the measured fractions and withdraws the promise. What remains is the
gap: the radiative and three-body datasets enter this model only in combination,
nothing in this thesis tests them apart, and a test that could would need a
density below the bottom of this grid, outside the conditions the thesis is
about.

%-----------------------------------------------------------------------------
\section{Reproduction, and what a document cannot certify}
\label{sec:reproducibility}
% QUESTION: can the numbers in this thesis be regenerated, and by whom?

\subsection{Regeneration from the canonical matrix}

Every quantitative result in Chapters~\ref{ch:theory} and~\ref{ch:results} was
regenerated from $\Lmat$ and compared against the stored output. The timescale
and step-error grids reproduce bit-identically, including the values that come
from numerical integration and not from linear algebra. The plateau and persistence
tables reproduce byte for byte, $680$ of $680$ rows, with $784$ superposition
residuals at exactly zero difference. The nineteen points corrected by the
eigenvalue filter reproduce exactly, including the ladder-shift identity.

This is reproduction in the narrow sense: the same code, run again, on the same
inputs, gives the same answers. It excludes a stale array and it excludes a
non-deterministic solver. It says nothing about whether the algebra is right.

\subsection{Independent re-derivation}

The stronger check is re-derivation. All six central results of
Chapter~\ref{ch:theory} were derived from the linear collisional--radiative
system alone by a party working from the physical statement of the problem,
implemented from scratch, with the predicted values written down before any
project file was opened.

In the table below, $\bdep$ is the ground-state departure coefficient of
Eq.~\eqref{eq:departure}, the ratio of the actual ground-state population to its
Saha--Boltzmann value, and $\bdep^{\rm CRE}$ is the value it takes in
collisional--radiative equilibrium (CRE), the steady state the lookup tables
assume.

\begin{center}
\begin{tabular}{lrr}
  \toprule
  quantity & \texttt{chapter3.tex} & independent re-derivation \\
  \midrule
  $\bdep^{\rm CRE}$ at the benchmark        & $956$                        & $956.509$ \\
  switching points $c_m/a_m$                & $2.73\times10^{3}$, $1.91\times10^{4}$ & $2727.3$, $19085.5$ \\
  peak $\bdep$                              & $7.2\times10^{3}$            & $7214.7$ \\
  $\ffrac{3}$, $\ffrac{4}$, difference       & $0.260$, $0.048$, $0.212$    & $0.259652$, $0.047725$, $0.211927$ \\
  maximum of $\ffrac{3}-\ffrac{4}$          & $0.4514$                     & $0.451357$ \\
  $\Robs$ at the benchmark                  & $0.7778$                     & $0.777768$ \\
  \bottomrule
\end{tabular}
\end{center}

The timescales reproduce to four digits, $\tauslow = \SI{22.73}{\micro\second}$,
$\taurel = \SI{2.277}{\nano\second}$, $\Msep = 9981.9$.

Two further invariances were measured in the same pass and belong with the
result they protect. The shell separation $\Delta$ that sets the bound of
Section~\ref{sec:bound} is invariant under rescaling the ground-to-ion
normalisation, verified exactly under factors of $10^{10}$ and $10^{-7}$, while
the individual channel fractions are not; and $\Delta$ varies by
\SI{0.07}{\percent} across five different $\ell$-weightings of the shell sum, so
the bound is not an artifact of how sublevels were weighted.

\subsection{A prediction that was falsified, and improved the claim}

The same re-derivation made one prediction that turned out to be wrong, and it
is more informative than the six that were right.

The prediction was that $\ffrac{3}-\ffrac{4}$ must vanish at high density, on
the reasoning that in local thermodynamic equilibrium both supply channels
produce the same Boltzmann distribution, so the two shells would draw
identically and the difference would close. It does not happen. At
$\Te = \benchTe$, $\Delta$ doubles over the lowest part of the grid and then
stops moving: $0.979$ at $\SI{e12}{\per\cubic\centi\metre}$, $1.946$ at
$\SI{1.4e14}{\per\cubic\centi\metre}$, $1.939$ at
$\SI{e15}{\per\cubic\centi\metre}$
(\texttt{validation/fujimoto\_r1\_consequence/}). Above about
$\SI{7e12}{\per\cubic\centi\metre}$ it is flat, and the bound $\tanh(\Delta/4)$
built on it spans only $0.394$ to $0.450$, a \SI{14}{\percent} range, across the
whole of the rest of the grid.

The reason is the open boundary again, this time in the model's favour: the
recombination source is imposed externally rather than tied to the ion density
by a Saha relation, so the two channels cannot merge by construction, and
Gate~C's \SI{1.5}{\percent} of Saha says the grid never approaches the limit in
which they would. The consequence for the density maximum of the diagnostic
error, a position effect rather than a span effect, is
Table~\ref{tab:position_effect} in Section~\ref{sec:density_mechanism}.

\subsection{One item demoted: the May/August cross-validation}

A cross-validation between two runs four months apart, on different matrices,
with different observables and different numerical methods, was recorded in this
project's register as verified and described as the strongest single check it
had. The agreement is close: $\tauslow$ to \SI{0.064}{\percent}, $\taurel$ to
\SI{0.200}{\percent}, $\Msep$ to \SI{0.264}{\percent}, $\epsstep$ to
\SI{0.014}{\percent}. Predicting how long the error stays above
\SI{10}{\percent} from the plateau-then-decay picture gives
\SI{12.478}{\micro\second} against the earlier run's full integration,
\SI{12.4794}{\micro\second}.

Two corrections apply, and both weaken it.

The apparent agreement to eight parts in $10^{5}$ on the duration is
cancellation. The inputs agree only to \SIrange{0.06}{0.6}{\percent}, and
carrying the earlier run's own internal numbers through gives
\SI{12.406}{\micro\second}, off by \SI{0.59}{\percent}. \textbf{The figure is
agreement at the \SI{0.6}{\percent} level}, which is still strong.

More seriously, \textbf{the earlier run has never been reproduced against the
canonical matrix}. This project's own derivation record says of these same
numbers that they are a reported result, not independently re-run, and should be
treated as strong documented evidence rather than as self-verified. A
cross-validation whose earlier half cannot be re-executed is a citation, not a
check. It is reported here as documented evidence and is not counted on the
ladder of Section~\ref{sec:validation_summary}. Under the ground-truth ordering
this thesis works to, physics over mathematics over code over documents,
a document is the weakest of the four, and it is where this particular claim
currently rests.

\subsection{Two provenance defects that are not repaired}

Two arrays of timescales, $\Msep$, $\tauslow$ and $\taurel$, are written to the
same three file paths by two different scripts. Whichever ran last wins, and six
downstream readers cannot tell whose numbers they are holding. This is exactly
the condition under which the eigenvalue filter propagated, since the filtered
and unfiltered versions of the same array have the same name.

Two things bound what that costs here, and both were checked rather than
assumed. The two writers are not two calculations: each sorts the eigenvalues of
the same canonical $\Lmat$ and keeps the negative ones, neither applies the
filter that caused the earlier incident, and the arrays on disk agree with the
independently produced \texttt{verify\_timescales.py} artifact to $2\times
10^{-16}$. And none of the six readers feeds this document: every figure and
table they write is absent from it. The benchmark $\Msep = 9982$ quoted in
Chapter~\ref{ch:theory} and in the abstract is not read from these arrays at
all; it is $\tauslow/\taurel$ computed from $\Lmat$ by
\texttt{verify\_timescales.py} and recorded in
\texttt{validation/spectrum\_testpoint.csv}. The defect is therefore a latent
hazard in the repository rather than an ambiguity in any number printed here,
and it is left in place and named rather than quietly repaired.

An audit script exists to detect single-writer violations, and it reports an
unfixed failure of this kind. It is also incomplete. Its test for a write reads the file mode from the second
argument or a keyword, so the idiom \texttt{path.open("w")} passes through it
unrecorded. At the tip of \texttt{main} on 23~August that idiom had four users;
in the tree this thesis is built
from it has $45$ call sites in $29$ scripts, $42$ of them writing
comma-separated output, among them the producer of the plateau map that
fifteen places in Chapter~\ref{ch:results} read from
(\texttt{verify\_writer\_census.py}, \texttt{validation/writer\_census/}). The
audit should not be cited as evidence that the single-writer property holds.
The census confirms that the property fails for exactly the three timescale
arrays, each written by both \texttt{qss\_analysis.py} and
\texttt{validate\_gates.py}, and holds for the plateau map, which has one
writer.

%-----------------------------------------------------------------------------
\section{What this chapter establishes}
\label{sec:validation_summary}
% QUESTION: after grading every check, what is the model licensed to be used
% for, and where?

Table~\ref{tab:ladder} collects the checks in the order a reader should weight
them, with the grade of Section~\ref{sec:gate_philosophy} attached to each.

\begin{table}[tbp]
\centering
\caption[The validation ladder, graded by severity and provenance]{Every check
this chapter grades, with what it tests, its severity, whether its result can be
regenerated, and the measured outcome. \emph{Tautological} means the check
cannot fail except on an arithmetic bug, because the quantity tested was
constructed from the relation being tested; \emph{weak} means it can fail but
only for an error far larger than any plausible one; \emph{severe} means a
plausible wrong version of this model fails it. The provenance column is the
more uncomfortable of the two: \emph{stamped} means a script in the repository
regenerates the number into an artifact under \texttt{validation/};
\emph{notes} means the number is recorded in this project's working files with
no producing script, so it cannot be re-executed and cannot fail anything today.
None of the checks graded severe are in that condition, and the one
\emph{notes} entry left is the $\tanh$ bound, which is graded tautological. The
counts stated in this caption are checked against the table's own rows by
\texttt{verify\_ladder\_counts.py}, which exits non-zero if they drift apart. Of the eighteen
internal checks, twelve are severe, four are tautological and two are weak; of the five external comparisons, two
pass, one passes only at the higher densities, one fails, and one is open.}
\label{tab:ladder}
\footnotesize
\begin{tabular}{p{3.9cm}p{1.8cm}p{1.3cm}p{5.6cm}}
\toprule
check & grade & record & measured \\
\midrule
Atomic-data provenance (Ch.~\ref{ch:atoms})  & severe & stamped & every coefficient traces to a named dated source; tallies exact \\
Excitation rates against RMPS (Ch.~\ref{ch:atoms}) & severe & stamped & \SI{81.4}{\percent} within \SI{20}{\percent} for $n_{\rm upper}\le4$, mean \SI{13.0}{\percent}; $n=5$ disagrees by factors of $2$--$5$, localised to the top shell of the RMPS basis \\
Raw cross sections, microscopic reversibility & severe & figure & ratio $0.9995 \pm 0.0032$, \SI{98.7}{\percent} within \SI{1}{\percent} \\
One energy ladder on both sides & severe & stamped & Saha ratio $0.999997$ for all $43$ states, spread $1.4\times10^{-15}$ \\
Superposition of the two channels & severe & stamped & $3.075\times10^{-14}$ over $784$ pairs; catches a $3\leftrightarrow4$ swap \\
Reduced against full $\Robs$ & severe & stamped & $0.777768$ by three routes; catches the same swap \\
Deliberate fault injection & severe & stamped & three of four injected faults invisible to the conservation residual; the fourth moves it by thirteen orders \\
Line-weighted against shell ratio & severe & stamped & ratio $0.948$ to $1.000$ over $784$ pairs, below unity at every one; census $44$ against $45$ \\
Escape-factor cross-check & severe & stamped & $\sigma_0$ against an independent module, \SI{0.059}{\percent} apart against a \SI{1}{\percent} raising tolerance \\
Pinned input hashes & severe & stamped & both SHA-256 match byte for byte \\
Recomputation guards on the figures & severe & code & every plotted row rebuilt from $\Lmat$ at $10^{-9}$, raises \\
Truncation convergence of $\ffrac{3}-\ffrac{4}$ & severe & stamped & increments one-signed and geometric at $245$ of $280$ defended points, ratio median $0.67$; extrapolated \SI{-0.46}{\percent} at the benchmark, \SI{+1.6}{\percent} at the cold corner, within \SIrange{-0.82}{+1.03}{\percent} on the defended set. Tests one quantity, not the state space \\
Gate~A, detailed balance & tautological & stamped & $<\SI{0.01}{\percent}$, but $K_{\rm deexc}$ was derived from $K_{\rm exc}$ \\
Column-sum conservation & tautological & stamped & $3.61\times10^{-11}$; see the fault-injection row for what it cannot see \\
The $\tanh$ bound & tautological & notes & holds at all $248$ points, and holds on corrupted input \\
Gate~E, timescale hierarchy & tautological & stamped & $\Msep \ge 86.8$ against a threshold of $1$ \\
Gate~B, coronal limit & weak & stamped & worst convergence $0.0034$ against a tolerance of $0.5$ \\
Gate~C, Saha ceiling & weak & stamped & monotone and below Saha; approach fraction \SI{1.5}{\percent}, excluded from the flag \\
\midrule
Fujimoto App.~4B timescales & external, passes & stamped & $0.34\times$ and $1.77\times$ at a grid point his example lands on \\
ACD96 against this model's recombination coefficient & external, passes & stamped & all $400$ points within \SI{11}{\percent}; $\eta$ from $0.894$ to $1.009$, median $0.982$, \SI{3.3}{\percent} at the benchmark (logarithmic read) \\
Fujimoto Table~4.1, $r_0$ & external, passes at $\lg\nel = 21$ & stamped & \SIrange{0.1}{1}{\percent} for $n\ge4$ at the higher density; \SIrange{12}{14}{\percent} at $p=3,4$ at $\lg\nel = 18$; $r_0(2) = 0.7392$ against $0.730$ \\
Fujimoto Table~4.1(b), $r_1$ & external, \textbf{open} & stamped & $8.3\times$ low at $p=3$. Truncation, Lyman trapping and the $\ell$~closure each tested and each eliminated; unexplained (\S\ref{sec:fujimoto}) \\
Gate~D, SCD96 against ADAS & \textbf{FAIL} & stamped & $\eta \in [5.73,\ 8346]$, $0$ of $400$ within a factor $2$. Rebuilt on the ground-fed channel alone and read logarithmically it reaches $400$ of $400$, $\eta$ from $0.702$ to $1.004$, median $0.828$; the gate itself is unchanged \\
\bottomrule
\end{tabular}
\end{table}

\subsection{What the model may be used for}

Three statements are supported by the evidence above, and they are the ones
Chapter~\ref{ch:results} relies on.

\textbf{The equations are solved correctly.} The superposition residual, the
independent re-derivation of all six central results, the bit-for-bit
regeneration, and the grid-wide conditioning measurement together establish that
the linear algebra does what Chapter~\ref{ch:theory} says it does, to a
precision many orders of magnitude finer than any quantity the thesis reports.

\textbf{The input rates are traceable and benchmarked where the observable
lives; one derived coefficient there is not.} The rates feeding $n \le 4$, which is where $H_\alpha$ and $H_\beta$
originate, agree with an independent R-matrix calculation at a mean absolute
error of \SI{13.0}{\percent}, and the two transitions this thesis leans on
hardest, $1s\to2p$ and $2p\to3d$, agree to between \SI{2.4}{\percent} and
\SI{13.3}{\percent} across the temperature range. Above $n=4$ the atomic data are
less certain and the truncation is measurable. Section~\ref{sec:convergence}
tests one component of the mechanism, the sensitivity $\ffrac{3}-\ffrac{4}$,
and finds it converged to within \SI{1.03}{\percent} at the defended points
where the geometric tail can be summed, in a stamped run
(\texttt{validation/nmax\_downward\_scan/}); under the same truncation
$u^{\mathrm{CRE}}$ is converged to within \SI{0.13}{\percent} at the $257$
defended points where its increments are geometric, $\Delta$ to within
\SI{0.57}{\percent} and $\tauslow$ to within \SI{1.6}{\percent} at all $280$,
and $\Delta\ln u$ across the defended heating steps moves by at most
$2.3\times10^{-5}$ between $n_{\max}=14$ and $15$ against a median step of
$0.20$ (Appendix~\ref{app:numerics}). What is measured
for the reservoir is a scan rather than a truncation: weakening the ionisation
and three-body exchange of the top six shells eightfold moves $\epsplat$ by
\SI{0.7}{\percent} at the median over the defended pairs
(Section~\ref{sec:ion_recomb}). All of that concerns the rates going in. The
coefficients that come out at the two shells the observable uses are a separate
matter, and one of them is in open disagreement: the ground-fed $r_1(3)$ and
$r_1(4)$ of this model sit factors $8.3$ and $4.4$ below Fujimoto's tabulation
at the lowest density, which is a deficit in $a_3$ and $a_4$ themselves and so
in the sensitivity built from them (Section~\ref{sec:fujimoto}). Benchmarked
inputs are therefore not the same thing as a benchmarked susceptibility, and
this thesis can claim the first but not the second.

\textbf{The structural results are robust to the choices that could have
manufactured them.} Closing the ion boundary changes the persistence census by
one part in $325$. Scaling $\ell$-mixing over two decades moves the observable
by under \SI{0.3}{\percent} where the thesis makes claims. The single-slow-mode
estimate of Eq.~\eqref{eq:elm_bound} reproduces the true time-averaged error to
within \SI{2.2}{\percent} at \SI{100}{\micro\second} over the $448$ pairs above
\SI{2}{\electronvolt}.

\subsection{What it may not be used for}

\textbf{No effective coefficient from this model should be compared against a
database without first separating its two supply channels.} Gate~D as written
fails at every point on the grid because it compares a coefficient that mixes
both channels against a table that separates them (Section~\ref{sec:gate_d});
rebuilt channel by channel it agrees within a factor two at all $400$ points
for SCD and ACD, and the model's total effective ionisation coefficient remains
a quantity ADAS does not tabulate.

\textbf{No claim rests on the terminal shell}, whose population is high by a
factor consistent with truncation, or on the $\ell$-distribution inside the
bundled shells, which is assumed rather than measured.
Section~\ref{sec:bundling_check} tests that assumption: $\ell$-mixing licenses
it over most of the grid, and at the two highest densities below
\SI{3}{\electronvolt}, where the locally evaluated mixing loses to collisional
loss, detailed balance among shells within $10^{-4}$ of Saha--Boltzmann
equilibrium licenses it instead; the adjacent-shell rates themselves are not
$\ell$-blind, and no grid point escapes both criteria.

\textbf{Nothing here validates the physics that is absent.} Every check in this
chapter tests a $43$-state atomic hydrogen model against itself, against
arithmetic, and against other atomic hydrogen models. None of them tests
transport, molecular channels, or radiation trapping, because none of those
appears in $\Lmat$ at all. A model can pass every check in
Table~\ref{tab:ladder} and still be the wrong model for a divertor.
Chapter~\ref{ch:limits} is where that question is asked, and it is asked with
numbers rather than with hedges.

\medskip
The instrument is now characterised: what it computes, how precisely, on which
region of the grid, and where the evidence for it runs out. What it says is the
subject of Chapter~\ref{ch:results}.

%     breaks no existing cross-reference.
%
%  2. OVERLAP, RESOLVED BY HOMES. Three passages exist in both chapters:
%       chapter6 sec:open_system   vs  Section~\ref{sec:attacks}   (home: ch4)
%       chapter6 sec:r1_deficit    vs  Section~\ref{sec:fujimoto}
%       chapter6 sec:nmax          vs  Section~\ref{sec:convergence}
%     The Chapter 6 versions carry the SCOPE consequence; the Chapter 4
%     versions carry the EVIDENTIAL standing. On the r_1 deficit: the table
%     is bundled-n (backlog I.1, 10 Sep 2026; findings_10 ADDENDUM I), the
%     l-closure was tested and does NOT account for the gap, and the deficit
%     DOES reach a_3 and a_4; its consequence for Delta, the cap and S is
%     quantified in the Status paragraph of sec:fujimoto
%     (validation/fujimoto_r1_consequence/). Do not reintroduce the earlier
%     "does not propagate" sentence.
%
%     [17-18 Sep 2026] verify_bundling_psm20.py HAS now been executed
%     (sec:bundling_check in chapter 6, outputs/bundling/); an earlier
%     version of this comment said the "never executed" sentence was still
%     true and must not be corrected. It was corrected. The n_max
%     convergence result carried in Section~\ref{sec:convergence} is now
%     stamped (verify_nmax_downward_scan.py, 18 Sep 2026); the working-note
%     values (ADDENDUM D.8) it replaced had the opposite sign at the
%     benchmark and the cold corner, and the text records that. Its \todo asking for provenance on the
%     4.9-6.2x terminal-shell excess is also closed there, by the internal
%     terminal-vs-interior measurement (a_p at [49,0]: 9.4x at p=8, 10.7x at
%     p=9, 6.3x at p=10), which does not bracket it (band 6.3 to 10.7 against
%     4.9 to 6.2) but lies within a factor 2.2 of it and shares the density
%     trend; stamped 21 Sep 2026 by verify_terminal_vs_interior.py.
%
%  3. chapter3.tex:222-225 asserts a severity for the conservation check that
%     Section~\ref{sec:conservation_severity} measures and refutes. Either
%     soften line 224 or let this chapter carry the correction; do not leave
%     both standing as written.
%=============================================================================

\chapter{Does the Model Work?}
\label{ch:validation}

Chapter~\ref{ch:theory} built a $43\times43$ rate matrix, eliminated the fast
variables, and extracted from the result a closed-form statement about how a
line ratio responds to its reservoir. Chapter~\ref{ch:results} will use that
machinery to make quantitative claims about a diagnostic in use on real
machines. Between the two sits a question that no amount of further algebra
answers: why should anyone believe a number this model produces?

The answer is not a single verdict. It is a graded one, because the
checks that have been run on this model are not of equal worth, and several of
the ones that look most impressive are worth almost nothing. This chapter sorts
them.

\medskip
\noindent\textbf{What this chapter takes as given.} The state space, the atomic
data and their sources, and the benchmark of the excitation rates against an
independent R-matrix-with-pseudostates (RMPS) calculation are established in
Chapter~\ref{ch:atoms}, Sections~\ref{sec:atomic_uncertainty}
and~\ref{sec:atoms_corrections}, and are not repeated. The construction of
$\Lmat$, the conservation identity, the quasi-steady-state (QSS) elimination,
and the two-channel decomposition of the excited populations come from
Chapter~\ref{ch:theory}. This chapter asks only what the checks run on that
construction are entitled to establish.

\medskip
\noindent\textbf{What this chapter does not do.} It reports no result about
divertor plasmas. Every number below is about the instrument, not about the
physics the instrument is pointed at. Keeping those separate matters here more
than anywhere else in the thesis: a model that predicts something interesting
is not thereby a correct model, and a check that uses the interesting
prediction as its evidence has verified nothing.

%-----------------------------------------------------------------------------
\section{What a check can and cannot show}
\label{sec:gate_philosophy}
% QUESTION: what distinguishes a check that could have failed from one that
% could not, and how is that distinction measured rather than asserted?

\subsection{Three things that get called validation}

Three activities are routinely reported under one heading, and they license
different conclusions.

\emph{Verification} asks whether the equations that were intended got solved
correctly. Conservation residuals, detailed balance, matrix identities,
convergence with tolerance, agreement between two independent code paths: all
of these compare the model against itself or against arithmetic. They can
detect a transposed index. They cannot detect a rate coefficient that is wrong
by a factor of three, because a uniformly wrong dataset satisfies every one of
them.

\emph{Validation} asks whether the model reproduces something established
independently of it. Published coefficients from another group's code, measured
cross sections, a tabulated limit. Only this class of check can reach the
atomic physics.

\emph{Analysis} asks what the model then says. That is Chapter~\ref{ch:results},
and none of it appears here.

\subsection{Severity, and how to measure it}

The useful property of a check is not whether it passes. It is the size of the
set of wrong models it would reject. Call that its severity.

The failure mode is the one a mass balance on a continuous stirred-tank
reactor shows. Compute the outlet stream as inlet minus consumption, then
verify that the balance closes. It closes to machine
precision, always, for any rate constant whatever, because the outlet was
constructed from the balance being tested. The check is real arithmetic and it
is entirely uninformative about the kinetics. A model can be wired correctly
and be wrong.

Severity is measurable, and two operations measure it.

\begin{enumerate}
\item \textbf{Provenance inspection.} Ask whether the quantity under test was
derived from the quantity doing the testing. If the de-excitation coefficients
were computed from the excitation coefficients through detailed balance, then a
detailed-balance check on them is arithmetic, not physics. Chapter~\ref{ch:atoms}
already states this for the rate dataset (Section~\ref{sec:two_energies}); this
chapter carries it through to the matrix.

\item \textbf{Fault injection.} Break the model deliberately, in a way that a
real mistake could plausibly break it, and see whether the check notices. A
check that survives a tenfold error in a rate coefficient has told you that a
tenfold error in that coefficient is invisible to it.
\end{enumerate}

Fault injection is the sharper of the two, because it produces a number rather
than an argument. It is used in Section~\ref{sec:conservation_severity} on the
conservation check, and its result is the least comfortable finding in this
chapter.

\subsection{The three grades used below}

Every check in this chapter is placed in one of three grades, and the grade is
stated wherever the check is reported.

\begin{description}
\item[Tautological.] Cannot fail except on an arithmetic bug, because the
quantity tested was constructed from the relation being tested. It confirms
that the code implements the construction. It is not evidence about the
physics, and reporting it as though it were makes the validation weaker, not
stronger, by displacing the checks that carry information.

\item[Weak.] Could fail in principle, but the tolerance is so much looser than
the measured value that only a gross error would trip it. A gate that passes
with a factor of $150$ in hand has told you the answer is not absurd.

\item[Severe.] A plausible wrong version of this model fails it. Either the
check has actually fired at some point in the project's history, or a specific
corruption is known to trip it.
\end{description}

Four quantities currently printed by this project's own scripts as validation
are tautological by this definition, and Table~\ref{tab:ladder} names them:
Gate~A, the column-sum conservation residual, the $\tanh$ bound, and Gate~E.
Sections~\ref{sec:conservation_severity} to~\ref{sec:severe_checks} give the
measurements behind each.

%-----------------------------------------------------------------------------
\section{The conservation check is exact by construction}
\label{sec:conservation_severity}
% QUESTION: does the column-sum identity detect a wrong rate coefficient?

Section~\ref{sec:conservation} established that each column of $\Lmat$ must sum
to the ionisation leak from the level that column represents, and reported the
measured residual of Eq.~\eqref{eq:conservation_measured} as
$3.61\times10^{-11}$, a relative residual (Section~\ref{sec:conservation}). The
$3.61\times10^{-11}$ is the maximum over every column at every one of the $400$
grid points; the fault-injection run below has its own baseline, $2.238\times
10^{-12}$, at the single point it was run at. The claim attached to it there is that ``an error in a
single off-diagonal element, a transposed index, or a missing back-reaction all
break this immediately.''

That claim is testable, and it is false.

\subsection{Four injected faults}

Four faults were introduced into the assembly, one at a time, and the residual
recomputed after each. The first three are the exact failure modes the claim
names. Table~\ref{tab:faultinject} gives the outcome.

\begin{table}[tbp]
  \centering
  \caption[Fault injection into the conservation check]{Four deliberate faults
  injected into the assembly of $\Lmat$, with the resulting change in the
  matrix and the conservation residual of
  Eq.~\eqref{eq:conservation_measured} recomputed after each. Three faults
  large enough to move individual matrix elements by up to
  $2.7\times10^{11}\,\si{\per\second}$ leave the residual unchanged at the
  baseline value; only an inconsistency between the Einstein coefficients $A$
  and the total radiative loss rates $\gamma$ is detected. Source:
  \texttt{outputs/CHANGE\_REPORT.md}~\S4.2.}
  \label{tab:faultinject}
  \begin{tabular}{lrl}
    \toprule
    injected fault & $\max|\Delta \Lmat|$ (\si{\per\second}) & residual \\
    \midrule
    none (baseline)                     & $0$                  & $2.238\times10^{-12}$ \\
    $K_{\rm exc}(1s\to2p)$ $\times$ ten   & $6.9\times10^{5}$    & $2.238\times10^{-12}$ \\
    all de-excitation deleted           & $1.9\times10^{11}$   & $2.171\times10^{-12}$ \\
    excitation array transposed         & $2.7\times10^{11}$   & $1.837\times10^{-12}$ \\
    $A_{pq}$ halved, $\gamma_p$ left alone & $3.1\times10^{8}$   & $3.63\times10^{1}$ \textbf{caught} \\
    \bottomrule
  \end{tabular}
\end{table}

A tenfold error in the single most important excitation coefficient in the
model, the one that supplies the entire ground-fed channel of
Section~\ref{sec:R_depends_on_b}, changes the residual in no digit. Deleting
every de-excitation process in the matrix, which is not a subtle mistake,
changes it in the third digit and leaves it three orders of magnitude below any
tolerance. Transposing the whole excitation input array, the classic
index-order mistake, does the same.

\subsection{Why, and what the check is actually worth}

The reason is in the assembly, not in the physics. The diagonal of $\Lmat$ is
built as minus the column sum of the same off-diagonal array that the check
then sums. Whatever numbers are in that array, correct or corrupted, the column
sum returns the ionisation leak. The identity of Eq.~\eqref{eq:conservation} is
imposed by the construction rather than tested by it. This is the reactor mass
balance of Section~\ref{sec:gate_philosophy} exactly.

The one fault that was caught is the one that breaks the construction. In the
notation of Eq.~\eqref{eq:level_balance}, $A_{pq}$ is the Einstein coefficient
for spontaneous emission from level $p$ to level $q$, in \si{\per\second}, and
$\gamma_p = \sum_{q<p} A_{pq}$ is the total radiative loss rate from $p$, also in
\si{\per\second}. The individual $A_{pq}$ enter the off-diagonal elements while
their sum $\gamma_p$ enters the diagonal, so halving one without the other makes
the two inconsistent, and the residual rises from $2.238\times10^{-12}$ to
$36.3$, thirteen orders of magnitude. The check raises rather than warns, it
costs nothing to run, and on this evidence that is the width of what it covers.

The rate magnitudes are reached only by Section~\ref{sec:atomic_uncertainty}'s
comparison against RMPS and by the external comparisons of
Section~\ref{sec:fujimoto}, and by nothing internal.

\paragraph{The grid point of the table.} Three of the four
$\max|\Delta\Lmat|$ values in Table~\ref{tab:faultinject} scale with $\nel$, so
the table fixes its own grid point. Sweeping all $400$ grid points and matching on those three values
identifies $[23,5]$, the benchmark, with a worst disagreement of
\SI{1.63}{\percent}, the rounding of the printed table to two significant
figures; every residual in Table~\ref{tab:faultinject} then reproduces to the
digits shown (\texttt{verify\_fault\_injection.py}, stamped in
\texttt{validation/fault\_injection/}). The script rebuilds $\Lmat$ through the
pipeline's own assembly rather than patching the stored matrix, and the rebuild
reproduces \texttt{L\_grid} exactly before anything is injected, which makes the
faults faults in the model rather than in a re-implementation of it.

%-----------------------------------------------------------------------------
\section{The gate suite, read one gate at a time}
\label{sec:gates_abc}
% QUESTION: for each automated gate, what quantity does it compare, against
% what tolerance, and what would have to be wrong for it to fail?

Five automated gates run over the grid in \texttt{src/validation/validate\_gates.py}
and write \texttt{validation/gate\_summary.txt}. Their headline outcome is four
passes and one failure. Read as a headline it says little, because the five
differ in severity by a wide margin and two of them cannot fail. This section
takes them one at a time; Gate~D, the failure, gets
Section~\ref{sec:gate_d} to itself.

\subsection{Gate A: detailed balance, and where its physics content actually
lives}

Gate~A takes every excitation rate coefficient $K_{\rm exc}(i\to j)$ in the
dataset, in \si{\cubic\centi\metre\per\second}, and predicts the reverse
coefficient from the Boltzmann relation between forward and reverse rates in a
Maxwellian electron distribution at temperature $\Te$,
\begin{equation}
  K_{\rm deexc}(j\to i)
  = K_{\rm exc}(i\to j)\,\frac{g_i}{g_j}\,
    \exp\!\left(\frac{E_j - E_i}{k_B\Te}\right),
  \label{eq:gateA}
\end{equation}
then compares the prediction against the stored reverse coefficient at three
temperatures spanning the grid. Here $g_i$ is the dimensionless statistical
weight of level $i$, $E_j-E_i$ is the excitation energy in
\si{\electronvolt}, and $k_B$ is Boltzmann's constant. The relation holds only
for a Maxwellian distribution, which Section~\ref{sec:maxwellian} established
and which is assumed throughout. The gate passes
if the maximum relative error is below \SI{0.01}{\percent}
(\texttt{validate\_gates.py}:118). Measured: below \SI{0.01}{\percent} over all
$819$ pairs at all three temperatures.

\textbf{Grade: tautological.} The stored de-excitation coefficients were
produced from the excitation coefficients through Eq.~\eqref{eq:gateA} in the
first place (\texttt{compute\_K\_CCC.py}, function \texttt{detailed\_balance}),
as Section~\ref{sec:two_energies} states. Gate~A confirms that the same formula
was applied twice with the same energies and weights. That is a wiring check,
and Chapter~\ref{ch:atoms} already flagged it as one.

The physics content of detailed balance in this project sits somewhere else and
is currently invisible in the outputs. \texttt{src/parsers/qc\_ccc.py}:87--168
tests whether Bray's \emph{raw} convergent-close-coupling (CCC) cross sections,
which this project did not produce, satisfy microscopic reversibility on their
own. The ratio of the two directions has mean $0.9995$ with standard deviation
$0.0032$, and \SI{98.7}{\percent} of comparisons fall within \SI{1}{\percent}
(\texttt{ccc\_qc\_report.png}). A \SI{0.05}{\percent} mean bias with
\SI{0.32}{\percent} scatter on external data is a genuine external check, and it
is the one to lead with.

\subsection{One energy ladder on both sides, which is severe}

There is a second detailed-balance result, and unlike Gate~A it could have
failed.

A model of this kind carries energies twice: once in the Boltzmann factors that
relate forward and reverse collisional rates, and once in the Saha factors that
relate ionisation to three-body recombination. If two different energy tables
were used, or one table were indexed differently from the other, the model would
still conserve particles, still satisfy Gate~A, and still produce plausible
populations. Nothing internal would notice.

Testing ionisation against three-body recombination through the Saha relation
at every level, with CODATA constants on the Saha side, gives a ratio of
$0.999997$ for every one of the $43$ states at every temperature, identical
across all $2150$ entries to $1.4\times10^{-15}$
(\texttt{verify\_saha\_balance.py}, \texttt{validation/saha\_balance/}); the
residual is the rate module's rounded values of $h$, $m_e$ and the
electron-volt against CODATA, and nothing else. The point is not the
closeness to unity, which detailed balance imposes. It is the identity of the residual across
levels whose Saha exponents span the full binding-energy ladder of hydrogen,
from $\mathrm{e}^{13.6\,\mathrm{eV}/k_B\Te}$ for the ground state down to
$\mathrm{e}^{0.06\,\mathrm{eV}/k_B\Te}$ for the top shell. At
$\Te = \SI{1}{\electronvolt}$ those two factors differ by
$\mathrm{e}^{13.54}$, a factor of $7.6\times10^{5}$. Two energy ladders that
differed would produce a residual drifting across that range; the measured
residual does not drift, and its common value is the eight-figure rounding of
the stored state table. \textbf{Grade: severe.} The check could have failed and
did not.

The same audit gives the collisional side over the $819$ excitation pairs at
each of three temperatures, $2457$ comparisons in total: maximum deviation
$7.31\times10^{-9}$, median $7.5\times10^{-10}$, again consistent with the state
table's rounding.

\subsection{Gate B: the coronal limit}

At low enough density an excited level is populated by electron impact from the
ground state and empties radiatively before a second collision reaches it, so
its population per unit ground-state density and per unit electron density
becomes independent of $\nel$. Gate~B tests that plateau: it computes
$n_{2p}/(\ngs\nel)$ and $n_{3d}/(\ngs\nel)$ at the two lowest densities on the
grid and asks whether they have converged.

The tolerance is \SI{50}{\percent} per temperature
(\texttt{validate\_gates.py}:175) and the gate passes if \SI{70}{\percent} of
temperatures pass (line 181). Measured: all $50$ temperatures pass, worst
convergence $0.0034$.

\textbf{Grade: weak.} Passing with a factor of $150$ between the measured
convergence and the tolerance means the gate rejects only a gross error. It
does establish that the low-density end of the grid is in the coronal regime,
which Chapter~\ref{ch:results} relies on when it reads the low-density column
as ionising, and that is worth having. It is not evidence about rate
magnitudes.

\subsection{Gate C: what the Saha gate tests, which is not the Saha limit}

Gate~C is named for the high-density limit in which collisions dominate
radiation and populations approach their Saha--Boltzmann values. It computes
three quantities per temperature: whether $n_{2p}/\ngs$ increases monotonically
with density, whether it stays below the Saha value (with \SI{5}{\percent}
slack), and whether it approaches that value, defined as reaching $10^{-4}$ of
it.

The pass flag uses the first two and \textbf{excludes the third}
(\texttt{validate\_gates.py}:250; the \texttt{approaching} column is computed at
line 239, stored, printed at line 261, and never consulted). The exclusion is
not what makes the gate weak, because the excluded column would pass: the
measured approach fraction runs $0.0123$ to $0.0155$, comfortably above the
$10^{-4}$ the criterion asks for. What makes it weak is that the criterion is
$150$ times slacker than the measurement. Reaching $10^{-4}$ of Saha is not
approaching Saha, and neither is reaching \SI{1.5}{\percent}.

The measurement itself is the useful part. At
$\nel = 10^{15}\,\si{\per\cubic\centi\metre}$, the top of the grid, the $2p$
population sits at about \SI{1.5}{\percent} of its Saha--Boltzmann value, so
that level is nowhere near local thermodynamic equilibrium (LTE) anywhere on
this grid. Only $2p$ was tested, so this is a statement about $2p$ and not about
the whole distribution. It is the expected behaviour for an optically thin
plasma whose ion population is an externally imposed reservoir rather than a
Saha partner, and Chapter~\ref{ch:results} depends on it: it is why the two
supply channels of Section~\ref{sec:R_depends_on_b} never merge, and why
$\ffrac{3}-\ffrac{4}$ does not vanish at high density.

\textbf{Grade: weak, and mis-named.} What Gate~C tests is monotonicity in
density and a Saha ceiling. Both would fail on a sign error in the recombination
source, which is worth something. Neither tests the limit the gate is named for.

%-----------------------------------------------------------------------------
\section{Gate E: the timescale hierarchy}
\label{sec:gate_e}
% QUESTION: is the two-clock structure of Chapter 3 present everywhere on the
% grid, and how demanding is the test that says so?

Gate~E diagonalises $\Lmat$ at every grid point, orders the negative
eigenvalues, and forms $\Msep = \tauslow/\taurel = |\lambda_1|/|\lambda_0|$. It
passes if $\Msep > 1$ at \SI{95}{\percent} of points
(\texttt{validate\_gates.py}:409--412).

Measured over all $400$ points: $\Msep$ ranges from $86.77$ at
$[49,7]$ ($\Te = \SI{10}{\electronvolt}$,
$\nel = \SI{e15}{\per\cubic\centi\metre}$) to $1.72928\times10^{9}$ at the cold
corner; $\tauslow$ spans \SI{75.4}{\nano\second} to \SI{67.2}{\second} and
$\taurel$ spans \SI{0.87}{\nano\second} to \SI{38.9}{\nano\second}
(project working note; Table~\ref{tab:ladder}). The gate asks for $\Msep>1$. The smallest value
anywhere is $86.8$.

\textbf{Grade: tautological as posed, and the real content is elsewhere.} The
threshold sits between two and seven orders of magnitude below the measured
values. No matrix this construction has produced comes within a factor of $86$
of failing, and none plausibly could, because the fast radiative rates and the
slow ionisation rate that set the two eigenvalues are separated by construction
of the physics rather than by a numerical accident.
The reportable quantity is the range, not the pass, together with the fact that
three independent implementations produce it bit-for-bit: \texttt{qss\_analysis.py},
Gate~E itself, and an unconditional recomputation stored in
\texttt{timescales\_unfiltered\_CHECK.npz}. That agreement is what caught the
eigenvalue-filter error of Section~\ref{sec:errors_found}, and a check that has
actually fired is worth more than a threshold nothing approaches.

Two properties of these numbers must travel with them, and both come from
Chapter~\ref{ch:theory}. $\tauslow$ is boundary-dependent, which
Section~\ref{sec:attacks} quantifies. And the slow eigenvalue is not an
equilibration rate: measured against the ground-state ionisation rate
coefficient $K_{\rm ion}(1s)$, in \si{\cubic\centi\metre\per\second}, it
satisfies $|\lambda_0| = 2.286\,K_{\rm ion}(1s)\,\nel$ across the grid
(project working note; Table~\ref{tab:ladder}), so $\tauslow$ is the collisional--radiative
ionisation lifetime of the ground-state reservoir and not the time for anything
to come to equilibrium.

\subsection{The gates do not stop anything}

No gate raises on failure. Each prints \texttt{PASS} or \texttt{FAIL} and
returns a flag; \texttt{main} sets \texttt{all\_pass = False}, prints a warning,
and exits normally. \textbf{The script returns exit code zero whether or not
Gate~D fails.} The residual gates elsewhere in the project do raise; this one
does not, and that asymmetry is a defect.

%-----------------------------------------------------------------------------
\section{The checks that could have failed}
\label{sec:severe_checks}
% QUESTION: which checks in this project would a plausible wrong model fail?

Everything so far has been a demotion. This section is the other half, and it
is short by design: a small number of severe checks is worth more than a long
list of consistency tests, and these are the ones the thesis leans on.

\subsection{The superposition residual}
\label{sec:superposition}

The entire mechanism of Chapter~\ref{ch:results} rests on one structural claim:
that the excited populations $\nF$ of Section~\ref{sec:qss} split exactly into a
ground-fed and a recombination-fed part,
$\nF = \bm{a}\,\ngs + \bm{c}\,\nion$, where the network responses $\bm{a}$ and
$\bm{c}$ of Section~\ref{sec:R_depends_on_b} are dimensionless and independent
of the two reservoir densities. If that split is not exact, the ground-fed
fractions $\ffrac{p}$ of Section~\ref{sec:ground_fed_fraction} are not well
defined and the closed form for the diagnostic error does not exist.

The check solves the full system three times at each grid point and each step
direction, once with each reservoir alone and once with both, and measures the
relative residual of the sum against the joint solve. Across the $784$
grid-point and step-direction pairs it covers, out of the $800$ the grid admits,
the largest residual is $3.075\times10^{-14}$
(\texttt{plateau\_gridmap.txt}, line 19).

\textbf{Grade: severe}, for a specific reason. Feeding the check a corrupted
input in which the $n=3$ and $n=4$ entries are swapped, which is the mistake
that would most easily go unnoticed in a two-shell observable, trips it
(project working note; Table~\ref{tab:ladder}). It is one of only two checks in the
project reported to catch that swap.

\subsection{The reduced-versus-full ratio test}

The second is its companion. The observable $\Robs = n_3/n_4$ can be computed
two ways: from the full $43$-state steady-state solve, and from the reduced
two-channel form of Eq.~\eqref{eq:R_of_b} with the ground-fed fractions
substituted in. The two routes share no code path beyond the matrix itself. At
the benchmark point $\Te = \benchTe$, $\nel = \benchne$, an independent
re-implementation obtained $\Robs = 0.777768$ by three routes against
Chapter~\ref{ch:theory}'s recorded $0.7778$
(project working note; Table~\ref{tab:ladder}). A $3\leftrightarrow4$ swap breaks the
agreement.

\subsection{Pinned hashes on the inputs}

The most common way for a computation like this to become wrong is not an
algebra error. It is a stale array: a matrix regenerated after a correction,
combined with an analysis output produced before it. This project has three
retractions in its history and at least one of them had that shape.

\texttt{src/validation/preflight.py} pins the SHA-256 hashes, fixed-length
fingerprints that change if any byte of a file changes, of the two canonical
inputs and compares them byte for byte before any verification script
runs. Recomputed for this chapter:
\begin{center}
\begin{tabular}{ll}
  \toprule
  file & SHA-256 (first 16 hex digits) \\
  \midrule
  \texttt{data/processed/cr\_matrix/L\_grid.npy} & \texttt{2d92b58e1693107d} \\
  \texttt{data/processed/cr\_matrix/S\_grid.npy} & \texttt{7822f536590c76ba} \\
  \bottomrule
\end{tabular}
\end{center}
Both match the pinned baselines at \texttt{preflight.py}:37--38 exactly, so the
matrix on disk at the time of writing is the one the pins were taken from; the
statement is only as good as the discipline of running the preflight before
each analysis.

\textbf{Grade: severe on the data, incomplete on the scripts.} Of the six
scripts the same preflight claims to check, \texttt{verify\_ch3\_groupB.py} and
\texttt{verify\_grid\_coverage.py} do not exist on disk under those names; a
missing script is downgraded to a non-blocking advisory
(\texttt{preflight.py}:258--262) while the summary still prints ``All checks
passed''. The hash check is sound; the sentence after it overstates what ran.

\subsection{Guards that recompute rather than trust}

\texttt{src/analysis/make\_ch5\_figures.py} does not plot the stored results. It
rebuilds every plotted quantity from $\Lmat$ and the recombination source
$\Svec$ of Eq.~\eqref{eq:source} and compares against
the stored comma-separated tables at a relative tolerance of $10^{-9}$, raising
with the offending grid point, column, and matrix hash if any row disagrees
(lines 393--432). It reads the temperature step size out of the file header
rather than redeclaring it, so a figure cannot be drawn against a step
convention different from the one that produced the data. It also verifies that
the untrapped run of the radiation-trapping sweep reproduces the canonical
matrix exactly, without which no trapped number in Chapter~\ref{ch:limits} would
be interpretable as a difference.

The same script carries a negative control: it computes the correlation between
$\Msep$ and the diagnostic error both raw and after controlling for temperature
and density, and refuses to draw the figure if the controlled correlation fails
to change sign (Section~\ref{sec:M_does_not_predict} draws the conclusion).

\subsection{Two checks that fired, and one that was withdrawn}

Three episodes say more than the pass/fail table, because in each the check did
something.

The Lyman-$\alpha$ line-centre cross-section gate, an independent implementation
in \texttt{escape\_factor.py} against the project's own Einstein coefficients,
fired at \SI{1156}{\percent} on its first run, located a $4\pi$ factor, and now
agrees to \SI{0.059}{\percent} against a raising \SI{1}{\percent} tolerance
(Section~\ref{sec:opacity}).

\texttt{verify\_fujimoto\_table41.py} carries a deliberately broken variant with
proton-impact $\ell$-mixing removed. It returns a recombination-fed population
coefficient $r_0(2) = 104.5$, and the script flags it itself as unphysical on
the ground that a level cannot exceed Saha equilibrium. Production gives
$0.7392$. A check that fires on a known-bad input is the definition of a severe
one.

Third, a semigroup propagator identity,
$\exp(\Lmat t_1)\exp(\Lmat t_2) = \exp(\Lmat(t_1+t_2))$ for a matrix exponential
computed by the same routine, returned $0.000\times10^{0}$ by construction and
was \emph{withdrawn by the author} (Section~\ref{sec:bridge}), the judgement
this chapter applies to Gate~A, applied earlier and unprompted.

\subsection{The bound that is a theorem}

Section~\ref{sec:bound} derives an upper bound on the diagnostic error of the
form $\tanh(\Delta/4)$, where $\Delta$ is the separation between the two shells'
switching points. A script reports that the bound is honoured at all $248$
points tested, and that line reads as validation.

It is not. Here $\Delta$ is the dimensionless separation between the two shells'
switching points, in units of the logarithm of the reservoir strength. Given
$\bm{a},\bm{c}\ge0$, which Section~\ref{sec:R_depends_on_b} proves from the
M-matrix structure of $-\Lmat_{FF}$, the bound is a theorem: it cannot fail
except on an arithmetic bug. Fed corrupted input it still passes. Swapping
shells $3\leftrightarrow4$ sends $\Delta\to-\Delta$ and leaves both $|\Delta|$
and the logarithmic change in $\Robs$ that it bounds unchanged, so the gate
cannot see it. Swapping $c_3\leftrightarrow c_4$ also
passes, but for a different reason: it moves $|\Delta|$ from $1.945614$ to
$0.939493$, a factor $2.07$, and the gate still passes because it compares
$|f_3-f_4|$ against $\tanh(|\Delta|/4)$ evaluated from the same corrupted
$\bm{c}$. The bound is satisfied by whatever $\bm{a},\bm{c}\ge0$ it is
handed.
\textbf{Grade: tautological.} The script's own docstring calls it a wiring
check; the line it prints does not, and the printed line is what a reader sees.

The bound is a result, and a useful one, because it holds for any atomic dataset
whatever. It is not evidence that this atomic dataset is right.

%-----------------------------------------------------------------------------
\section{Conditioning, precision, and state-space convergence}
\label{sec:conditioning}
% QUESTION: could the answers be artifacts of finite-precision arithmetic or of
% where the level ladder was cut off?

Two numerical questions sit underneath every result in this thesis. The
elimination of Section~\ref{sec:qss} inverts a $42\times42$ matrix whose entries
span many orders of magnitude, which invites the question of whether the inverse
means anything. And the level ladder was cut at $n_{\max}=15$, which invites the
question of what the discarded levels would have changed.

\subsection{The inverse is well behaved, and here is how well}

The relevant quantity is the condition number of $\Lmat_{FF}$, the block of
$\Lmat$ coupling the $42$ excited levels to each other
(Section~\ref{sec:qss}), and it is the factor by which a relative perturbation
of the matrix can be amplified in the solution.
Computed from the canonical $\Lmat$ at all $400$ grid points by
\texttt{src/validation/verify\_operator\_conditioning.py} and stamped to
\texttt{validation/operator\_conditioning/operator\_conditioning.csv}: minimum
$1.4821\times10^{3}$ at $[0,0]$ ($\Te = \SI{1.000}{\electronvolt}$), maximum
$1.7436\times10^{5}$ at $[33,7]$ ($\Te = \SI{4.715}{\electronvolt}$), median
$2.1044\times10^{4}$. In double precision, which carries about sixteen
significant digits, a condition number of $10^{5}$ costs five of them. Eleven
remain, and every conclusion in this thesis turns on the first three.

The range was first asserted in Chapter~\ref{ch:theory} on an unscripted
working note, which is not a check; the stamped script reproduces it to
\SI{0.14}{\percent} and \SI{0.21}{\percent}. The number was right the whole
time and was not evidence until it could be re-run.

Three further measurements bound the arithmetic. The eigenvalue condition
number is a different quantity from the matrix condition number above: it is
$1/|\bm{y}^H\bm{x}|$ for a left and right eigenvector pair, and it measures how
far a single eigenvalue moves under a perturbation of the matrix. For the two
eigenvalues the thesis uses it is of order unity
($1.71$ and $1.98$ at the benchmark point, $3.26$ and $2.59$ at the cold
corner), despite a largest absolute column sum of
$2.0\times10^{14}\,\si{\per\second}$; that large norm is a
statement about the fast radiative modes and does not propagate into
$\lambda_0$ or $\lambda_1$. Recomputing the two supply channels in single
precision moves them by $1.3\times10^{-5}$, so double precision is nowhere near
a cliff. And tightening the integrator changes nothing that matters: a relative
tolerance of $10^{-9}$ against $10^{-12}$ shifts the result by
$1.20\times10^{-8}$, and computing the propagator by a dense matrix exponential
against a Krylov action differs by $1.13\times10^{-10}$.

The eigenvalue condition numbers and the single-precision test cover three of
$400$ points; the tolerance comparisons cover two. The condition numbers and the
displacement are recomputed under the definition $1/|\bm{y}^H\bm{x}|$ with
unit-norm left and right eigenvectors by \texttt{verify\_residual\_claims.py}
(\texttt{validation/residual\_claims/}). They are point measurements, not grid
scans, and are quoted as such.

\subsection{Where the ladder was cut}
\label{sec:convergence}

The quantity that has to converge is not a population but the difference in
ground-fed fraction between the two Balmer shells, $\ffrac{3}-\ffrac{4}$, since
Chapter~\ref{ch:results} shows the diagnostic error to be proportional to it.

Rebuilding $\Lmat$ with successively lower $n_{\max}$, from $15$ down to $9$,
with the loss channels of every removed state returned to the diagonal of the
survivors, and recomputing $\ffrac{3}-\ffrac{4}$ at each truncation's own
$u^{\mathrm{CRE}}$ (\texttt{verify\_nmax\_downward\_scan.py},
\texttt{validation/nmax\_downward\_scan/}) gives increments of one sign at the
three named points, with the last measured ratio $d(15)/d(14)$ equal to
$0.88$ at the benchmark, $0.72$ at the cold corner $[0,0]$ and $0.33$ at the
ridge $[15,3]$; at the $245$ of the $280$ defended points where the increments
are one-signed and geometric the ratio has median $0.67$. Summing the geometric
tail with that ratio gives an extrapolated truncation error of
\SI{-0.46}{\percent} at the benchmark, \SI{+1.6}{\percent} at the cold corner
and \SI{+0.005}{\percent} at the ridge; holding $u$ at its $n_{\max}=15$ value
instead gives \SI{-0.38}{\percent}, \SI{+1.4}{\percent} and
\SI{+0.005}{\percent}, so the sign does not depend on that choice. A historical
project value with no producing script, not used as evidence here and superseded
by \texttt{verify\_nmax\_downward\_scan.py}, records
\SI{+0.9}{\percent}, \SI{-1.4}{\percent} and \SI{-0.55}{\percent} with a
common ratio of about $0.75$; the stamped run reverses both signs and the
conclusion that rested on them. Truncation makes
$\ffrac{3}-\ffrac{4}$ slightly too large at the benchmark and too small at the
cold corner, so it is not a bias that can be corrected by a single factor.
Over the defended points where the tail can be summed the extrapolated error
has median \SI{+0.002}{\percent} and runs from \SI{-0.82}{\percent} to
\SI{+1.03}{\percent}, smaller than the spread between the two step conventions
of Section~\ref{sec:step_convention} at every such point. At the $35$ defended
points of the density column $\nel = \SI{5.2e13}{\per\cubic\centi\metre}$ the
increments change sign among $n_{\max} = 12$ to $15$ and no tail can be
summed; the last measured increment there is below \SI{0.041}{\percent} of
$\ffrac{3}-\ffrac{4}$ at every point.

\paragraph{A trap in this test.} The first attempt dropped the top shells from
the state vector without removing their loss channels from the diagonal of the
surviving states, which then lost population to states that no longer existed,
and $\ffrac{3}-\ffrac{4}$ jumped. The stamped scan repeats the mistake
deliberately: dropping $n=15$ alone without the correction moves
$\ffrac{3}-\ffrac{4}$ by \SI{11.5}{\percent} at the benchmark,
\SI{12.6}{\percent} at the cold corner and \SI{5.0}{\percent} at the ridge,
where the correct truncation to $n_{\max}=14$ moves it by \SI{0.06}{\percent},
\SI{0.64}{\percent} and \SI{0.01}{\percent}; dropping six shells without it
gives \SI{17.9}{\percent}, \SI{25.9}{\percent} and \SI{13.3}{\percent} (a
working note recorded \SIrange{22}{37}{\percent}). The jump is an artifact of
the truncation procedure, not a dependence on $n_{\max}$. Removing state $k$
requires adding $\sum_{k\ \rm dropped} L_{ki}$ back to $L_{ii}$ for every
surviving $i$, as the production matrix does. Anyone repeating this test will
reproduce the spurious number first.

\paragraph{The terminal shell is over-populated, and by a measured amount.}
Truncating the ladder at $n_{\max}$ removes, from the balance of the terminal
shell, the collisional excitation out of it to $n > n_{\max}$ and the
de-excitation and cascade into it from those levels (cascade is downward, so
what is lost is a way out, not something ``to cascade out to''). The loss removed is not small. Continuing the $12\to13$, $13\to14$
and $14\to15$ excitation coefficients geometrically, the missing $15\to16$
channel would carry $1.30$ times the whole retained loss out of $n=15$ at the
benchmark and $1.20$ times at $[0,4]$, between $1.20$ and $1.35$ times at
every grid point, while radiative decay is below $10^{-6}$ of that loss and
ionisation is $7$ to \SI{10}{\percent} of it
(\texttt{validation/terminal\_shell\_budget/}). The top shell sits within
$6\times10^{-5}$ of Saha--Boltzmann, so for the total population the removed
exchange is nearly balanced and the sign of the bias is decided channel by
channel: the ground-fed channel, whose net flux through the top of the ladder
is upward, loses only an exit and can only be biased high, and it is the
ground-fed coefficient $r_1(15)$ that the external comparison finds high while
$r_0(15)$ agrees (Section~\ref{sec:fujimoto}). Comparing each shell's ground-fed coefficient $a_p$ as terminal
shell (operator rebuilt with $n_{\max}=p$, the loss channels into the removed states kept on the diagonal) with its
value in the full $n_{\max}=15$ operator gives at the Fujimoto point $[49,0]$ (\SI{10}{\electronvolt},
\SI{e12}{\per\cubic\centi\metre}) excesses of $9.4$ at $p=8$, $10.7$ at $p=9$ and $6.3$ at $p=10$, then $4.4$,
$3.0$, $2.1$ and $1.4$ at $p=11$ to $14$ (\texttt{verify\_terminal\_vs\_interior.py},
\texttt{validation/terminal\_vs\_interior/}). At the benchmark $p=8$ to $10$ give $14.2$, $12.3$ and $6.5$; over
the $280$ defended points the $p=8$ excess spans $8.2$ to $17.7$ and the $p=10$ excess $5.4$ to $7.5$. The excess
is confined to the ground-fed channel: the shell's total population and recombination-fed coefficient move under
\SI{1}{\percent} at the benchmark and at most \SI{32}{\percent} and \SI{30}{\percent} at the low-density named
points. A historical project value with no producing script, not used as evidence here and superseded by
\texttt{verify\_terminal\_vs\_interior.py}, records $9.4$, $10.7$ and
$6.3$ with neither quantity nor grid point named; the stamped run reproduces them to \SI{0.4}{\percent} as $a_p$
at $[49,0]$. The absence of a monotone trend over $p=8$ to $10$ holds there and at $47$ of the $280$ points;
from $p=9$ upward the excess falls at every grid point, at the benchmark from $p=8$. The external excess at the
terminal shell, $4.9$ to $6.2$ against published values (Section~\ref{sec:fujimoto}), lies outside the $6.3$ to
$10.7$ band, not inside it as the working note states, but within a factor $2.2$ of every value in it, and both rise
with density ($4.9$ to $6.2$ externally, $9.4$ to $17.7$ at $p=8$ internally at \SI{10}{\electronvolt}). That
agreement in magnitude and in the sign of the density trend makes truncation a sufficient explanation for the
$n=15$ discrepancy without invoking the atomic data. The fall with $p$ cannot be extrapolated to $p=15$: the
reference for $p=14$ is itself a ladder with one shell above it, so the internal measurement neither predicts nor
excludes $4.9$ to $6.2$ there.
No result in this thesis rests on the
top shell: the observable is built from $n=3$ and $n=4$, both fully resolved in
orbital angular momentum $\ell$ and eleven shells below the cut at
$n_{\max}=15$.

\paragraph{What has now been checked.} Above $n=8$ the shells are bundled and
their internal $\ell$-distribution is assumed statistical rather than computed.
The script that tests that assumption, \texttt{verify\_bundling\_psm20.py}, had
never been executed; two of the three defects found on its first run, the
synthesised temperature grid and the identically zero $\ell$-mixing array for
the bundled indices, were predicted here. Section~\ref{sec:bundling_check}
reports all three defects and the result.

%-----------------------------------------------------------------------------
\section{Is the shell ratio the line ratio?}
\label{sec:balmer_equivalence}
% QUESTION: the analysis computes n_3/n_4, but a spectrometer measures
% Halpha/Hbeta. What does the substitution cost?

Every analytic result in Chapter~\ref{ch:theory} is written for the shell ratio
$\Robs = n_3/n_4$, because the two-channel decomposition applies to shell
populations. A spectrometer does not measure shell populations. It measures the
$H_\alpha$ and $H_\beta$ line intensities, each of which is a sum over the
electric-dipole-allowed sublevel transitions, weighted by their Einstein
coefficients. Section~\ref{sec:radiative} established that these are not
proportional: $4f$ holds \SI{43.75}{\percent} of the $n=4$ shell by statistical
weight and contributes nothing to $H_\beta$, since $4f\to2p$ would change $\ell$
by two and is forbidden.

The substitution therefore has to be measured, not assumed.

The line intensities are built as $A$-weighted sums over resolved channels,
$H_\alpha$ from $3s\to2p$, $3p\to2s$ and $3d\to2p$, and $H_\beta$ from
$4s\to2p$, $4p\to2s$ and $4d\to2p$, with populations taken from the solve.
Nothing statistical is assumed anywhere in that construction: searching the two
producing scripts for statistical-weight expressions returns no matches, and the
$A$-value lookup raises unless exactly one row matches each transition.

Two measurements follow. The $4f$ fraction of the $n=4$ shell runs from
$0.4361$ to $0.4375$ across the grid, against the statistical value
$14/32 = 0.4375$; proton-impact $\ell$-mixing drives the shell to a statistical
distribution to better than \SI{0.3}{\percent} everywhere, with the departure
concentrated at $\nel = \SI{e12}{\per\cubic\centi\metre}$ where $\ell$-mixing
weakens and radiative decay competes. That fraction is a property of the full
equilibrium population; the quantity that sets the line-versus-shell difference
below is the $\ell$-distribution within each supply channel separately, which
is not statistical at the low-density edge. And the ratio of the transient error
computed on the line ratio to the same error computed on the shell ratio runs
from $0.9482$ to $0.9999$ over the $784$ point-and-direction pairs, below $1$ at
every one: \textbf{worst case \SI{5.2}{\percent}, at the heating step from
$9.54$ to \SI{10.0}{\electronvolt} at the lowest density}
(\texttt{verify\_weighted\_census.py}, stamped to
\texttt{validation/weighted\_census/}). At
$[0,4]$, the point carrying the largest error anywhere, the two agree to
\SI{0.06}{\percent} ($0.386683$ against $0.386903$). Because the ratio is below
$1$ everywhere, the shell census of Chapter~\ref{ch:results} is conservative
for the line observable: counted on the line ratio, $44$ instead of $45$ of the
$448$ pairs above \SI{2}{\electronvolt} exceed \SI{10}{\percent} at
\SI{100}{\micro\second}, with the worst case $0.1746$ instead of $0.1748$ at
the same point. The \SI{5.2}{\percent} scales inversely with the proton
$\ell$-mixing rate: with the within-shell $\ell$-changing rates scaled by a
common factor and the column sums of the operator preserved, the corner
departure is \SI{10.0}{\percent} at half the rate, \SI{2.6}{\percent} at double
and \SI{1.1}{\percent} at five times, while the benchmark line-to-shell ratio
stays between $0.9987$ and $0.9999$ (\texttt{validation/weighted\_census/},
$\ell$-mixing sensitivity block and \texttt{weighted\_census\_lmix.csv}); and
it sits on the grid boundary.

That \SI{0.06}{\percent} is a single-point figure and must not be quoted as a
grid-wide one. The grid-wide statement is the \SI{5.2}{\percent}.

\subsection{And is the plasma transparent to its own Balmer light?}

The shell-to-line substitution assumes the emitted photons escape. For the
Lyman series that assumption fails at the cold edge and
Chapter~\ref{ch:limits} treats it at length. For the Balmer lines themselves it
closes here, because it bears on the observable rather than on the matrix.

Section~\ref{sec:opacity} defines the optical depth $\tau$ and the escape factor
and treats the Lyman series, where transparency fails. For the Balmer lines the
absorbers are the $n=2$ atoms, $n_{2s}+n_{2p}$ in collisional--radiative
equilibrium with $\nion=\nel$, the Doppler width is taken at $\Tn=\Te$, and the
multiplet oscillator strengths are $f(2\to3)=0.641$ and
$f(2\to4)=0.119$~\cite{WieseFuhr2009}. The line-centre optical depth of
$H_\alpha$ across a \SI{10}{\centi\metre} slab is then $4.9\times10^{-7}$ at the
cold corner, $8.2\times10^{-4}$ at $[0,4]$ and $0.214$ at $[0,7]$, the single
worst cell on the grid; resolving the absorbers into $2s$ and $2p$ with the
model's own $A$ coefficients changes these by at most \SI{4}{\percent}. At
$[0,7]$ the population escape factor for a photon born at the slab centre
($\kappa_0 D/2 = 0.107$) is $0.928$ for $H_\alpha$ and $0.990$ for $H_\beta$, so
the observed ratio is lowered by \SI{6.3}{\percent} in that one cell; over the
full path ($\tau=0.214$) the figures are $0.861$ and \SI{12}{\percent}. Above
\SI{2}{\electronvolt} the largest $H_\alpha$ optical depth is $0.010$ at
$[15,7]$ and the ratio moves by less than \SI{0.7}{\percent} anywhere
(\texttt{verify\_balmer\_optical\_depth.py},
\texttt{validation/balmer\_optical\_depth/}). The Balmer lines are optically
thin over the region where this thesis makes quantitative claims. A historical
project value with no producing script, not used as evidence here and superseded
by \texttt{verify\_balmer\_optical\_depth.py}, records
$6.3\times10^{-7}$, $1.0\times10^{-3}$ and $0.263$ for the same three cells;
those values are $1.23$ to $1.28$ times the stamped ones and no convention
tested recovers them, and the escape factor above $0.85$ and ratio shift of at
most about \SI{10}{\percent} recorded with them are half-slab figures attached
to a full-slab optical depth. The conclusion holds with either set of values.

%-----------------------------------------------------------------------------
\section{Two tests of the reduction, and their expected outcomes}
\label{sec:attacks}
% QUESTION: the two objections most likely to destroy the result were tested;
% what was the falsifier in each case, and did it appear?

The checks so far ask whether the model computes what it intends. This section
asks something harder: whether two specific structural choices manufacture the
result. In each case the size of effect that would have destroyed the claim was
written down before the test was run, which is what makes the outcome
informative rather than reassuring.

\subsection{Does the open boundary manufacture the slow clock?}

$\Lmat$ describes an open system: atoms that ionise leave the accounting, the
recombination feed $\Svec\nion$ enters as a constant source rather than a matrix
element, and there is no forty-fourth state for ions to accumulate in. Since the
persistence claim of Section~\ref{sec:persistence} rests entirely on $\tauslow$
exceeding the duration of the disturbance, the first objection is that
$\tauslow$ is an artifact of the open boundary.

\textbf{The falsifier, named first.} Write $\taudrive$ for the duration of the
disturbance, fixed at \SI{100}{\micro\second} throughout this thesis as a
representative edge-localised-mode (ELM) time
(Section~\ref{sec:persistence}). At $[0,4]$, the worst point on the grid,
$\tauslow/\taudrive \approx 2300$. If closing the system shortened $\tauslow$ by
that factor, the slow clock would fall to $\taudrive$ itself and the persistence
would vanish. A factor of $2300$ refutes the claim.

Two grid labels appear in the table for the same worst case: the error is
evaluated at $[0,4]$, but the operator governing the decay after the step is the
post-step one, $\Lmat[1,4]$, and every timescale in this thesis names the
operator it belongs to (Chapter~\ref{ch:results}).

The manifold was closed by appending an ion state,
\begin{equation}
  \Lmat_c =
  \begin{pmatrix}
    \Lmat & \Svec \\
    -\operatorname{colsum}(\Lmat) & -\textstyle\sum\Svec
  \end{pmatrix},
  \label{eq:closure_ch4}
\end{equation}
so that every ionising atom lands in the new state instead of being deleted.
Each column of $\Lmat_c$ sums to zero to machine precision, and the construction
reproduces the closed-system value at the cold corner recorded in this
project's working notes, \SI{1.6226}{\second}, to five digits, which confirms it is the same
closure rather than a new one. The measured factors are given in
Table~\ref{tab:closure}.

\begin{table}[tbp]
  \centering
  \caption[Effect of closing the ion boundary on the slow timescale]{The slow
  timescale $\tauslow$ computed with the ion treated as an external reservoir
  (open) and with a forty-fourth ion state appended (closed), at four grid
  points. Over the whole grid the ratio has minimum $1.00$, median $1.00$ and
  maximum $41.4$: the boundary condition matters only where $\tauslow$ already
  exceeds the disturbance duration by three orders of magnitude. Every entry is
  reduced from \texttt{validation/ion\_closure/ion\_closure\_summary.csv}
  (\texttt{verify\_ion\_closure.py}) by
  \texttt{verify\_headline\_provenance.py}; the open column is the float64
  eigensolver, which at the cold corner differs from the artifact's 40-digit
  path by one part in $10^{4}$.}
  \label{tab:closure}
  \begin{tabular}{lrrr}
    \toprule
    point & $\tauslow$ open & $\tauslow$ closed & factor \\
    \midrule
    cold corner $[0,0]$ & \SI{67.23}{\second} & \SI{1.6226}{\second} & $41.4$ \\
    post-step operator $\Lmat[1,4]$ & \SI{0.23324}{\second} & \SI{0.015173}{\second} & $15.4$ \\
    ridge $[15,3]$ & \SI{2.027}{\milli\second} & \SI{1.999}{\milli\second} & $1.01$ \\
    benchmark $[23,5]$ & \SI{22.728}{\micro\second} & \SI{22.708}{\micro\second} & $1.00$ \\
    \bottomrule
  \end{tabular}
\end{table}

The shape of the distribution carries the argument: the factor is large only
where $\tauslow$ already exceeds $\taudrive$ by about a thousand, so the error
has already saturated, and it is unity everywhere $\tauslow$ approaches
$\taudrive$. Recomputing the entire persistence census with closed-system
$\tauslow$ moves the breakdown count from $202$ to $200$ of $680$ and the worst
case from $0.38682$ to $0.38563$, about one part in $325$; restricted to
$\Te \ge \SI{2}{\electronvolt}$, where Chapter~\ref{ch:results} makes
quantitative claims, the count is $45$ before and $45$ after.

The falsifier did not appear. A factor of $2300$ was needed; the measured factor
at that operator is $15$.

\textbf{The caveat, which must not be softened.} The closure adds an ion
reservoir; it does not close the electron balance, since $\nel$ is held fixed
while the ion population moves, and it adds no transport. It answers the
boundary condition on the atomic manifold and nothing else:
Section~\ref{sec:transport} states the objection it cannot answer, and
Section~\ref{sec:quasineutrality} the one it makes sharper rather than softer.

\subsection{Is the single-slow-mode estimate a bound on the time-averaged error?}

Chapter~\ref{ch:results} evaluates the error at the partial-equilibrium plateau
and builds from it the single-slow-mode estimate of Eq.~\eqref{eq:elm_bound},
which multiplies the plateau error $\epsplat$ of
Section~\ref{sec:eps_definitions} by the fraction of $\taudrive$ over which an
exponential decay on $\tauslow$ keeps it alive. It was first written as a lower
bound, on the assumption that the observable cannot decay faster than the
slowest mode. That assumption is testable, and an estimate quoted as though it
were the answer is a claim about accuracy that has to be measured.

Propagating the full $43$-state system through the step by direct
eigen-decomposition and integrating the error over \SIrange{0}{100}{\micro\second}
gives the comparison in Table~\ref{tab:lowerbound}.

\begin{table}[tbp]
  \centering
  \caption[Accuracy of the single-slow-mode estimate]{The plateau error, the
  true error averaged over \SIrange{0}{100}{\micro\second} from direct
  eigen-propagation of the full $43$-state system, and the single-slow-mode
  estimate of Eq.~\eqref{eq:elm_bound}, at five grid points. The averages are
  integrated at $400$ time samples per decade; at $1600$ per decade the five
  change by at most $4.7\times10^{-6}$ relative, and the $[0,4]$ average by
  $2.7\times10^{-10}$. Integration on a uniform \SI{25}{\nano\second} mesh, which
  under-resolves the nanosecond rise, reads low by $6\times10^{-5}$ to
  $5\times10^{-4}$ relative. The estimate reproduces
  the true average to better than \SI{0.8}{\percent} at the five points. The
  value below unity at $[0,4]$ is real: a deficit of $2.2\times10^{-5}$ that
  survives the denser grid. Source: \texttt{verify\_trajectory\_census.py},
  \texttt{validation/trajectory\_census/}.}
  \label{tab:lowerbound}
  \begin{tabular}{lrrrr}
    \toprule
    point & $\epsplat$ & true \SI{100}{\micro\second} average & estimate, Eq.~\eqref{eq:elm_bound} & true/estimate \\
    \midrule
    $[0,4]$ worst case & $0.386903$ & $0.386811$ & $0.386820$ & $0.99998$ \\
    $[15,3]$ ridge     & $0.180728$ & $0.175283$ & $0.174812$ & $1.00270$ \\
    $[15,5]$           & $0.103550$ & $0.072185$ & $0.071722$ & $1.00646$ \\
    $[23,5]$ benchmark & $0.063612$ & $0.011796$ & $0.011712$ & $1.00722$ \\
    $[0,0]$            & $0.083004$ & $0.083084$ & $0.083004$ & $1.00096$ \\
    \bottomrule
  \end{tabular}
\end{table}

At the five points tested the estimate sits within \SI{0.8}{\percent} of the
directly propagated average, and at $[0,4]$ it sits above it, so it is an
estimate and not a bound. The accuracy at these five points does not
generalise: over the $448$ pairs above \SI{2}{\electronvolt} the ratio of true
average to estimate departs from unity by up to \SI{2.2}{\percent} at
\SI{100}{\micro\second} and \SI{3.8}{\percent} at \SI{506}{\micro\second}
(Section~\ref{sec:persistence}). What matters here is that the assumption
underneath it was tested rather than asserted.

The reason is structural, and it is a property the shell ratio has and another
observable in this same project does not. The spectrum of $\Lmat$ has one slow
eigenvalue and then a band, with nothing between: at $\Lmat[1,4]$ the three
slowest eigenvalues are $-4.29$, $-1.60\times10^{8}$ and
$-3.44\times10^{8}\,\si{\per\second}$, a gap of $3.7\times10^{7}$. An observable
that depends on the excited populations only through their ratio cannot decay
faster than $\tauslow$: once the manifold is quasi-steady the excited vector is
a function of $u$ alone, so
$\ln\Robs - \ln\Robs^{\mathrm{CRE}} = (S/u^{*})\,\delta u + O(\delta u^{2})$
with $\delta u \propto \mathrm{e}^{-t/\tauslow}$, and the leading term
vanishes only where $S = \ffrac{3}-\ffrac{4} = 0$. It does not: $S$ is positive
at all $400$ grid points, with minimum $0.046$ at $[44,7]$, and the slow
eigenvector's projection onto $\ln\Robs$ equals $S/u^{*}$ to
$1.2\times10^{-4}$ at the benchmark and $1.6\times10^{-8}$ at $[0,4]$, the
deviation being the first-order quasi-steady correction
$\lambda_{\mathrm{slow}}/\lambda_{2}$ itself. It exceeds \SI{1}{\percent} only
at the two hot, dense points $[48,7]$ and $[49,7]$, where the separation is
$87$-fold and which the window criterion excludes from the analysed set
(\texttt{validation/slow\_mode\_projection/}). That decay on an intermediate
mode was the failure mode being hunted, and it does occur for the
ground-normalised error measure, which does not cancel the reservoir and whose
$1/\mathrm{e}$ time is \SI{2.58}{\second} against
$\tauslow = \SI{67.2}{\second}$ at the cold corner $[0,0]$. The estimate is
accurate for the ratio and not for that measure; why it is nevertheless not a
bound, the denominator relaxing on the same $\tauslow$, is set out in
Section~\ref{sec:persistence}.

%-----------------------------------------------------------------------------
\section{External comparison: Fujimoto's population coefficients}
\label{sec:fujimoto}
% QUESTION: does an independently computed collisional-radiative model, built
% decades earlier from different atomic data, agree with this one?

Everything to this point compares the model against itself or against
arithmetic. Only comparison against an independent calculation reaches the
atomic physics, and this project has one such comparison in depth: against
Fujimoto's collisional--radiative treatment of atomic
hydrogen~\cite{Fujimoto2004}, computed decades earlier, with different atomic
data, by a different method.

\subsection{Timescales, at a point that lies exactly on this grid}

Fujimoto works a numerical example at $\Te = \SI{10}{\electronvolt}$ and
$\nel = \SI{e18}{\per\cubic\metre}$, which is
$\SI{e12}{\per\cubic\centi\metre}$ and lands exactly on grid point $[49,0]$. No
interpolation is involved.

\begin{center}
\begin{tabular}{lrrr}
  \toprule
  quantity & Fujimoto App.~4B & this model, $[49,0]$ & ratio \\
  \midrule
  excited-state response & ${\sim}\SI{1e-7}{\second}$ & \SI{3.4057e-8}{\second} & $0.34$ \\
  ground-state depletion & ${\sim}\SI{1e-4}{\second}$ & \SI{1.7657e-4}{\second} & $1.77$ \\
  \bottomrule
\end{tabular}
\end{center}

The two-clock structure appears in both, separated by three orders of
magnitude in both, with the fast time short and the slow time long in the same
sense.

\textbf{What this comparison is not.} Fujimoto's response time is
$\max_p t_{\rm tr}(p)$, the \SI{63}{\percent} rise time of the slowest-responding
level, which is not $1/|\lambda_1|$. The two are different functionals of the
same spectrum, so agreement to a factor of a few is the meaningful comparison
and exact agreement would be suspicious. Quoting these ratios as though they
were a precision test would overstate them.

One qualification on the independence of this check, found on reading the
printed appendix rather than a summary of it. Fujimoto's Figs.~4B.1 and 4B.2 are
captioned as quoted from Sawada and Fujimoto (1994)
\cite{SawadaFujimoto1994}, so the appendix and the paper
Chapter~\ref{ch:conclusions} discusses are the same calculation, not two.
\textbf{The comparison is therefore independent of this thesis but not two
independent checks}, and it is counted once. The appendix also states the same
construction the paper does, that $n(1)$ is held constant, which matters for
Chapter~\ref{ch:conclusions} and is taken up there.

A second, independent route corroborates the same physics. Fujimoto's boundary
level, the principal quantum number above which collisions dominate radiative
decay, sits at $p_G \approx 4$--$5$ at this density. Tracing the boundary
descent in this model puts it at $5g$ at the low-density end of the grid, which
is the crossing Chapter~\ref{ch:atoms} deferred to this chapter
(Section~\ref{sec:radiative}). Two independently computed models place the same
structural feature within one shell of each other, at one density.

\subsection{The recombining coefficients agree}

Fujimoto's Table~4.1 tabulates the two population coefficients of the
decomposition that Section~\ref{sec:R_depends_on_b} derives: $r_0(p)$, how
strongly level $p$ is fed from the ion continuum, and $r_1(p)$, how strongly it
is fed from the ground state. Both are dimensionless, and both are normalised to
the Saha--Boltzmann population of level $p$, which is why $r_0$ is of order unity
while $r_1$ at low density is of order $10^{-5}$, and why a value of $r_0$ above
one is unphysical.

The recombining coefficients agree to \SIrange{0.1}{1}{\percent} for $n \ge 4$.
The discriminating case is $r_0(2)$, where cascade and $\ell$-mixing structure
matter most and a hydrogenic approximation would go wrong: this model gives
$0.7392$ against the tabulated $0.730$, a difference of \SI{1.3}{\percent}, a
test the model could have failed and did not (the broken variant of
Section~\ref{sec:severe_checks} drives the same quantity to $104.5$).

\subsection{The ionising coefficients at low $p$ disagree, and the target is
not verified}

At $\nel = \SI{e12}{\per\cubic\centi\metre}$ this model's $r_1$ sits a factor
$8.3$ low at $p=3$ and $4.4$ low at $p=4$ against Fujimoto's Table~4.1(b).
Those are precisely the two shells whose difference is the mechanism of
Chapter~\ref{ch:results}, so the discrepancy cannot be dismissed as peripheral.

Reading it as a failure of the model is not supportable, for six reasons, and
the status is an open comparison whose \emph{target} is unverified.

\begin{enumerate}
\item \textbf{The model's $r_1$ is reproduced by hand.} In the coronal limit
$r_1(2)$ follows from the single excitation rate coefficient
$K_{\rm exc}(1s\to n{=}2) = \SI{9.2796e-9}{\cubic\centi\metre\per\second}$ and
gives $1.37\times10^{-5}$, against the full $43$-state solve's
$1.66\times10^{-5}$. The two agree to \SI{21}{\percent}, which is the accuracy a
coronal estimate that ignores cascades should have. Whatever produces a factor
of $8$ is not inside the solve.

\item \textbf{The atomic stack passes three independent external checks.} The
radiative recombination coefficient into the ground state is within
\SI{2}{\percent} of the scaled standard value; the sum over excited states,
$\SI{2.08e-13}{\cubic\centi\metre\per\second}$, sits \SI{10}{\percent} below the
standard case-B coefficient $\alpha_B \approx
\SI{2.31e-13}{\cubic\centi\metre\per\second}$, the accepted total for capture
into every excited level, and that \SI{10}{\percent} shortfall is exactly the
$n>15$ truncation of Section~\ref{sec:convergence}; and the ground-state ionisation
coefficient at \SI{10}{\electronvolt}, $\SI{4.864e-9}{\cubic\centi\metre\per\second}$,
sits \SI{10.4}{\percent} below ADAS SCD96 read in its low-density limit,
$\SI{5.43e-9}{\cubic\centi\metre\per\second}$ at
$\nel = \SI{5e7}{\per\cubic\centi\metre}$ \cite{Summers2006, ADAS}, where the
effective coefficient reduces to ionisation out of the ground state because
every excited level decays before it can be ionised. The comparison is at the
low-density limit deliberately: SCD96 varies by \SI{0.9}{\percent} over a
factor four in density there, so the limit is reached. \textbf{This
\SI{10}{\percent} is not the truncation of the item above.} Direct ionisation
from $1s$ is a single cross section and does not know where the level ladder
stops; the difference is between this model's ionisation cross-section source
and the one behind SCD96, and it is reported rather than reconciled.

\item \textbf{The tabulated asymptote demands an excitation coefficient fifteen
times the accepted value.}
Reproducing Fujimoto's quoted coronal asymptote requires
$K_{\rm exc}(1s\to n{=}2) = \SI{1.44e-7}{\cubic\centi\metre\per\second}$ at
\SI{11.03}{\electronvolt}, fifteen times the accepted value. No error in this
model's implementation produces that number; only a different input does.

\item \textbf{Three candidate explanations are eliminated quantitatively.}
Truncation moves $r_1(3)$ by \SI{1.2}{\percent} between $n_{\max}=10$ and
$n_{\max}=15$, which is a factor $8$ short. Lyman trapping at an escape factor
of $0.1$ brings $r_1(2)$ to $0.92$ of the table but leaves $r_1(3)$ at $0.26$ of
it and drives $r_0(2)$ to $8.2$, a population above Saha equilibrium and
therefore unphysical. The third candidate, the $\ell$~closure, is eliminated in
the item below.

\item \textbf{The $\ell$~closure is a real difference, and it does not account
for the gap.} Table~4.1(b) is indexed by $p = 2, 3, 4, 5, 7, 10, 15$, principal
quantum numbers, and carries no $\ell$ label: the published coefficients are
shell quantities computed with the $\ell$~substructure bundled, while this model
resolves $\ell$ below $n=9$. That difference is real, and it is not the
explanation. Imposing the source's own
closure directly, by bundling this model's operator under statistical
$\ell$~weights as $\Lmat_{\rm bundle} = \mathsf{P}\,\Lmat\,\mathsf{R}$ with
$\mathsf{P}_{ki} = 1$ for $i$ in shell $k$ and $\mathsf{R}_{ik} = g_i/g_k$,
moves $r_1(3)$ by \SI{0.42}{\percent}: the factor $8.29$ becomes $8.26$
(\texttt{verify\_fujimoto\_bundle.py}, stamped in
\texttt{validation/fujimoto\_bundle/}). At
$p=2$ the agreement becomes \emph{worse}, from a factor $10.8$ low to a factor
$11.8$ low, and $r_0(2)$ moves from $0.7392$, which agrees with the tabulated
$0.730$, to $0.6459$, which does not. The reason is structural: the bundling
projector annihilates the $\ell$-mixing operator identically, $\max|\mathsf{P}
\mathsf{B} \mathsf{R}| = 8.5\times10^{-7}$ against a matrix scale of
$1.0\times10^{11}$, so bundling cannot express what $\ell$-mixing does; and the
production model is already at the statistical-$\ell$ limit for $p \ge 3$, its
sublevel populations sitting within \SI{5}{\percent} of $g_{n\ell}/g_n$. The
$\ell$-resolved calculation was therefore already the like-for-like comparison.
Removing proton $\ell$-mixing altogether moves $r_1(2)$ by a factor $118$, but
that is the opposite limit from a statistical closure rather than a proxy for
it, and the two do not bracket the tabulated value.

\item \textbf{The transcription and the row labels are correct.} Both were
checked against the printed table. The $\log_{10}\nel = 18$ row reads
$0.730$, $0.835$, $0.947$ and $0.983$ for $r_0$ at $p = 2, 3, 4, 5$, and
$1.79\times10^{-4}$, $5.72\times10^{-5}$, $1.66\times10^{-5}$ and
$5.08\times10^{-6}$ for $r_1$, matching the values used here exactly. The row
labels are $\log_{10}(\nel/\si{\per\cubic\metre})$, which the table establishes
internally: $r_0(2)$ reaches $0.981$ at $\log_{10}\nel = 21$, and Griem's
criterion places local thermodynamic equilibrium for $n=2$ at
\SI{1.74e16}{\per\cubic\centi\metre}, which is $10^{21}\,\si{\per\cubic\metre}$
to within the criterion's own precision. Read as \si{\per\cubic\centi\metre} the
onset would sit five orders above Griem, which is impossible; and the rows run
from $\log_{10}\nel = 12$ to $24$, a span of $10^{6}$ to
$10^{18}\,\si{\per\cubic\centi\metre}$ read as \si{\per\cubic\metre} and no
plasma at all at the top read as \si{\per\cubic\centi\metre}. A one-decade
offset in the density rows is not a tenable reading either.
\end{enumerate}

\textbf{Status.} The $r_1$ deficit at low $p$ is an open external
disagreement, and this thesis does not explain it; the three candidate causes
above each fail to account for a factor of $8$. The $\ell$-closure test carries
a severity check: bundling an operator whose $\ell$~distribution is deliberately
non-statistical moves $r_1(2)$ by a factor $129$, so the null result at $p=3$ is
a measurement and not a blind instrument. What the section does \emph{not}
establish is that the disagreement leaves the quantity
Chapter~\ref{ch:results} uses untouched.
Equation~\eqref{eq:fujimoto_correspondence} makes
$a_p = r_1(p)\,\Zsb{p}/\Zsb{1}$: $r_1$ is the ground-fed coefficient from which
$\ffrac{p}$ is built, so a deficit in $r_1(3)$ and $r_1(4)$ is a deficit in
$a_3$ and $a_4$ themselves. Bundling is not the protection, since neither shell
passes through a bundled coefficient; the only protection is that $\Delta$ and
the cap depend on the ratio $a_3/a_4$, in which a deficit common to both shells
cancels. At \SI{e12}{\per\cubic\centi\metre} the deficit is not common, $8.3$
at $p=3$ against $4.4$ at $p=4$; at \SI{e15}{\per\cubic\centi\metre} it nearly
is, $1.8$ against $1.9$. Dividing the model's $a_3$ and $a_4$ by those ratios,
so that they become what Fujimoto's coefficients would give, and recomputing
at every grid point (\texttt{verify\_fujimoto\_r1\_consequence.py},
\texttt{validation/fujimoto\_r1\_consequence/}) moves $\Delta$ by $+0.64$
under the first pair and by $-0.08$ under the second. Over the $280$ points
above \SI{2}{\electronvolt} the cap then moves by \SI{+29}{\percent} at the
median under the first and \SI{-4}{\percent} under the second; the crest
$u^{\mathrm{peak}}$ falls by a factor $6.0$ and $1.9$; and the sensitivity at
the operating point, $\ffrac{3}-\ffrac{4}$ evaluated at $u^{\mathrm{CRE}}$,
moves by \SI{+30}{\percent} and \SI{+11}{\percent} at the median, with ranges
across the points of $-58$ to \SI{+524}{\percent} and $-38$ to
\SI{+68}{\percent}. Both exceed the \SIrange{8}{9}{\percent} that the R-matrix
excitation substitution of Chapter~\ref{ch:atoms} produces. The disagreement
is therefore direct evidence about the coefficients Chapter~\ref{ch:results}
uses; what is open is whether the deficit lies in this model or in the
tabulation, and the scenario applies a deficit measured at
\SI{10}{\electronvolt} at every temperature, which is its stated assumption.
An unexplained disagreement in an absolute population coefficient is reported
here as unexplained, with its stake stated.

\textbf{What the benchmark does establish.} The recombination-fed channel
agrees, and because $r_0$ and $r_1$ carry the same $p$-dependent factor $Z(p)$,
a normalisation error large enough to explain the $r_1$ deficit is excluded by
$r_0$ agreeing. That is a pass, and it should be read as one.

\textbf{What it does not establish, and two limits on it.} It does not fix the
absolute normalisation of the ground-fed channel against an independent
calculation, because no $\ell$-resolved tabulation was available to compare
against. Table~4.1(b) is computed at $T_\mathrm{e} = \SI{1.28e5}{\kelvin}$,
which is \SI{11.03}{\electronvolt}, while this grid reaches
\SI{10}{\electronvolt}, so the comparison sits at the grid edge, with the
tabulated temperature \SI{10.3}{\percent} above it. The companion
Table~4.1(a) is at \SI{e3}{\kelvin} =
\SI{0.0862}{\electronvolt}, two orders below this grid, so there is no second
usable case. The external check is one temperature at one edge.

%-----------------------------------------------------------------------------
\section{Gate D fails, and the failure is in the gate}
\label{sec:gate_d}
% QUESTION: the one gate that compares against an external atomic database
% disagrees at every point. What does that mean, and what does it not mean?

The Atomic Data and Analysis Structure (ADAS) distributes effective
stage-to-stage rate coefficients for hydrogen: SCD, the effective ionisation
coefficient, and ACD, the effective recombination coefficient~\cite{Summers2006}.
These are the coefficients a transport code uses when it does not resolve
excited states, and they are the most widely used external reference this model
could be compared against.

Gate~D forms the model's counterpart by taking the steady-state population
$n_p^{\rm ss}$ of each level at each grid point, weighting it by that level's
ionisation rate coefficient $Q_p$ of Eq.~\eqref{eq:level_balance}, and dividing
by the total neutral population:
\begin{equation}
  {\rm SCD}_{\rm model}
  = \frac{\sum_p Q_p\,n_p^{\rm ss}}{\sum_p n_p^{\rm ss}},
  \label{eq:scd_model}
\end{equation}
then compares against ADAS interpolated to the same temperature within a
factor-of-two density band. The gate passes a point if the ratio
$\eta = {\rm SCD}_{\rm model}/{\rm SCD}_{\rm ADAS}$ lies between $0.5$ and $2.0$.

\begin{center}
\texttt{Gate D: FAIL  (0\% of points)} \hfill \texttt{validation/gate\_summary.txt}
\end{center}

Measured, $\eta$ ranges from $5.73$ to $8346$ over the full grid, and no point
agrees within a factor of two. The extremes are not typical: over most of the
grid the model sits high by a factor between $9$ and $65$, and that systematic
offset peaks near $\Te \approx \SI{2}{\electronvolt}$ (these are the gate's
own numbers, under its linear read of the table; the read is discussed
below). The disagreement has a
shape, which is what makes it diagnosable rather than random.

\subsection{The diagnosis, and the test that settles it}
\label{sec:gate_d_diagnosis}

Two readings were open. One of them has now been tested and it is the
explanation.

The first is that Eq.~\eqref{eq:scd_model} is not the quantity ADAS tabulates.
ADAS defines SCD as the \emph{ionising} coefficient, the ionisation flux
attributable to the ground-state population. The sum in
Eq.~\eqref{eq:scd_model} runs over the full steady state, which includes the
recombination-fed Rydberg levels, the $\nzero = \bm{c}\,\nion$ channel of
Section~\ref{sec:R_depends_on_b}. Ionisation out of
those levels is a recombination-driven flux and belongs under ACD, not SCD.
Since the high levels have small ionisation potentials and large ionisation
coefficients, moving them to the wrong side of the ledger inflates the model's
SCD by a large factor, and inflation is what is measured. The test that would
settle it is to recompute Eq.~\eqref{eq:scd_model} over the ground-fed channel
$\none$ alone, and it has now been run
(\texttt{diagnose\_gate\_d.py}).

The two channels separate exactly, with no approximation:
$\nvec_E = \nzero + \none$ to a relative residual below $10^{-8}$ at every grid
point, which the script checks before using either. Rebuilding the coefficient
on the ground-fed channel alone,
\begin{equation}
  {\rm SCD}_{\rm model} = Q_{1s}
  + \sum_{p \neq 1s} Q_p\,\frac{n^{(1)}_p}{\ngs},
  \label{eq:scd_groundfed}
\end{equation}

\noindent and comparing against the same ADAS table with the same factor-of-two
tolerance moves the agreement from $0$ of $400$ points to $400$ of $400$. The
ratio $\eta$ falls from a range of $5.73$ to $8346$ to a range of $0.702$ to
$1.004$, with median $0.828$; above \SI{2}{\electronvolt} it runs from $0.721$
to $0.937$, median $0.831$, and at the benchmark it is $0.768$
(\texttt{validation/reservoir\_vs\_adas/}). The read matters. ADAS96
tabulates seven temperatures between $1$ and \SI{10}{\electronvolt}, and
\texttt{diagnose\_gate\_d.py}, which first made this comparison, interpolates
linearly in value between them (\texttt{adas\_at}, lines 114 to 124); across
$1$ to \SI{1.5}{\electronvolt}, where SCD rises by nearly two orders of
magnitude, that overstates the table by up to a factor $6.65$. That read gives
$323$ of $400$, a range of $0.130$ to $0.931$ and a median of $0.696$, with a deficit deepening toward low temperature; read logarithmically
in both axes, as the table's spacing requires, the tail disappears, the row
median being $0.72$ at \SI{1}{\electronvolt} and $0.90$ by
\SI{1.2}{\electronvolt}. The diagnosis script is left in place; every ADAS
number in this section comes from the logarithmic read.

The mechanism is visible directly. At collisional--radiative equilibrium the
recombination-fed channel carries between \SI{89.5}{\percent} and
\SI{99.99}{\percent} of $\sum_p Q_p n_p$, and its share rises monotonically with
density, from \SI{97.8}{\percent} at $\SI{e12}{\per\cubic\centi\metre}$ to
\SI{99.99}{\percent} at $\SI{e15}{\per\cubic\centi\metre}$ at
\SI{1}{\electronvolt}. Equation~\eqref{eq:scd_model} is therefore reporting the
recombination channel with a small ground-fed correction, which is why $\eta$
tracks $\nel$: three-body recombination carries an extra factor of $\nel$ and
the mislabelled flux carries it into the coefficient.

\paragraph{The recombination half, built and compared for the first time.}
The same decomposition supplies the coefficient the gate never computed. The net
flux into the neutral ground state carried by the recombination-fed channel,
per unit $\nel\nion$, is
\begin{equation}
  {\rm ACD}_{\rm model}
  = \frac{S_{1s}\,\nion + \sum_{p \neq 1s} L_{1s,p}\,n^{(0)}_p}{\nel\,\nion},
  \label{eq:acd_groundfed}
\end{equation}

\noindent direct radiative recombination into $1s$ plus everything that cascades
or is collisionally transferred down into it. Against ACD96 this agrees at
\emph{all} $400$ points, with $\eta$ between $0.894$ and $1.009$ and median
$0.982$; at the benchmark the model gives
$\SI{2.788e-13}{\cubic\centi\metre\per\second}$ against ADAS's
$\SI{2.883e-13}{}$, a ratio of $0.967$.

\textbf{That is the strongest external check in this thesis.} It is an absolute
coefficient, compared against an independent database built from independent
atomic data, agreeing to better than \SI{10}{\percent} at the benchmark and
within a factor of two everywhere. It was available all along and was not taken,
because the half of the gate that would have taken it was never written.

The second reading was that the two quantities are not defined on the same
system at all. ADAS coefficients describe a closed atom-plus-ion system; this
model holds the ion as a fixed external reservoir with no forty-fourth state
(Section~\ref{sec:two_references}). A stage-to-stage coefficient extracted from
an open system need not equal one extracted from a closed one, and
Section~\ref{sec:attacks} has already shown that the boundary condition moves
$\tauslow$ by up to a factor of $41$. \textbf{The ACD result argues against this
being the dominant effect}: ACD is extracted from the same open system through
the same fixed-ion reservoir and agrees with ADAS to within \SI{3.3}{\percent}
at the benchmark, \SI{1.8}{\percent} at the grid median and \SI{11}{\percent}
at every point; an open-system artifact large enough to produce $\eta$ of
$8346$ in SCD would not leave ACD intact.

\paragraph{The residual SCD deficit, stated rather than dismissed.} Rebuilding
SCD on the ground-fed channel does not make the agreement perfect, and the
remaining disagreement is systematic in a way worth reporting: the model's
ionisation sits below ADAS at all but one of the $400$ points, by up to
\SI{30}{\percent} and by \SI{17}{\percent} at the median, with the maximum
$\eta = 1.004$ at $[4,7]$. Under the linear read of the table the gap closes
steeply with temperature, from a median of $0.302$ over $1$ to
\SI{1.5}{\electronvolt} to $0.817$ over $5$ to \SI{10}{\electronvolt}, a shape
that reads as the signature of truncation; the steep part is the linear read,
and under the logarithmic read the deficit is between $0.72$ and $0.94$ at every
temperature above \SI{2}{\electronvolt}. A deficit of that size is still what
truncation would produce, since ADAS96 carries levels above $n_{\max} = 15$
and the levels nearest the continuum ionise most easily, but its temperature
shape no longer singles that cause out (the $n \ge 12$ flux share the diagnosis
script printed against the linear read is not carried forward). The test that
would settle the
cause is the $n_{\max}$ convergence study of Chapter~\ref{ch:limits} rather
than this one.

\paragraph{The reservoir itself, against the same table.} ADAS's two
coefficients fix the equilibrium reservoir: at ionisation balance
$n_{1s}\,\mathrm{SCD} = \nion\,\mathrm{ACD}$, so
$u^{\mathrm{CRE}} = \mathrm{ACD}/\mathrm{SCD}$, and the model obeys the same
identity with its own ground-fed SCD and ACD to $4\times10^{-13}$ at every grid
point. That makes the quantity Chapter~\ref{ch:results} is built on directly
comparable, and the comparison has two parts
(\texttt{verify\_reservoir\_vs\_adas.py},
\texttt{validation/reservoir\_vs\_adas/}). The \emph{level} is high:
$u^{\mathrm{CRE}}_{\mathrm{model}}/u^{\mathrm{CRE}}_{\mathrm{ADAS}}$ has median
$1.18$ over the $400$ points, runs from $0.89$ at $[4,7]$ to $1.40$ at $[0,1]$,
and is $1.26$ at the benchmark, structured about equally in $\Te$ and in
$\nel$. The model's ground reservoir per ion sits a fifth above ADAS's, which
moves the operating point beneath the cap of Section~\ref{sec:bound} without
touching the cap. The \emph{displacement}, which is what the plateau error is
built from, agrees. Over the $392$ one-interval temperature steps the model's
$\Delta\ln u$ has median magnitude $0.26$ and differs from ADAS's by a median
$0.014$; above \SI{2}{\electronvolt} the median difference is $0.011$ and the
worst $0.050$, at $[15,0]$; at the benchmark step the two are $-0.270$ and
$-0.258$. The reservoir gain $G$ this thesis tabulates is therefore within
about \SI{5}{\percent} of the one ADAS implies, at the median over the
defended range, and the refuter set before the run, a difference above $0.1$
at any interior step above \SI{2}{\electronvolt}, appeared at none of $264$.
Every number in this paragraph comes from the logarithmic read; on the linear
read of \texttt{diagnose\_gate\_d.py}, which overstates SCD by up to a factor
$6.65$ over $1$ to \SI{1.5}{\electronvolt}, the refuter would have appeared at
$9$ of the $264$ steps.

\subsection{Three further defects in the gate itself}

The gate is also incompletely implemented, and this must be said alongside the
failure rather than instead of it.

The recombination half does not exist \emph{in the gate}. The ACD table is
loaded once, at line 299 of \texttt{validate\_gates.py}, and the variable
holding it is never read again anywhere in the file; only the ionisation
comparison runs, despite a docstring that describes both. That half is computed
in \texttt{diagnose\_gate\_d.py} (Section~\ref{sec:gate_d_diagnosis}).

The interpolation is wrapped in a catch-all exception handler that sets the
comparison value to not-a-number and continues, so a failure to find the right
density band is indistinguishable from a point that was legitimately skipped.

And the gate's status is recorded three incompatible ways: \texttt{README.md}
line 119 calls it documented, the docstring at line 47 predicts a pass above
$\Te = \SI{5}{\electronvolt}$, and \texttt{gate\_summary.txt} reports failure
at \SI{0}{\percent} of points. The correct entry is the measured one: Gate~D
fails.

\subsection{What is claimed, and what is not}

\textbf{Gate~D itself still fails, and it has not been changed.}
\texttt{validate\_gates.py} is untouched: no tolerance was widened, no subset
was selected, and \texttt{gate\_summary.txt} still records
\SI{0}{\percent}. What has changed is that the failure now has a demonstrated
cause rather than two undemonstrated candidates. The gate compares a mixed
quantity against a table that keeps the two channels apart by design, which is
a defect in the comparison and not in the matrix.

Deciding what Gate~D should test belongs with whoever owns the gate. What is
claimed here is narrower: rebuilt on the two channels separately, the
coefficients agree with ADAS within a factor two at all $400$ points for SCD,
the model $0$ to \SI{30}{\percent} low with a median ratio of $0.83$, consistent
with the $n_{\max}$ truncation but not uniquely so, and at all $400$ for ACD
with a median ratio of $0.98$.

The failure does not cost the thesis its only comparison of absolute effective
coefficients against an external database; that comparison exists, hidden by
the definition. The failure does not touch the structural results
of Chapter~\ref{ch:theory}, which are statements about whatever $\Lmat$ is
supplied and would hold for the ADAS-consistent matrix equally. It does not
touch the level-resolved comparison of Section~\ref{sec:fujimoto}, which
compares the same decomposition against tabulated coefficients rather than
against a contraction of $42$ coefficients into one. And the contraction is the
reason the two comparisons carry different weight: many incorrect $\bm{a}$ would
contract to the correct SCD, so a passing Gate~D would have been weak evidence
even had it passed.

%-----------------------------------------------------------------------------
\section{Errors found by our own audit}
\label{sec:errors_found}
% QUESTION: what did this project get wrong, how was each error found, and how
% large was it?

Section~\ref{sec:atoms_corrections} reports two implementation errors that this
project found in itself, with their sizes, and observed that both were invisible
to every consistency test the pipeline ran. This section states what each
changes here, adds two more, and keeps the four together rather than in an
appendix: a model whose authors found and bounded four of its own errors is
better characterised than one reporting none, and in three of the four cases
the bound is itself a result.

\subsection{The $\ell$-mixing Coulomb logarithm}

Proton-impact $\ell$-mixing redistributes population between sublevels of the
same $n$ and, as Section~\ref{sec:lmixing} showed, is the dominant loss channel
for the metastable $2s$ level. The rate coefficient follows Badnell's
formulation~\cite{Badnell2021}, whose Eq.~9 carries a factor
\begin{equation}
  F(U_m) = \frac{\sqrt{\pi}}{2}U_m^{-3/2}\operatorname{erf}(\sqrt{U_m})
           - \frac{\mathrm{e}^{-U_m}}{U_m} + E_1(U_m),
  \label{eq:FUm}
\end{equation}
in which $\operatorname{erf}$ is the error function, $E_1$ the exponential
integral, and $U_m$ the dimensionless ratio of the smallest energy transfer the
Debye cutoff permits to $k_B\Te$; $F(U_m)$ is the Coulomb logarithm of the
interaction, which diverges as $U_m \to 0$.

The implementation set $F = 1$, justified in its docstring as the low-density
limit $U_m \to 0$; the limit was inverted, and
Section~\ref{sec:atoms_corrections} gives the finding, the factor of three to
seven by which the rates were low, and the measured effect on the spectrum.
Evaluated with the local Debye length at the benchmark point $[23,5]$,
$F = 6.62$ for $2p$--$2s$, $5.71$ for $3d$--$3p$, $4.26$ for $5g$--$5f$ and
$2.89$ for $8k$--$8i$ (\texttt{validation/lmix\_cutoffs/}); with the
production table's frozen density of \SI{e14}{\per\cubic\centi\metre} the
$2p$--$2s$ value is $6.95$ (Section~\ref{sec:atoms_corrections}). A derivation
note records $6.65$, $5.73$, $4.29$ and $2.92$ for the same four pairs.

The insensitivity of the slow eigenvalues to that error, which
Section~\ref{sec:atoms_corrections} measures, is specific to a process already
far faster than $\taurel$; it is not a general licence to be careless about
fast rates, and the $\ell$-mixing scan below is what establishes the range over
which it holds.

\subsection{And $\ell$-mixing is saturated}

The same insensitivity was then tested directly rather than argued. Scaling
every $\ell$-mixing rate in the matrix over a range of $\times0.1$ to $\times10$
moves $\ffrac{3}-\ffrac{4}$ by less than \SI{0.3}{\percent} at the benchmark
point and at the ridge, and by less than \SI{5}{\percent} at the cold corner.
Only switching $\ell$-mixing off entirely changes anything. Rebuilding the
matrix with a density-consistent Debye cutoff in place of a frozen one moves
$\ffrac{3}-\ffrac{4}$ by at most \SI{0.42}{\percent}.

The claim that the frozen Debye density varies the cutoff factor by about
\SI{30}{\percent} across the grid is wrong: the module's own functions give
$\times3.7$ for $n=2$ and at least $\times50$ for $n=8$
(Section~\ref{sec:lmixing}); the conclusion survives because the observable does
not depend on the input over that range.

\subsection{The eigenvalue magnitude filter}

The spectrum of $\Lmat$ was extracted by a routine containing the line
\texttt{eigs = eigs[eigs < -1.0]}, which discards every eigenvalue smaller in
magnitude than $1\,\si{\per\second}$. Where the slow eigenvalue is smaller than
that, at the cold end of the grid, the filter deleted it and promoted the whole
ladder by one step, so $\tauslow$, $\taurel$ and $\Msep$ were wrong together
while all three remained finite, negative and correctly ordered. It was found by
recomputing the spectrum unconditionally and comparing bit-for-bit against the
stored arrays, the three-implementation agreement of Section~\ref{sec:gate_e}.
Section~\ref{sec:atoms_corrections} gives the extent of the damage, the ladder
shift that confirmed the mechanism, the preserved filtered outputs and the
removal of the literal from two dormant files. What belongs here is the shape
of the wrong answer: at the cold corner the filter turned
$\Msep = 1.729\times10^{9}$ into $\Msep = 1.467$, and the grid minimum from
$86.77$ into $1.341$, values that read as a physics finding, since
near-breakdown of timescale separation in the coldest, thinnest plasma is
exactly what a reader would expect to see. The later correction of the two
dormant files changes no number in this thesis: at the benchmark the two filters
return the same $\tauslow$ to every digit, and they differ only where the slow
mode is slower than $1\,\si{\per\second}$.

\subsection{The central quantity was misnamed}

The largest of the four errors changed no number and inverted a sentence.

The quantity that had been reported as the quasi-steady-state error was
defined as the difference between the tabulated ratio and the true
one, with the tabulated ratio evaluated from a two-argument function of $\Te$
and $\nel$. Section~\ref{sec:two_references} shows that the excited-state ratio
is a function of three arguments, not two, the third being the state of the
ground-state reservoir. Fixing the first two therefore silently imposes
equilibrium ionisation balance, so what was measured was the distance to the
equilibrium answer, not the error of the quasi-steady-state closure.

Integrating the full system and evaluating both errors along the same trajectory
separates them:

\begin{center}
\begin{tabular}{lrr}
  \toprule
  point & error against equilibrium & error of the QSS closure \\
  \midrule
  benchmark $[23,5]$ & $6.34\times10^{-2}$ & $8.66\times10^{-6}$ \\
  worst case $[0,4]$ & $3.869\times10^{-1}$ & $6.73\times10^{-9}$ \\
  \bottomrule
\end{tabular}
\end{center}

They differ by four to eight orders of magnitude, most at the point carrying the
largest error. The physics was intact and every number reproduced; what was
wrong was the name and the sentence built on it, and Chapter~\ref{ch:results}
states the consequence, that the approximation everyone worries about is not
the one that fails. The error was found by an adversarial reading of the thesis text
itself, not of the code, and it is why Section~\ref{sec:eps_definitions}
defines four error measures separately.

\subsection{A statement in Chapter~\ref{ch:theory} that measurement contradicted}
\label{sec:promise_broken}

The statement, in Section~\ref{sec:source_quadratic}, was that at
$\nel = \SI{e12}{\per\cubic\centi\metre}$ three-body recombination supplies
under \SI{10}{\percent} of the feed into any level, and it was used to propose
testing the radiative and three-body datasets separately at the two ends of the
density grid. Measured against the matrix, $30$ of $36$ levels exceed
\SI{10}{\percent} three-body feed at both \SI{1.00}{\electronvolt} and
\SI{2.95}{\electronvolt}, and $26$ of $36$ at \SI{10}{\electronvolt}. Into the
bundled shells it is essentially the whole feed.

The density grid therefore does not separate the two recombination datasets, and
\textbf{no such separate test was run or can be}; Section~\ref{sec:source_quadratic}
now states the measured fractions and withdraws the promise. What remains is the
gap: the radiative and three-body datasets enter this model only in combination,
nothing in this thesis tests them apart, and a test that could would need a
density below the bottom of this grid, outside the conditions the thesis is
about.

%-----------------------------------------------------------------------------
\section{Reproduction, and what a document cannot certify}
\label{sec:reproducibility}
% QUESTION: can the numbers in this thesis be regenerated, and by whom?

\subsection{Regeneration from the canonical matrix}

Every quantitative result in Chapters~\ref{ch:theory} and~\ref{ch:results} was
regenerated from $\Lmat$ and compared against the stored output. The timescale
and step-error grids reproduce bit-identically, including the values that come
from numerical integration and not from linear algebra. The plateau and persistence
tables reproduce byte for byte, $680$ of $680$ rows, with $784$ superposition
residuals at exactly zero difference. The nineteen points corrected by the
eigenvalue filter reproduce exactly, including the ladder-shift identity.

This is reproduction in the narrow sense: the same code, run again, on the same
inputs, gives the same answers. It excludes a stale array and it excludes a
non-deterministic solver. It says nothing about whether the algebra is right.

\subsection{Independent re-derivation}

The stronger check is re-derivation. All six central results of
Chapter~\ref{ch:theory} were derived from the linear collisional--radiative
system alone by a party working from the physical statement of the problem,
implemented from scratch, with the predicted values written down before any
project file was opened.

In the table below, $\bdep$ is the ground-state departure coefficient of
Eq.~\eqref{eq:departure}, the ratio of the actual ground-state population to its
Saha--Boltzmann value, and $\bdep^{\rm CRE}$ is the value it takes in
collisional--radiative equilibrium (CRE), the steady state the lookup tables
assume.

\begin{center}
\begin{tabular}{lrr}
  \toprule
  quantity & \texttt{chapter3.tex} & independent re-derivation \\
  \midrule
  $\bdep^{\rm CRE}$ at the benchmark        & $956$                        & $956.509$ \\
  switching points $c_m/a_m$                & $2.73\times10^{3}$, $1.91\times10^{4}$ & $2727.3$, $19085.5$ \\
  peak $\bdep$                              & $7.2\times10^{3}$            & $7214.7$ \\
  $\ffrac{3}$, $\ffrac{4}$, difference       & $0.260$, $0.048$, $0.212$    & $0.259652$, $0.047725$, $0.211927$ \\
  maximum of $\ffrac{3}-\ffrac{4}$          & $0.4514$                     & $0.451357$ \\
  $\Robs$ at the benchmark                  & $0.7778$                     & $0.777768$ \\
  \bottomrule
\end{tabular}
\end{center}

The timescales reproduce to four digits, $\tauslow = \SI{22.73}{\micro\second}$,
$\taurel = \SI{2.277}{\nano\second}$, $\Msep = 9981.9$.

Two further invariances were measured in the same pass and belong with the
result they protect. The shell separation $\Delta$ that sets the bound of
Section~\ref{sec:bound} is invariant under rescaling the ground-to-ion
normalisation, verified exactly under factors of $10^{10}$ and $10^{-7}$, while
the individual channel fractions are not; and $\Delta$ varies by
\SI{0.07}{\percent} across five different $\ell$-weightings of the shell sum, so
the bound is not an artifact of how sublevels were weighted.

\subsection{A prediction that was falsified, and improved the claim}

The same re-derivation made one prediction that turned out to be wrong, and it
is more informative than the six that were right.

The prediction was that $\ffrac{3}-\ffrac{4}$ must vanish at high density, on
the reasoning that in local thermodynamic equilibrium both supply channels
produce the same Boltzmann distribution, so the two shells would draw
identically and the difference would close. It does not happen. At
$\Te = \benchTe$, $\Delta$ doubles over the lowest part of the grid and then
stops moving: $0.979$ at $\SI{e12}{\per\cubic\centi\metre}$, $1.946$ at
$\SI{1.4e14}{\per\cubic\centi\metre}$, $1.939$ at
$\SI{e15}{\per\cubic\centi\metre}$
(\texttt{validation/fujimoto\_r1\_consequence/}). Above about
$\SI{7e12}{\per\cubic\centi\metre}$ it is flat, and the bound $\tanh(\Delta/4)$
built on it spans only $0.394$ to $0.450$, a \SI{14}{\percent} range, across the
whole of the rest of the grid.

The reason is the open boundary again, this time in the model's favour: the
recombination source is imposed externally rather than tied to the ion density
by a Saha relation, so the two channels cannot merge by construction, and
Gate~C's \SI{1.5}{\percent} of Saha says the grid never approaches the limit in
which they would. The consequence for the density maximum of the diagnostic
error, a position effect rather than a span effect, is
Table~\ref{tab:position_effect} in Section~\ref{sec:density_mechanism}.

\subsection{One item demoted: the May/August cross-validation}

A cross-validation between two runs four months apart, on different matrices,
with different observables and different numerical methods, was recorded in this
project's register as verified and described as the strongest single check it
had. The agreement is close: $\tauslow$ to \SI{0.064}{\percent}, $\taurel$ to
\SI{0.200}{\percent}, $\Msep$ to \SI{0.264}{\percent}, $\epsstep$ to
\SI{0.014}{\percent}. Predicting how long the error stays above
\SI{10}{\percent} from the plateau-then-decay picture gives
\SI{12.478}{\micro\second} against the earlier run's full integration,
\SI{12.4794}{\micro\second}.

Two corrections apply, and both weaken it.

The apparent agreement to eight parts in $10^{5}$ on the duration is
cancellation. The inputs agree only to \SIrange{0.06}{0.6}{\percent}, and
carrying the earlier run's own internal numbers through gives
\SI{12.406}{\micro\second}, off by \SI{0.59}{\percent}. \textbf{The figure is
agreement at the \SI{0.6}{\percent} level}, which is still strong.

More seriously, \textbf{the earlier run has never been reproduced against the
canonical matrix}. This project's own derivation record says of these same
numbers that they are a reported result, not independently re-run, and should be
treated as strong documented evidence rather than as self-verified. A
cross-validation whose earlier half cannot be re-executed is a citation, not a
check. It is reported here as documented evidence and is not counted on the
ladder of Section~\ref{sec:validation_summary}. Under the ground-truth ordering
this thesis works to, physics over mathematics over code over documents,
a document is the weakest of the four, and it is where this particular claim
currently rests.

\subsection{Two provenance defects that are not repaired}

Two arrays of timescales, $\Msep$, $\tauslow$ and $\taurel$, are written to the
same three file paths by two different scripts. Whichever ran last wins, and six
downstream readers cannot tell whose numbers they are holding. This is exactly
the condition under which the eigenvalue filter propagated, since the filtered
and unfiltered versions of the same array have the same name.

Two things bound what that costs here, and both were checked rather than
assumed. The two writers are not two calculations: each sorts the eigenvalues of
the same canonical $\Lmat$ and keeps the negative ones, neither applies the
filter that caused the earlier incident, and the arrays on disk agree with the
independently produced \texttt{verify\_timescales.py} artifact to $2\times
10^{-16}$. And none of the six readers feeds this document: every figure and
table they write is absent from it. The benchmark $\Msep = 9982$ quoted in
Chapter~\ref{ch:theory} and in the abstract is not read from these arrays at
all; it is $\tauslow/\taurel$ computed from $\Lmat$ by
\texttt{verify\_timescales.py} and recorded in
\texttt{validation/spectrum\_testpoint.csv}. The defect is therefore a latent
hazard in the repository rather than an ambiguity in any number printed here,
and it is left in place and named rather than quietly repaired.

An audit script exists to detect single-writer violations, and it reports an
unfixed failure of this kind. It is also incomplete. Its test for a write reads the file mode from the second
argument or a keyword, so the idiom \texttt{path.open("w")} passes through it
unrecorded. At the tip of \texttt{main} on 23~August that idiom had four users;
in the tree this thesis is built
from it has $45$ call sites in $29$ scripts, $42$ of them writing
comma-separated output, among them the producer of the plateau map that
fifteen places in Chapter~\ref{ch:results} read from
(\texttt{verify\_writer\_census.py}, \texttt{validation/writer\_census/}). The
audit should not be cited as evidence that the single-writer property holds.
The census confirms that the property fails for exactly the three timescale
arrays, each written by both \texttt{qss\_analysis.py} and
\texttt{validate\_gates.py}, and holds for the plateau map, which has one
writer.

%-----------------------------------------------------------------------------
\section{What this chapter establishes}
\label{sec:validation_summary}
% QUESTION: after grading every check, what is the model licensed to be used
% for, and where?

Table~\ref{tab:ladder} collects the checks in the order a reader should weight
them, with the grade of Section~\ref{sec:gate_philosophy} attached to each.

\begin{table}[tbp]
\centering
\caption[The validation ladder, graded by severity and provenance]{Every check
this chapter grades, with what it tests, its severity, whether its result can be
regenerated, and the measured outcome. \emph{Tautological} means the check
cannot fail except on an arithmetic bug, because the quantity tested was
constructed from the relation being tested; \emph{weak} means it can fail but
only for an error far larger than any plausible one; \emph{severe} means a
plausible wrong version of this model fails it. The provenance column is the
more uncomfortable of the two: \emph{stamped} means a script in the repository
regenerates the number into an artifact under \texttt{validation/};
\emph{notes} means the number is recorded in this project's working files with
no producing script, so it cannot be re-executed and cannot fail anything today.
None of the checks graded severe are in that condition, and the one
\emph{notes} entry left is the $\tanh$ bound, which is graded tautological. The
counts stated in this caption are checked against the table's own rows by
\texttt{verify\_ladder\_counts.py}, which exits non-zero if they drift apart. Of the eighteen
internal checks, twelve are severe, four are tautological and two are weak; of the five external comparisons, two
pass, one passes only at the higher densities, one fails, and one is open.}
\label{tab:ladder}
\footnotesize
\begin{tabular}{p{3.9cm}p{1.8cm}p{1.3cm}p{5.6cm}}
\toprule
check & grade & record & measured \\
\midrule
Atomic-data provenance (Ch.~\ref{ch:atoms})  & severe & stamped & every coefficient traces to a named dated source; tallies exact \\
Excitation rates against RMPS (Ch.~\ref{ch:atoms}) & severe & stamped & \SI{81.4}{\percent} within \SI{20}{\percent} for $n_{\rm upper}\le4$, mean \SI{13.0}{\percent}; $n=5$ disagrees by factors of $2$--$5$, localised to the top shell of the RMPS basis \\
Raw cross sections, microscopic reversibility & severe & figure & ratio $0.9995 \pm 0.0032$, \SI{98.7}{\percent} within \SI{1}{\percent} \\
One energy ladder on both sides & severe & stamped & Saha ratio $0.999997$ for all $43$ states, spread $1.4\times10^{-15}$ \\
Superposition of the two channels & severe & stamped & $3.075\times10^{-14}$ over $784$ pairs; catches a $3\leftrightarrow4$ swap \\
Reduced against full $\Robs$ & severe & stamped & $0.777768$ by three routes; catches the same swap \\
Deliberate fault injection & severe & stamped & three of four injected faults invisible to the conservation residual; the fourth moves it by thirteen orders \\
Line-weighted against shell ratio & severe & stamped & ratio $0.948$ to $1.000$ over $784$ pairs, below unity at every one; census $44$ against $45$ \\
Escape-factor cross-check & severe & stamped & $\sigma_0$ against an independent module, \SI{0.059}{\percent} apart against a \SI{1}{\percent} raising tolerance \\
Pinned input hashes & severe & stamped & both SHA-256 match byte for byte \\
Recomputation guards on the figures & severe & code & every plotted row rebuilt from $\Lmat$ at $10^{-9}$, raises \\
Truncation convergence of $\ffrac{3}-\ffrac{4}$ & severe & stamped & increments one-signed and geometric at $245$ of $280$ defended points, ratio median $0.67$; extrapolated \SI{-0.46}{\percent} at the benchmark, \SI{+1.6}{\percent} at the cold corner, within \SIrange{-0.82}{+1.03}{\percent} on the defended set. Tests one quantity, not the state space \\
Gate~A, detailed balance & tautological & stamped & $<\SI{0.01}{\percent}$, but $K_{\rm deexc}$ was derived from $K_{\rm exc}$ \\
Column-sum conservation & tautological & stamped & $3.61\times10^{-11}$; see the fault-injection row for what it cannot see \\
The $\tanh$ bound & tautological & notes & holds at all $248$ points, and holds on corrupted input \\
Gate~E, timescale hierarchy & tautological & stamped & $\Msep \ge 86.8$ against a threshold of $1$ \\
Gate~B, coronal limit & weak & stamped & worst convergence $0.0034$ against a tolerance of $0.5$ \\
Gate~C, Saha ceiling & weak & stamped & monotone and below Saha; approach fraction \SI{1.5}{\percent}, excluded from the flag \\
\midrule
Fujimoto App.~4B timescales & external, passes & stamped & $0.34\times$ and $1.77\times$ at a grid point his example lands on \\
ACD96 against this model's recombination coefficient & external, passes & stamped & all $400$ points within \SI{11}{\percent}; $\eta$ from $0.894$ to $1.009$, median $0.982$, \SI{3.3}{\percent} at the benchmark (logarithmic read) \\
Fujimoto Table~4.1, $r_0$ & external, passes at $\lg\nel = 21$ & stamped & \SIrange{0.1}{1}{\percent} for $n\ge4$ at the higher density; \SIrange{12}{14}{\percent} at $p=3,4$ at $\lg\nel = 18$; $r_0(2) = 0.7392$ against $0.730$ \\
Fujimoto Table~4.1(b), $r_1$ & external, \textbf{open} & stamped & $8.3\times$ low at $p=3$. Truncation, Lyman trapping and the $\ell$~closure each tested and each eliminated; unexplained (\S\ref{sec:fujimoto}) \\
Gate~D, SCD96 against ADAS & \textbf{FAIL} & stamped & $\eta \in [5.73,\ 8346]$, $0$ of $400$ within a factor $2$. Rebuilt on the ground-fed channel alone and read logarithmically it reaches $400$ of $400$, $\eta$ from $0.702$ to $1.004$, median $0.828$; the gate itself is unchanged \\
\bottomrule
\end{tabular}
\end{table}

\subsection{What the model may be used for}

Three statements are supported by the evidence above, and they are the ones
Chapter~\ref{ch:results} relies on.

\textbf{The equations are solved correctly.} The superposition residual, the
independent re-derivation of all six central results, the bit-for-bit
regeneration, and the grid-wide conditioning measurement together establish that
the linear algebra does what Chapter~\ref{ch:theory} says it does, to a
precision many orders of magnitude finer than any quantity the thesis reports.

\textbf{The input rates are traceable and benchmarked where the observable
lives; one derived coefficient there is not.} The rates feeding $n \le 4$, which is where $H_\alpha$ and $H_\beta$
originate, agree with an independent R-matrix calculation at a mean absolute
error of \SI{13.0}{\percent}, and the two transitions this thesis leans on
hardest, $1s\to2p$ and $2p\to3d$, agree to between \SI{2.4}{\percent} and
\SI{13.3}{\percent} across the temperature range. Above $n=4$ the atomic data are
less certain and the truncation is measurable. Section~\ref{sec:convergence}
tests one component of the mechanism, the sensitivity $\ffrac{3}-\ffrac{4}$,
and finds it converged to within \SI{1.03}{\percent} at the defended points
where the geometric tail can be summed, in a stamped run
(\texttt{validation/nmax\_downward\_scan/}); under the same truncation
$u^{\mathrm{CRE}}$ is converged to within \SI{0.13}{\percent} at the $257$
defended points where its increments are geometric, $\Delta$ to within
\SI{0.57}{\percent} and $\tauslow$ to within \SI{1.6}{\percent} at all $280$,
and $\Delta\ln u$ across the defended heating steps moves by at most
$2.3\times10^{-5}$ between $n_{\max}=14$ and $15$ against a median step of
$0.20$ (Appendix~\ref{app:numerics}). What is measured
for the reservoir is a scan rather than a truncation: weakening the ionisation
and three-body exchange of the top six shells eightfold moves $\epsplat$ by
\SI{0.7}{\percent} at the median over the defended pairs
(Section~\ref{sec:ion_recomb}). All of that concerns the rates going in. The
coefficients that come out at the two shells the observable uses are a separate
matter, and one of them is in open disagreement: the ground-fed $r_1(3)$ and
$r_1(4)$ of this model sit factors $8.3$ and $4.4$ below Fujimoto's tabulation
at the lowest density, which is a deficit in $a_3$ and $a_4$ themselves and so
in the sensitivity built from them (Section~\ref{sec:fujimoto}). Benchmarked
inputs are therefore not the same thing as a benchmarked susceptibility, and
this thesis can claim the first but not the second.

\textbf{The structural results are robust to the choices that could have
manufactured them.} Closing the ion boundary changes the persistence census by
one part in $325$. Scaling $\ell$-mixing over two decades moves the observable
by under \SI{0.3}{\percent} where the thesis makes claims. The single-slow-mode
estimate of Eq.~\eqref{eq:elm_bound} reproduces the true time-averaged error to
within \SI{2.2}{\percent} at \SI{100}{\micro\second} over the $448$ pairs above
\SI{2}{\electronvolt}.

\subsection{What it may not be used for}

\textbf{No effective coefficient from this model should be compared against a
database without first separating its two supply channels.} Gate~D as written
fails at every point on the grid because it compares a coefficient that mixes
both channels against a table that separates them (Section~\ref{sec:gate_d});
rebuilt channel by channel it agrees within a factor two at all $400$ points
for SCD and ACD, and the model's total effective ionisation coefficient remains
a quantity ADAS does not tabulate.

\textbf{No claim rests on the terminal shell}, whose population is high by a
factor consistent with truncation, or on the $\ell$-distribution inside the
bundled shells, which is assumed rather than measured.
Section~\ref{sec:bundling_check} tests that assumption: $\ell$-mixing licenses
it over most of the grid, and at the two highest densities below
\SI{3}{\electronvolt}, where the locally evaluated mixing loses to collisional
loss, detailed balance among shells within $10^{-4}$ of Saha--Boltzmann
equilibrium licenses it instead; the adjacent-shell rates themselves are not
$\ell$-blind, and no grid point escapes both criteria.

\textbf{Nothing here validates the physics that is absent.} Every check in this
chapter tests a $43$-state atomic hydrogen model against itself, against
arithmetic, and against other atomic hydrogen models. None of them tests
transport, molecular channels, or radiation trapping, because none of those
appears in $\Lmat$ at all. A model can pass every check in
Table~\ref{tab:ladder} and still be the wrong model for a divertor.
Chapter~\ref{ch:limits} is where that question is asked, and it is asked with
numbers rather than with hedges.

\medskip
The instrument is now characterised: what it computes, how precisely, on which
region of the grid, and where the evidence for it runs out. What it says is the
subject of Chapter~\ref{ch:results}.

%     breaks no existing cross-reference.
%
%  2. OVERLAP, RESOLVED BY HOMES. Three passages exist in both chapters:
%       chapter6 sec:open_system   vs  Section~\ref{sec:attacks}   (home: ch4)
%       chapter6 sec:r1_deficit    vs  Section~\ref{sec:fujimoto}
%       chapter6 sec:nmax          vs  Section~\ref{sec:convergence}
%     The Chapter 6 versions carry the SCOPE consequence; the Chapter 4
%     versions carry the EVIDENTIAL standing. On the r_1 deficit: the table
%     is bundled-n (backlog I.1, 10 Sep 2026; findings_10 ADDENDUM I), the
%     l-closure was tested and does NOT account for the gap, and the deficit
%     DOES reach a_3 and a_4; its consequence for Delta, the cap and S is
%     quantified in the Status paragraph of sec:fujimoto
%     (validation/fujimoto_r1_consequence/). Do not reintroduce the earlier
%     "does not propagate" sentence.
%
%     [17-18 Sep 2026] verify_bundling_psm20.py HAS now been executed
%     (sec:bundling_check in chapter 6, outputs/bundling/); an earlier
%     version of this comment said the "never executed" sentence was still
%     true and must not be corrected. It was corrected. The n_max
%     convergence result carried in Section~\ref{sec:convergence} is now
%     stamped (verify_nmax_downward_scan.py, 18 Sep 2026); the working-note
%     values (ADDENDUM D.8) it replaced had the opposite sign at the
%     benchmark and the cold corner, and the text records that. Its \todo asking for provenance on the
%     4.9-6.2x terminal-shell excess is also closed there, by the internal
%     terminal-vs-interior measurement (a_p at [49,0]: 9.4x at p=8, 10.7x at
%     p=9, 6.3x at p=10), which does not bracket it (band 6.3 to 10.7 against
%     4.9 to 6.2) but lies within a factor 2.2 of it and shares the density
%     trend; stamped 21 Sep 2026 by verify_terminal_vs_interior.py.
%
%  3. chapter3.tex:222-225 asserts a severity for the conservation check that
%     Section~\ref{sec:conservation_severity} measures and refutes. Either
%     soften line 224 or let this chapter carry the correction; do not leave
%     both standing as written.
%=============================================================================

\chapter{Does the Model Work?}
\label{ch:validation}

Chapter~\ref{ch:theory} built a $43\times43$ rate matrix, eliminated the fast
variables, and extracted from the result a closed-form statement about how a
line ratio responds to its reservoir. Chapter~\ref{ch:results} will use that
machinery to make quantitative claims about a diagnostic in use on real
machines. Between the two sits a question that no amount of further algebra
answers: why should anyone believe a number this model produces?

The answer is not a single verdict. It is a graded one, because the
checks that have been run on this model are not of equal worth, and several of
the ones that look most impressive are worth almost nothing. This chapter sorts
them.

\medskip
\noindent\textbf{What this chapter takes as given.} The state space, the atomic
data and their sources, and the benchmark of the excitation rates against an
independent R-matrix-with-pseudostates (RMPS) calculation are established in
Chapter~\ref{ch:atoms}, Sections~\ref{sec:atomic_uncertainty}
and~\ref{sec:atoms_corrections}, and are not repeated. The construction of
$\Lmat$, the conservation identity, the quasi-steady-state (QSS) elimination,
and the two-channel decomposition of the excited populations come from
Chapter~\ref{ch:theory}. This chapter asks only what the checks run on that
construction are entitled to establish.

\medskip
\noindent\textbf{What this chapter does not do.} It reports no result about
divertor plasmas. Every number below is about the instrument, not about the
physics the instrument is pointed at. Keeping those separate matters here more
than anywhere else in the thesis: a model that predicts something interesting
is not thereby a correct model, and a check that uses the interesting
prediction as its evidence has verified nothing.

%-----------------------------------------------------------------------------
\section{What a check can and cannot show}
\label{sec:gate_philosophy}
% QUESTION: what distinguishes a check that could have failed from one that
% could not, and how is that distinction measured rather than asserted?

\subsection{Three things that get called validation}

Three activities are routinely reported under one heading, and they license
different conclusions.

\emph{Verification} asks whether the equations that were intended got solved
correctly. Conservation residuals, detailed balance, matrix identities,
convergence with tolerance, agreement between two independent code paths: all
of these compare the model against itself or against arithmetic. They can
detect a transposed index. They cannot detect a rate coefficient that is wrong
by a factor of three, because a uniformly wrong dataset satisfies every one of
them.

\emph{Validation} asks whether the model reproduces something established
independently of it. Published coefficients from another group's code, measured
cross sections, a tabulated limit. Only this class of check can reach the
atomic physics.

\emph{Analysis} asks what the model then says. That is Chapter~\ref{ch:results},
and none of it appears here.

\subsection{Severity, and how to measure it}

The useful property of a check is not whether it passes. It is the size of the
set of wrong models it would reject. Call that its severity.

The failure mode is the one a mass balance on a continuous stirred-tank
reactor shows. Compute the outlet stream as inlet minus consumption, then
verify that the balance closes. It closes to machine
precision, always, for any rate constant whatever, because the outlet was
constructed from the balance being tested. The check is real arithmetic and it
is entirely uninformative about the kinetics. A model can be wired correctly
and be wrong.

Severity is measurable, and two operations measure it.

\begin{enumerate}
\item \textbf{Provenance inspection.} Ask whether the quantity under test was
derived from the quantity doing the testing. If the de-excitation coefficients
were computed from the excitation coefficients through detailed balance, then a
detailed-balance check on them is arithmetic, not physics. Chapter~\ref{ch:atoms}
already states this for the rate dataset (Section~\ref{sec:two_energies}); this
chapter carries it through to the matrix.

\item \textbf{Fault injection.} Break the model deliberately, in a way that a
real mistake could plausibly break it, and see whether the check notices. A
check that survives a tenfold error in a rate coefficient has told you that a
tenfold error in that coefficient is invisible to it.
\end{enumerate}

Fault injection is the sharper of the two, because it produces a number rather
than an argument. It is used in Section~\ref{sec:conservation_severity} on the
conservation check, and its result is the least comfortable finding in this
chapter.

\subsection{The three grades used below}

Every check in this chapter is placed in one of three grades, and the grade is
stated wherever the check is reported.

\begin{description}
\item[Tautological.] Cannot fail except on an arithmetic bug, because the
quantity tested was constructed from the relation being tested. It confirms
that the code implements the construction. It is not evidence about the
physics, and reporting it as though it were makes the validation weaker, not
stronger, by displacing the checks that carry information.

\item[Weak.] Could fail in principle, but the tolerance is so much looser than
the measured value that only a gross error would trip it. A gate that passes
with a factor of $150$ in hand has told you the answer is not absurd.

\item[Severe.] A plausible wrong version of this model fails it. Either the
check has actually fired at some point in the project's history, or a specific
corruption is known to trip it.
\end{description}

Four quantities currently printed by this project's own scripts as validation
are tautological by this definition, and Table~\ref{tab:ladder} names them:
Gate~A, the column-sum conservation residual, the $\tanh$ bound, and Gate~E.
Sections~\ref{sec:conservation_severity} to~\ref{sec:severe_checks} give the
measurements behind each.

%-----------------------------------------------------------------------------
\section{The conservation check is exact by construction}
\label{sec:conservation_severity}
% QUESTION: does the column-sum identity detect a wrong rate coefficient?

Section~\ref{sec:conservation} established that each column of $\Lmat$ must sum
to the ionisation leak from the level that column represents, and reported the
measured residual of Eq.~\eqref{eq:conservation_measured} as
$3.61\times10^{-11}$, a relative residual (Section~\ref{sec:conservation}). The
$3.61\times10^{-11}$ is the maximum over every column at every one of the $400$
grid points; the fault-injection run below has its own baseline, $2.238\times
10^{-12}$, at the single point it was run at. The claim attached to it there is that ``an error in a
single off-diagonal element, a transposed index, or a missing back-reaction all
break this immediately.''

That claim is testable, and it is false.

\subsection{Four injected faults}

Four faults were introduced into the assembly, one at a time, and the residual
recomputed after each. The first three are the exact failure modes the claim
names. Table~\ref{tab:faultinject} gives the outcome.

\begin{table}[tbp]
  \centering
  \caption[Fault injection into the conservation check]{Four deliberate faults
  injected into the assembly of $\Lmat$, with the resulting change in the
  matrix and the conservation residual of
  Eq.~\eqref{eq:conservation_measured} recomputed after each. Three faults
  large enough to move individual matrix elements by up to
  $2.7\times10^{11}\,\si{\per\second}$ leave the residual unchanged at the
  baseline value; only an inconsistency between the Einstein coefficients $A$
  and the total radiative loss rates $\gamma$ is detected. Source:
  \texttt{outputs/CHANGE\_REPORT.md}~\S4.2.}
  \label{tab:faultinject}
  \begin{tabular}{lrl}
    \toprule
    injected fault & $\max|\Delta \Lmat|$ (\si{\per\second}) & residual \\
    \midrule
    none (baseline)                     & $0$                  & $2.238\times10^{-12}$ \\
    $K_{\rm exc}(1s\to2p)$ $\times$ ten   & $6.9\times10^{5}$    & $2.238\times10^{-12}$ \\
    all de-excitation deleted           & $1.9\times10^{11}$   & $2.171\times10^{-12}$ \\
    excitation array transposed         & $2.7\times10^{11}$   & $1.837\times10^{-12}$ \\
    $A_{pq}$ halved, $\gamma_p$ left alone & $3.1\times10^{8}$   & $3.63\times10^{1}$ \textbf{caught} \\
    \bottomrule
  \end{tabular}
\end{table}

A tenfold error in the single most important excitation coefficient in the
model, the one that supplies the entire ground-fed channel of
Section~\ref{sec:R_depends_on_b}, changes the residual in no digit. Deleting
every de-excitation process in the matrix, which is not a subtle mistake,
changes it in the third digit and leaves it three orders of magnitude below any
tolerance. Transposing the whole excitation input array, the classic
index-order mistake, does the same.

\subsection{Why, and what the check is actually worth}

The reason is in the assembly, not in the physics. The diagonal of $\Lmat$ is
built as minus the column sum of the same off-diagonal array that the check
then sums. Whatever numbers are in that array, correct or corrupted, the column
sum returns the ionisation leak. The identity of Eq.~\eqref{eq:conservation} is
imposed by the construction rather than tested by it. This is the reactor mass
balance of Section~\ref{sec:gate_philosophy} exactly.

The one fault that was caught is the one that breaks the construction. In the
notation of Eq.~\eqref{eq:level_balance}, $A_{pq}$ is the Einstein coefficient
for spontaneous emission from level $p$ to level $q$, in \si{\per\second}, and
$\gamma_p = \sum_{q<p} A_{pq}$ is the total radiative loss rate from $p$, also in
\si{\per\second}. The individual $A_{pq}$ enter the off-diagonal elements while
their sum $\gamma_p$ enters the diagonal, so halving one without the other makes
the two inconsistent, and the residual rises from $2.238\times10^{-12}$ to
$36.3$, thirteen orders of magnitude. The check raises rather than warns, it
costs nothing to run, and on this evidence that is the width of what it covers.

The rate magnitudes are reached only by Section~\ref{sec:atomic_uncertainty}'s
comparison against RMPS and by the external comparisons of
Section~\ref{sec:fujimoto}, and by nothing internal.

\paragraph{The grid point of the table.} Three of the four
$\max|\Delta\Lmat|$ values in Table~\ref{tab:faultinject} scale with $\nel$, so
the table fixes its own grid point. Sweeping all $400$ grid points and matching on those three values
identifies $[23,5]$, the benchmark, with a worst disagreement of
\SI{1.63}{\percent}, the rounding of the printed table to two significant
figures; every residual in Table~\ref{tab:faultinject} then reproduces to the
digits shown (\texttt{verify\_fault\_injection.py}, stamped in
\texttt{validation/fault\_injection/}). The script rebuilds $\Lmat$ through the
pipeline's own assembly rather than patching the stored matrix, and the rebuild
reproduces \texttt{L\_grid} exactly before anything is injected, which makes the
faults faults in the model rather than in a re-implementation of it.

%-----------------------------------------------------------------------------
\section{The gate suite, read one gate at a time}
\label{sec:gates_abc}
% QUESTION: for each automated gate, what quantity does it compare, against
% what tolerance, and what would have to be wrong for it to fail?

Five automated gates run over the grid in \texttt{src/validation/validate\_gates.py}
and write \texttt{validation/gate\_summary.txt}. Their headline outcome is four
passes and one failure. Read as a headline it says little, because the five
differ in severity by a wide margin and two of them cannot fail. This section
takes them one at a time; Gate~D, the failure, gets
Section~\ref{sec:gate_d} to itself.

\subsection{Gate A: detailed balance, and where its physics content actually
lives}

Gate~A takes every excitation rate coefficient $K_{\rm exc}(i\to j)$ in the
dataset, in \si{\cubic\centi\metre\per\second}, and predicts the reverse
coefficient from the Boltzmann relation between forward and reverse rates in a
Maxwellian electron distribution at temperature $\Te$,
\begin{equation}
  K_{\rm deexc}(j\to i)
  = K_{\rm exc}(i\to j)\,\frac{g_i}{g_j}\,
    \exp\!\left(\frac{E_j - E_i}{k_B\Te}\right),
  \label{eq:gateA}
\end{equation}
then compares the prediction against the stored reverse coefficient at three
temperatures spanning the grid. Here $g_i$ is the dimensionless statistical
weight of level $i$, $E_j-E_i$ is the excitation energy in
\si{\electronvolt}, and $k_B$ is Boltzmann's constant. The relation holds only
for a Maxwellian distribution, which Section~\ref{sec:maxwellian} established
and which is assumed throughout. The gate passes
if the maximum relative error is below \SI{0.01}{\percent}
(\texttt{validate\_gates.py}:118). Measured: below \SI{0.01}{\percent} over all
$819$ pairs at all three temperatures.

\textbf{Grade: tautological.} The stored de-excitation coefficients were
produced from the excitation coefficients through Eq.~\eqref{eq:gateA} in the
first place (\texttt{compute\_K\_CCC.py}, function \texttt{detailed\_balance}),
as Section~\ref{sec:two_energies} states. Gate~A confirms that the same formula
was applied twice with the same energies and weights. That is a wiring check,
and Chapter~\ref{ch:atoms} already flagged it as one.

The physics content of detailed balance in this project sits somewhere else and
is currently invisible in the outputs. \texttt{src/parsers/qc\_ccc.py}:87--168
tests whether Bray's \emph{raw} convergent-close-coupling (CCC) cross sections,
which this project did not produce, satisfy microscopic reversibility on their
own. The ratio of the two directions has mean $0.9995$ with standard deviation
$0.0032$, and \SI{98.7}{\percent} of comparisons fall within \SI{1}{\percent}
(\texttt{ccc\_qc\_report.png}). A \SI{0.05}{\percent} mean bias with
\SI{0.32}{\percent} scatter on external data is a genuine external check, and it
is the one to lead with.

\subsection{One energy ladder on both sides, which is severe}

There is a second detailed-balance result, and unlike Gate~A it could have
failed.

A model of this kind carries energies twice: once in the Boltzmann factors that
relate forward and reverse collisional rates, and once in the Saha factors that
relate ionisation to three-body recombination. If two different energy tables
were used, or one table were indexed differently from the other, the model would
still conserve particles, still satisfy Gate~A, and still produce plausible
populations. Nothing internal would notice.

Testing ionisation against three-body recombination through the Saha relation
at every level, with CODATA constants on the Saha side, gives a ratio of
$0.999997$ for every one of the $43$ states at every temperature, identical
across all $2150$ entries to $1.4\times10^{-15}$
(\texttt{verify\_saha\_balance.py}, \texttt{validation/saha\_balance/}); the
residual is the rate module's rounded values of $h$, $m_e$ and the
electron-volt against CODATA, and nothing else. The point is not the
closeness to unity, which detailed balance imposes. It is the identity of the residual across
levels whose Saha exponents span the full binding-energy ladder of hydrogen,
from $\mathrm{e}^{13.6\,\mathrm{eV}/k_B\Te}$ for the ground state down to
$\mathrm{e}^{0.06\,\mathrm{eV}/k_B\Te}$ for the top shell. At
$\Te = \SI{1}{\electronvolt}$ those two factors differ by
$\mathrm{e}^{13.54}$, a factor of $7.6\times10^{5}$. Two energy ladders that
differed would produce a residual drifting across that range; the measured
residual does not drift, and its common value is the eight-figure rounding of
the stored state table. \textbf{Grade: severe.} The check could have failed and
did not.

The same audit gives the collisional side over the $819$ excitation pairs at
each of three temperatures, $2457$ comparisons in total: maximum deviation
$7.31\times10^{-9}$, median $7.5\times10^{-10}$, again consistent with the state
table's rounding.

\subsection{Gate B: the coronal limit}

At low enough density an excited level is populated by electron impact from the
ground state and empties radiatively before a second collision reaches it, so
its population per unit ground-state density and per unit electron density
becomes independent of $\nel$. Gate~B tests that plateau: it computes
$n_{2p}/(\ngs\nel)$ and $n_{3d}/(\ngs\nel)$ at the two lowest densities on the
grid and asks whether they have converged.

The tolerance is \SI{50}{\percent} per temperature
(\texttt{validate\_gates.py}:175) and the gate passes if \SI{70}{\percent} of
temperatures pass (line 181). Measured: all $50$ temperatures pass, worst
convergence $0.0034$.

\textbf{Grade: weak.} Passing with a factor of $150$ between the measured
convergence and the tolerance means the gate rejects only a gross error. It
does establish that the low-density end of the grid is in the coronal regime,
which Chapter~\ref{ch:results} relies on when it reads the low-density column
as ionising, and that is worth having. It is not evidence about rate
magnitudes.

\subsection{Gate C: what the Saha gate tests, which is not the Saha limit}

Gate~C is named for the high-density limit in which collisions dominate
radiation and populations approach their Saha--Boltzmann values. It computes
three quantities per temperature: whether $n_{2p}/\ngs$ increases monotonically
with density, whether it stays below the Saha value (with \SI{5}{\percent}
slack), and whether it approaches that value, defined as reaching $10^{-4}$ of
it.

The pass flag uses the first two and \textbf{excludes the third}
(\texttt{validate\_gates.py}:250; the \texttt{approaching} column is computed at
line 239, stored, printed at line 261, and never consulted). The exclusion is
not what makes the gate weak, because the excluded column would pass: the
measured approach fraction runs $0.0123$ to $0.0155$, comfortably above the
$10^{-4}$ the criterion asks for. What makes it weak is that the criterion is
$150$ times slacker than the measurement. Reaching $10^{-4}$ of Saha is not
approaching Saha, and neither is reaching \SI{1.5}{\percent}.

The measurement itself is the useful part. At
$\nel = 10^{15}\,\si{\per\cubic\centi\metre}$, the top of the grid, the $2p$
population sits at about \SI{1.5}{\percent} of its Saha--Boltzmann value, so
that level is nowhere near local thermodynamic equilibrium (LTE) anywhere on
this grid. Only $2p$ was tested, so this is a statement about $2p$ and not about
the whole distribution. It is the expected behaviour for an optically thin
plasma whose ion population is an externally imposed reservoir rather than a
Saha partner, and Chapter~\ref{ch:results} depends on it: it is why the two
supply channels of Section~\ref{sec:R_depends_on_b} never merge, and why
$\ffrac{3}-\ffrac{4}$ does not vanish at high density.

\textbf{Grade: weak, and mis-named.} What Gate~C tests is monotonicity in
density and a Saha ceiling. Both would fail on a sign error in the recombination
source, which is worth something. Neither tests the limit the gate is named for.

%-----------------------------------------------------------------------------
\section{Gate E: the timescale hierarchy}
\label{sec:gate_e}
% QUESTION: is the two-clock structure of Chapter 3 present everywhere on the
% grid, and how demanding is the test that says so?

Gate~E diagonalises $\Lmat$ at every grid point, orders the negative
eigenvalues, and forms $\Msep = \tauslow/\taurel = |\lambda_1|/|\lambda_0|$. It
passes if $\Msep > 1$ at \SI{95}{\percent} of points
(\texttt{validate\_gates.py}:409--412).

Measured over all $400$ points: $\Msep$ ranges from $86.77$ at
$[49,7]$ ($\Te = \SI{10}{\electronvolt}$,
$\nel = \SI{e15}{\per\cubic\centi\metre}$) to $1.72928\times10^{9}$ at the cold
corner; $\tauslow$ spans \SI{75.4}{\nano\second} to \SI{67.2}{\second} and
$\taurel$ spans \SI{0.87}{\nano\second} to \SI{38.9}{\nano\second}
(project working note; Table~\ref{tab:ladder}). The gate asks for $\Msep>1$. The smallest value
anywhere is $86.8$.

\textbf{Grade: tautological as posed, and the real content is elsewhere.} The
threshold sits between two and seven orders of magnitude below the measured
values. No matrix this construction has produced comes within a factor of $86$
of failing, and none plausibly could, because the fast radiative rates and the
slow ionisation rate that set the two eigenvalues are separated by construction
of the physics rather than by a numerical accident.
The reportable quantity is the range, not the pass, together with the fact that
three independent implementations produce it bit-for-bit: \texttt{qss\_analysis.py},
Gate~E itself, and an unconditional recomputation stored in
\texttt{timescales\_unfiltered\_CHECK.npz}. That agreement is what caught the
eigenvalue-filter error of Section~\ref{sec:errors_found}, and a check that has
actually fired is worth more than a threshold nothing approaches.

Two properties of these numbers must travel with them, and both come from
Chapter~\ref{ch:theory}. $\tauslow$ is boundary-dependent, which
Section~\ref{sec:attacks} quantifies. And the slow eigenvalue is not an
equilibration rate: measured against the ground-state ionisation rate
coefficient $K_{\rm ion}(1s)$, in \si{\cubic\centi\metre\per\second}, it
satisfies $|\lambda_0| = 2.286\,K_{\rm ion}(1s)\,\nel$ across the grid
(project working note; Table~\ref{tab:ladder}), so $\tauslow$ is the collisional--radiative
ionisation lifetime of the ground-state reservoir and not the time for anything
to come to equilibrium.

\subsection{The gates do not stop anything}

No gate raises on failure. Each prints \texttt{PASS} or \texttt{FAIL} and
returns a flag; \texttt{main} sets \texttt{all\_pass = False}, prints a warning,
and exits normally. \textbf{The script returns exit code zero whether or not
Gate~D fails.} The residual gates elsewhere in the project do raise; this one
does not, and that asymmetry is a defect.

%-----------------------------------------------------------------------------
\section{The checks that could have failed}
\label{sec:severe_checks}
% QUESTION: which checks in this project would a plausible wrong model fail?

Everything so far has been a demotion. This section is the other half, and it
is short by design: a small number of severe checks is worth more than a long
list of consistency tests, and these are the ones the thesis leans on.

\subsection{The superposition residual}
\label{sec:superposition}

The entire mechanism of Chapter~\ref{ch:results} rests on one structural claim:
that the excited populations $\nF$ of Section~\ref{sec:qss} split exactly into a
ground-fed and a recombination-fed part,
$\nF = \bm{a}\,\ngs + \bm{c}\,\nion$, where the network responses $\bm{a}$ and
$\bm{c}$ of Section~\ref{sec:R_depends_on_b} are dimensionless and independent
of the two reservoir densities. If that split is not exact, the ground-fed
fractions $\ffrac{p}$ of Section~\ref{sec:ground_fed_fraction} are not well
defined and the closed form for the diagnostic error does not exist.

The check solves the full system three times at each grid point and each step
direction, once with each reservoir alone and once with both, and measures the
relative residual of the sum against the joint solve. Across the $784$
grid-point and step-direction pairs it covers, out of the $800$ the grid admits,
the largest residual is $3.075\times10^{-14}$
(\texttt{plateau\_gridmap.txt}, line 19).

\textbf{Grade: severe}, for a specific reason. Feeding the check a corrupted
input in which the $n=3$ and $n=4$ entries are swapped, which is the mistake
that would most easily go unnoticed in a two-shell observable, trips it
(project working note; Table~\ref{tab:ladder}). It is one of only two checks in the
project reported to catch that swap.

\subsection{The reduced-versus-full ratio test}

The second is its companion. The observable $\Robs = n_3/n_4$ can be computed
two ways: from the full $43$-state steady-state solve, and from the reduced
two-channel form of Eq.~\eqref{eq:R_of_b} with the ground-fed fractions
substituted in. The two routes share no code path beyond the matrix itself. At
the benchmark point $\Te = \benchTe$, $\nel = \benchne$, an independent
re-implementation obtained $\Robs = 0.777768$ by three routes against
Chapter~\ref{ch:theory}'s recorded $0.7778$
(project working note; Table~\ref{tab:ladder}). A $3\leftrightarrow4$ swap breaks the
agreement.

\subsection{Pinned hashes on the inputs}

The most common way for a computation like this to become wrong is not an
algebra error. It is a stale array: a matrix regenerated after a correction,
combined with an analysis output produced before it. This project has three
retractions in its history and at least one of them had that shape.

\texttt{src/validation/preflight.py} pins the SHA-256 hashes, fixed-length
fingerprints that change if any byte of a file changes, of the two canonical
inputs and compares them byte for byte before any verification script
runs. Recomputed for this chapter:
\begin{center}
\begin{tabular}{ll}
  \toprule
  file & SHA-256 (first 16 hex digits) \\
  \midrule
  \texttt{data/processed/cr\_matrix/L\_grid.npy} & \texttt{2d92b58e1693107d} \\
  \texttt{data/processed/cr\_matrix/S\_grid.npy} & \texttt{7822f536590c76ba} \\
  \bottomrule
\end{tabular}
\end{center}
Both match the pinned baselines at \texttt{preflight.py}:37--38 exactly, so the
matrix on disk at the time of writing is the one the pins were taken from; the
statement is only as good as the discipline of running the preflight before
each analysis.

\textbf{Grade: severe on the data, incomplete on the scripts.} Of the six
scripts the same preflight claims to check, \texttt{verify\_ch3\_groupB.py} and
\texttt{verify\_grid\_coverage.py} do not exist on disk under those names; a
missing script is downgraded to a non-blocking advisory
(\texttt{preflight.py}:258--262) while the summary still prints ``All checks
passed''. The hash check is sound; the sentence after it overstates what ran.

\subsection{Guards that recompute rather than trust}

\texttt{src/analysis/make\_ch5\_figures.py} does not plot the stored results. It
rebuilds every plotted quantity from $\Lmat$ and the recombination source
$\Svec$ of Eq.~\eqref{eq:source} and compares against
the stored comma-separated tables at a relative tolerance of $10^{-9}$, raising
with the offending grid point, column, and matrix hash if any row disagrees
(lines 393--432). It reads the temperature step size out of the file header
rather than redeclaring it, so a figure cannot be drawn against a step
convention different from the one that produced the data. It also verifies that
the untrapped run of the radiation-trapping sweep reproduces the canonical
matrix exactly, without which no trapped number in Chapter~\ref{ch:limits} would
be interpretable as a difference.

The same script carries a negative control: it computes the correlation between
$\Msep$ and the diagnostic error both raw and after controlling for temperature
and density, and refuses to draw the figure if the controlled correlation fails
to change sign (Section~\ref{sec:M_does_not_predict} draws the conclusion).

\subsection{Two checks that fired, and one that was withdrawn}

Three episodes say more than the pass/fail table, because in each the check did
something.

The Lyman-$\alpha$ line-centre cross-section gate, an independent implementation
in \texttt{escape\_factor.py} against the project's own Einstein coefficients,
fired at \SI{1156}{\percent} on its first run, located a $4\pi$ factor, and now
agrees to \SI{0.059}{\percent} against a raising \SI{1}{\percent} tolerance
(Section~\ref{sec:opacity}).

\texttt{verify\_fujimoto\_table41.py} carries a deliberately broken variant with
proton-impact $\ell$-mixing removed. It returns a recombination-fed population
coefficient $r_0(2) = 104.5$, and the script flags it itself as unphysical on
the ground that a level cannot exceed Saha equilibrium. Production gives
$0.7392$. A check that fires on a known-bad input is the definition of a severe
one.

Third, a semigroup propagator identity,
$\exp(\Lmat t_1)\exp(\Lmat t_2) = \exp(\Lmat(t_1+t_2))$ for a matrix exponential
computed by the same routine, returned $0.000\times10^{0}$ by construction and
was \emph{withdrawn by the author} (Section~\ref{sec:bridge}), the judgement
this chapter applies to Gate~A, applied earlier and unprompted.

\subsection{The bound that is a theorem}

Section~\ref{sec:bound} derives an upper bound on the diagnostic error of the
form $\tanh(\Delta/4)$, where $\Delta$ is the separation between the two shells'
switching points. A script reports that the bound is honoured at all $248$
points tested, and that line reads as validation.

It is not. Here $\Delta$ is the dimensionless separation between the two shells'
switching points, in units of the logarithm of the reservoir strength. Given
$\bm{a},\bm{c}\ge0$, which Section~\ref{sec:R_depends_on_b} proves from the
M-matrix structure of $-\Lmat_{FF}$, the bound is a theorem: it cannot fail
except on an arithmetic bug. Fed corrupted input it still passes. Swapping
shells $3\leftrightarrow4$ sends $\Delta\to-\Delta$ and leaves both $|\Delta|$
and the logarithmic change in $\Robs$ that it bounds unchanged, so the gate
cannot see it. Swapping $c_3\leftrightarrow c_4$ also
passes, but for a different reason: it moves $|\Delta|$ from $1.945614$ to
$0.939493$, a factor $2.07$, and the gate still passes because it compares
$|f_3-f_4|$ against $\tanh(|\Delta|/4)$ evaluated from the same corrupted
$\bm{c}$. The bound is satisfied by whatever $\bm{a},\bm{c}\ge0$ it is
handed.
\textbf{Grade: tautological.} The script's own docstring calls it a wiring
check; the line it prints does not, and the printed line is what a reader sees.

The bound is a result, and a useful one, because it holds for any atomic dataset
whatever. It is not evidence that this atomic dataset is right.

%-----------------------------------------------------------------------------
\section{Conditioning, precision, and state-space convergence}
\label{sec:conditioning}
% QUESTION: could the answers be artifacts of finite-precision arithmetic or of
% where the level ladder was cut off?

Two numerical questions sit underneath every result in this thesis. The
elimination of Section~\ref{sec:qss} inverts a $42\times42$ matrix whose entries
span many orders of magnitude, which invites the question of whether the inverse
means anything. And the level ladder was cut at $n_{\max}=15$, which invites the
question of what the discarded levels would have changed.

\subsection{The inverse is well behaved, and here is how well}

The relevant quantity is the condition number of $\Lmat_{FF}$, the block of
$\Lmat$ coupling the $42$ excited levels to each other
(Section~\ref{sec:qss}), and it is the factor by which a relative perturbation
of the matrix can be amplified in the solution.
Computed from the canonical $\Lmat$ at all $400$ grid points by
\texttt{src/validation/verify\_operator\_conditioning.py} and stamped to
\texttt{validation/operator\_conditioning/operator\_conditioning.csv}: minimum
$1.4821\times10^{3}$ at $[0,0]$ ($\Te = \SI{1.000}{\electronvolt}$), maximum
$1.7436\times10^{5}$ at $[33,7]$ ($\Te = \SI{4.715}{\electronvolt}$), median
$2.1044\times10^{4}$. In double precision, which carries about sixteen
significant digits, a condition number of $10^{5}$ costs five of them. Eleven
remain, and every conclusion in this thesis turns on the first three.

The range was first asserted in Chapter~\ref{ch:theory} on an unscripted
working note, which is not a check; the stamped script reproduces it to
\SI{0.14}{\percent} and \SI{0.21}{\percent}. The number was right the whole
time and was not evidence until it could be re-run.

Three further measurements bound the arithmetic. The eigenvalue condition
number is a different quantity from the matrix condition number above: it is
$1/|\bm{y}^H\bm{x}|$ for a left and right eigenvector pair, and it measures how
far a single eigenvalue moves under a perturbation of the matrix. For the two
eigenvalues the thesis uses it is of order unity
($1.71$ and $1.98$ at the benchmark point, $3.26$ and $2.59$ at the cold
corner), despite a largest absolute column sum of
$2.0\times10^{14}\,\si{\per\second}$; that large norm is a
statement about the fast radiative modes and does not propagate into
$\lambda_0$ or $\lambda_1$. Recomputing the two supply channels in single
precision moves them by $1.3\times10^{-5}$, so double precision is nowhere near
a cliff. And tightening the integrator changes nothing that matters: a relative
tolerance of $10^{-9}$ against $10^{-12}$ shifts the result by
$1.20\times10^{-8}$, and computing the propagator by a dense matrix exponential
against a Krylov action differs by $1.13\times10^{-10}$.

The eigenvalue condition numbers and the single-precision test cover three of
$400$ points; the tolerance comparisons cover two. The condition numbers and the
displacement are recomputed under the definition $1/|\bm{y}^H\bm{x}|$ with
unit-norm left and right eigenvectors by \texttt{verify\_residual\_claims.py}
(\texttt{validation/residual\_claims/}). They are point measurements, not grid
scans, and are quoted as such.

\subsection{Where the ladder was cut}
\label{sec:convergence}

The quantity that has to converge is not a population but the difference in
ground-fed fraction between the two Balmer shells, $\ffrac{3}-\ffrac{4}$, since
Chapter~\ref{ch:results} shows the diagnostic error to be proportional to it.

Rebuilding $\Lmat$ with successively lower $n_{\max}$, from $15$ down to $9$,
with the loss channels of every removed state returned to the diagonal of the
survivors, and recomputing $\ffrac{3}-\ffrac{4}$ at each truncation's own
$u^{\mathrm{CRE}}$ (\texttt{verify\_nmax\_downward\_scan.py},
\texttt{validation/nmax\_downward\_scan/}) gives increments of one sign at the
three named points, with the last measured ratio $d(15)/d(14)$ equal to
$0.88$ at the benchmark, $0.72$ at the cold corner $[0,0]$ and $0.33$ at the
ridge $[15,3]$; at the $245$ of the $280$ defended points where the increments
are one-signed and geometric the ratio has median $0.67$. Summing the geometric
tail with that ratio gives an extrapolated truncation error of
\SI{-0.46}{\percent} at the benchmark, \SI{+1.6}{\percent} at the cold corner
and \SI{+0.005}{\percent} at the ridge; holding $u$ at its $n_{\max}=15$ value
instead gives \SI{-0.38}{\percent}, \SI{+1.4}{\percent} and
\SI{+0.005}{\percent}, so the sign does not depend on that choice. A historical
project value with no producing script, not used as evidence here and superseded
by \texttt{verify\_nmax\_downward\_scan.py}, records
\SI{+0.9}{\percent}, \SI{-1.4}{\percent} and \SI{-0.55}{\percent} with a
common ratio of about $0.75$; the stamped run reverses both signs and the
conclusion that rested on them. Truncation makes
$\ffrac{3}-\ffrac{4}$ slightly too large at the benchmark and too small at the
cold corner, so it is not a bias that can be corrected by a single factor.
Over the defended points where the tail can be summed the extrapolated error
has median \SI{+0.002}{\percent} and runs from \SI{-0.82}{\percent} to
\SI{+1.03}{\percent}, smaller than the spread between the two step conventions
of Section~\ref{sec:step_convention} at every such point. At the $35$ defended
points of the density column $\nel = \SI{5.2e13}{\per\cubic\centi\metre}$ the
increments change sign among $n_{\max} = 12$ to $15$ and no tail can be
summed; the last measured increment there is below \SI{0.041}{\percent} of
$\ffrac{3}-\ffrac{4}$ at every point.

\paragraph{A trap in this test.} The first attempt dropped the top shells from
the state vector without removing their loss channels from the diagonal of the
surviving states, which then lost population to states that no longer existed,
and $\ffrac{3}-\ffrac{4}$ jumped. The stamped scan repeats the mistake
deliberately: dropping $n=15$ alone without the correction moves
$\ffrac{3}-\ffrac{4}$ by \SI{11.5}{\percent} at the benchmark,
\SI{12.6}{\percent} at the cold corner and \SI{5.0}{\percent} at the ridge,
where the correct truncation to $n_{\max}=14$ moves it by \SI{0.06}{\percent},
\SI{0.64}{\percent} and \SI{0.01}{\percent}; dropping six shells without it
gives \SI{17.9}{\percent}, \SI{25.9}{\percent} and \SI{13.3}{\percent} (a
working note recorded \SIrange{22}{37}{\percent}). The jump is an artifact of
the truncation procedure, not a dependence on $n_{\max}$. Removing state $k$
requires adding $\sum_{k\ \rm dropped} L_{ki}$ back to $L_{ii}$ for every
surviving $i$, as the production matrix does. Anyone repeating this test will
reproduce the spurious number first.

\paragraph{The terminal shell is over-populated, and by a measured amount.}
Truncating the ladder at $n_{\max}$ removes, from the balance of the terminal
shell, the collisional excitation out of it to $n > n_{\max}$ and the
de-excitation and cascade into it from those levels (cascade is downward, so
what is lost is a way out, not something ``to cascade out to''). The loss removed is not small. Continuing the $12\to13$, $13\to14$
and $14\to15$ excitation coefficients geometrically, the missing $15\to16$
channel would carry $1.30$ times the whole retained loss out of $n=15$ at the
benchmark and $1.20$ times at $[0,4]$, between $1.20$ and $1.35$ times at
every grid point, while radiative decay is below $10^{-6}$ of that loss and
ionisation is $7$ to \SI{10}{\percent} of it
(\texttt{validation/terminal\_shell\_budget/}). The top shell sits within
$6\times10^{-5}$ of Saha--Boltzmann, so for the total population the removed
exchange is nearly balanced and the sign of the bias is decided channel by
channel: the ground-fed channel, whose net flux through the top of the ladder
is upward, loses only an exit and can only be biased high, and it is the
ground-fed coefficient $r_1(15)$ that the external comparison finds high while
$r_0(15)$ agrees (Section~\ref{sec:fujimoto}). Comparing each shell's ground-fed coefficient $a_p$ as terminal
shell (operator rebuilt with $n_{\max}=p$, the loss channels into the removed states kept on the diagonal) with its
value in the full $n_{\max}=15$ operator gives at the Fujimoto point $[49,0]$ (\SI{10}{\electronvolt},
\SI{e12}{\per\cubic\centi\metre}) excesses of $9.4$ at $p=8$, $10.7$ at $p=9$ and $6.3$ at $p=10$, then $4.4$,
$3.0$, $2.1$ and $1.4$ at $p=11$ to $14$ (\texttt{verify\_terminal\_vs\_interior.py},
\texttt{validation/terminal\_vs\_interior/}). At the benchmark $p=8$ to $10$ give $14.2$, $12.3$ and $6.5$; over
the $280$ defended points the $p=8$ excess spans $8.2$ to $17.7$ and the $p=10$ excess $5.4$ to $7.5$. The excess
is confined to the ground-fed channel: the shell's total population and recombination-fed coefficient move under
\SI{1}{\percent} at the benchmark and at most \SI{32}{\percent} and \SI{30}{\percent} at the low-density named
points. A historical project value with no producing script, not used as evidence here and superseded by
\texttt{verify\_terminal\_vs\_interior.py}, records $9.4$, $10.7$ and
$6.3$ with neither quantity nor grid point named; the stamped run reproduces them to \SI{0.4}{\percent} as $a_p$
at $[49,0]$. The absence of a monotone trend over $p=8$ to $10$ holds there and at $47$ of the $280$ points;
from $p=9$ upward the excess falls at every grid point, at the benchmark from $p=8$. The external excess at the
terminal shell, $4.9$ to $6.2$ against published values (Section~\ref{sec:fujimoto}), lies outside the $6.3$ to
$10.7$ band, not inside it as the working note states, but within a factor $2.2$ of every value in it, and both rise
with density ($4.9$ to $6.2$ externally, $9.4$ to $17.7$ at $p=8$ internally at \SI{10}{\electronvolt}). That
agreement in magnitude and in the sign of the density trend makes truncation a sufficient explanation for the
$n=15$ discrepancy without invoking the atomic data. The fall with $p$ cannot be extrapolated to $p=15$: the
reference for $p=14$ is itself a ladder with one shell above it, so the internal measurement neither predicts nor
excludes $4.9$ to $6.2$ there.
No result in this thesis rests on the
top shell: the observable is built from $n=3$ and $n=4$, both fully resolved in
orbital angular momentum $\ell$ and eleven shells below the cut at
$n_{\max}=15$.

\paragraph{What has now been checked.} Above $n=8$ the shells are bundled and
their internal $\ell$-distribution is assumed statistical rather than computed.
The script that tests that assumption, \texttt{verify\_bundling\_psm20.py}, had
never been executed; two of the three defects found on its first run, the
synthesised temperature grid and the identically zero $\ell$-mixing array for
the bundled indices, were predicted here. Section~\ref{sec:bundling_check}
reports all three defects and the result.

%-----------------------------------------------------------------------------
\section{Is the shell ratio the line ratio?}
\label{sec:balmer_equivalence}
% QUESTION: the analysis computes n_3/n_4, but a spectrometer measures
% Halpha/Hbeta. What does the substitution cost?

Every analytic result in Chapter~\ref{ch:theory} is written for the shell ratio
$\Robs = n_3/n_4$, because the two-channel decomposition applies to shell
populations. A spectrometer does not measure shell populations. It measures the
$H_\alpha$ and $H_\beta$ line intensities, each of which is a sum over the
electric-dipole-allowed sublevel transitions, weighted by their Einstein
coefficients. Section~\ref{sec:radiative} established that these are not
proportional: $4f$ holds \SI{43.75}{\percent} of the $n=4$ shell by statistical
weight and contributes nothing to $H_\beta$, since $4f\to2p$ would change $\ell$
by two and is forbidden.

The substitution therefore has to be measured, not assumed.

The line intensities are built as $A$-weighted sums over resolved channels,
$H_\alpha$ from $3s\to2p$, $3p\to2s$ and $3d\to2p$, and $H_\beta$ from
$4s\to2p$, $4p\to2s$ and $4d\to2p$, with populations taken from the solve.
Nothing statistical is assumed anywhere in that construction: searching the two
producing scripts for statistical-weight expressions returns no matches, and the
$A$-value lookup raises unless exactly one row matches each transition.

Two measurements follow. The $4f$ fraction of the $n=4$ shell runs from
$0.4361$ to $0.4375$ across the grid, against the statistical value
$14/32 = 0.4375$; proton-impact $\ell$-mixing drives the shell to a statistical
distribution to better than \SI{0.3}{\percent} everywhere, with the departure
concentrated at $\nel = \SI{e12}{\per\cubic\centi\metre}$ where $\ell$-mixing
weakens and radiative decay competes. That fraction is a property of the full
equilibrium population; the quantity that sets the line-versus-shell difference
below is the $\ell$-distribution within each supply channel separately, which
is not statistical at the low-density edge. And the ratio of the transient error
computed on the line ratio to the same error computed on the shell ratio runs
from $0.9482$ to $0.9999$ over the $784$ point-and-direction pairs, below $1$ at
every one: \textbf{worst case \SI{5.2}{\percent}, at the heating step from
$9.54$ to \SI{10.0}{\electronvolt} at the lowest density}
(\texttt{verify\_weighted\_census.py}, stamped to
\texttt{validation/weighted\_census/}). At
$[0,4]$, the point carrying the largest error anywhere, the two agree to
\SI{0.06}{\percent} ($0.386683$ against $0.386903$). Because the ratio is below
$1$ everywhere, the shell census of Chapter~\ref{ch:results} is conservative
for the line observable: counted on the line ratio, $44$ instead of $45$ of the
$448$ pairs above \SI{2}{\electronvolt} exceed \SI{10}{\percent} at
\SI{100}{\micro\second}, with the worst case $0.1746$ instead of $0.1748$ at
the same point. The \SI{5.2}{\percent} scales inversely with the proton
$\ell$-mixing rate: with the within-shell $\ell$-changing rates scaled by a
common factor and the column sums of the operator preserved, the corner
departure is \SI{10.0}{\percent} at half the rate, \SI{2.6}{\percent} at double
and \SI{1.1}{\percent} at five times, while the benchmark line-to-shell ratio
stays between $0.9987$ and $0.9999$ (\texttt{validation/weighted\_census/},
$\ell$-mixing sensitivity block and \texttt{weighted\_census\_lmix.csv}); and
it sits on the grid boundary.

That \SI{0.06}{\percent} is a single-point figure and must not be quoted as a
grid-wide one. The grid-wide statement is the \SI{5.2}{\percent}.

\subsection{And is the plasma transparent to its own Balmer light?}

The shell-to-line substitution assumes the emitted photons escape. For the
Lyman series that assumption fails at the cold edge and
Chapter~\ref{ch:limits} treats it at length. For the Balmer lines themselves it
closes here, because it bears on the observable rather than on the matrix.

Section~\ref{sec:opacity} defines the optical depth $\tau$ and the escape factor
and treats the Lyman series, where transparency fails. For the Balmer lines the
absorbers are the $n=2$ atoms, $n_{2s}+n_{2p}$ in collisional--radiative
equilibrium with $\nion=\nel$, the Doppler width is taken at $\Tn=\Te$, and the
multiplet oscillator strengths are $f(2\to3)=0.641$ and
$f(2\to4)=0.119$~\cite{WieseFuhr2009}. The line-centre optical depth of
$H_\alpha$ across a \SI{10}{\centi\metre} slab is then $4.9\times10^{-7}$ at the
cold corner, $8.2\times10^{-4}$ at $[0,4]$ and $0.214$ at $[0,7]$, the single
worst cell on the grid; resolving the absorbers into $2s$ and $2p$ with the
model's own $A$ coefficients changes these by at most \SI{4}{\percent}. At
$[0,7]$ the population escape factor for a photon born at the slab centre
($\kappa_0 D/2 = 0.107$) is $0.928$ for $H_\alpha$ and $0.990$ for $H_\beta$, so
the observed ratio is lowered by \SI{6.3}{\percent} in that one cell; over the
full path ($\tau=0.214$) the figures are $0.861$ and \SI{12}{\percent}. Above
\SI{2}{\electronvolt} the largest $H_\alpha$ optical depth is $0.010$ at
$[15,7]$ and the ratio moves by less than \SI{0.7}{\percent} anywhere
(\texttt{verify\_balmer\_optical\_depth.py},
\texttt{validation/balmer\_optical\_depth/}). The Balmer lines are optically
thin over the region where this thesis makes quantitative claims. A historical
project value with no producing script, not used as evidence here and superseded
by \texttt{verify\_balmer\_optical\_depth.py}, records
$6.3\times10^{-7}$, $1.0\times10^{-3}$ and $0.263$ for the same three cells;
those values are $1.23$ to $1.28$ times the stamped ones and no convention
tested recovers them, and the escape factor above $0.85$ and ratio shift of at
most about \SI{10}{\percent} recorded with them are half-slab figures attached
to a full-slab optical depth. The conclusion holds with either set of values.

%-----------------------------------------------------------------------------
\section{Two tests of the reduction, and their expected outcomes}
\label{sec:attacks}
% QUESTION: the two objections most likely to destroy the result were tested;
% what was the falsifier in each case, and did it appear?

The checks so far ask whether the model computes what it intends. This section
asks something harder: whether two specific structural choices manufacture the
result. In each case the size of effect that would have destroyed the claim was
written down before the test was run, which is what makes the outcome
informative rather than reassuring.

\subsection{Does the open boundary manufacture the slow clock?}

$\Lmat$ describes an open system: atoms that ionise leave the accounting, the
recombination feed $\Svec\nion$ enters as a constant source rather than a matrix
element, and there is no forty-fourth state for ions to accumulate in. Since the
persistence claim of Section~\ref{sec:persistence} rests entirely on $\tauslow$
exceeding the duration of the disturbance, the first objection is that
$\tauslow$ is an artifact of the open boundary.

\textbf{The falsifier, named first.} Write $\taudrive$ for the duration of the
disturbance, fixed at \SI{100}{\micro\second} throughout this thesis as a
representative edge-localised-mode (ELM) time
(Section~\ref{sec:persistence}). At $[0,4]$, the worst point on the grid,
$\tauslow/\taudrive \approx 2300$. If closing the system shortened $\tauslow$ by
that factor, the slow clock would fall to $\taudrive$ itself and the persistence
would vanish. A factor of $2300$ refutes the claim.

Two grid labels appear in the table for the same worst case: the error is
evaluated at $[0,4]$, but the operator governing the decay after the step is the
post-step one, $\Lmat[1,4]$, and every timescale in this thesis names the
operator it belongs to (Chapter~\ref{ch:results}).

The manifold was closed by appending an ion state,
\begin{equation}
  \Lmat_c =
  \begin{pmatrix}
    \Lmat & \Svec \\
    -\operatorname{colsum}(\Lmat) & -\textstyle\sum\Svec
  \end{pmatrix},
  \label{eq:closure_ch4}
\end{equation}
so that every ionising atom lands in the new state instead of being deleted.
Each column of $\Lmat_c$ sums to zero to machine precision, and the construction
reproduces the closed-system value at the cold corner recorded in this
project's working notes, \SI{1.6226}{\second}, to five digits, which confirms it is the same
closure rather than a new one. The measured factors are given in
Table~\ref{tab:closure}.

\begin{table}[tbp]
  \centering
  \caption[Effect of closing the ion boundary on the slow timescale]{The slow
  timescale $\tauslow$ computed with the ion treated as an external reservoir
  (open) and with a forty-fourth ion state appended (closed), at four grid
  points. Over the whole grid the ratio has minimum $1.00$, median $1.00$ and
  maximum $41.4$: the boundary condition matters only where $\tauslow$ already
  exceeds the disturbance duration by three orders of magnitude. Every entry is
  reduced from \texttt{validation/ion\_closure/ion\_closure\_summary.csv}
  (\texttt{verify\_ion\_closure.py}) by
  \texttt{verify\_headline\_provenance.py}; the open column is the float64
  eigensolver, which at the cold corner differs from the artifact's 40-digit
  path by one part in $10^{4}$.}
  \label{tab:closure}
  \begin{tabular}{lrrr}
    \toprule
    point & $\tauslow$ open & $\tauslow$ closed & factor \\
    \midrule
    cold corner $[0,0]$ & \SI{67.23}{\second} & \SI{1.6226}{\second} & $41.4$ \\
    post-step operator $\Lmat[1,4]$ & \SI{0.23324}{\second} & \SI{0.015173}{\second} & $15.4$ \\
    ridge $[15,3]$ & \SI{2.027}{\milli\second} & \SI{1.999}{\milli\second} & $1.01$ \\
    benchmark $[23,5]$ & \SI{22.728}{\micro\second} & \SI{22.708}{\micro\second} & $1.00$ \\
    \bottomrule
  \end{tabular}
\end{table}

The shape of the distribution carries the argument: the factor is large only
where $\tauslow$ already exceeds $\taudrive$ by about a thousand, so the error
has already saturated, and it is unity everywhere $\tauslow$ approaches
$\taudrive$. Recomputing the entire persistence census with closed-system
$\tauslow$ moves the breakdown count from $202$ to $200$ of $680$ and the worst
case from $0.38682$ to $0.38563$, about one part in $325$; restricted to
$\Te \ge \SI{2}{\electronvolt}$, where Chapter~\ref{ch:results} makes
quantitative claims, the count is $45$ before and $45$ after.

The falsifier did not appear. A factor of $2300$ was needed; the measured factor
at that operator is $15$.

\textbf{The caveat, which must not be softened.} The closure adds an ion
reservoir; it does not close the electron balance, since $\nel$ is held fixed
while the ion population moves, and it adds no transport. It answers the
boundary condition on the atomic manifold and nothing else:
Section~\ref{sec:transport} states the objection it cannot answer, and
Section~\ref{sec:quasineutrality} the one it makes sharper rather than softer.

\subsection{Is the single-slow-mode estimate a bound on the time-averaged error?}

Chapter~\ref{ch:results} evaluates the error at the partial-equilibrium plateau
and builds from it the single-slow-mode estimate of Eq.~\eqref{eq:elm_bound},
which multiplies the plateau error $\epsplat$ of
Section~\ref{sec:eps_definitions} by the fraction of $\taudrive$ over which an
exponential decay on $\tauslow$ keeps it alive. It was first written as a lower
bound, on the assumption that the observable cannot decay faster than the
slowest mode. That assumption is testable, and an estimate quoted as though it
were the answer is a claim about accuracy that has to be measured.

Propagating the full $43$-state system through the step by direct
eigen-decomposition and integrating the error over \SIrange{0}{100}{\micro\second}
gives the comparison in Table~\ref{tab:lowerbound}.

\begin{table}[tbp]
  \centering
  \caption[Accuracy of the single-slow-mode estimate]{The plateau error, the
  true error averaged over \SIrange{0}{100}{\micro\second} from direct
  eigen-propagation of the full $43$-state system, and the single-slow-mode
  estimate of Eq.~\eqref{eq:elm_bound}, at five grid points. The averages are
  integrated at $400$ time samples per decade; at $1600$ per decade the five
  change by at most $4.7\times10^{-6}$ relative, and the $[0,4]$ average by
  $2.7\times10^{-10}$. Integration on a uniform \SI{25}{\nano\second} mesh, which
  under-resolves the nanosecond rise, reads low by $6\times10^{-5}$ to
  $5\times10^{-4}$ relative. The estimate reproduces
  the true average to better than \SI{0.8}{\percent} at the five points. The
  value below unity at $[0,4]$ is real: a deficit of $2.2\times10^{-5}$ that
  survives the denser grid. Source: \texttt{verify\_trajectory\_census.py},
  \texttt{validation/trajectory\_census/}.}
  \label{tab:lowerbound}
  \begin{tabular}{lrrrr}
    \toprule
    point & $\epsplat$ & true \SI{100}{\micro\second} average & estimate, Eq.~\eqref{eq:elm_bound} & true/estimate \\
    \midrule
    $[0,4]$ worst case & $0.386903$ & $0.386811$ & $0.386820$ & $0.99998$ \\
    $[15,3]$ ridge     & $0.180728$ & $0.175283$ & $0.174812$ & $1.00270$ \\
    $[15,5]$           & $0.103550$ & $0.072185$ & $0.071722$ & $1.00646$ \\
    $[23,5]$ benchmark & $0.063612$ & $0.011796$ & $0.011712$ & $1.00722$ \\
    $[0,0]$            & $0.083004$ & $0.083084$ & $0.083004$ & $1.00096$ \\
    \bottomrule
  \end{tabular}
\end{table}

At the five points tested the estimate sits within \SI{0.8}{\percent} of the
directly propagated average, and at $[0,4]$ it sits above it, so it is an
estimate and not a bound. The accuracy at these five points does not
generalise: over the $448$ pairs above \SI{2}{\electronvolt} the ratio of true
average to estimate departs from unity by up to \SI{2.2}{\percent} at
\SI{100}{\micro\second} and \SI{3.8}{\percent} at \SI{506}{\micro\second}
(Section~\ref{sec:persistence}). What matters here is that the assumption
underneath it was tested rather than asserted.

The reason is structural, and it is a property the shell ratio has and another
observable in this same project does not. The spectrum of $\Lmat$ has one slow
eigenvalue and then a band, with nothing between: at $\Lmat[1,4]$ the three
slowest eigenvalues are $-4.29$, $-1.60\times10^{8}$ and
$-3.44\times10^{8}\,\si{\per\second}$, a gap of $3.7\times10^{7}$. An observable
that depends on the excited populations only through their ratio cannot decay
faster than $\tauslow$: once the manifold is quasi-steady the excited vector is
a function of $u$ alone, so
$\ln\Robs - \ln\Robs^{\mathrm{CRE}} = (S/u^{*})\,\delta u + O(\delta u^{2})$
with $\delta u \propto \mathrm{e}^{-t/\tauslow}$, and the leading term
vanishes only where $S = \ffrac{3}-\ffrac{4} = 0$. It does not: $S$ is positive
at all $400$ grid points, with minimum $0.046$ at $[44,7]$, and the slow
eigenvector's projection onto $\ln\Robs$ equals $S/u^{*}$ to
$1.2\times10^{-4}$ at the benchmark and $1.6\times10^{-8}$ at $[0,4]$, the
deviation being the first-order quasi-steady correction
$\lambda_{\mathrm{slow}}/\lambda_{2}$ itself. It exceeds \SI{1}{\percent} only
at the two hot, dense points $[48,7]$ and $[49,7]$, where the separation is
$87$-fold and which the window criterion excludes from the analysed set
(\texttt{validation/slow\_mode\_projection/}). That decay on an intermediate
mode was the failure mode being hunted, and it does occur for the
ground-normalised error measure, which does not cancel the reservoir and whose
$1/\mathrm{e}$ time is \SI{2.58}{\second} against
$\tauslow = \SI{67.2}{\second}$ at the cold corner $[0,0]$. The estimate is
accurate for the ratio and not for that measure; why it is nevertheless not a
bound, the denominator relaxing on the same $\tauslow$, is set out in
Section~\ref{sec:persistence}.

%-----------------------------------------------------------------------------
\section{External comparison: Fujimoto's population coefficients}
\label{sec:fujimoto}
% QUESTION: does an independently computed collisional-radiative model, built
% decades earlier from different atomic data, agree with this one?

Everything to this point compares the model against itself or against
arithmetic. Only comparison against an independent calculation reaches the
atomic physics, and this project has one such comparison in depth: against
Fujimoto's collisional--radiative treatment of atomic
hydrogen~\cite{Fujimoto2004}, computed decades earlier, with different atomic
data, by a different method.

\subsection{Timescales, at a point that lies exactly on this grid}

Fujimoto works a numerical example at $\Te = \SI{10}{\electronvolt}$ and
$\nel = \SI{e18}{\per\cubic\metre}$, which is
$\SI{e12}{\per\cubic\centi\metre}$ and lands exactly on grid point $[49,0]$. No
interpolation is involved.

\begin{center}
\begin{tabular}{lrrr}
  \toprule
  quantity & Fujimoto App.~4B & this model, $[49,0]$ & ratio \\
  \midrule
  excited-state response & ${\sim}\SI{1e-7}{\second}$ & \SI{3.4057e-8}{\second} & $0.34$ \\
  ground-state depletion & ${\sim}\SI{1e-4}{\second}$ & \SI{1.7657e-4}{\second} & $1.77$ \\
  \bottomrule
\end{tabular}
\end{center}

The two-clock structure appears in both, separated by three orders of
magnitude in both, with the fast time short and the slow time long in the same
sense.

\textbf{What this comparison is not.} Fujimoto's response time is
$\max_p t_{\rm tr}(p)$, the \SI{63}{\percent} rise time of the slowest-responding
level, which is not $1/|\lambda_1|$. The two are different functionals of the
same spectrum, so agreement to a factor of a few is the meaningful comparison
and exact agreement would be suspicious. Quoting these ratios as though they
were a precision test would overstate them.

One qualification on the independence of this check, found on reading the
printed appendix rather than a summary of it. Fujimoto's Figs.~4B.1 and 4B.2 are
captioned as quoted from Sawada and Fujimoto (1994)
\cite{SawadaFujimoto1994}, so the appendix and the paper
Chapter~\ref{ch:conclusions} discusses are the same calculation, not two.
\textbf{The comparison is therefore independent of this thesis but not two
independent checks}, and it is counted once. The appendix also states the same
construction the paper does, that $n(1)$ is held constant, which matters for
Chapter~\ref{ch:conclusions} and is taken up there.

A second, independent route corroborates the same physics. Fujimoto's boundary
level, the principal quantum number above which collisions dominate radiative
decay, sits at $p_G \approx 4$--$5$ at this density. Tracing the boundary
descent in this model puts it at $5g$ at the low-density end of the grid, which
is the crossing Chapter~\ref{ch:atoms} deferred to this chapter
(Section~\ref{sec:radiative}). Two independently computed models place the same
structural feature within one shell of each other, at one density.

\subsection{The recombining coefficients agree}

Fujimoto's Table~4.1 tabulates the two population coefficients of the
decomposition that Section~\ref{sec:R_depends_on_b} derives: $r_0(p)$, how
strongly level $p$ is fed from the ion continuum, and $r_1(p)$, how strongly it
is fed from the ground state. Both are dimensionless, and both are normalised to
the Saha--Boltzmann population of level $p$, which is why $r_0$ is of order unity
while $r_1$ at low density is of order $10^{-5}$, and why a value of $r_0$ above
one is unphysical.

The recombining coefficients agree to \SIrange{0.1}{1}{\percent} for $n \ge 4$.
The discriminating case is $r_0(2)$, where cascade and $\ell$-mixing structure
matter most and a hydrogenic approximation would go wrong: this model gives
$0.7392$ against the tabulated $0.730$, a difference of \SI{1.3}{\percent}, a
test the model could have failed and did not (the broken variant of
Section~\ref{sec:severe_checks} drives the same quantity to $104.5$).

\subsection{The ionising coefficients at low $p$ disagree, and the target is
not verified}

At $\nel = \SI{e12}{\per\cubic\centi\metre}$ this model's $r_1$ sits a factor
$8.3$ low at $p=3$ and $4.4$ low at $p=4$ against Fujimoto's Table~4.1(b).
Those are precisely the two shells whose difference is the mechanism of
Chapter~\ref{ch:results}, so the discrepancy cannot be dismissed as peripheral.

Reading it as a failure of the model is not supportable, for six reasons, and
the status is an open comparison whose \emph{target} is unverified.

\begin{enumerate}
\item \textbf{The model's $r_1$ is reproduced by hand.} In the coronal limit
$r_1(2)$ follows from the single excitation rate coefficient
$K_{\rm exc}(1s\to n{=}2) = \SI{9.2796e-9}{\cubic\centi\metre\per\second}$ and
gives $1.37\times10^{-5}$, against the full $43$-state solve's
$1.66\times10^{-5}$. The two agree to \SI{21}{\percent}, which is the accuracy a
coronal estimate that ignores cascades should have. Whatever produces a factor
of $8$ is not inside the solve.

\item \textbf{The atomic stack passes three independent external checks.} The
radiative recombination coefficient into the ground state is within
\SI{2}{\percent} of the scaled standard value; the sum over excited states,
$\SI{2.08e-13}{\cubic\centi\metre\per\second}$, sits \SI{10}{\percent} below the
standard case-B coefficient $\alpha_B \approx
\SI{2.31e-13}{\cubic\centi\metre\per\second}$, the accepted total for capture
into every excited level, and that \SI{10}{\percent} shortfall is exactly the
$n>15$ truncation of Section~\ref{sec:convergence}; and the ground-state ionisation
coefficient at \SI{10}{\electronvolt}, $\SI{4.864e-9}{\cubic\centi\metre\per\second}$,
sits \SI{10.4}{\percent} below ADAS SCD96 read in its low-density limit,
$\SI{5.43e-9}{\cubic\centi\metre\per\second}$ at
$\nel = \SI{5e7}{\per\cubic\centi\metre}$ \cite{Summers2006, ADAS}, where the
effective coefficient reduces to ionisation out of the ground state because
every excited level decays before it can be ionised. The comparison is at the
low-density limit deliberately: SCD96 varies by \SI{0.9}{\percent} over a
factor four in density there, so the limit is reached. \textbf{This
\SI{10}{\percent} is not the truncation of the item above.} Direct ionisation
from $1s$ is a single cross section and does not know where the level ladder
stops; the difference is between this model's ionisation cross-section source
and the one behind SCD96, and it is reported rather than reconciled.

\item \textbf{The tabulated asymptote demands an excitation coefficient fifteen
times the accepted value.}
Reproducing Fujimoto's quoted coronal asymptote requires
$K_{\rm exc}(1s\to n{=}2) = \SI{1.44e-7}{\cubic\centi\metre\per\second}$ at
\SI{11.03}{\electronvolt}, fifteen times the accepted value. No error in this
model's implementation produces that number; only a different input does.

\item \textbf{Three candidate explanations are eliminated quantitatively.}
Truncation moves $r_1(3)$ by \SI{1.2}{\percent} between $n_{\max}=10$ and
$n_{\max}=15$, which is a factor $8$ short. Lyman trapping at an escape factor
of $0.1$ brings $r_1(2)$ to $0.92$ of the table but leaves $r_1(3)$ at $0.26$ of
it and drives $r_0(2)$ to $8.2$, a population above Saha equilibrium and
therefore unphysical. The third candidate, the $\ell$~closure, is eliminated in
the item below.

\item \textbf{The $\ell$~closure is a real difference, and it does not account
for the gap.} Table~4.1(b) is indexed by $p = 2, 3, 4, 5, 7, 10, 15$, principal
quantum numbers, and carries no $\ell$ label: the published coefficients are
shell quantities computed with the $\ell$~substructure bundled, while this model
resolves $\ell$ below $n=9$. That difference is real, and it is not the
explanation. Imposing the source's own
closure directly, by bundling this model's operator under statistical
$\ell$~weights as $\Lmat_{\rm bundle} = \mathsf{P}\,\Lmat\,\mathsf{R}$ with
$\mathsf{P}_{ki} = 1$ for $i$ in shell $k$ and $\mathsf{R}_{ik} = g_i/g_k$,
moves $r_1(3)$ by \SI{0.42}{\percent}: the factor $8.29$ becomes $8.26$
(\texttt{verify\_fujimoto\_bundle.py}, stamped in
\texttt{validation/fujimoto\_bundle/}). At
$p=2$ the agreement becomes \emph{worse}, from a factor $10.8$ low to a factor
$11.8$ low, and $r_0(2)$ moves from $0.7392$, which agrees with the tabulated
$0.730$, to $0.6459$, which does not. The reason is structural: the bundling
projector annihilates the $\ell$-mixing operator identically, $\max|\mathsf{P}
\mathsf{B} \mathsf{R}| = 8.5\times10^{-7}$ against a matrix scale of
$1.0\times10^{11}$, so bundling cannot express what $\ell$-mixing does; and the
production model is already at the statistical-$\ell$ limit for $p \ge 3$, its
sublevel populations sitting within \SI{5}{\percent} of $g_{n\ell}/g_n$. The
$\ell$-resolved calculation was therefore already the like-for-like comparison.
Removing proton $\ell$-mixing altogether moves $r_1(2)$ by a factor $118$, but
that is the opposite limit from a statistical closure rather than a proxy for
it, and the two do not bracket the tabulated value.

\item \textbf{The transcription and the row labels are correct.} Both were
checked against the printed table. The $\log_{10}\nel = 18$ row reads
$0.730$, $0.835$, $0.947$ and $0.983$ for $r_0$ at $p = 2, 3, 4, 5$, and
$1.79\times10^{-4}$, $5.72\times10^{-5}$, $1.66\times10^{-5}$ and
$5.08\times10^{-6}$ for $r_1$, matching the values used here exactly. The row
labels are $\log_{10}(\nel/\si{\per\cubic\metre})$, which the table establishes
internally: $r_0(2)$ reaches $0.981$ at $\log_{10}\nel = 21$, and Griem's
criterion places local thermodynamic equilibrium for $n=2$ at
\SI{1.74e16}{\per\cubic\centi\metre}, which is $10^{21}\,\si{\per\cubic\metre}$
to within the criterion's own precision. Read as \si{\per\cubic\centi\metre} the
onset would sit five orders above Griem, which is impossible; and the rows run
from $\log_{10}\nel = 12$ to $24$, a span of $10^{6}$ to
$10^{18}\,\si{\per\cubic\centi\metre}$ read as \si{\per\cubic\metre} and no
plasma at all at the top read as \si{\per\cubic\centi\metre}. A one-decade
offset in the density rows is not a tenable reading either.
\end{enumerate}

\textbf{Status.} The $r_1$ deficit at low $p$ is an open external
disagreement, and this thesis does not explain it; the three candidate causes
above each fail to account for a factor of $8$. The $\ell$-closure test carries
a severity check: bundling an operator whose $\ell$~distribution is deliberately
non-statistical moves $r_1(2)$ by a factor $129$, so the null result at $p=3$ is
a measurement and not a blind instrument. What the section does \emph{not}
establish is that the disagreement leaves the quantity
Chapter~\ref{ch:results} uses untouched.
Equation~\eqref{eq:fujimoto_correspondence} makes
$a_p = r_1(p)\,\Zsb{p}/\Zsb{1}$: $r_1$ is the ground-fed coefficient from which
$\ffrac{p}$ is built, so a deficit in $r_1(3)$ and $r_1(4)$ is a deficit in
$a_3$ and $a_4$ themselves. Bundling is not the protection, since neither shell
passes through a bundled coefficient; the only protection is that $\Delta$ and
the cap depend on the ratio $a_3/a_4$, in which a deficit common to both shells
cancels. At \SI{e12}{\per\cubic\centi\metre} the deficit is not common, $8.3$
at $p=3$ against $4.4$ at $p=4$; at \SI{e15}{\per\cubic\centi\metre} it nearly
is, $1.8$ against $1.9$. Dividing the model's $a_3$ and $a_4$ by those ratios,
so that they become what Fujimoto's coefficients would give, and recomputing
at every grid point (\texttt{verify\_fujimoto\_r1\_consequence.py},
\texttt{validation/fujimoto\_r1\_consequence/}) moves $\Delta$ by $+0.64$
under the first pair and by $-0.08$ under the second. Over the $280$ points
above \SI{2}{\electronvolt} the cap then moves by \SI{+29}{\percent} at the
median under the first and \SI{-4}{\percent} under the second; the crest
$u^{\mathrm{peak}}$ falls by a factor $6.0$ and $1.9$; and the sensitivity at
the operating point, $\ffrac{3}-\ffrac{4}$ evaluated at $u^{\mathrm{CRE}}$,
moves by \SI{+30}{\percent} and \SI{+11}{\percent} at the median, with ranges
across the points of $-58$ to \SI{+524}{\percent} and $-38$ to
\SI{+68}{\percent}. Both exceed the \SIrange{8}{9}{\percent} that the R-matrix
excitation substitution of Chapter~\ref{ch:atoms} produces. The disagreement
is therefore direct evidence about the coefficients Chapter~\ref{ch:results}
uses; what is open is whether the deficit lies in this model or in the
tabulation, and the scenario applies a deficit measured at
\SI{10}{\electronvolt} at every temperature, which is its stated assumption.
An unexplained disagreement in an absolute population coefficient is reported
here as unexplained, with its stake stated.

\textbf{What the benchmark does establish.} The recombination-fed channel
agrees, and because $r_0$ and $r_1$ carry the same $p$-dependent factor $Z(p)$,
a normalisation error large enough to explain the $r_1$ deficit is excluded by
$r_0$ agreeing. That is a pass, and it should be read as one.

\textbf{What it does not establish, and two limits on it.} It does not fix the
absolute normalisation of the ground-fed channel against an independent
calculation, because no $\ell$-resolved tabulation was available to compare
against. Table~4.1(b) is computed at $T_\mathrm{e} = \SI{1.28e5}{\kelvin}$,
which is \SI{11.03}{\electronvolt}, while this grid reaches
\SI{10}{\electronvolt}, so the comparison sits at the grid edge, with the
tabulated temperature \SI{10.3}{\percent} above it. The companion
Table~4.1(a) is at \SI{e3}{\kelvin} =
\SI{0.0862}{\electronvolt}, two orders below this grid, so there is no second
usable case. The external check is one temperature at one edge.

%-----------------------------------------------------------------------------
\section{Gate D fails, and the failure is in the gate}
\label{sec:gate_d}
% QUESTION: the one gate that compares against an external atomic database
% disagrees at every point. What does that mean, and what does it not mean?

The Atomic Data and Analysis Structure (ADAS) distributes effective
stage-to-stage rate coefficients for hydrogen: SCD, the effective ionisation
coefficient, and ACD, the effective recombination coefficient~\cite{Summers2006}.
These are the coefficients a transport code uses when it does not resolve
excited states, and they are the most widely used external reference this model
could be compared against.

Gate~D forms the model's counterpart by taking the steady-state population
$n_p^{\rm ss}$ of each level at each grid point, weighting it by that level's
ionisation rate coefficient $Q_p$ of Eq.~\eqref{eq:level_balance}, and dividing
by the total neutral population:
\begin{equation}
  {\rm SCD}_{\rm model}
  = \frac{\sum_p Q_p\,n_p^{\rm ss}}{\sum_p n_p^{\rm ss}},
  \label{eq:scd_model}
\end{equation}
then compares against ADAS interpolated to the same temperature within a
factor-of-two density band. The gate passes a point if the ratio
$\eta = {\rm SCD}_{\rm model}/{\rm SCD}_{\rm ADAS}$ lies between $0.5$ and $2.0$.

\begin{center}
\texttt{Gate D: FAIL  (0\% of points)} \hfill \texttt{validation/gate\_summary.txt}
\end{center}

Measured, $\eta$ ranges from $5.73$ to $8346$ over the full grid, and no point
agrees within a factor of two. The extremes are not typical: over most of the
grid the model sits high by a factor between $9$ and $65$, and that systematic
offset peaks near $\Te \approx \SI{2}{\electronvolt}$ (these are the gate's
own numbers, under its linear read of the table; the read is discussed
below). The disagreement has a
shape, which is what makes it diagnosable rather than random.

\subsection{The diagnosis, and the test that settles it}
\label{sec:gate_d_diagnosis}

Two readings were open. One of them has now been tested and it is the
explanation.

The first is that Eq.~\eqref{eq:scd_model} is not the quantity ADAS tabulates.
ADAS defines SCD as the \emph{ionising} coefficient, the ionisation flux
attributable to the ground-state population. The sum in
Eq.~\eqref{eq:scd_model} runs over the full steady state, which includes the
recombination-fed Rydberg levels, the $\nzero = \bm{c}\,\nion$ channel of
Section~\ref{sec:R_depends_on_b}. Ionisation out of
those levels is a recombination-driven flux and belongs under ACD, not SCD.
Since the high levels have small ionisation potentials and large ionisation
coefficients, moving them to the wrong side of the ledger inflates the model's
SCD by a large factor, and inflation is what is measured. The test that would
settle it is to recompute Eq.~\eqref{eq:scd_model} over the ground-fed channel
$\none$ alone, and it has now been run
(\texttt{diagnose\_gate\_d.py}).

The two channels separate exactly, with no approximation:
$\nvec_E = \nzero + \none$ to a relative residual below $10^{-8}$ at every grid
point, which the script checks before using either. Rebuilding the coefficient
on the ground-fed channel alone,
\begin{equation}
  {\rm SCD}_{\rm model} = Q_{1s}
  + \sum_{p \neq 1s} Q_p\,\frac{n^{(1)}_p}{\ngs},
  \label{eq:scd_groundfed}
\end{equation}

\noindent and comparing against the same ADAS table with the same factor-of-two
tolerance moves the agreement from $0$ of $400$ points to $400$ of $400$. The
ratio $\eta$ falls from a range of $5.73$ to $8346$ to a range of $0.702$ to
$1.004$, with median $0.828$; above \SI{2}{\electronvolt} it runs from $0.721$
to $0.937$, median $0.831$, and at the benchmark it is $0.768$
(\texttt{validation/reservoir\_vs\_adas/}). The read matters. ADAS96
tabulates seven temperatures between $1$ and \SI{10}{\electronvolt}, and
\texttt{diagnose\_gate\_d.py}, which first made this comparison, interpolates
linearly in value between them (\texttt{adas\_at}, lines 114 to 124); across
$1$ to \SI{1.5}{\electronvolt}, where SCD rises by nearly two orders of
magnitude, that overstates the table by up to a factor $6.65$. That read gives
$323$ of $400$, a range of $0.130$ to $0.931$ and a median of $0.696$, with a deficit deepening toward low temperature; read logarithmically
in both axes, as the table's spacing requires, the tail disappears, the row
median being $0.72$ at \SI{1}{\electronvolt} and $0.90$ by
\SI{1.2}{\electronvolt}. The diagnosis script is left in place; every ADAS
number in this section comes from the logarithmic read.

The mechanism is visible directly. At collisional--radiative equilibrium the
recombination-fed channel carries between \SI{89.5}{\percent} and
\SI{99.99}{\percent} of $\sum_p Q_p n_p$, and its share rises monotonically with
density, from \SI{97.8}{\percent} at $\SI{e12}{\per\cubic\centi\metre}$ to
\SI{99.99}{\percent} at $\SI{e15}{\per\cubic\centi\metre}$ at
\SI{1}{\electronvolt}. Equation~\eqref{eq:scd_model} is therefore reporting the
recombination channel with a small ground-fed correction, which is why $\eta$
tracks $\nel$: three-body recombination carries an extra factor of $\nel$ and
the mislabelled flux carries it into the coefficient.

\paragraph{The recombination half, built and compared for the first time.}
The same decomposition supplies the coefficient the gate never computed. The net
flux into the neutral ground state carried by the recombination-fed channel,
per unit $\nel\nion$, is
\begin{equation}
  {\rm ACD}_{\rm model}
  = \frac{S_{1s}\,\nion + \sum_{p \neq 1s} L_{1s,p}\,n^{(0)}_p}{\nel\,\nion},
  \label{eq:acd_groundfed}
\end{equation}

\noindent direct radiative recombination into $1s$ plus everything that cascades
or is collisionally transferred down into it. Against ACD96 this agrees at
\emph{all} $400$ points, with $\eta$ between $0.894$ and $1.009$ and median
$0.982$; at the benchmark the model gives
$\SI{2.788e-13}{\cubic\centi\metre\per\second}$ against ADAS's
$\SI{2.883e-13}{}$, a ratio of $0.967$.

\textbf{That is the strongest external check in this thesis.} It is an absolute
coefficient, compared against an independent database built from independent
atomic data, agreeing to better than \SI{10}{\percent} at the benchmark and
within a factor of two everywhere. It was available all along and was not taken,
because the half of the gate that would have taken it was never written.

The second reading was that the two quantities are not defined on the same
system at all. ADAS coefficients describe a closed atom-plus-ion system; this
model holds the ion as a fixed external reservoir with no forty-fourth state
(Section~\ref{sec:two_references}). A stage-to-stage coefficient extracted from
an open system need not equal one extracted from a closed one, and
Section~\ref{sec:attacks} has already shown that the boundary condition moves
$\tauslow$ by up to a factor of $41$. \textbf{The ACD result argues against this
being the dominant effect}: ACD is extracted from the same open system through
the same fixed-ion reservoir and agrees with ADAS to within \SI{3.3}{\percent}
at the benchmark, \SI{1.8}{\percent} at the grid median and \SI{11}{\percent}
at every point; an open-system artifact large enough to produce $\eta$ of
$8346$ in SCD would not leave ACD intact.

\paragraph{The residual SCD deficit, stated rather than dismissed.} Rebuilding
SCD on the ground-fed channel does not make the agreement perfect, and the
remaining disagreement is systematic in a way worth reporting: the model's
ionisation sits below ADAS at all but one of the $400$ points, by up to
\SI{30}{\percent} and by \SI{17}{\percent} at the median, with the maximum
$\eta = 1.004$ at $[4,7]$. Under the linear read of the table the gap closes
steeply with temperature, from a median of $0.302$ over $1$ to
\SI{1.5}{\electronvolt} to $0.817$ over $5$ to \SI{10}{\electronvolt}, a shape
that reads as the signature of truncation; the steep part is the linear read,
and under the logarithmic read the deficit is between $0.72$ and $0.94$ at every
temperature above \SI{2}{\electronvolt}. A deficit of that size is still what
truncation would produce, since ADAS96 carries levels above $n_{\max} = 15$
and the levels nearest the continuum ionise most easily, but its temperature
shape no longer singles that cause out (the $n \ge 12$ flux share the diagnosis
script printed against the linear read is not carried forward). The test that
would settle the
cause is the $n_{\max}$ convergence study of Chapter~\ref{ch:limits} rather
than this one.

\paragraph{The reservoir itself, against the same table.} ADAS's two
coefficients fix the equilibrium reservoir: at ionisation balance
$n_{1s}\,\mathrm{SCD} = \nion\,\mathrm{ACD}$, so
$u^{\mathrm{CRE}} = \mathrm{ACD}/\mathrm{SCD}$, and the model obeys the same
identity with its own ground-fed SCD and ACD to $4\times10^{-13}$ at every grid
point. That makes the quantity Chapter~\ref{ch:results} is built on directly
comparable, and the comparison has two parts
(\texttt{verify\_reservoir\_vs\_adas.py},
\texttt{validation/reservoir\_vs\_adas/}). The \emph{level} is high:
$u^{\mathrm{CRE}}_{\mathrm{model}}/u^{\mathrm{CRE}}_{\mathrm{ADAS}}$ has median
$1.18$ over the $400$ points, runs from $0.89$ at $[4,7]$ to $1.40$ at $[0,1]$,
and is $1.26$ at the benchmark, structured about equally in $\Te$ and in
$\nel$. The model's ground reservoir per ion sits a fifth above ADAS's, which
moves the operating point beneath the cap of Section~\ref{sec:bound} without
touching the cap. The \emph{displacement}, which is what the plateau error is
built from, agrees. Over the $392$ one-interval temperature steps the model's
$\Delta\ln u$ has median magnitude $0.26$ and differs from ADAS's by a median
$0.014$; above \SI{2}{\electronvolt} the median difference is $0.011$ and the
worst $0.050$, at $[15,0]$; at the benchmark step the two are $-0.270$ and
$-0.258$. The reservoir gain $G$ this thesis tabulates is therefore within
about \SI{5}{\percent} of the one ADAS implies, at the median over the
defended range, and the refuter set before the run, a difference above $0.1$
at any interior step above \SI{2}{\electronvolt}, appeared at none of $264$.
Every number in this paragraph comes from the logarithmic read; on the linear
read of \texttt{diagnose\_gate\_d.py}, which overstates SCD by up to a factor
$6.65$ over $1$ to \SI{1.5}{\electronvolt}, the refuter would have appeared at
$9$ of the $264$ steps.

\subsection{Three further defects in the gate itself}

The gate is also incompletely implemented, and this must be said alongside the
failure rather than instead of it.

The recombination half does not exist \emph{in the gate}. The ACD table is
loaded once, at line 299 of \texttt{validate\_gates.py}, and the variable
holding it is never read again anywhere in the file; only the ionisation
comparison runs, despite a docstring that describes both. That half is computed
in \texttt{diagnose\_gate\_d.py} (Section~\ref{sec:gate_d_diagnosis}).

The interpolation is wrapped in a catch-all exception handler that sets the
comparison value to not-a-number and continues, so a failure to find the right
density band is indistinguishable from a point that was legitimately skipped.

And the gate's status is recorded three incompatible ways: \texttt{README.md}
line 119 calls it documented, the docstring at line 47 predicts a pass above
$\Te = \SI{5}{\electronvolt}$, and \texttt{gate\_summary.txt} reports failure
at \SI{0}{\percent} of points. The correct entry is the measured one: Gate~D
fails.

\subsection{What is claimed, and what is not}

\textbf{Gate~D itself still fails, and it has not been changed.}
\texttt{validate\_gates.py} is untouched: no tolerance was widened, no subset
was selected, and \texttt{gate\_summary.txt} still records
\SI{0}{\percent}. What has changed is that the failure now has a demonstrated
cause rather than two undemonstrated candidates. The gate compares a mixed
quantity against a table that keeps the two channels apart by design, which is
a defect in the comparison and not in the matrix.

Deciding what Gate~D should test belongs with whoever owns the gate. What is
claimed here is narrower: rebuilt on the two channels separately, the
coefficients agree with ADAS within a factor two at all $400$ points for SCD,
the model $0$ to \SI{30}{\percent} low with a median ratio of $0.83$, consistent
with the $n_{\max}$ truncation but not uniquely so, and at all $400$ for ACD
with a median ratio of $0.98$.

The failure does not cost the thesis its only comparison of absolute effective
coefficients against an external database; that comparison exists, hidden by
the definition. The failure does not touch the structural results
of Chapter~\ref{ch:theory}, which are statements about whatever $\Lmat$ is
supplied and would hold for the ADAS-consistent matrix equally. It does not
touch the level-resolved comparison of Section~\ref{sec:fujimoto}, which
compares the same decomposition against tabulated coefficients rather than
against a contraction of $42$ coefficients into one. And the contraction is the
reason the two comparisons carry different weight: many incorrect $\bm{a}$ would
contract to the correct SCD, so a passing Gate~D would have been weak evidence
even had it passed.

%-----------------------------------------------------------------------------
\section{Errors found by our own audit}
\label{sec:errors_found}
% QUESTION: what did this project get wrong, how was each error found, and how
% large was it?

Section~\ref{sec:atoms_corrections} reports two implementation errors that this
project found in itself, with their sizes, and observed that both were invisible
to every consistency test the pipeline ran. This section states what each
changes here, adds two more, and keeps the four together rather than in an
appendix: a model whose authors found and bounded four of its own errors is
better characterised than one reporting none, and in three of the four cases
the bound is itself a result.

\subsection{The $\ell$-mixing Coulomb logarithm}

Proton-impact $\ell$-mixing redistributes population between sublevels of the
same $n$ and, as Section~\ref{sec:lmixing} showed, is the dominant loss channel
for the metastable $2s$ level. The rate coefficient follows Badnell's
formulation~\cite{Badnell2021}, whose Eq.~9 carries a factor
\begin{equation}
  F(U_m) = \frac{\sqrt{\pi}}{2}U_m^{-3/2}\operatorname{erf}(\sqrt{U_m})
           - \frac{\mathrm{e}^{-U_m}}{U_m} + E_1(U_m),
  \label{eq:FUm}
\end{equation}
in which $\operatorname{erf}$ is the error function, $E_1$ the exponential
integral, and $U_m$ the dimensionless ratio of the smallest energy transfer the
Debye cutoff permits to $k_B\Te$; $F(U_m)$ is the Coulomb logarithm of the
interaction, which diverges as $U_m \to 0$.

The implementation set $F = 1$, justified in its docstring as the low-density
limit $U_m \to 0$; the limit was inverted, and
Section~\ref{sec:atoms_corrections} gives the finding, the factor of three to
seven by which the rates were low, and the measured effect on the spectrum.
Evaluated with the local Debye length at the benchmark point $[23,5]$,
$F = 6.62$ for $2p$--$2s$, $5.71$ for $3d$--$3p$, $4.26$ for $5g$--$5f$ and
$2.89$ for $8k$--$8i$ (\texttt{validation/lmix\_cutoffs/}); with the
production table's frozen density of \SI{e14}{\per\cubic\centi\metre} the
$2p$--$2s$ value is $6.95$ (Section~\ref{sec:atoms_corrections}). A derivation
note records $6.65$, $5.73$, $4.29$ and $2.92$ for the same four pairs.

The insensitivity of the slow eigenvalues to that error, which
Section~\ref{sec:atoms_corrections} measures, is specific to a process already
far faster than $\taurel$; it is not a general licence to be careless about
fast rates, and the $\ell$-mixing scan below is what establishes the range over
which it holds.

\subsection{And $\ell$-mixing is saturated}

The same insensitivity was then tested directly rather than argued. Scaling
every $\ell$-mixing rate in the matrix over a range of $\times0.1$ to $\times10$
moves $\ffrac{3}-\ffrac{4}$ by less than \SI{0.3}{\percent} at the benchmark
point and at the ridge, and by less than \SI{5}{\percent} at the cold corner.
Only switching $\ell$-mixing off entirely changes anything. Rebuilding the
matrix with a density-consistent Debye cutoff in place of a frozen one moves
$\ffrac{3}-\ffrac{4}$ by at most \SI{0.42}{\percent}.

The claim that the frozen Debye density varies the cutoff factor by about
\SI{30}{\percent} across the grid is wrong: the module's own functions give
$\times3.7$ for $n=2$ and at least $\times50$ for $n=8$
(Section~\ref{sec:lmixing}); the conclusion survives because the observable does
not depend on the input over that range.

\subsection{The eigenvalue magnitude filter}

The spectrum of $\Lmat$ was extracted by a routine containing the line
\texttt{eigs = eigs[eigs < -1.0]}, which discards every eigenvalue smaller in
magnitude than $1\,\si{\per\second}$. Where the slow eigenvalue is smaller than
that, at the cold end of the grid, the filter deleted it and promoted the whole
ladder by one step, so $\tauslow$, $\taurel$ and $\Msep$ were wrong together
while all three remained finite, negative and correctly ordered. It was found by
recomputing the spectrum unconditionally and comparing bit-for-bit against the
stored arrays, the three-implementation agreement of Section~\ref{sec:gate_e}.
Section~\ref{sec:atoms_corrections} gives the extent of the damage, the ladder
shift that confirmed the mechanism, the preserved filtered outputs and the
removal of the literal from two dormant files. What belongs here is the shape
of the wrong answer: at the cold corner the filter turned
$\Msep = 1.729\times10^{9}$ into $\Msep = 1.467$, and the grid minimum from
$86.77$ into $1.341$, values that read as a physics finding, since
near-breakdown of timescale separation in the coldest, thinnest plasma is
exactly what a reader would expect to see. The later correction of the two
dormant files changes no number in this thesis: at the benchmark the two filters
return the same $\tauslow$ to every digit, and they differ only where the slow
mode is slower than $1\,\si{\per\second}$.

\subsection{The central quantity was misnamed}

The largest of the four errors changed no number and inverted a sentence.

The quantity that had been reported as the quasi-steady-state error was
defined as the difference between the tabulated ratio and the true
one, with the tabulated ratio evaluated from a two-argument function of $\Te$
and $\nel$. Section~\ref{sec:two_references} shows that the excited-state ratio
is a function of three arguments, not two, the third being the state of the
ground-state reservoir. Fixing the first two therefore silently imposes
equilibrium ionisation balance, so what was measured was the distance to the
equilibrium answer, not the error of the quasi-steady-state closure.

Integrating the full system and evaluating both errors along the same trajectory
separates them:

\begin{center}
\begin{tabular}{lrr}
  \toprule
  point & error against equilibrium & error of the QSS closure \\
  \midrule
  benchmark $[23,5]$ & $6.34\times10^{-2}$ & $8.66\times10^{-6}$ \\
  worst case $[0,4]$ & $3.869\times10^{-1}$ & $6.73\times10^{-9}$ \\
  \bottomrule
\end{tabular}
\end{center}

They differ by four to eight orders of magnitude, most at the point carrying the
largest error. The physics was intact and every number reproduced; what was
wrong was the name and the sentence built on it, and Chapter~\ref{ch:results}
states the consequence, that the approximation everyone worries about is not
the one that fails. The error was found by an adversarial reading of the thesis text
itself, not of the code, and it is why Section~\ref{sec:eps_definitions}
defines four error measures separately.

\subsection{A statement in Chapter~\ref{ch:theory} that measurement contradicted}
\label{sec:promise_broken}

The statement, in Section~\ref{sec:source_quadratic}, was that at
$\nel = \SI{e12}{\per\cubic\centi\metre}$ three-body recombination supplies
under \SI{10}{\percent} of the feed into any level, and it was used to propose
testing the radiative and three-body datasets separately at the two ends of the
density grid. Measured against the matrix, $30$ of $36$ levels exceed
\SI{10}{\percent} three-body feed at both \SI{1.00}{\electronvolt} and
\SI{2.95}{\electronvolt}, and $26$ of $36$ at \SI{10}{\electronvolt}. Into the
bundled shells it is essentially the whole feed.

The density grid therefore does not separate the two recombination datasets, and
\textbf{no such separate test was run or can be}; Section~\ref{sec:source_quadratic}
now states the measured fractions and withdraws the promise. What remains is the
gap: the radiative and three-body datasets enter this model only in combination,
nothing in this thesis tests them apart, and a test that could would need a
density below the bottom of this grid, outside the conditions the thesis is
about.

%-----------------------------------------------------------------------------
\section{Reproduction, and what a document cannot certify}
\label{sec:reproducibility}
% QUESTION: can the numbers in this thesis be regenerated, and by whom?

\subsection{Regeneration from the canonical matrix}

Every quantitative result in Chapters~\ref{ch:theory} and~\ref{ch:results} was
regenerated from $\Lmat$ and compared against the stored output. The timescale
and step-error grids reproduce bit-identically, including the values that come
from numerical integration and not from linear algebra. The plateau and persistence
tables reproduce byte for byte, $680$ of $680$ rows, with $784$ superposition
residuals at exactly zero difference. The nineteen points corrected by the
eigenvalue filter reproduce exactly, including the ladder-shift identity.

This is reproduction in the narrow sense: the same code, run again, on the same
inputs, gives the same answers. It excludes a stale array and it excludes a
non-deterministic solver. It says nothing about whether the algebra is right.

\subsection{Independent re-derivation}

The stronger check is re-derivation. All six central results of
Chapter~\ref{ch:theory} were derived from the linear collisional--radiative
system alone by a party working from the physical statement of the problem,
implemented from scratch, with the predicted values written down before any
project file was opened.

In the table below, $\bdep$ is the ground-state departure coefficient of
Eq.~\eqref{eq:departure}, the ratio of the actual ground-state population to its
Saha--Boltzmann value, and $\bdep^{\rm CRE}$ is the value it takes in
collisional--radiative equilibrium (CRE), the steady state the lookup tables
assume.

\begin{center}
\begin{tabular}{lrr}
  \toprule
  quantity & \texttt{chapter3.tex} & independent re-derivation \\
  \midrule
  $\bdep^{\rm CRE}$ at the benchmark        & $956$                        & $956.509$ \\
  switching points $c_m/a_m$                & $2.73\times10^{3}$, $1.91\times10^{4}$ & $2727.3$, $19085.5$ \\
  peak $\bdep$                              & $7.2\times10^{3}$            & $7214.7$ \\
  $\ffrac{3}$, $\ffrac{4}$, difference       & $0.260$, $0.048$, $0.212$    & $0.259652$, $0.047725$, $0.211927$ \\
  maximum of $\ffrac{3}-\ffrac{4}$          & $0.4514$                     & $0.451357$ \\
  $\Robs$ at the benchmark                  & $0.7778$                     & $0.777768$ \\
  \bottomrule
\end{tabular}
\end{center}

The timescales reproduce to four digits, $\tauslow = \SI{22.73}{\micro\second}$,
$\taurel = \SI{2.277}{\nano\second}$, $\Msep = 9981.9$.

Two further invariances were measured in the same pass and belong with the
result they protect. The shell separation $\Delta$ that sets the bound of
Section~\ref{sec:bound} is invariant under rescaling the ground-to-ion
normalisation, verified exactly under factors of $10^{10}$ and $10^{-7}$, while
the individual channel fractions are not; and $\Delta$ varies by
\SI{0.07}{\percent} across five different $\ell$-weightings of the shell sum, so
the bound is not an artifact of how sublevels were weighted.

\subsection{A prediction that was falsified, and improved the claim}

The same re-derivation made one prediction that turned out to be wrong, and it
is more informative than the six that were right.

The prediction was that $\ffrac{3}-\ffrac{4}$ must vanish at high density, on
the reasoning that in local thermodynamic equilibrium both supply channels
produce the same Boltzmann distribution, so the two shells would draw
identically and the difference would close. It does not happen. At
$\Te = \benchTe$, $\Delta$ doubles over the lowest part of the grid and then
stops moving: $0.979$ at $\SI{e12}{\per\cubic\centi\metre}$, $1.946$ at
$\SI{1.4e14}{\per\cubic\centi\metre}$, $1.939$ at
$\SI{e15}{\per\cubic\centi\metre}$
(\texttt{validation/fujimoto\_r1\_consequence/}). Above about
$\SI{7e12}{\per\cubic\centi\metre}$ it is flat, and the bound $\tanh(\Delta/4)$
built on it spans only $0.394$ to $0.450$, a \SI{14}{\percent} range, across the
whole of the rest of the grid.

The reason is the open boundary again, this time in the model's favour: the
recombination source is imposed externally rather than tied to the ion density
by a Saha relation, so the two channels cannot merge by construction, and
Gate~C's \SI{1.5}{\percent} of Saha says the grid never approaches the limit in
which they would. The consequence for the density maximum of the diagnostic
error, a position effect rather than a span effect, is
Table~\ref{tab:position_effect} in Section~\ref{sec:density_mechanism}.

\subsection{One item demoted: the May/August cross-validation}

A cross-validation between two runs four months apart, on different matrices,
with different observables and different numerical methods, was recorded in this
project's register as verified and described as the strongest single check it
had. The agreement is close: $\tauslow$ to \SI{0.064}{\percent}, $\taurel$ to
\SI{0.200}{\percent}, $\Msep$ to \SI{0.264}{\percent}, $\epsstep$ to
\SI{0.014}{\percent}. Predicting how long the error stays above
\SI{10}{\percent} from the plateau-then-decay picture gives
\SI{12.478}{\micro\second} against the earlier run's full integration,
\SI{12.4794}{\micro\second}.

Two corrections apply, and both weaken it.

The apparent agreement to eight parts in $10^{5}$ on the duration is
cancellation. The inputs agree only to \SIrange{0.06}{0.6}{\percent}, and
carrying the earlier run's own internal numbers through gives
\SI{12.406}{\micro\second}, off by \SI{0.59}{\percent}. \textbf{The figure is
agreement at the \SI{0.6}{\percent} level}, which is still strong.

More seriously, \textbf{the earlier run has never been reproduced against the
canonical matrix}. This project's own derivation record says of these same
numbers that they are a reported result, not independently re-run, and should be
treated as strong documented evidence rather than as self-verified. A
cross-validation whose earlier half cannot be re-executed is a citation, not a
check. It is reported here as documented evidence and is not counted on the
ladder of Section~\ref{sec:validation_summary}. Under the ground-truth ordering
this thesis works to, physics over mathematics over code over documents,
a document is the weakest of the four, and it is where this particular claim
currently rests.

\subsection{Two provenance defects that are not repaired}

Two arrays of timescales, $\Msep$, $\tauslow$ and $\taurel$, are written to the
same three file paths by two different scripts. Whichever ran last wins, and six
downstream readers cannot tell whose numbers they are holding. This is exactly
the condition under which the eigenvalue filter propagated, since the filtered
and unfiltered versions of the same array have the same name.

Two things bound what that costs here, and both were checked rather than
assumed. The two writers are not two calculations: each sorts the eigenvalues of
the same canonical $\Lmat$ and keeps the negative ones, neither applies the
filter that caused the earlier incident, and the arrays on disk agree with the
independently produced \texttt{verify\_timescales.py} artifact to $2\times
10^{-16}$. And none of the six readers feeds this document: every figure and
table they write is absent from it. The benchmark $\Msep = 9982$ quoted in
Chapter~\ref{ch:theory} and in the abstract is not read from these arrays at
all; it is $\tauslow/\taurel$ computed from $\Lmat$ by
\texttt{verify\_timescales.py} and recorded in
\texttt{validation/spectrum\_testpoint.csv}. The defect is therefore a latent
hazard in the repository rather than an ambiguity in any number printed here,
and it is left in place and named rather than quietly repaired.

An audit script exists to detect single-writer violations, and it reports an
unfixed failure of this kind. It is also incomplete. Its test for a write reads the file mode from the second
argument or a keyword, so the idiom \texttt{path.open("w")} passes through it
unrecorded. At the tip of \texttt{main} on 23~August that idiom had four users;
in the tree this thesis is built
from it has $45$ call sites in $29$ scripts, $42$ of them writing
comma-separated output, among them the producer of the plateau map that
fifteen places in Chapter~\ref{ch:results} read from
(\texttt{verify\_writer\_census.py}, \texttt{validation/writer\_census/}). The
audit should not be cited as evidence that the single-writer property holds.
The census confirms that the property fails for exactly the three timescale
arrays, each written by both \texttt{qss\_analysis.py} and
\texttt{validate\_gates.py}, and holds for the plateau map, which has one
writer.

%-----------------------------------------------------------------------------
\section{What this chapter establishes}
\label{sec:validation_summary}
% QUESTION: after grading every check, what is the model licensed to be used
% for, and where?

Table~\ref{tab:ladder} collects the checks in the order a reader should weight
them, with the grade of Section~\ref{sec:gate_philosophy} attached to each.

\begin{table}[tbp]
\centering
\caption[The validation ladder, graded by severity and provenance]{Every check
this chapter grades, with what it tests, its severity, whether its result can be
regenerated, and the measured outcome. \emph{Tautological} means the check
cannot fail except on an arithmetic bug, because the quantity tested was
constructed from the relation being tested; \emph{weak} means it can fail but
only for an error far larger than any plausible one; \emph{severe} means a
plausible wrong version of this model fails it. The provenance column is the
more uncomfortable of the two: \emph{stamped} means a script in the repository
regenerates the number into an artifact under \texttt{validation/};
\emph{notes} means the number is recorded in this project's working files with
no producing script, so it cannot be re-executed and cannot fail anything today.
None of the checks graded severe are in that condition, and the one
\emph{notes} entry left is the $\tanh$ bound, which is graded tautological. The
counts stated in this caption are checked against the table's own rows by
\texttt{verify\_ladder\_counts.py}, which exits non-zero if they drift apart. Of the eighteen
internal checks, twelve are severe, four are tautological and two are weak; of the five external comparisons, two
pass, one passes only at the higher densities, one fails, and one is open.}
\label{tab:ladder}
\footnotesize
\begin{tabular}{p{3.9cm}p{1.8cm}p{1.3cm}p{5.6cm}}
\toprule
check & grade & record & measured \\
\midrule
Atomic-data provenance (Ch.~\ref{ch:atoms})  & severe & stamped & every coefficient traces to a named dated source; tallies exact \\
Excitation rates against RMPS (Ch.~\ref{ch:atoms}) & severe & stamped & \SI{81.4}{\percent} within \SI{20}{\percent} for $n_{\rm upper}\le4$, mean \SI{13.0}{\percent}; $n=5$ disagrees by factors of $2$--$5$, localised to the top shell of the RMPS basis \\
Raw cross sections, microscopic reversibility & severe & figure & ratio $0.9995 \pm 0.0032$, \SI{98.7}{\percent} within \SI{1}{\percent} \\
One energy ladder on both sides & severe & stamped & Saha ratio $0.999997$ for all $43$ states, spread $1.4\times10^{-15}$ \\
Superposition of the two channels & severe & stamped & $3.075\times10^{-14}$ over $784$ pairs; catches a $3\leftrightarrow4$ swap \\
Reduced against full $\Robs$ & severe & stamped & $0.777768$ by three routes; catches the same swap \\
Deliberate fault injection & severe & stamped & three of four injected faults invisible to the conservation residual; the fourth moves it by thirteen orders \\
Line-weighted against shell ratio & severe & stamped & ratio $0.948$ to $1.000$ over $784$ pairs, below unity at every one; census $44$ against $45$ \\
Escape-factor cross-check & severe & stamped & $\sigma_0$ against an independent module, \SI{0.059}{\percent} apart against a \SI{1}{\percent} raising tolerance \\
Pinned input hashes & severe & stamped & both SHA-256 match byte for byte \\
Recomputation guards on the figures & severe & code & every plotted row rebuilt from $\Lmat$ at $10^{-9}$, raises \\
Truncation convergence of $\ffrac{3}-\ffrac{4}$ & severe & stamped & increments one-signed and geometric at $245$ of $280$ defended points, ratio median $0.67$; extrapolated \SI{-0.46}{\percent} at the benchmark, \SI{+1.6}{\percent} at the cold corner, within \SIrange{-0.82}{+1.03}{\percent} on the defended set. Tests one quantity, not the state space \\
Gate~A, detailed balance & tautological & stamped & $<\SI{0.01}{\percent}$, but $K_{\rm deexc}$ was derived from $K_{\rm exc}$ \\
Column-sum conservation & tautological & stamped & $3.61\times10^{-11}$; see the fault-injection row for what it cannot see \\
The $\tanh$ bound & tautological & notes & holds at all $248$ points, and holds on corrupted input \\
Gate~E, timescale hierarchy & tautological & stamped & $\Msep \ge 86.8$ against a threshold of $1$ \\
Gate~B, coronal limit & weak & stamped & worst convergence $0.0034$ against a tolerance of $0.5$ \\
Gate~C, Saha ceiling & weak & stamped & monotone and below Saha; approach fraction \SI{1.5}{\percent}, excluded from the flag \\
\midrule
Fujimoto App.~4B timescales & external, passes & stamped & $0.34\times$ and $1.77\times$ at a grid point his example lands on \\
ACD96 against this model's recombination coefficient & external, passes & stamped & all $400$ points within \SI{11}{\percent}; $\eta$ from $0.894$ to $1.009$, median $0.982$, \SI{3.3}{\percent} at the benchmark (logarithmic read) \\
Fujimoto Table~4.1, $r_0$ & external, passes at $\lg\nel = 21$ & stamped & \SIrange{0.1}{1}{\percent} for $n\ge4$ at the higher density; \SIrange{12}{14}{\percent} at $p=3,4$ at $\lg\nel = 18$; $r_0(2) = 0.7392$ against $0.730$ \\
Fujimoto Table~4.1(b), $r_1$ & external, \textbf{open} & stamped & $8.3\times$ low at $p=3$. Truncation, Lyman trapping and the $\ell$~closure each tested and each eliminated; unexplained (\S\ref{sec:fujimoto}) \\
Gate~D, SCD96 against ADAS & \textbf{FAIL} & stamped & $\eta \in [5.73,\ 8346]$, $0$ of $400$ within a factor $2$. Rebuilt on the ground-fed channel alone and read logarithmically it reaches $400$ of $400$, $\eta$ from $0.702$ to $1.004$, median $0.828$; the gate itself is unchanged \\
\bottomrule
\end{tabular}
\end{table}

\subsection{What the model may be used for}

Three statements are supported by the evidence above, and they are the ones
Chapter~\ref{ch:results} relies on.

\textbf{The equations are solved correctly.} The superposition residual, the
independent re-derivation of all six central results, the bit-for-bit
regeneration, and the grid-wide conditioning measurement together establish that
the linear algebra does what Chapter~\ref{ch:theory} says it does, to a
precision many orders of magnitude finer than any quantity the thesis reports.

\textbf{The input rates are traceable and benchmarked where the observable
lives; one derived coefficient there is not.} The rates feeding $n \le 4$, which is where $H_\alpha$ and $H_\beta$
originate, agree with an independent R-matrix calculation at a mean absolute
error of \SI{13.0}{\percent}, and the two transitions this thesis leans on
hardest, $1s\to2p$ and $2p\to3d$, agree to between \SI{2.4}{\percent} and
\SI{13.3}{\percent} across the temperature range. Above $n=4$ the atomic data are
less certain and the truncation is measurable. Section~\ref{sec:convergence}
tests one component of the mechanism, the sensitivity $\ffrac{3}-\ffrac{4}$,
and finds it converged to within \SI{1.03}{\percent} at the defended points
where the geometric tail can be summed, in a stamped run
(\texttt{validation/nmax\_downward\_scan/}); under the same truncation
$u^{\mathrm{CRE}}$ is converged to within \SI{0.13}{\percent} at the $257$
defended points where its increments are geometric, $\Delta$ to within
\SI{0.57}{\percent} and $\tauslow$ to within \SI{1.6}{\percent} at all $280$,
and $\Delta\ln u$ across the defended heating steps moves by at most
$2.3\times10^{-5}$ between $n_{\max}=14$ and $15$ against a median step of
$0.20$ (Appendix~\ref{app:numerics}). What is measured
for the reservoir is a scan rather than a truncation: weakening the ionisation
and three-body exchange of the top six shells eightfold moves $\epsplat$ by
\SI{0.7}{\percent} at the median over the defended pairs
(Section~\ref{sec:ion_recomb}). All of that concerns the rates going in. The
coefficients that come out at the two shells the observable uses are a separate
matter, and one of them is in open disagreement: the ground-fed $r_1(3)$ and
$r_1(4)$ of this model sit factors $8.3$ and $4.4$ below Fujimoto's tabulation
at the lowest density, which is a deficit in $a_3$ and $a_4$ themselves and so
in the sensitivity built from them (Section~\ref{sec:fujimoto}). Benchmarked
inputs are therefore not the same thing as a benchmarked susceptibility, and
this thesis can claim the first but not the second.

\textbf{The structural results are robust to the choices that could have
manufactured them.} Closing the ion boundary changes the persistence census by
one part in $325$. Scaling $\ell$-mixing over two decades moves the observable
by under \SI{0.3}{\percent} where the thesis makes claims. The single-slow-mode
estimate of Eq.~\eqref{eq:elm_bound} reproduces the true time-averaged error to
within \SI{2.2}{\percent} at \SI{100}{\micro\second} over the $448$ pairs above
\SI{2}{\electronvolt}.

\subsection{What it may not be used for}

\textbf{No effective coefficient from this model should be compared against a
database without first separating its two supply channels.} Gate~D as written
fails at every point on the grid because it compares a coefficient that mixes
both channels against a table that separates them (Section~\ref{sec:gate_d});
rebuilt channel by channel it agrees within a factor two at all $400$ points
for SCD and ACD, and the model's total effective ionisation coefficient remains
a quantity ADAS does not tabulate.

\textbf{No claim rests on the terminal shell}, whose population is high by a
factor consistent with truncation, or on the $\ell$-distribution inside the
bundled shells, which is assumed rather than measured.
Section~\ref{sec:bundling_check} tests that assumption: $\ell$-mixing licenses
it over most of the grid, and at the two highest densities below
\SI{3}{\electronvolt}, where the locally evaluated mixing loses to collisional
loss, detailed balance among shells within $10^{-4}$ of Saha--Boltzmann
equilibrium licenses it instead; the adjacent-shell rates themselves are not
$\ell$-blind, and no grid point escapes both criteria.

\textbf{Nothing here validates the physics that is absent.} Every check in this
chapter tests a $43$-state atomic hydrogen model against itself, against
arithmetic, and against other atomic hydrogen models. None of them tests
transport, molecular channels, or radiation trapping, because none of those
appears in $\Lmat$ at all. A model can pass every check in
Table~\ref{tab:ladder} and still be the wrong model for a divertor.
Chapter~\ref{ch:limits} is where that question is asked, and it is asked with
numbers rather than with hedges.

\medskip
The instrument is now characterised: what it computes, how precisely, on which
region of the grid, and where the evidence for it runs out. What it says is the
subject of Chapter~\ref{ch:results}.
